\documentclass[11pt,a4paper]{article}

% Packages
\usepackage[utf8]{inputenc}
\usepackage[T1]{fontenc}
\usepackage[ngerman]{babel}
\usepackage{amsmath,amssymb,amsthm}
\usepackage{mathtools}
\usepackage{booktabs}
\usepackage{longtable}
\usepackage{array}
\usepackage{graphicx}
\usepackage{hyperref}
\usepackage{cleveref}
\usepackage{enumitem}
\usepackage{geometry}
\usepackage{fancyhdr}
\usepackage{tikz}
\usepackage{epigraph}
\usepackage{csquotes}
\usepackage{tcolorbox}
\usepackage{xcolor}
\usepackage{listings}
\usetikzlibrary{arrows.meta, positioning, shapes.geometric, fit, backgrounds}

% Code Listings Style für BEATRIX
\definecolor{codegreen}{RGB}{0,128,0}
\definecolor{codegray}{RGB}{128,128,128}
\definecolor{codepurple}{RGB}{128,0,128}
\definecolor{backcolour}{RGB}{245,245,245}
\lstdefinestyle{pythonstyle}{
    backgroundcolor=\color{backcolour},
    commentstyle=\color{codegreen},
    keywordstyle=\color{codepurple},
    numberstyle=\tiny\color{codegray},
    stringstyle=\color{codepurple},
    basicstyle=\ttfamily\footnotesize,
    breakatwhitespace=false,
    breaklines=true,
    captionpos=b,
    keepspaces=true,
    numbers=left,
    numbersep=5pt,
    showspaces=false,
    showstringspaces=false,
    showtabs=false,
    tabsize=2,
    language=Python
}
\lstset{style=pythonstyle}

\geometry{margin=2.5cm}

% Epigraph styling
\setlength{\epigraphwidth}{0.85\textwidth}
\setlength{\epigraphrule}{0pt}

% Theorem environments
\theoremstyle{definition}
\newtheorem{axiom}{Axiom}[section]
\newtheorem{definition}{Definition}[section]
\theoremstyle{plain}
\newtheorem{theorem}{Theorem}[section]
\newtheorem{corollary}{Korollar}[theorem]

% Custom commands
\newcommand{\BCM}{\textsc{BCM}}
\newcommand{\META}{\textsc{META}}
\newcommand{\BEATRIX}{\textsc{BEATRIX}}
\newcommand{\module}[1]{\textsc{#1}}

% Header
\pagestyle{fancy}
\fancyhf{}
\rhead{META-Modul im BCM 2.2}
\lhead{FehrAdvice}
\rfoot{\thepage}

\title{%
    \Large \textbf{Das META-Modul: Axiomatische Grundlagen des Behavioral Change Model} \\[0.5em]
    \large Formale Spezifikation der ontologischen, epistemologischen, governancebezogenen und dynamischen Fundamente
}

\author{%
    BEATRIX Research Team \\
    FehrAdvice \& Partners AG \\[0.5em]
    \small Version 2.2 | Dezember 2024
}

\date{}

\begin{document}

\maketitle

\begin{abstract}
Das \META-Modul bildet das axiomatische Fundament des Behavioral Change Model (\BCM) 2.2. Es definiert 115 Axiome in vier Kategorien: Ontologie (ONT), Epistemologie (EPI), Governance (GOV) und Dynamik (DYN). Diese Axiome konstituieren die philosophischen und methodologischen Grundlagen, auf denen alle operativen BCM-Module (INU, KNU, IDN, LAMBDA) aufbauen. Das vorliegende Paper dokumentiert systematisch die Struktur des \META-Moduls, erläutert die Funktion jeder Axiom-Kategorie und zeigt die Verknüpfung mit dem \BEATRIX-System zur LLM-basierten Verhaltensmodellierung. Besonderer Fokus liegt auf der Integration von 38 kognitiven und methodologischen Axiomen (EPI-Kategorie), der 5-Ebenen-Hierarchie (META-MAKRO-MESO-MIKRO-IND), der Evidenzhierarchie für BCM-Parameter (Labor-Feld-Triangulation) sowie den Governance-Prinzipien für ethische Interventionsgestaltung.
\end{abstract}

\tableofcontents
\newpage

%==============================================================================
% KAPITEL
%==============================================================================

% Kapitel 1: Einführung
%==============================================================================
\section{Einführung}
%==============================================================================

\subsection{Motivation: Warum ein META-Modul?}

\epigraph{\textit{``The formulation of a problem is often more essential than its solution, which may be merely a matter of mathematical or experimental skill.''}}{--- Albert Einstein \& Leopold Infeld, \textit{The Evolution of Physics} (1938)}

Das Behavioral Change Model (\BCM) ist ein formal-mathematisches Framework zur Modellierung, Prognose und Beeinflussung menschlichen Verhaltens. Im Gegensatz zu rein empirischen Ansätzen basiert das \BCM{} auf einem axiomatischen Fundament, das explizite Annahmen über die Natur von Verhalten, Wissen, Erlaubtem und Veränderung macht.

Diese Entscheidung für Axiomatik ist nicht selbstverständlich. Die Verhaltensökonomie hat sich seit Kahneman und Tversky (1979) als empirische Disziplin verstanden -- als Sammlung von Anomalien, Effekten und Heuristiken. Doch genau diese empirische Stärke ist zugleich ihre theoretische Schwäche: Ohne explizite Fundamente bleibt die Disziplin eine lose Kollektion von Phänomenen, die sich weder systematisch integrieren noch maschinell operationalisieren lässt.

Das \META-Modul adressiert diese Lücke. Es macht explizit, was bisher implizit blieb -- und schafft damit die Voraussetzung für eine neue Generation verhaltensökonomischer Anwendungen.

%------------------------------------------------------------------------------
\subsubsection{Das Erbe impliziter Annahmen}
%------------------------------------------------------------------------------

\epigraph{\textit{``Economics is haunted by more fallacies than any other study known to man.''}}{--- Henry Hazlitt, \textit{Economics in One Lesson} (1946)}

Die Geschichte der Wirtschaftswissenschaften ist eine Geschichte impliziter Annahmen. Adam Smith sprach 1776 von der ``invisible hand'', ohne je zu definieren, unter welchen Bedingungen sie funktioniert. Léon Walras formalisierte 1874 das allgemeine Gleichgewicht, ohne die psychologischen Grundlagen seiner Akteure zu explizieren. Und selbst Kenneth Arrow und Gérard Debreu, die 1954 die Existenz eines Gleichgewichts bewiesen, taten dies auf Basis von Annahmen, die sie selbst als ``heroisch'' bezeichneten.

Der \textit{Homo Oeconomicus} -- jene Fiktion eines vollständig rationalen, nutzenmaximierenden Agenten -- war nie als empirische Beschreibung gedacht. Er war ein \textit{Als-Ob-Konstrukt}, eine methodologische Vereinfachung. Doch über Jahrzehnte verwechselten Ökonomen das Modell mit der Realität. Milton Friedman verteidigte diese Praxis 1953 mit seinem berühmten ``As-If''-Argument:

\begin{quote}
\textit{``The relevant question to ask about the `assumptions' of a theory is not whether they are descriptively `realistic,' for they never are, but whether they are sufficiently good approximations for the purpose in hand.''}
\end{quote}

Das Problem: Friedman gab nie an, \textit{welche} Annahmen der Homo Oeconomicus eigentlich macht. Sie blieben implizit -- und damit unkritisierbar.

Die Behavioral Economics Revolution ab 1979 war eine Reaktion auf diese impliziten Annahmen. Kahneman und Tversky dokumentierten systematische Abweichungen vom Rationalmodell: Loss Aversion, Framing Effects, Anchoring, Overconfidence. Doch -- und das ist der entscheidende Punkt -- sie \textit{ersetzten} das implizite Axiomensystem nicht. Sie \textit{kritisierten} es lediglich.

Richard Thaler, der diese Kritik in die Ökonomie trug, beschrieb das Dilemma in seinem Nobelpreis-Vortrag 2017:

\begin{quote}
\textit{``We have learned a lot about how people deviate from the predictions of economic theory, but we have not yet built a coherent alternative theory.''}
\end{quote}

Dieses ``noch nicht'' dauert nun fast 50 Jahre. Die Verhaltensökonomie hat Anomalien gesammelt wie ein Kuriositätenkabinett -- aber kein neues Fundament gebaut. Das \META-Modul ist der Versuch, dieses Fundament nachzuliefern.

\paragraph{Das Problem der Parametervariation.}

Ein konkretes Beispiel illustriert die Konsequenzen impliziter Annahmen: Loss Aversion. Kahneman und Tversky (1979) schätzten den Loss-Aversion-Koeffizienten $\lambda$ auf etwa 2.0 -- Verluste wiegen doppelt so schwer wie Gewinne. Doch nachfolgende Studien fanden erhebliche Variation:

\begin{table}[h]
\centering
\caption{Variation des Loss-Aversion-Koeffizienten $\lambda$ in der Literatur}
\begin{tabular}{@{}llc@{}}
\toprule
\textbf{Studie} & \textbf{Kontext} & \textbf{$\lambda$} \\
\midrule
Kahneman \& Tversky (1979) & Hypothetische Lotterien & 2.25 \\
Tversky \& Kahneman (1992) & Cumulative Prospect Theory & 2.25 \\
Abdellaoui et al. (2007) & Reale Auszahlungen & 1.43 \\
Gächter et al. (2022) & Meta-Analyse & 1.5--2.5 \\
Yechiam (2019) & Kritische Analyse & 1.0--3.0 \\
\bottomrule
\end{tabular}
\end{table}

Ist $\lambda = 2.0$ nun ein Naturgesetz, ein Durchschnitt oder ein Artefakt? Die Literatur schweigt, weil die \textit{ontologische Frage} nie gestellt wurde: Was \textit{ist} Loss Aversion eigentlich -- ein stabiler Persönlichkeitstrait, ein kontextabhängiger Zustand oder ein statistisches Aggregat? Ohne explizite Ontologie bleibt diese Frage unbeantwortbar.

Das \META-Modul adressiert dies in Axiom MA-EPI-15:

\begin{quote}
\textbf{Loss Aversion (MA-EPI-15):} $|U(\text{Loss})| > |U(\text{Gain})|$ für $|\text{Loss}| = |\text{Gain}|$. Der Koeffizient $\lambda$ ist kontextabhängig und liegt typischerweise im Bereich $[1.5, 2.5]$.
\end{quote}

Der entscheidende Unterschied: Das Axiom macht die Kontextabhängigkeit \textit{explizit}. Es behauptet nicht, dass $\lambda = 2.0$ universal gilt, sondern dass es einen \textit{Bereich} gibt -- und dass die \BEATRIX-Schätzung diesen Bereich respektieren muss.

%------------------------------------------------------------------------------
\subsubsection{Die axiomatische Tradition in der Wissenschaft}
%------------------------------------------------------------------------------

\epigraph{\textit{``The method of `postulating' what we want has many advantages; they are the same as the advantages of theft over honest toil.''}}{--- Bertrand Russell, \textit{Introduction to Mathematical Philosophy} (1919)}

Die Idee, Wissen auf expliziten Grundannahmen aufzubauen, ist nicht neu. Sie bildet den Kern des wissenschaftlichen Denkens seit der Antike.

\paragraph{Euklids Elemente (ca. 300 v. Chr.).}

Euklid begann seine \textit{Elemente} mit fünf Postulaten -- Annahmen, die nicht bewiesen, sondern gesetzt werden. Das berühmteste ist das Parallelenpostulat: \textit{Durch einen Punkt außerhalb einer Geraden gibt es genau eine Parallele.} Aus diesen fünf Postulaten leitete Euklid 465 Sätze ab -- ein Gebäude von atemberaubender Kohärenz.

Für zwei Jahrtausende galt Euklid als Modell wissenschaftlicher Strenge. Erst im 19. Jahrhundert zeigten Lobatschewski, Bolyai und Riemann, dass man das Parallelenpostulat \textit{ändern} kann -- und konsistente, aber radikal andere Geometrien erhält. Diese Entdeckung war nur möglich, \textit{weil} Euklid seine Annahmen explizit gemacht hatte.

Die Lektion für die Verhaltensökonomie: Implizite Annahmen können nicht systematisch variiert werden. Man weiß nicht einmal, dass man sie macht.

\paragraph{Die Peano-Axiome (1889).}

Giuseppe Peano formalisierte die natürlichen Zahlen in fünf Axiomen. Das erste: \textit{0 ist eine natürliche Zahl.} Das zweite: \textit{Jede natürliche Zahl hat einen Nachfolger.} Aus diesen minimalen Annahmen folgt die gesamte Arithmetik.

Die Eleganz der Peano-Axiome liegt in ihrer \textit{Minimalität}: Sie machen genau die Annahmen, die nötig sind -- nicht mehr und nicht weniger. Das \META-Modul strebt dasselbe Ideal an: 115 Axiome, die das Minimum an Annahmen darstellen, um menschliches Verhalten konsistent zu modellieren.

\paragraph{Kolmogorovs Wahrscheinlichkeitstheorie (1933).}

Andrei Kolmogorov axiomatisierte die Wahrscheinlichkeitstheorie in drei Grundannahmen:
\begin{enumerate}
    \item Die Wahrscheinlichkeit eines Ereignisses ist nicht-negativ.
    \item Die Wahrscheinlichkeit des sicheren Ereignisses ist 1.
    \item Die Wahrscheinlichkeit der Vereinigung disjunkter Ereignisse ist die Summe ihrer Einzelwahrscheinlichkeiten.
\end{enumerate}

Vor Kolmogorov war ``Wahrscheinlichkeit'' ein vager Begriff -- manchmal Häufigkeit, manchmal Plausibilität, manchmal Wettquote. Nach Kolmogorov war klar, welche Operationen erlaubt sind und welche nicht.

Die Verhaltensökonomie befindet sich 2024 in der Situation der Wahrscheinlichkeitstheorie vor 1933: Ein Sammelsurium von intuitiven Begriffen -- ``Nutzen'', ``Bias'', ``Heuristik'' -- ohne klare axiomatische Fundierung.

\paragraph{Von Neumann-Morgenstern Expected Utility (1944).}

John von Neumann und Oskar Morgenstern axiomatisierten in ihrer \textit{Theory of Games and Economic Behavior} die Entscheidungstheorie. Ihre Axiome -- Vollständigkeit, Transitivität, Kontinuität, Unabhängigkeit -- definierten erstmals präzise, was ``rationale Präferenzen'' bedeuten.

Kahneman und Tversky (1979) zeigten, dass Menschen das Unabhängigkeitsaxiom systematisch verletzen (Allais-Paradox). Doch sie formulierten keine \textit{alternativen} Axiome. Prospect Theory blieb eine deskriptive Theorie ohne axiomatische Fundierung -- ein ``Addon'' zur Expected Utility, kein Ersatz.

Daniel Ellsberg bemerkte in seinem wegweisenden Paper von 1961:

\begin{quote}
\textit{``If people's actual behavior is persistently inconsistent with [the axioms], then the axioms must be modified or abandoned.''}
\end{quote}

Das ``Modifizieren oder Aufgeben'' ist nie geschehen. Die Verhaltensökonomie lebt seit 60 Jahren mit gebrochenen Axiomen, ohne neue zu formulieren. Das \META-Modul schließt diese Lücke.

\paragraph{Arrows Unmöglichkeitstheorem (1951).}

Kenneth Arrow bewies, dass es keine Wahlregel gibt, die fünf scheinbar vernünftige Axiome gleichzeitig erfüllt (Pareto-Effizienz, Unabhängigkeit irrelevanter Alternativen, Nicht-Diktatur, universelle Domäne, Transitivität). Dieses ``Unmöglichkeitstheorem'' war ein intellektueller Schock -- aber nur möglich, \textit{weil} Arrow seine Annahmen explizit gemacht hatte.

Die Lektion: Axiomatisierung führt manchmal zu unbequemen Erkenntnissen. Aber ohne sie bleibt man in behaglicher Unwissenheit.

%------------------------------------------------------------------------------
\subsubsection{Von der Philosophie zur Infrastruktur: Die drei Zeitalter der Ontologie}
%------------------------------------------------------------------------------

\epigraph{\textit{``The limits of my language mean the limits of my world.''}}{--- Ludwig Wittgenstein, \textit{Tractatus Logico-Philosophicus} (1921)}

Die Frage drängt sich auf: Warum entsteht das \META-Modul ausgerechnet jetzt -- 2024, nicht 1984 oder 2004? Die Antwort liegt in einem fundamentalen Wandel der Rolle von Ontologie: \textbf{Was früher akademischer Luxus war, ist heute operative Infrastruktur.}

Um diesen Wandel zu verstehen, müssen wir drei distinkte Epochen unterscheiden: das vor-digitale Zeitalter (bis ca. 1990), das digitale Zeitalter vor Machine Learning (1990--2020), und die LLM-Ära (ab 2022).

%..............................................................................
\paragraph{Epoche I: Das vor-digitale Zeitalter (bis ca. 1990)}
%..............................................................................

\subparagraph{Wie Verhaltensökonomie praktiziert wurde.}

In der vor-digitalen Ära war verhaltensökonomische Arbeit ein zutiefst \textit{menschlicher} Prozess. Die gesamte Wertschöpfungskette -- von der Problemidentifikation über die Analyse bis zur Intervention -- lief über menschliche Köpfe, menschliche Gespräche, menschliche Intuition.

\textbf{Fallbeispiel 1: Richard Thaler und die Pensionsfonds (1980er Jahre)}

\begin{quote}
\textit{Richard Thaler berät ein Unternehmen zur Erhöhung der Pensionsfonds-Teilnahme. Sein Vorgehen:}
\begin{enumerate}
    \item \textit{Er führt 20 Interviews mit Mitarbeitern -- persönlich, mit Notizblock}
    \item \textit{Er beobachtet das Anmeldeverfahren vor Ort}
    \item \textit{Er diskutiert mit dem HR-Team über Barrieren}
    \item \textit{Er entwickelt eine Hypothese: ``Mitarbeiter sind nicht irrational -- sie sind busy und das Formular ist zu kompliziert''}
    \item \textit{Er schlägt vor: ``Macht die Einschreibung automatisch, mit Opt-Out''}
    \item \textit{Er präsentiert die Idee dem Management -- in einem einstündigen Meeting, mit Folien auf Overhead-Projektor}
\end{enumerate}
\end{quote}

Beachten Sie, was in diesem Prozess \textit{nicht} vorkam:
\begin{itemize}
    \item Keine formale Definition von ``Gegenwartsorientierung'' oder ``Present Bias''
    \item Keine Quantifizierung des $\beta$-Parameters
    \item Keine Spezifikation der Randbedingungen (für welche Mitarbeiter? in welchem kulturellen Kontext?)
    \item Keine maschinenlesbare Dokumentation
\end{itemize}

Der Erfolg hing vollständig von Thalers \textit{tacit knowledge} ab -- seinem Gespür, seiner Erfahrung, seiner Überzeugungskraft.

\textbf{Fallbeispiel 2: Kahneman \& Tversky im Labor (1970er Jahre)}

Amos Tversky und Daniel Kahneman entwickelten die Prospect Theory in jahrelanger Zusammenarbeit. Ihr Arbeitsprozess:
\begin{itemize}
    \item \textbf{Gedankenexperimente}: Sie erfanden Szenarien (``Stellen Sie sich vor, Sie könnten €3.000 sicher bekommen oder mit 80\% Chance €4.000...'') und diskutierten ihre eigenen Intuitionen
    \item \textbf{Kleine Stichproben}: Typische Studien hatten 50-100 Teilnehmer, oft Studenten
    \item \textbf{Handschriftliche Analyse}: Daten wurden auf Papier erfasst und mit Taschenrechner ausgewertet
    \item \textbf{Konzeptuelle Diskussion}: Stundenlange Gespräche über ``Was bedeutet es, risikoavers zu sein?''
\end{itemize}

Das Ergebnis -- die Prospect Theory (1979) -- war ein Meilenstein. Aber sie enthielt \textit{keine explizite Ontologie}:
\begin{itemize}
    \item Der ``Entscheidungsträger'' wurde nicht formal definiert (Individuum? Haushalt? Firma?)
    \item ``Referenzpunkt'' blieb vage (aktueller Besitz? Erwartung? sozialer Vergleich?)
    \item Die Parameter $\alpha$, $\beta$, $\lambda$ wurden geschätzt, aber ihre Stabilität (Trait vs. State?) blieb offen
    \item ``Framing'' wurde beschrieben, aber nicht taxonomiert
\end{itemize}

\textbf{Warum das funktionierte}: Die Paper wurden von Menschen gelesen -- Ökonomen, Psychologen, Praktikern -- die den fehlenden Kontext aus ihrem eigenen Wissen ergänzten. Wenn Kahneman ``Referenzpunkt'' schrieb, verstand der kompetente Leser \textit{ungefähr}, was gemeint war.

\textbf{Fallbeispiel 3: Wissenstransfer durch Meister-Lehrling-Beziehung}

Wie wurde verhaltensökonomisches Wissen in dieser Epoche weitergegeben?
\begin{itemize}
    \item \textbf{Doktoranden-Betreuung}: Ein Doktorand arbeitete 4-6 Jahre mit einem Professor. Er lernte nicht nur Methoden, sondern auch ``wie man denkt'' -- durch Osmose
    \item \textbf{Konferenzen}: Jährliche Treffen (AEA, SJDM) ermöglichten persönlichen Austausch. Missverständnisse wurden in Kaffeepausen geklärt
    \item \textbf{Bücher}: Thalers ``Quasi Rational Economics'' (1991) oder Kahneman, Slovic \& Tversky's ``Judgment under Uncertainty'' (1982) waren \textit{Prosa} -- keine Axiomensysteme
    \item \textbf{Informelle Netzwerke}: ``Ruf Richard an und frag ihn'' war eine legitime Methode der Wissensgewinnung
\end{itemize}

\textbf{Die Rolle der Ontologie in Epoche I:}

Ontologie war ein Thema für Philosophen -- nicht für Praktiker. Wer die Frage stellte ``Was \textit{ist} ein Referenzpunkt ontologisch?'' wurde als übertrieben theoretisch angesehen. Die pragmatische Antwort: ``Wir wissen es, wenn wir es sehen.''

\textbf{Zusammenfassung Epoche I:}
\begin{table}[h]
\centering
\begin{tabular}{@{}ll@{}}
\toprule
\textbf{Merkmal} & \textbf{Ausprägung} \\
\midrule
Wissensträger & Menschen (tacit knowledge) \\
Dokumentation & Prosa, Papers, Bücher \\
Wissenstransfer & Meister-Lehrling, Konferenzen \\
Skala & Dutzende Fälle pro Experte \\
Ontologie & Implizit, ``wir wissen es wenn wir es sehen'' \\
Funktionierte weil... & Menschen interpretierten und ergänzten \\
\bottomrule
\end{tabular}
\end{table}

%..............................................................................
\paragraph{Epoche II: Das digitale Zeitalter vor Machine Learning (1990--2020)}
%..............................................................................

\subparagraph{Der Übergang zu datengetriebener Arbeit.}

Die Digitalisierung veränderte die Verhaltensökonomie graduell aber fundamental. Neue Werkzeuge ermöglichten neue Arbeitsweisen -- aber die grundlegende Abhängigkeit von menschlicher Interpretation blieb.

\textbf{Was sich änderte:}
\begin{itemize}
    \item \textbf{Größere Datensätze}: Statt 50 Interviews nun 50.000 Datenpunkte
    \item \textbf{Statistische Software}: SPSS (1968), Stata (1985), R (1993) ermöglichten komplexe Analysen
    \item \textbf{Online-Experimente}: Amazon Mechanical Turk (2005) revolutionierte die Datenerhebung
    \item \textbf{Institutionalisierung}: Nudge Units entstanden (BIT UK 2010, US SBST 2014, hunderte weltweit)
    \item \textbf{A/B-Testing}: Tech-Firmen testeten Verhaltensinterventionen in Echtzeit
\end{itemize}

\textbf{Fallbeispiel 4: Das Behavioural Insights Team UK (2010-2020)}

Das BIT -- die weltweit erste staatliche Nudge Unit -- revolutionierte die Anwendung von Verhaltensökonomie. Ein typisches Projekt:

\begin{quote}
\textit{Das BIT will die Steuerzahlungsquote erhöhen. Vorgehen:}
\begin{enumerate}
    \item \textit{Hypothese: Soziale Normen motivieren Compliance (``Die meisten Ihrer Nachbarn haben bereits bezahlt'')}
    \item \textit{Design: Randomized Controlled Trial mit 100.000 Steuerzahlern}
    \item \textit{Treatment-Gruppen: Standard-Brief vs. Brief mit Soziale-Norm-Botschaft}
    \item \textit{Analyse: Regressionsanalyse in Stata}
    \item \textit{Ergebnis: 15\% höhere Zahlungsrate in Treatment-Gruppe}
    \item \textit{Publikation: Working Paper mit p-Werten und Konfidenzintervallen}
\end{enumerate}
\end{quote}

\textbf{Was fehlte -- auch beim BIT:}
\begin{itemize}
    \item \textbf{Mechanismus-Spezifikation}: \textit{Warum} wirkte die soziale Norm? Deskriptive Norm? Injunktive Norm? Gruppendruck? Informationsupdate?
    \item \textbf{Moderator-Analyse}: Für \textit{wen} wirkte sie? (Alter? Einkommen? Persönlichkeit? Vorgeschichte?)
    \item \textbf{Kontextbedingungen}: \textit{Wann} funktioniert Social Proof? (In welchen Kulturen? Bei welchen Verhaltenstypen?)
    \item \textbf{Transfer-Regeln}: Kann man das Ergebnis auf Müllentsorgung, Energiesparen, Impfungen übertragen?
\end{itemize}

Das BIT produzierte \textit{Fakten} (``dieser Nudge wirkte hier''), aber kein \textit{transferierbares Wissen}.

\textbf{Fallbeispiel 5: Google und die 41 Blautöne (2009)}

Ein berühmtes Beispiel für datengetriebene Verhaltensoptimierung:

\begin{quote}
\textit{Google testete 41 verschiedene Blautöne für Links, um die Click-Through-Rate zu optimieren. Ergebnis: Der optimale Blauton brachte \$200 Millionen zusätzlichen Jahresumsatz.}
\end{quote}

Das Beispiel illustriert die Stärken \textit{und} Grenzen der Epoche II:
\begin{itemize}
    \item \textbf{Stärke}: Präzise Messung, große Stichprobe, klares Ergebnis
    \item \textbf{Schwäche 1}: Kein Verständnis des Mechanismus -- \textit{warum} ist dieser Blauton besser?
    \item \textbf{Schwäche 2}: Keine Generalisierbarkeit -- gilt das auch für Bing? Für eine Bank-Website? In Japan?
    \item \textbf{Schwäche 3}: Keine Theorie -- das Ergebnis trägt nicht zum kumulativen Wissen bei
\end{itemize}

\textbf{Die Krise wird sichtbar: Replikationsprobleme}

Ab 2015 wurde deutlich, dass etwas grundlegend nicht stimmte:

\begin{enumerate}
    \item \textbf{DellaVigna \& Linos (2022)}: Analysierten 126 RCTs von Nudge Units. Ergebnis: Durchschnittliche Effektgröße in der Praxis = \textbf{1.4 Prozentpunkte} -- verglichen mit 8.7 Prozentpunkten in akademischen Publikationen. Die Praxis-Effekte waren nur \textbf{1/6} der publizierten Effekte.

    \item \textbf{Warum?} Die Autoren vermuten: Publication Bias, aber auch \textit{unterschiedliche Kontextbedingungen}. Ein Nudge, der in einem Paper mit Studenten funktionierte, scheitert bei einer anderen Bevölkerung.

    \item \textbf{Das fundamentale Problem}: Ohne explizite Ontologie fehlte die Sprache, um zu analysieren, \textit{warum} ein Nudge in Kontext A funktioniert und in Kontext B nicht.
\end{enumerate}

\textbf{Zusammenfassung Epoche II:}
\begin{table}[h]
\centering
\begin{tabular}{@{}ll@{}}
\toprule
\textbf{Merkmal} & \textbf{Ausprägung} \\
\midrule
Wissensträger & Menschen + statistische Software \\
Dokumentation & Papers mit p-Werten, RCT-Reports \\
Wissenstransfer & Publikationen, Toolkits, Cookbooks \\
Skala & Tausende bis Millionen Datenpunkte \\
Ontologie & Implizit, versteckt in Variablennamen \\
Funktionierte weil... & Menschen blieben in der Interpretations-Schleife \\
Scheiterte weil... & Kein Transfer, keine Kumulation, Replikationskrise \\
\bottomrule
\end{tabular}
\end{table}

\textbf{Die zentrale Erkenntnis aus Epoche II:}

Man kann mit Big Data und RCTs sehr präzise messen, \textit{dass} etwas wirkt. Aber ohne explizite Ontologie bleibt unklar, \textit{warum} es wirkt, \textit{für wen} es wirkt, und \textit{wann} es wirkt. Jede Studie ist eine Insel. Kumulatives Wissen entsteht nicht.

%..............................................................................
\paragraph{Epoche III: Das LLM-Zeitalter (ab 2022)}
%..............................................................................

\subparagraph{Der Paradigmenwechsel.}

Die Veröffentlichung von GPT-3 (2020), ChatGPT (2022) und GPT-4 (2023) markiert eine Diskontinuität. Zum ersten Mal sind Maschinen verfügbar, die:
\begin{itemize}
    \item Natürliche Sprache verstehen und generieren
    \item Komplexe Reasoning-Aufgaben durchführen
    \item Mit unstrukturierten Daten (Text, Beschreibungen) arbeiten
    \item In Echtzeit auf neue Anfragen reagieren
\end{itemize}

Diese Fähigkeiten eröffnen revolutionäre Möglichkeiten für die Verhaltensökonomie -- aber sie erzwingen auch eine fundamentale Transformation: \textbf{Was implizit war, muss explizit werden.}

\textbf{Fallbeispiel 9: ChatGPT ohne Ontologie -- Ein Experiment}

Wir haben folgendes Experiment durchgeführt (März 2024):

\begin{quote}
\textbf{Prompt an ChatGPT-4:} ``Anna ist 35 Jahre alt, arbeitet als Marketing-Managerin, verdient €80.000/Jahr, ist verheiratet, hat zwei Kinder, und beschreibt sich als `umweltbewusst aber pragmatisch'. Schätze ihren Loss Aversion Parameter $\lambda$.''
\end{quote}

\textbf{Antwort ChatGPT (Run 1):}
\begin{quote}
``Basierend auf Annas Profil würde ich ihren Loss Aversion Parameter auf etwa $\lambda = 2.3$ schätzen. Als Marketing-Managerin ist sie wahrscheinlich risikoaverser als der Durchschnitt, und mit zwei Kindern hat sie mehr zu verlieren...''
\end{quote}

\textbf{Antwort ChatGPT (Run 2, identischer Prompt):}
\begin{quote}
``Annas $\lambda$ dürfte bei etwa $\lambda = 1.8$ liegen. Als berufstätige Mutter hat sie gelernt, pragmatische Entscheidungen zu treffen, was auf moderate Loss Aversion hindeutet...''
\end{quote}

\textbf{Antwort ChatGPT (Run 3, identischer Prompt):}
\begin{quote}
``Ich schätze $\lambda \approx 2.5$ für Anna. Ihre Umweltbewusstheit deutet auf Zukunftsorientierung hin, was oft mit höherer Loss Aversion korreliert...''
\end{quote}

\textbf{Was ist problematisch?}
\begin{enumerate}
    \item \textbf{Inkonsistenz}: Drei Runs, drei verschiedene Werte (1.8, 2.3, 2.5)
    \item \textbf{Keine Base Rate}: Keine der Antworten erwähnt die empirische Verteilung von $\lambda$ in der Population
    \item \textbf{Pseudo-Reasoning}: Die ``Begründungen'' klingen plausibel, sind aber post-hoc erfunden
    \item \textbf{Keine Unsicherheitsquantifizierung}: Keine Konfidenzintervalle, keine Varianzangabe
    \item \textbf{Keine Mechanismus-Spezifikation}: Warum sollte ``Umweltbewusstheit'' mit Loss Aversion korrelieren?
\end{enumerate}

\textbf{Fallbeispiel 11: BEATRIX mit META -- Der Unterschied}

Derselbe Fall mit \BEATRIX{} und expliziter \META-Ontologie:

\begin{quote}
\textbf{Input}: Persona-Beschreibung von Anna (wie oben)

\textbf{BEATRIX-Prozess}:
\begin{enumerate}
    \item \textbf{PreFlight Check}: Ist die Beschreibung vollständig? (FEPSDE-Dimensionen abgedeckt?)
    \item \textbf{Base Rate Anchoring}: Starte mit empirischen Priors:
    \begin{itemize}
        \item $\lambda \sim \mathcal{N}(2.25, 0.5)$ basierend auf MA-EPI-15
        \item P(Conditional Cooperator) = 0.50 basierend auf MA-EPI-26
    \end{itemize}
    \item \textbf{Social Desirability Korrektur}: ``umweltbewusst'' → prüfe \textit{Verhaltensevidenz}, nicht nur Selbstbeschreibung
    \item \textbf{Ensemble Generation}: 5 parallele LLM-Runs
    \item \textbf{Aggregation}: $\hat{\lambda} = \text{Median}(\lambda_1, ..., \lambda_5) = 2.2$
    \item \textbf{Unsicherheitsquantifizierung}: $\lambda \in [1.9, 2.5]$ mit 80\% Konfidenz
    \item \textbf{Bimodality Check}: Sind die 5 Schätzungen unimodal (Unsicherheit) oder bimodal (Ambivalenz)?
\end{enumerate}
\end{quote}

\textbf{Der Unterschied}:
\begin{itemize}
    \item \textbf{Konsistenz}: Derselbe Input ergibt denselben Output
    \item \textbf{Verankerung}: Empirische Base Rates statt willkürliche Schätzungen
    \item \textbf{Unsicherheit}: Explizite Konfidenzintervalle
    \item \textbf{Bias-Korrektur}: Social Desirability wird aktiv adressiert
    \item \textbf{Transparenz}: Der gesamte Reasoning-Prozess ist dokumentiert
\end{itemize}

\textbf{Warum explizite Ontologie jetzt NOTWENDIG ist:}

\begin{enumerate}
    \item \textbf{LLMs haben kein Tacit Knowledge}

    Ein LLM hat das gesamte Internet gelesen -- aber es hat nicht das ``Gefühl'' eines erfahrenen Praktikers. \textbf{Was nicht expliziert wurde, kann das LLM nicht systematisch anwenden.}

    \item \textbf{Code-Interpreter erfordern Typen und Schemas}

    Moderne LLM-Systeme wie \BEATRIX{} integrieren Code-Ausführung. Das erfordert formale Definitionen. \textbf{Ohne ontologische Festlegung kann kein Code geschrieben werden.}

    \item \textbf{JSON-Schemas erfordern Entitäts-Definition}

    \BEATRIX{} arbeitet mit strukturierten Outputs (Pydantic-Schemas). Das erfordert Antworten auf ontologische Fragen:
    \begin{itemize}
        \item Welche Entitäten existieren? → MA-ONT-01 bis MA-ONT-06
        \item Welche Attribute haben sie? → MA-ONT-08 (FEPSDE)
        \item Welche Beziehungen bestehen? → MA-ONT-07 (Drei-Komponenten-Struktur)
        \item Welche Wertebereiche sind gültig? → MA-EPI-15 ($\lambda \in [1.5, 2.5]$)
    \end{itemize}

    \textbf{Das ist nichts anderes als angewandte Ontologie.}

    \item \textbf{Prompt Engineering ist implizite Ontologie}

    Jeder System-Prompt macht ontologische Annahmen -- ob der Autor es weiß oder nicht. \textbf{Wer keine explizite Ontologie hat, macht trotzdem ontologische Annahmen -- nur unkontrolliert.}
\end{enumerate}

\textbf{Zusammenfassung Epoche III:}
\begin{table}[h]
\centering
\begin{tabular}{@{}ll@{}}
\toprule
\textbf{Merkmal} & \textbf{Ausprägung} \\
\midrule
Wissensträger & LLMs + Code + Menschen (Aufsicht) \\
Dokumentation & Strukturierte Schemas, APIs, Axiome \\
Wissenstransfer & Codierte Systeme, trainierbare Modelle \\
Skala & Millionen Interaktionen in Echtzeit \\
Ontologie & \textbf{Operative Infrastruktur} \\
Funktioniert nur wenn... & Explizite Ontologie vorhanden \\
Scheitert wenn... & LLM ``frei'' operiert ohne Constraints \\
\bottomrule
\end{tabular}
\end{table}

%..............................................................................
\paragraph{Die historische Analogie: Kolmogorovs Axiomatisierung}
%..............................................................................

Ein instruktiver Vergleich: Vor 1933 war ``Wahrscheinlichkeit'' ein vager Begriff. Verschiedene Autoren meinten Verschiedenes:
\begin{itemize}
    \item \textbf{Frequentisten}: Wahrscheinlichkeit = Grenzwert relativer Häufigkeit
    \item \textbf{Bayesianer}: Wahrscheinlichkeit = Grad der Überzeugung
    \item \textbf{Logizisten}: Wahrscheinlichkeit = logische Relation zwischen Propositionen
\end{itemize}

Menschen konnten trotzdem mit Wahrscheinlichkeit arbeiten, weil sie intuitiv wussten, was in welchem Kontext gemeint war.

\textbf{Aber: Computer konnten nicht mit vager Wahrscheinlichkeit rechnen.}

Andrei Kolmogorov (1933) axiomatisierte die Wahrscheinlichkeitstheorie in drei Grundannahmen:
\begin{enumerate}
    \item $P(A) \geq 0$ für alle Ereignisse $A$
    \item $P(\Omega) = 1$ für den Ereignisraum $\Omega$
    \item $P(A \cup B) = P(A) + P(B)$ für disjunkte $A, B$
\end{enumerate}

Diese Axiome sagten nicht, was Wahrscheinlichkeit ``wirklich ist'' -- sie definierten, wie man mit ihr \textit{rechnen} kann. Erst danach war probabilistische Software möglich.

\textbf{Dieselbe Transformation passiert jetzt mit Verhaltensökonomie:}
\begin{itemize}
    \item Vor LLMs: Intuitive, implizite Behavioral Science
    \item Mit LLMs: Axiomatische, explizite Behavioral Science (= BCM META)
\end{itemize}

Das \META-Modul ist für Verhaltensökonomie, was Kolmogorovs Axiome für Wahrscheinlichkeit waren: die Voraussetzung für maschinelle Operationalisierung.

%..............................................................................
\paragraph{Zusammenfassung: Der Statuswandel der Ontologie}
%..............................................................................

\begin{table}[h]
\centering
\caption{Die drei Epochen der Ontologie in der Verhaltensökonomie}
\begin{tabular}{@{}p{3cm}p{4cm}p{4cm}p{4cm}@{}}
\toprule
\textbf{Aspekt} & \textbf{Vor-digital (--1990)} & \textbf{Digital (1990--2020)} & \textbf{LLM-Ära (2022--)} \\
\midrule
\textbf{Wissensträger} & Menschen mit tacit knowledge & Menschen + statistische Software & Menschen + LLMs + Code \\
\textbf{Skala} & Dutzende Fälle & Tausende Datenpunkte & Millionen Interaktionen \\
\textbf{Ambiguitäts-toleranz} & Hoch (Menschen interpretieren) & Mittel (Zahlen sind präzise, Bedeutung vage) & Niedrig (Maschinen brauchen Präzision) \\
\textbf{Wissenstransfer} & Meister-Lehrling, Konferenzen & Papers, Cookbooks & Codierte Systeme, APIs \\
\textbf{Status der Ontologie} & Philosophischer Luxus & Methodologische Option & \textbf{Operative Notwendigkeit} \\
\textbf{Konsequenz ohne Ontologie} & Ineffizienz, aber funktional & Replikationsprobleme, Transfer-Schwierigkeiten & \textbf{System-Dysfunktion} \\
\bottomrule
\end{tabular}
\end{table}

\textbf{Die zentrale Einsicht:}

Wittgenstein schrieb: ``Die Grenzen meiner Sprache bedeuten die Grenzen meiner Welt.'' Für Menschen ist das eine philosophische Aussage. Für LLMs ist es \textbf{wörtlich wahr}:
\begin{itemize}
    \item Was nicht in Sprache (Text) expliziert wurde, existiert für das LLM nicht
    \item Das LLM kann nur über Entitäten reasonen, die benannt wurden
    \item Beziehungen, die nicht spezifiziert wurden, kann es nicht inferieren
\end{itemize}

Das \META-Modul ist der Versuch, die ``Welt'' der Verhaltensökonomie vollständig in Sprache zu fassen -- nicht als philosophische Übung, sondern als \textbf{operative Infrastruktur} für die LLM-basierte Verhaltensmodellierung.

%------------------------------------------------------------------------------
\subsubsection{Die vier Funktionen des META-Moduls}
%------------------------------------------------------------------------------

Das \META-Modul erfüllt vier distinkte, aber komplementäre Funktionen im \BCM-Ökosystem. Jede verdient eine detaillierte Betrachtung.

\paragraph{(i) Konsistenzgarantie.}

\epigraph{\textit{``A foolish consistency is the hobgoblin of little minds.''}}{--- Ralph Waldo Emerson, \textit{Self-Reliance} (1841)}

Emerson mochte Konsistenz verachten -- für formale Systeme ist sie existentiell. Ein inkonsistentes Axiomensystem ist wertlos: Aus einem Widerspruch folgt alles (ex falso quodlibet).

Das \META-Modul definiert die Konsistenzanforderungen für alle operativen Module:

\begin{itemize}
    \item Das \textbf{INU-Modul} (Individual Utility) muss konsistent sein mit den ontologischen Axiomen über Subjekte (MA-ONT-01), Nutzen (MA-ONT-03) und FEPSDE-Dimensionen (MA-ONT-08).

    \item Das \textbf{KNU-Modul} (Collective Utility) muss konsistent sein mit den Axiomen über Gruppen (MA-ONT-10), Normen (MA-ONT-11) und soziale Präferenzen.

    \item Das \textbf{IDN-Modul} (Identity) muss konsistent sein mit den Axiomen über multiple Identitäten (MA-ONT-09), Identitätsaktivierung (MA-ONT-17) und Identitätsevolution (MA-ONT-23).

    \item Das \textbf{LAMBDA-Modul} muss konsistent sein mit den dynamischen Axiomen über $\lambda$-Dynamik (MA-DYN-23) und Feedback-Schleifen (MA-ONT-21).
\end{itemize}

\textbf{Beispiel für Konsistenzprüfung:} Angenommen, ein Praktiker möchte das INU-Modul um eine siebte Nutzen-Dimension erweitern -- etwa ``Spiritual'' (S). Diese Erweiterung muss konsistent sein mit:
\begin{itemize}
    \item MA-ONT-08: Die Gewichtungssumme muss weiterhin $\sum \omega_d = 1$ ergeben.
    \item MA-EPI-01: Unvollständiges Wissen gilt auch für die neue Dimension.
    \item MA-DYN-19: Hedonische Adaptation gilt auch für spirituelle Erfahrungen.
\end{itemize}

Ohne explizite Axiome wäre eine solche Konsistenzprüfung unmöglich -- man wüsste nicht, \textit{wogegen} man prüfen soll.

\paragraph{(ii) Transparenz und Kritisierbarkeit.}

\epigraph{\textit{``A theory which is not refutable by any conceivable event is non-scientific.''}}{--- Karl Popper, \textit{Conjectures and Refutations} (1963)}

Karl Popper argumentierte, dass Wissenschaftlichkeit an Falsifizierbarkeit hängt. Eine Theorie, die mit \textit{jeder} möglichen Beobachtung kompatibel ist, sagt nichts über die Welt.

Implizite Annahmen sind per Definition unfalsifizierbar: Man kann sie nicht testen, weil man nicht weiß, dass man sie macht. Explizite Axiome hingegen sind kritisierbar. Man kann fragen:
\begin{itemize}
    \item Ist MA-EPI-15 (Loss Aversion) empirisch haltbar?
    \item Ist MA-GOV-06 (Subsidiarität) ethisch gerechtfertigt?
    \item Ist MA-DYN-11 (Kritische Masse bei 25\%) korrekt kalibriert?
\end{itemize}

Die Antworten mögen kontrovers sein -- aber die \textit{Fragen} können gestellt werden. Das ist Wissenschaft.

Imre Lakatos, Poppers Schüler und Kritiker, unterschied zwischen dem ``harten Kern'' einer Forschungsprogramms und seinem ``Schutzgürtel'' von Hilfshypothesen. Das \META-Modul macht diese Unterscheidung explizit:
\begin{itemize}
    \item \textbf{Harter Kern}: Die ontologischen Grundaxiome (Subjekte existieren, Nutzen existiert, Zeit existiert).
    \item \textbf{Schutzgürtel}: Die spezifischen Parameter (Loss Aversion $\lambda \approx 2$, kritische Masse $P^* \approx 25\%$).
\end{itemize}

Der Schutzgürtel darf -- und soll -- durch empirische Forschung modifiziert werden. Der harte Kern bleibt stabil.

\paragraph{(iii) Systematische Erweiterbarkeit.}

\epigraph{\textit{``If I have seen further, it is by standing on the shoulders of giants.''}}{--- Isaac Newton, Brief an Robert Hooke (1675)}

Wissenschaftlicher Fortschritt ist kumulativ. Newton baute auf Galilei und Kepler. Einstein baute auf Newton und Maxwell. Jede Generation steht auf den Schultern der vorigen.

Die Verhaltensökonomie hat dieses Ideal bisher nicht erreicht. Jedes neue Paper startet quasi bei Null: eigene Definitionen, eigene Operationalisierungen, eigene implizite Annahmen. Replikationsstudien scheitern oft nicht an methodischen Fehlern, sondern an begrifflichen Inkompatibilitäten.

Das \META-Modul schafft eine gemeinsame Grundlage. Neue Module können systematisch integriert werden, indem sie ihre Konsistenz mit den META-Axiomen nachweisen.

\textbf{Beispiel: Integration eines AFFECT-Moduls.} Angenommen, zukünftige BCM-Versionen sollen Affekt-Dynamik modellieren -- die schnellen emotionalen Reaktionen, die Entscheidungen beeinflussen. Ein solches AFFECT-Modul müsste:
\begin{enumerate}
    \item Neue ontologische Axiome definieren (``Affekte existieren'', ``Affekte sind kurzlebig'').
    \item Konsistenz mit bestehenden Axiomen nachweisen (z.B. Interaktion mit MA-EPI-05, der Zwei-Systeme-Theorie).
    \item Epistemologische Constraints respektieren (z.B. MA-EPI-37, die Affect Heuristic).
    \item Dynamische Axiome für Affekt-Verläufe formulieren.
\end{enumerate}

Das Ergebnis wäre eine \textit{kontrollierte} Erweiterung: Neues Wissen wird integriert, ohne die Kohärenz des Gesamtsystems zu zerstören.

\paragraph{(iv) Interdisziplinäre Verankerung.}

\epigraph{\textit{``Philosophy is written in this grand book, the universe... It is written in the language of mathematics.''}}{--- Galileo Galilei, \textit{Il Saggiatore} (1623)}

Die vier Kategorien des \META-Moduls -- Ontologie, Epistemologie, Governance, Dynamik -- sind keine Erfindungen der Verhaltensökonomie. Sie entsprechen klassischen philosophischen Disziplinen, die seit der Antike reflektiert werden.

\textbf{Ontologie} (griech. \textit{on} = seiend): Die Lehre vom Sein. Was existiert? Welche Entitäten und Beziehungen sind fundamental?

Die META-Ontologie definiert: Subjekte ($\sigma$), Verhalten ($V$), Nutzen ($U$), Kontext ($\rho$), Zeit ($\tau$), Gruppen ($G$), Identitäten ($I$), Normen ($N$). Diese Entitäten sind die ``Möbel'' des BCM-Universums.

Die philosophische Tradition -- von Aristoteles' Kategorien über Leibniz' Monadologie bis zu Quines ``On What There Is'' (1948) -- liefert das begriffliche Rüstzeug für diese Fragen.

\textbf{Epistemologie} (griech. \textit{episteme} = Wissen): Die Lehre vom Wissen. Was können wir wissen? Welche Grenzen hat Erkenntnis?

Die META-Epistemologie definiert: Wissen ist unvollständig (MA-EPI-01), Rationalität ist begrenzt (MA-EPI-02), Biases sind systematisch (MA-EPI-06 bis MA-EPI-37). Diese Axiome stehen in direkter Linie zur philosophischen Tradition -- von Platons Höhlengleichnis über Kants ``Grenzen der Vernunft'' bis zu Quines ``Two Dogmas of Empiricism'' (1951).

\textbf{Governance} (lat. \textit{gubernare} = steuern): Die Lehre vom Erlaubten. Was darf man tun? Welche Interventionen sind ethisch?

Die META-Governance definiert: Autonomie hat Primat (MA-GOV-03), Transparenz ist geboten (MA-GOV-05), Subsidiarität gilt (MA-GOV-06). Diese Prinzipien verbinden die Ethik -- von Kants kategorischem Imperativ über Mills Utilitarismus bis zu Rawls' Theorie der Gerechtigkeit -- mit der praktischen Interventionsgestaltung.

\textbf{Dynamik} (griech. \textit{dynamis} = Kraft): Die Lehre von der Veränderung. Wie entwickeln sich Systeme? Was sind die Gesetze des Wandels?

Die META-Dynamik definiert: Zeit existiert (MA-DYN-01), Pfadabhängigkeit gilt (MA-DYN-03), Gleichgewichte können multipel sein (MA-DYN-04). Diese Axiome verbinden die Metaphysik der Zeit -- von Heraklits ``Alles fließt'' über Hegels Dialektik bis zu Whiteheads Prozessphilosophie -- mit der Systemtheorie und den dynamischen Modellen der modernen Sozialwissenschaft.

Die interdisziplinäre Verankerung hat praktische Konsequenzen: Das \BCM{} kann von Philosophen geprüft, von Ökonomen kalibriert, von Psychologen validiert und von Ingenieuren implementiert werden. Jede Disziplin findet Anknüpfungspunkte.

%------------------------------------------------------------------------------
\subsubsection{Abgrenzung: Was META nicht ist}
%------------------------------------------------------------------------------

Um Missverständnisse zu vermeiden, sei explizit gesagt, was das \META-Modul \textit{nicht} ist:

\paragraph{Kein Bias-Katalog.}

Das \META-Modul enthält in der EPI-Kategorie 37 Axiome zu kognitiven Biases. Aber es ist kein ``Katalog aller Biases''. Es gibt Hunderte dokumentierter Biases in der Literatur -- von Anchoring über Dunning-Kruger bis zum Zeigarnik-Effekt. Das \META-Modul wählt \textit{die für BCM relevanten} Biases aus und formalisiert sie. Vollständigkeit ist nicht das Ziel; Kohärenz ist es.

\paragraph{Keine Heuristik-Sammlung.}

Heuristiken (wie ``Take-the-Best'' oder ``Recognition Heuristic'') sind \textit{Beschreibungen} menschlicher Entscheidungsstrategien. Axiome sind \textit{Annahmen} über die Struktur der Welt. Das \META-Modul enthält Axiome \textit{über} Heuristiken (z.B. MA-EPI-04: ``Heuristiken sind adaptiv''), aber nicht die Heuristiken selbst.

\paragraph{Kein Interventions-Cookbook.}

Das \META-Modul definiert acht Interventionstypen (MA-GOV-08) und ethische Prinzipien für ihren Einsatz. Aber es enthält keine konkreten Rezepte (``Wenn X, dann nutze Nudge Y''). Solche Rezepte gehören in die Anwendungsschicht (\BEATRIX), nicht in die Fundamente.

\paragraph{Keine empirische Theorie.}

Axiome sind \textit{Setzungen}, keine Hypothesen. Man kann MA-ONT-01 (``Subjekte existieren'') nicht empirisch widerlegen -- es ist eine definitorische Annahme, keine Prognose. Die empirischen Gehalte des \BCM{} entstehen erst in den operativen Modulen, die auf den META-Axiomen aufbauen.

%------------------------------------------------------------------------------
\subsubsection{Praktische Konsequenzen fehlender Axiomatik}
%------------------------------------------------------------------------------

Die Notwendigkeit eines \META-Moduls wird am deutlichsten, wenn man die Konsequenzen seines Fehlens betrachtet. Drei Fallbeispiele illustrieren das Problem.

\paragraph{Fallbeispiel 1: Inkommensurable Effektschätzungen.}

Die Behavioral Insights Team (BIT) in UK schätzte 2015 die Effektgröße eines Tax-Nudges auf etwa 15 Prozentpunkte. Eine Replikation durch die australische Nudge-Unit fand nur 3 Prozentpunkte. Beide Teams arbeiteten seriös -- aber mit unterschiedlichen impliziten Annahmen über Baseline-Verhalten, Kontexteffekte und Messzeitpunkte.

Ohne gemeinsame Sprache ist unklar, ob der Unterschied an der Intervention, der Population, dem Kontext oder der Messung liegt. Das \META-Modul liefert diese Sprache: MA-ONT-04 definiert, was ``Kontext'' bedeutet; MA-DYN-19 definiert, wie Effekte über Zeit abklingen; MA-EPI-26 fordert Berücksichtigung von Base Rates.

\paragraph{Fallbeispiel 2: LLM-Halluzination bei Persona-Schätzung.}

Ein Praktiker beschreibt eine Persona: ``Anna ist eine 35-jährige Managerin, die sich für Nachhaltigkeit engagiert und Fair-Trade-Produkte kauft.'' Ein naives LLM schätzt: ``Anna ist hochgradig altruistisch, $\lambda \approx 0.9$.''

Das Problem: Die Beschreibung enthält \textit{stated values} (was Anna sagt), nicht \textit{revealed behavior} (was Anna tut). Social Desirability Bias (MA-EPI-32) legt nahe, dass die wahren Werte niedriger liegen. Base Rate Neglect (MA-EPI-26) erinnert daran, dass nur etwa 10\% der Bevölkerung ``echte Altruisten'' sind.

Ein durch \META{} informiertes System wie \BEATRIX{} würde:
\begin{enumerate}
    \item Die Beschreibung als potenziell verzerrt flaggen.
    \item Die Schätzung an empirischen Base Rates verankern.
    \item Unsicherheit explizit quantifizieren (``$\lambda$ zwischen 0.5 und 0.8 mit 70\% Konfidenz'').
\end{enumerate}

\paragraph{Fallbeispiel 3: Kulturblindheit bei Interventions-Transfer.}

Eine Nudge-Intervention funktioniert in Dänemark: Opt-Out für Organspende erhöht die Spenderquote von 4\% auf 85\%. Dieselbe Intervention wird nach Japan übertragen -- und scheitert.

Warum? Ohne die META-Ebene (MA-ONT-26: 5-Ebenen-Hierarchie) fehlt die Sprache, um den Unterschied zu beschreiben. Mit \META{} ist klar:
\begin{itemize}
    \item \textbf{META-Ebene} (Kultur): Dänemark ist individualistisch, Japan ist kollektivistisch.
    \item \textbf{MAKRO-Ebene} (Institution): Unterschiedliche rechtliche Rahmenbedingungen für Organspende.
    \item \textbf{MESO-Ebene} (Organisation): Unterschiedliche Rolle von Krankenhäusern.
    \item \textbf{MIKRO-Ebene} (Familie): In Japan entscheidet die Familie, nicht das Individuum.
\end{itemize}

Die 5-Ebenen-Hierarchie erklärt, warum ``derselbe'' Nudge in verschiedenen Kontexten radikal unterschiedlich wirkt -- und warum naive Übertragungen scheitern.

%------------------------------------------------------------------------------
\subsubsection{Zusammenfassung: Die Notwendigkeit des META}
%------------------------------------------------------------------------------

Das \META-Modul ist kein akademisches Luxusgut. Es ist eine praktische Notwendigkeit für jedes System, das menschliches Verhalten \textit{konsistent}, \textit{transparent}, \textit{erweiterbar} und \textit{interdisziplinär verankert} modellieren will.

Die Alternative -- implizite Annahmen, fragmentierte Terminologie, nicht-kumulative Forschung -- hat die Verhaltensökonomie 50 Jahre lang behindert. Die LLM-Revolution erzwingt nun die Explizitmachung: Maschinen können nicht mit Intuition arbeiten; sie brauchen Axiome.

Das \BCM{} liefert diese Axiome. Die Frage ist nicht, \textit{ob} ein solches Framework entstehen wird, sondern \textit{wer} es zuerst etabliert -- und damit die Standards für eine ganze Generation von AI-gestützter Verhaltensmodellierung setzt.

%------------------------------------------------------------------------------

\subsection{Forschungslücke: Warum existiert 2025 kein öffentliches META?}

Trotz sieben Jahrzehnten Forschung zu Bounded Rationality (Simon, 1955), fünf Jahrzehnten Behavioral Economics (Kahneman \& Tversky, 1979) und zwei Jahrzehnten institutionalisierter Nudge-Praxis existiert 2025 kein allgemein verfügbares, öffentliches axiomatisches Framework für Verhaltensmodellierung. Diese bemerkenswerte Lücke -- ein fehlendes ``Periodensystem'' der Verhaltenswissenschaften -- hat tiefgreifende strukturelle Ursachen, die wir im Folgenden analysieren.

\subsubsection{Das gescheiterte Projekt der Bounded Rationality}

Herbert Simon führte 1955 den Begriff \textit{Bounded Rationality} ein als ``shorthand for his proposal to replace the perfect rationality assumptions of homo economicus with a concept of rationality better suited to cognitively limited agents'' (Stanford Encyclopedia of Philosophy). Simon identifizierte drei fundamentale Beschränkungen menschlicher Rationalität:

\begin{enumerate}[label=(\roman*)]
    \item \textbf{Limitierte Information}: Entscheidungsträger haben selten vollständige Information über alle Möglichkeiten
    \item \textbf{Kognitive Constraints}: Der menschliche Geist hat begrenzte Rechenkapazität
    \item \textbf{Zeitdruck}: Die meisten realen Entscheidungen müssen unter Zeitbeschränkungen getroffen werden, die erschöpfende Analyse verhindern
\end{enumerate}

Obwohl Simon 1978 den Nobelpreis für Wirtschaftswissenschaften erhielt ``for his pioneering research into the decision-making process within economic organizations'', blieb sein Ansatz ein Fremdkörper im ökonomischen Mainstream. Wie Barros (2010) analysiert: ``His approach was probably too far from the mainstream framework -- neither optimization, nor formal representations, nor axiomatic foundation -- to attract the attention it deserved.''

Simon selbst beklagte 1978, dass der Einfluss seiner Bounded-Rationality-Konzeption auf die Mainstream-Ökonomie noch immer begrenzt sei. Das Problem ist fundamental: Während für einen idealen Optimierer nur eine Handlungsoption offen steht, existieren für prozedurale Satisficer so viele Ergebnisse wie Satisficing-Prozeduren (Rubinstein, 1998). Mit anderen Worten: Es gibt viele Wege zur Bounded Rationality, aber keinen kanonischen.

\textbf{Die Konsequenz}: Bounded Rationality wurde nie zu einem konstruktiven alternativen Framework, sondern blieb eine kritische Perspektive auf die Mainstream-Ökonomie -- wichtig, aber nicht operationalisierbar.

\subsubsection{Akademische Fragmentierung: Die Sammeldisziplin ohne Synthese}

Die Verhaltensökonomie hat sich nicht als axiomatisches System entwickelt, sondern als \textit{Sammeldisziplin} empirisch dokumentierter Anomalien. David Gal hat argumentiert, dass viele Probleme der Verhaltensökonomie daher rühren, dass sie ``too concerned with understanding how behavior deviates from standard economic models rather than with understanding why people behave the way they do'' ist. Das Verständnis des \textit{Warum} ist jedoch notwendig für die Schaffung generalisierbaren Wissens -- dem Ziel der Wissenschaft.

\paragraph{Die Effekt-Katalog-Mentalität.}
Kahneman und Tversky dokumentierten ab 1972 systematisch kognitive Biases: Verfügbarkeitsheuristik (1973), Anchoring (1974), Framing (1981), Loss Aversion (1979). Thaler erweiterte den Katalog um den Endowment Effect (1980), Mental Accounting (1985) und Hyperbolic Discounting. Doch diese Phänomene wurden nie in ein kohärentes Axiomensystem integriert. Wie die Wikipedia-Analyse zusammenfasst: ``Despite a great deal of rhetoric, no unified behavioral theory has yet been espoused: behavioral economists have proposed no alternative unified theory of their own to replace neoclassical economics with.''

\paragraph{Die ``New Behavioral Economics'' und ihre Grenzen.}
Die ``New Behavioral Economists'' versuchten, empirisch entdeckte Anomalien in die bekannten formalen Modelle zu integrieren, mit axiomatischen Grundlagen und nutzenmaximierenden Funktionen. Dies geschah durch Lockerung strenger Annahmen neoklassischer Modelle und Einführung neuer Parameter für spezifische Verhaltensphänomene. Doch wie Primrose (2017) kritisiert, führte dies zu einer ``incongruous methodological configuration, simultaneously promulgating the need for more realistic accounts of behavior while retaining unrealistic assumptions associated with hyper-rationality.''

Forscher wie Gigerenzer argumentieren, dass Modelle, die menschliche Entscheidungsfindung als Optimierungsprozess darstellen -- statt als Anwendung vereinfachter Entscheidungsregeln wie Heuristiken -- die wahre Bedeutung der Behavioral Economics verfehlen. Diese Kritiker sehen sich als die wahren Erben von Simons Theorie, nicht Kahneman und Tversky.

\paragraph{Existierende Frameworks und ihre Limitationen.}
Es existieren durchaus Versuche, Verhaltensänderung systematisch zu modellieren. Das bekannteste ist das \textbf{COM-B Model} und das zugehörige \textbf{Behaviour Change Wheel} (BCW) von Michie, van Stralen und West (2011). Das COM-B-System postuliert drei essentielle Bedingungen für Verhalten:

\begin{itemize}
    \item \textbf{Capability}: Psychologische und physische Fähigkeit
    \item \textbf{Opportunity}: Externe Faktoren, die Verhalten ermöglichen
    \item \textbf{Motivation}: Bewusste und unbewusste kognitive Prozesse
\end{itemize}

Das BCW wurde aus 19 verschiedenen Verhaltensänderungs-Frameworks synthetisiert und bietet 9 Interventionsfunktionen sowie 7 Policy-Kategorien. Es ist jedoch primär ein \textit{Klassifikationssystem} für Interventionen, kein axiomatisches Fundament. Es fehlen:

\begin{itemize}
    \item Formale mathematische Spezifikation
    \item Explizite epistemologische Annahmen
    \item Integration kognitiver Biases
    \item Dynamische Modellierung über Zeit
    \item LLM-kompatible Operationalisierung
\end{itemize}

Ähnliche Limitationen betreffen die \textbf{Theory of Planned Behavior} (Ajzen, 1991), das \textbf{Health Belief Model} (Rosenstock, 1966) und das \textbf{Transtheoretical Model} (Prochaska \& DiClemente, 1983). Keines dieser Frameworks bietet eine axiomatische Grundlage im Sinne formaler Wissenschaftstheorie.

\subsubsection{Die Replikationskrise: Empirische Fundamente in Frage}

Die Verhaltensökonomie wird durch eine tiefgreifende \textbf{Replikationskrise} erschüttert, die das empirische Fundament vieler Kernbefunde in Frage stellt.

\paragraph{Das DellaVigna-Linos-Paradox.}
Die wegweisende Studie von DellaVigna und Linos (2022), publiziert in \textit{Econometrica}, analysierte 126 RCTs mit 23 Millionen Individuen aus den zwei größten Nudge Units der USA. Die Ergebnisse sind ernüchternd:

\begin{quote}
``In papers published in academic journals, the average impact of a nudge is very large -- an 8.7 percentage point take-up effect, which is a 33.4\% increase over the average control. In the Nudge Units sample, the average impact is still sizable and highly statistically significant, but smaller at 1.4 percentage points, an 8.0\% increase.''
\end{quote}

Die Effektgrößen in der Praxis betragen also nur \textbf{ein Sechstel} der publizierten akademischen Ergebnisse. DellaVigna und Linos erklären etwa 70\% dieser Differenz mit \textbf{Publication Bias}, verschärft durch niedrige statistische Power.

\paragraph{Systematische Replikationsfehler.}
Ein 2024 Wall Street Journal Review berichtet, dass weniger als die Hälfte der Nudging-Experimente langfristige Wirkung zeigen, wenn sie erneut getestet werden -- was ihre externe Validität unter realen Bedingungen in Frage stellt. In manchen Studien wurden Effektgrößen von Cohen's $d = -0.01$ bis $d = 0.05$ gemessen, wobei Äquivalenztests zeigten, dass die Treatment-Effekte in vier von fünf Experimenten statistisch äquivalent zu Null waren.

Chatziathanasiou (2022) argumentiert in ``Nudging after the Replication Crisis'' sogar, dass ``when publication bias is appropriately corrected for, no evidence for the effectiveness of nudges remains'' -- eine radikale Position, die jedoch die Dringlichkeit eines theoretisch fundierteren Ansatzes unterstreicht.

\paragraph{Konsequenzen für die Feldentwicklung.}
Die Replikationskrise offenbart ein fundamentales Problem: Ohne axiomatische Grundlage ist es unmöglich zu unterscheiden, welche Effekte robust sind und welche Artefakte von Kontextfaktoren, Publication Bias oder statistischer Fluktuation. Ein META-Framework würde ermöglichen:

\begin{itemize}
    \item Theoriegeleitete Vorhersagen statt explorativer Effekt-Jagd
    \item Klare Spezifikation von Boundary Conditions
    \item Kumulative Wissensakkumulation statt isolierter Befunde
\end{itemize}

\subsubsection{Kommerzielle Fragmentierung: Das Proprietary-Dilemma}

Behavioral Insights Units weltweit haben durchaus operative Frameworks entwickelt -- doch diese bleiben systematisch proprietär.

\paragraph{Die globale Proliferation von Nudge Units.}
Das UK Behavioural Insights Team (BIT), gegründet 2010, hat die Gründung von über 200 ähnlichen Nudge Units weltweit inspiriert. Internationale Institutionen wie Weltbank, UN-Agenturen, OECD und EU haben ebenfalls Behavioral Insights Units etabliert. Die Verbreitung umfasst:

\begin{itemize}
    \item \textbf{Nordamerika}: US Social and Behavioral Sciences Team (Executive Order 2015), Canadian Behavioural Insights Units (Ontario, British Columbia)
    \item \textbf{Europa}: BIT UK, Nudge Units in Deutschland, Niederlande, Skandinavien
    \item \textbf{Asien-Pazifik}: Singapur, Indien, Indonesien, Australien (BIT Sydney), Neuseeland (BIT Wellington seit 2016)
    \item \textbf{Lateinamerika}: Peru, Kolumbien, Brasilien
\end{itemize}

\paragraph{Das Wissens-Paradox.}
Trotz dieser Proliferation existiert kein öffentlicher Standard. Jede Unit entwickelt eigene Methoden, Taxonomien und Best Practices. Die Gründe sind ökonomisch rational:

\begin{enumerate}[label=(\arabic*)]
    \item \textbf{Wettbewerbsvorteil}: Proprietäre Frameworks differenzieren Anbieter
    \item \textbf{Consulting-Modell}: Wissen als wiederkehrende Einnahmequelle
    \item \textbf{Lock-in-Effekte}: Kunden werden an spezifische Methoden gebunden
    \item \textbf{Fehlende Incentives}: Keine Organisation profitiert von Standardisierung
\end{enumerate}

\paragraph{Evidence Adoption Gap.}
DellaVigna, Kim und Linos (2024) untersuchten in ``Bottlenecks for Evidence Adoption'' (\textit{Journal of Political Economy}), wie Regierungen RCT-Ergebnisse in Policy umsetzen. In einer Studie von 30 US-Städten mit 73 RCTs adoptierten Städte ein Nudge-Treatment nur in \textbf{27\% der Fälle}. Überraschenderweise prädizierte die Stärke der Evidenz die Adoption nicht -- der stärkste Prädiktor war, ob das RCT in bestehende Kommunikation implementiert wurde. Die Autoren identifizieren \textbf{organisatorische Trägheit} als Haupterklärung.

Dies unterstreicht: Selbst wenn Nudge Units effektive Interventionen entwickeln, fehlt ein Framework für systematischen Wissenstransfer und kumulative Verbesserung.

\subsubsection{Philosophie-Aversion: Das epistemologische Vakuum}

Die Ökonomie als Disziplin vermeidet traditionell explizite philosophische Positionierungen -- eine Haltung, die die Entwicklung eines META-Frameworks strukturell behindert.

\paragraph{Implizite Axiome des Homo Oeconomicus.}
Die neoklassische Ökonomie basiert auf Axiomen, die selten expliziert werden: vollständige Präferenzen, Transitivität, Unabhängigkeit von irrelevanten Alternativen (Arrow-Debreu). Diese axiomatische Grundlage wurde zum ``cornerstone of economic science'' (Barros, 2010), doch die Axiome selbst wurden als selbstevident behandelt, nicht als falsifizierbare Annahmen.

\paragraph{Die methodologische Inkongruenz der Behavioral Economics.}
Die epistemologischen, ontologischen und methodologischen Komponenten der Behavioral Economics werden zunehmend debattiert, insbesondere von Wirtschaftshistorikern und Methodologen. Eine kritische Literatur hat gezeigt, dass Behavioral Economics durch eine ``incongruous methodological configuration'' gekennzeichnet ist: Sie propagiert die Notwendigkeit realistischerer Verhaltensannahmen, während sie unrealistische Annahmen der Hyper-Rationalität beibehält.

\paragraph{Das Normativitätsproblem.}
Von Anfang an war die ordinale und erwartete Nutzentheorie ein \textit{normatives} Unterfangen. Die Rationalitätskonzepte, die durch die verschiedenen Axiome konstituiert werden, sind evaluativ und damit potenziell präskriptiv. Solange man annimmt, dass das tatsächliche Verhalten der Menschen die Axiome zumindest approximiert, ist keine normative Inferenz darüber nötig, wie sich Menschen verhalten \textit{sollten}. Doch wenn Entscheidungen keine kohärenten Präferenzen offenbaren, ist es nicht mehr möglich, Entscheidungen mit ``wahren'' Präferenzen zu identifizieren -- was die Grundlagen sowohl der Wohlfahrts- als auch der Wahlfreiheitsinterpretation untergräbt.

\paragraph{Konsequenz für die Feldentwicklung.}
Ein META-Framework erfordert genau diese Explizitmachung: Was nehmen wir an über die Natur von Entscheidungen (Ontologie)? Was können wir wissen über Präferenzen und Verhalten (Epistemologie)? Was ist erlaubt bei Interventionen (Governance)? Wie verändern sich Systeme über Zeit (Dynamik)? Diese philosophische Arbeit liegt außerhalb des disziplinären Selbstverständnisses der Ökonomie -- aber sie ist unumgänglich für ein konsistentes Framework.

\subsubsection{Der technologische Imperativ: LLMs als Katalysator}

Der entscheidende Wendepunkt ist die \textbf{LLM-Revolution ab 2022}. Vor GPT-4, Claude und vergleichbaren Modellen wäre ein axiomatisches Verhaltens-Framework ``nur Philosophie'' gewesen -- interessant für Methodologen, aber ohne operative Relevanz.

\paragraph{LLMs und Behavioral Economics: Konvergenz zweier Felder.}
Aktuelle Forschung zeigt, dass LLMs selbst verhaltensökonomisch relevante Eigenschaften aufweisen. Wie Forschung auf arXiv dokumentiert:

\begin{quote}
``Humans exhibit systematic behavioral biases such as loss aversion, anchoring, framing, etc., which lead to suboptimal economic decisions. Insofar as such biases may be embedded in text data on which LLMs are trained, LLMs are prone to the same behavioral biases.''
\end{quote}

LLMs zeigen sogar ``gambling-like risk-taking behaviors analogous to those observed in gambling psychology, including overconfidence bias, loss-chasing tendencies, and probability misjudgment'' (arXiv, 2024). Dies macht ein axiomatisches Framework nicht nur wünschenswert, sondern notwendig -- sowohl für das Verständnis von LLM-Verhalten als auch für dessen Steuerung.

\paragraph{LLMs als Experimental-Plattform.}
Die Integration von generativer KI in Behavioral und Experimental Economics stellt einen transformativen wissenschaftlichen Fortschritt dar. Wie Shapira (2024) argumentiert:

\begin{quote}
``LLMs can revolutionize experimental economics by simulating human behavior in foundational economic scenarios. Traditional approaches in economics are limited by the costs and complexities of human-based experiments, and replacing humans with LLMs enables large-scale simulations.''
\end{quote}

Ein axiomatisches META-Framework würde es ermöglichen, LLMs nicht nur als Simulatoren menschlichen Verhaltens einzusetzen, sondern als \textit{verhaltensökonomisch informierte Agenten}, die systematisch Biases berücksichtigen, Interventionen planen und personalisierte Beratung bieten.

\paragraph{Agent-Based Modeling mit LLMs.}
LLMs haben vielversprechende Ergebnisse bei der Verbesserung agentenbasierter Simulationen gezeigt, indem sie ``more nuanced and realistic representations of agents' decision-making processes, communication, and adaptation within simulated environments'' ermöglichen (Nature Humanities and Social Sciences Communications, 2024). Ein Memory-Modul in LLMs kann entscheidend sein für das Verstehen und Adaptieren an dynamische soziale Umgebungen.

\paragraph{Die neue Notwendigkeit: Maschinenlesbare Verhaltenstheorie.}
Die technologische Möglichkeit erzeugt erstmals einen echten \textit{Bedarf} nach einem maschinenlesbaren META. Ein solches Framework muss:

\begin{enumerate}[label=(\roman*)]
    \item \textbf{Formal spezifiziert} sein, damit LLMs es parsen und anwenden können
    \item \textbf{Parametrisiert} sein mit empirisch validierten Base Rates
    \item \textbf{Modular} sein, um verschiedene Anwendungskontexte zu unterstützen
    \item \textbf{Erweiterbar} sein, um neue Befunde zu integrieren
    \item \textbf{Validierbar} sein gegen empirische Daten
\end{enumerate}

Das \BCM{} META-Modul wurde genau für diese Anforderungen entwickelt -- als erstes öffentliches, axiomatisches, LLM-implementiertes Framework für Verhaltensmodellierung.

\subsection{Das Rennen um das erste Modell}

Die Konvergenz von Behavioral Science und LLM-Technologie hat ein \textbf{strategisches Wettrennen} ausgelöst. Mehrere Akteursgruppen arbeiten derzeit -- mit unterschiedlichen Ansätzen -- an verhaltensökonomischen Foundations für KI-Systeme:

\begin{table}[h]
\centering
\caption{Akteure im Wettbewerb um verhaltensökonomische AI-Frameworks}
\label{tab:competitors}
\begin{tabular}{@{}p{3cm}p{4cm}p{5cm}@{}}
\toprule
\textbf{Akteursgruppe} & \textbf{Ansatz} & \textbf{Limitationen} \\
\midrule
Big Tech (Google DeepMind, OpenAI, Anthropic) & LLM-native Verhaltensmodellierung & Fokus auf Alignment, nicht auf Behavioral Economics; keine explizite Axiomatik \\
\addlinespace
Behavioral Insights Units (BIT, ideas42) & Praxis-getriebene Frameworks & Proprietär; nicht formalisiert; nicht LLM-kompatibel \\
\addlinespace
Akademische Institutionen & Theoretische Integration & Langsam; fragmentiert; kein Implementierungsfokus \\
\addlinespace
Management-Beratungen & Client-spezifische Lösungen & Nicht generalisierbar; keine Veröffentlichung \\
\addlinespace
Regierungs-Nudge-Units & Policy-fokussierte Tools & Enger Scope; keine universelle Axiomatik \\
\bottomrule
\end{tabular}
\end{table}

\subsubsection{Unsere Prognose: Wer wird sich durchsetzen?}

Wir argumentieren, dass \textbf{derjenige Akteur das Rennen gewinnt}, der als erster ein \textit{offenes}, \textit{axiomatisch fundiertes}, \textit{LLM-implementiertes} und \textit{empirisch validiertes} Framework etabliert. Die Kriterien sind kumulativ:

\begin{enumerate}[label=(\alph*)]
    \item \textbf{Offen}: Proprietäre Systeme werden langfristig von offenen Standards verdrängt (vgl. Linux vs. proprietäre Unix-Varianten, TCP/IP vs. OSI)
    \item \textbf{Axiomatisch}: Ad-hoc-Sammlungen von Biases skalieren nicht; nur ein konsistentes System ermöglicht Erweiterung und Validierung
    \item \textbf{LLM-implementiert}: Theorie ohne Implementierung bleibt akademisch; erst die Operationalisierung schafft Netzwerkeffekte
    \item \textbf{Empirisch validiert}: Base Rates und Effektgrößen aus der Literatur müssen integriert sein (vgl. \texttt{baserates.json} im BCM 2.2)
\end{enumerate}

Das \BCM{} META-Modul erfüllt diese vier Kriterien. Die Frage ist nicht, \textit{ob} ein solches Framework entstehen wird, sondern \textit{wer} es zuerst etabliert.

\subsection{First-Mover-Advantage: Warum der Erste die Pfründe erntet}

Die Dynamik des Wettbewerbs um das erste verhaltensökonomische META-Framework folgt klassischen \textbf{Pfadabhängigkeits-Mechanismen}, die in einer umfangreichen wirtschaftswissenschaftlichen Literatur dokumentiert sind.

\subsubsection{Theoretische Grundlagen: Lieberman \& Montgomery}

Die systematische Analyse von First-Mover-Advantages (FMA) beginnt mit dem Seminalpaper von Lieberman und Montgomery (1988) im \textit{Strategic Management Journal}, das 1996 den Strategic Management Society Award für den herausragenden Beitrag zur Entwicklung des Feldes erhielt und über 5.000 Zitationen akkumulierte.

Lieberman und Montgomery definieren FMA als ``the ability of pioneering firms to earn positive economic profits (i.e., profits in excess of the cost of capital)'' und identifizieren drei primäre Quellen:

\begin{enumerate}[label=(\arabic*)]
    \item \textbf{Technological Leadership}: Lernkurven-Effekte und Patent-/F\&E-Races
    \item \textbf{Preemption of Assets}: Vorwegnahme knapper Ressourcen (Talente, Daten, Standards)
    \item \textbf{Buyer Switching Costs}: Aufbau von Wechselkosten beim Kunden
\end{enumerate}

Empirische Studien zeigen, dass Pioniere im Durchschnitt 10--20\% höhere Marktanteile erzielen als spätere Marktteilnehmer (Lieberman \& Montgomery, 1998). Allerdings sind die Ergebnisse kontextabhängig: In Märkten mit starken Netzwerkeffekten, patentierbarer Technologie oder hohen Switching Costs sind FMA ausgeprägter.

\subsubsection{Pfadabhängigkeit und Lock-In: Arthur und David}

Die theoretische Fundierung von Pfadabhängigkeit stammt von W. Brian Arthur (1989) und Paul David (1985). Arthurs wegweisendes Paper ``Competing Technologies, Increasing Returns, and Lock-In by Historical Events'' im \textit{Economic Journal} argumentiert:

\begin{quote}
In einer Welt mit \textit{increasing returns to scale} kann es multiple mögliche Gleichgewichtslösungen geben -- und die dominante Lösung kann suboptimal sein. Das vorherrschende ökonomische Ergebnis kann selbst das Produkt eines scheinbar unbedeutenden und zufälligen Ereignisses sein, das durch den Prozess steigender Erträge letztlich einen First-Mover-Vorteil gegenüber anderen möglichen Ergebnissen erhält, selbst wenn letztere ökonomisch effizienter sind.
\end{quote}

Arthur identifiziert vier Lock-In-Mechanismen:
\begin{itemize}
    \item \textbf{Hohe initiale Setup-Kosten} mit steigenden Erträgen bei mehr Nutzern
    \item \textbf{Lerneffekte} mit steigenden Erträgen für etablierte Nutzer
    \item \textbf{Koordinations- und Netzwerkeffekte}
    \item \textbf{Adaptive Erwartungen} der Marktteilnehmer
\end{itemize}

Paul Davids (1985) berühmte QWERTY-Studie illustriert, wie eine suboptimale Tastaturanordnung zum Standard wurde und trotz effizienterer Alternativen persistierte -- ein Paradebeispiel für pfadabhängiges Lock-In.

\subsubsection{Netzwerkeffekte: Katz \& Shapiro}

Die akademische Literatur zu Netzwerkeffekten wurde von Katz und Shapiro (1985, 1986) sowie Farrell und Saloner (1985, 1986) begründet. Katz und Shapiro formalisierten, wie \textit{network externalities} in Märkten mit kompatiblen Systemen funktionieren:

\begin{quote}
``Markets with network effects can exhibit inefficient outcomes where winners may be determined by historical accident or path dependence.'' (Katz \& Shapiro, 1986)
\end{quote}

Für ein META-Framework bedeutet dies: Je mehr Nutzer das BCM-System verwenden, desto wertvoller wird es durch:
\begin{itemize}
    \item Akkumulation von Trainingsdaten und Validierungsfällen
    \item Dokumentation von Anwendungsfällen und Best Practices
    \item Entstehung von Community-Wissen und Praktiker-Netzwerken
    \item Entwicklung komplementärer Tools und Erweiterungen
\end{itemize}

Ein etabliertes META wird zum \textbf{Schelling-Punkt} (Schelling, 1960) der Branche -- dem natürlichen Fokus für Koordination ohne explizite Kommunikation.

\subsubsection{Switching Costs und Kundenbindung}

Klemperer (1987, 1995) formalisierte die Rolle von Switching Costs in Wettbewerbsmärkten. Organisationen, die ihre Interventionsprogramme auf einem Framework aufbauen, haben erhebliche Wechselkosten:

\begin{itemize}
    \item \textbf{Technische Kosten}: Re-Training von Modellen, Re-Kalibrierung von Parametern
    \item \textbf{Organisatorische Kosten}: Re-Education von Mitarbeitern, Prozessanpassungen
    \item \textbf{Kognitive Kosten}: Umlernen von Terminologie und Denkmodellen
    \item \textbf{Beziehungskosten}: Verlust von aufgebautem Netzwerk und Support
\end{itemize}

Empirische Studien im E-Commerce zeigen, dass Early Mover Advantages durch ``consumer demand inertia'' verstärkt werden können, insbesondere wenn Customer Relationship Management-Fähigkeiten systematisch aufgebaut werden (Panel-Daten von 7.309 E-Tailern bestätigen diese Hypothese).

\subsubsection{Standardsetzung in Technologiemärkten}

Wer zuerst eine kritische Masse erreicht, definiert den \textit{de facto}-Standard. Die Forschung zu Technologiestandards (Shapiro \& Varian, 1999; Suarez, 2004) zeigt:

\begin{quote}
``A stronger advantage from technology leadership arises when the first mover can establish their product as the industry standard, making it more difficult for followers to gain customer acceptance.''
\end{quote}

Historische Beispiele illustrieren diese Dynamik:
\begin{itemize}
    \item \textbf{VHS vs. Betamax}: Technisch unterlegener Standard dominiert durch Netzwerkeffekte
    \item \textbf{QWERTY}: Persistenz trotz Ineffizienz (David, 1985)
    \item \textbf{Windows}: Dominanz trotz technischer Alternativen
    \item \textbf{TCP/IP vs. OSI}: Offener Standard verdrängt proprietäre Alternative
\end{itemize}

Die Lektion für BCM: Ein offenes, gut dokumentiertes Framework hat langfristig bessere Chancen als proprietäre Alternativen -- genau die Strategie, die das \BCM{} verfolgt.

\subsubsection{Akademische Zitationsökonomie und Matthew-Effekt}

Robert Merton (1968) beschrieb den \textbf{Matthew-Effekt} in der Wissenschaft: ``For unto every one that hath shall be given'' -- Zitationen akkumulieren exponentiell beim Erstpublizierer. Das erste publizierte, umfassende Framework wird zum Referenzpunkt aller Folgearbeiten.

Dieser Effekt ist besonders stark bei:
\begin{itemize}
    \item Methodologischen Frameworks (die Methode wird zitiert, nicht nur die Ergebnisse)
    \item Interdisziplinären Ansätzen (breitere Rezeption)
    \item Praktisch anwendbaren Systemen (Zitation durch Praktiker und Akademiker)
\end{itemize}

Das \BCM{} META-Modul erfüllt alle drei Kriterien und ist damit für akademische Impactmaximierung positioniert.

\subsubsection{Empirische Evidenz: Gemischte Befunde und Kontextfaktoren}

Die empirische Literatur zu FMA zeigt gemischte Ergebnisse. Suarez und Lanzolla (2005) argumentieren in ``The Half-Truth of First-Mover Advantage'':

\begin{quote}
``Despite claims that `earlier is better,' results on early-mover advantages are inconsistent. Pioneers do not consistently outperform followers in longevity, with failure rates around 47\% for true first entrants.''
\end{quote}

Entscheidend sind \textbf{Kontextfaktoren}:
\begin{itemize}
    \item \textbf{Marktevolution}: Graduelle Evolution begünstigt First Mover
    \item \textbf{Technologieevolution}: Radikale Disruption kann Incumbents benachteiligen
    \item \textbf{Ecosystem Economics}: Komplementäre Assets verstärken FMA
    \item \textbf{Organizational Will and Skill}: Timing ist notwendig, aber nicht hinreichend
\end{itemize}

Für verhaltensökonomische AI-Frameworks argumentieren wir, dass die Kontextfaktoren FMA \textit{begünstigen}:
\begin{itemize}
    \item Der Markt ist im Entstehen (graduelle Evolution)
    \item Die Technologie (LLMs) ist stabilisiert (keine radikale Disruption erwartet)
    \item Ecosystem-Effekte sind stark (Daten, Tools, Community)
    \item Die Kombination von Behavioral Science und AI-Expertise ist selten
\end{itemize}

\subsubsection{Wirtschaftliche Implikationen}

Die ökonomischen Pfründe eines etablierten META-Frameworks umfassen:
\begin{itemize}
    \item \textbf{API-Ökonomie}: Lizenzierung der Axiom-Engine für Drittanbieter
    \item \textbf{Beratungsrenten}: Expertise in der Implementierung und Anpassung
    \item \textbf{Datenhoheit}: Aggregierte Verhaltensdaten aus Anwendungen
    \item \textbf{Trainingsmarkt}: Zertifizierungen und Schulungen für BCM-Praktiker
    \item \textbf{Tooling-Ecosystem}: Add-ons, Visualisierungen, Branchenadaptionen
\end{itemize}

Der \textit{Total Addressable Market} für verhaltensökonomische AI-Anwendungen -- von Healthcare über Finance bis Public Policy -- wird auf dreistellige Milliardenbeträge geschätzt. Wer den Standard setzt, kontrolliert einen signifikanten Anteil dieser Wertschöpfung.

\subsection{Die Zürich-Konstellation: Warum gerade hier?}

\epigraph{\textit{``The best way to predict the future is to invent it.''}}{--- Alan Kay, Xerox PARC (1971)}

Die vorhergehenden Abschnitte haben gezeigt: Ein axiomatisches Verhaltens-Framework wird entstehen. Die Frage ist, \textit{wer} es etabliert. Wir argumentieren, dass eine einzigartige Konstellation in Zürich -- die Verbindung von Universität Zürich und FehrAdvice \& Partners -- die besten Voraussetzungen hat, dieses Rennen zu gewinnen.

Diese Behauptung mag unbescheiden klingen. Sie ist es nicht. Sie folgt aus einer nüchternen Analyse der erforderlichen Kompetenzen und der Frage, wo diese Kompetenzen \textit{gemeinsam} vorhanden sind.

%------------------------------------------------------------------------------
\subsubsection{Das Kompetenz-Quartett}
%------------------------------------------------------------------------------

Die Entwicklung und Skalierung eines META-Verhaltensmodells erfordert vier distinkte Kompetenzfelder, die selten in einer Hand liegen:

\begin{enumerate}
    \item \textbf{Verhaltensökonomische Theorie}: Die wissenschaftlichen Grundlagen menschlicher Entscheidungen
    \item \textbf{Software Engineering \& KI}: Die technische Implementierung in LLM-Systemen
    \item \textbf{Evidenzbasierte Strategie}: Die methodologische Brücke zwischen Forschung und Praxis
    \item \textbf{Anwendungsexpertise}: Das Wissen, was in der realen Welt funktioniert
\end{enumerate}

Die Zürich-Konstellation vereint diese vier Felder in einer einzigartigen Konfiguration:

\paragraph{Säule 1: Ernst Fehr und die Behavioral Economics.}

Ernst Fehr, Professor für Mikroökonomie und experimentelle Wirtschaftsforschung an der Universität Zürich, gehört zu den einflussreichsten Ökonomen der Gegenwart. Mit einem h-Index von über 150 und mehr als 200.000 Zitationen ist er einer der meistzitierten Wirtschaftswissenschaftler weltweit.

Entscheidender als die Metriken ist der \textit{Inhalt} seiner Arbeit. Fehr hat -- gemeinsam mit Klaus Schmidt, Urs Fischbacher, Simon Gächter und anderen -- die theoretischen Fundamente entwickelt, auf denen das \BCM{} aufbaut:

\begin{itemize}
    \item \textbf{Fehr \& Schmidt (1999)}: Das Modell der Ungleichheitsaversion, das erklärt, warum Menschen unfaire Verteilungen ablehnen -- selbst wenn sie dadurch schlechter gestellt werden. Dieses Paper ist die theoretische Grundlage des KNU-Moduls.

    \item \textbf{Fischbacher, Gächter \& Fehr (2001)}: Die empirische Identifikation von Kooperationstypen -- Conditional Cooperators (ca. 50\%), Self-Interested (ca. 30\%), Altruisten (ca. 10\%). Diese Base Rates sind direkt in das \BEATRIX-System integriert.

    \item \textbf{Fehr \& Fischbacher (2002)}: Die systematische Analyse sozialer Präferenzen und ihrer Auswirkungen auf Wettbewerb, Kooperation und Anreize.

    \item \textbf{Fehr \& Gächter (2000)}: Die Entdeckung der ``altruistischen Bestrafung'' -- Menschen bestrafen Trittbrettfahrer, auch wenn es sie selbst kostet. Ein Mechanismus, der für das Verständnis von Normdurchsetzung zentral ist.
\end{itemize}

Fehr steht in der Tradition von Kahneman und Tversky, aber mit einem entscheidenden Unterschied: Seine Arbeit ist \textit{formaler}. Wo Kahneman Effekte beschrieb, entwickelte Fehr mathematische Modelle. Diese Orientierung an formaler Theorie macht seine Arbeit besonders geeignet als Grundlage für ein axiomatisches System.

\paragraph{Säule 2: Harald Gall und die Informatik.}

Professor Harald Gall leitet das Institut für Informatik an der Universität Zürich und ist einer der führenden Forscher im Bereich Software Engineering, Software Evolution und AI for Software Engineering.

Seine Relevanz für das \BCM-Projekt liegt in mehreren Dimensionen:

\begin{itemize}
    \item \textbf{Software-Architektur}: Die Frage, wie man ein System wie \BEATRIX{} skalierbar, wartbar und erweiterbar baut, ist keine triviale. Galls Expertise in Software-Evolution adressiert genau diese Herausforderung.

    \item \textbf{AI Engineering}: Die Integration von LLMs in produktionsreife Systeme erfordert spezifisches Know-how -- von Prompt Engineering über Fine-Tuning bis zu Evaluation Pipelines. Die Informatik-Abteilung der UZH bringt diese Kompetenzen ein.

    \item \textbf{Qualitätssicherung}: Wie validiert man, dass ein LLM-basiertes System zuverlässige Outputs produziert? Die Methoden der empirischen Software-Forschung -- Metriken, Benchmarks, statistische Tests -- sind direkt übertragbar.
\end{itemize}

Die Verbindung von Behavioral Economics und Informatik ist nicht offensichtlich -- und genau deshalb so wertvoll. Die meisten Behavioral-Science-Gruppen haben keinen Zugang zu erstklassiger Software-Engineering-Expertise. Die meisten Informatik-Gruppen haben kein Interesse an Verhaltenstheorie. In Zürich existiert beides unter einem institutionellen Dach.

\paragraph{Säule 3: Johannes Luger und die Evidence-Based Strategy.}

Professor Johannes Luger verbindet als Experte für evidenzbasierte Strategie die akademische Forschung mit der strategischen Umsetzung in Organisationen.

Seine Rolle im Zürich-Ökosystem ist die eines \textit{Übersetzers}:

\begin{itemize}
    \item \textbf{Von Axiomen zu Hypothesen}: Wie übersetzt man theoretische Axiome in testbare Vorhersagen? Lugers methodologische Expertise ermöglicht rigorose empirische Validierung.

    \item \textbf{Von Effekten zu Business Outcomes}: Akademische Papers berichten Effektgrößen in Cohen's $d$ oder Prozentpunkten. Unternehmen wollen wissen: Was bringt das in Euro? Die Brücke zwischen diesen Welten zu bauen, erfordert spezifische Expertise.

    \item \textbf{Skalierungs-Strategie}: Wie geht man von einem funktionierenden Prototyp zu einem Produkt, das tausende Nutzer hat? Die Fragen der Skalierung sind strategischer Natur -- nicht nur technisch.

    \item \textbf{RCT-Design für die Praxis}: Randomized Controlled Trials sind der Goldstandard der Evaluation. Aber wie designt man RCTs in realen Organisationen, wo perfekte Randomisierung oft unmöglich ist? Lugers Expertise in ``messy real-world evidence'' ist hier zentral.
\end{itemize}

\paragraph{Säule 4: FehrAdvice \& Partners und die Praxis.}

FehrAdvice \& Partners, gegründet 2008, ist die direkte Brücke zwischen der akademischen Forschung der Fehr-Tradition und der Anwendung in der realen Welt.

Die Firma bringt ein, was keine Universität bieten kann:

\begin{itemize}
    \item \textbf{15+ Jahre Feldexpertise}: Hunderte von Projekten in Finance, Healthcare, Public Policy, Energy und Retail. Dieses Erfahrungswissen -- was funktioniert, was nicht, welche Fallstricke es gibt -- ist nicht publiziert, aber essentiell.

    \item \textbf{Validierte Interventionsmuster}: FehrAdvice hat nicht nur Effekte gemessen, sondern Interventionen \textit{implementiert}. Das Wissen, wie man einen Nudge in einem Schweizer Krankenhaus oder einer deutschen Bank umsetzt, existiert nirgendwo sonst in dieser Dichte.

    \item \textbf{Kunden-Pipeline}: Jedes neue Framework braucht Early Adopters. FehrAdvice hat bestehende Kundenbeziehungen, die als Testfeld für \BEATRIX{} dienen können -- ein enormer Vorteil gegenüber akademischen Projekten, die Anwendungspartner erst suchen müssen.

    \item \textbf{Kommerzialisierungs-Know-how}: Die Frage, wie man ein wissenschaftliches Framework in ein Produkt verwandelt -- Pricing, Lizenzierung, Support, Weiterentwicklung --, erfordert unternehmerische Kompetenz.
\end{itemize}

%------------------------------------------------------------------------------
\subsubsection{Die Kompetenz-Matrix}
%------------------------------------------------------------------------------

Die folgende Matrix visualisiert, wie die vier Säulen die erforderlichen Kompetenzen abdecken:

\begin{table}[h]
\centering
\caption{Kompetenzverteilung in der Zürich-Konstellation}
\label{tab:kompetenzmatrix}
\begin{tabular}{@{}lccccc@{}}
\toprule
\textbf{Kompetenzfeld} & \textbf{Fehr} & \textbf{Gall} & \textbf{Luger} & \textbf{FehrAdvice} & \textbf{Summe} \\
\midrule
Verhaltensökonomische Theorie & $\bullet\bullet\bullet$ & -- & $\bullet$ & $\bullet\bullet$ & stark \\
Software/AI Engineering & -- & $\bullet\bullet\bullet$ & -- & $\bullet$ & stark \\
Empirische Methodik & $\bullet\bullet$ & $\bullet$ & $\bullet\bullet\bullet$ & $\bullet\bullet$ & stark \\
Strategische Skalierung & -- & $\bullet$ & $\bullet\bullet\bullet$ & $\bullet\bullet$ & stark \\
Praxis/Kommerzialisierung & -- & -- & $\bullet$ & $\bullet\bullet\bullet$ & stark \\
\bottomrule
\end{tabular}
\end{table}

Die Tabelle zeigt: Jedes Kompetenzfeld ist mindestens durch eine Säule mit höchster Expertise ($\bullet\bullet\bullet$) abgedeckt. Gleichzeitig gibt es Überlappungen, die Redundanz und damit Robustheit schaffen.

%------------------------------------------------------------------------------
\subsubsection{Warum andere Akteure diese Kombination nicht haben}
%------------------------------------------------------------------------------

Um die Einzigartigkeit der Zürich-Konstellation zu verstehen, lohnt ein Blick auf die Alternativen:

\paragraph{Big Tech (Google DeepMind, OpenAI, Anthropic).}

Die großen KI-Labore haben erstklassige technische Kompetenz. Was ihnen fehlt:
\begin{itemize}
    \item Kein Fokus auf Behavioral Economics -- ihre Agenda ist ``AI Alignment'', nicht Verhaltensmodellierung
    \item Keine akademische Incentive-Struktur für grundlegende Theoriearbeit
    \item Plattform-Interessen, die Neutralität gefährden
    \item Keine Praxis-Expertise in Verhaltensinterventionen
\end{itemize}

\paragraph{Reine Akademie (Stanford, MIT, Oxford).}

Elite-Universitäten haben exzellente Forscher. Was ihnen fehlt:
\begin{itemize}
    \item Kein Zugang zu Praxis-Partnern für Validierung
    \item Langsame Publikationszyklen (Jahre statt Monate)
    \item Incentive-Strukturen, die Integration bestrafen (``Salamitaktik'' bei Publikationen)
    \item Keine Kommerzialisierungs-Kompetenz
\end{itemize}

\paragraph{Beratungen (McKinsey, BCG, Accenture).}

Management-Beratungen haben Marktzugang und Praxis-Know-how. Was ihnen fehlt:
\begin{itemize}
    \item Keine wissenschaftliche Tiefe -- ihre ``Frameworks'' sind oft oberflächlich
    \item Keine Incentive für Publikation und Transparenz
    \item Client-spezifische Lösungen statt generalisierbarer Systeme
    \item IP bleibt bei Kunden oder ist proprietär
\end{itemize}

\paragraph{Behavioral Insights Units (BIT, ideas42).}

Die spezialisierten Nudge-Units haben Praxis-Expertise. Was ihnen fehlt:
\begin{itemize}
    \item Fokus auf Policy statt auf universelle Theorie
    \item Keine LLM-Kompetenz
    \item Fragmentierte, nicht-formalisierte Frameworks
    \item Abhängigkeit von staatlicher Finanzierung
\end{itemize}

\begin{table}[h]
\centering
\caption{Wettbewerbsvergleich: Wer hat welche Kompetenzen?}
\label{tab:wettbewerb}
\begin{tabular}{@{}lccccc@{}}
\toprule
\textbf{Akteur} & \textbf{Theorie} & \textbf{Tech} & \textbf{Methodik} & \textbf{Praxis} & \textbf{Gesamt} \\
\midrule
Big Tech & $\circ$ & $\bullet\bullet\bullet$ & $\bullet$ & $\circ$ & 4/12 \\
Elite-Universitäten & $\bullet\bullet\bullet$ & $\bullet$ & $\bullet\bullet\bullet$ & $\circ$ & 7/12 \\
Management-Beratungen & $\circ$ & $\bullet$ & $\bullet$ & $\bullet\bullet\bullet$ & 5/12 \\
Nudge-Units & $\bullet\bullet$ & $\circ$ & $\bullet\bullet$ & $\bullet\bullet$ & 6/12 \\
\textbf{Zürich-Konstellation} & $\bullet\bullet\bullet$ & $\bullet\bullet\bullet$ & $\bullet\bullet\bullet$ & $\bullet\bullet\bullet$ & \textbf{12/12} \\
\bottomrule
\end{tabular}
\end{table}

%------------------------------------------------------------------------------
\subsubsection{Die ``Valley of Death'' und ihre Überwindung}
%------------------------------------------------------------------------------

\epigraph{\textit{``Basic research is what I am doing when I don't know what I am doing.''}}{--- Wernher von Braun}

In der Innovationsforschung ist das ``Valley of Death'' ein bekanntes Phänomen: Die Lücke zwischen erfolgreicher Grundlagenforschung und erfolgreichem Produkt, in der vielversprechende Ideen sterben.

Die Verhaltensökonomie leidet besonders unter diesem Problem. Es gibt brillante Forschung -- Nobelpreise für Kahneman (2002), Thaler (2017), bald vielleicht Fehr -- aber die Übersetzung in skalierbare Produkte bleibt aus.

Warum? Weil die Übersetzung \textit{vier} Brücken erfordert:

\begin{enumerate}
    \item \textbf{Theorie → Formalisierung}: Verbale Beschreibungen müssen in mathematische Modelle überführt werden.
    \item \textbf{Formalisierung → Software}: Mathematische Modelle müssen in Code implementiert werden.
    \item \textbf{Software → Validierung}: Software muss empirisch getestet werden.
    \item \textbf{Validierung → Markt}: Validierte Software muss kommerzialisiert werden.
\end{enumerate}

Die Zürich-Konstellation hat für \textit{jede} dieser Brücken einen designierten Experten:

\begin{table}[h]
\centering
\caption{Die vier Brücken über das Valley of Death}
\begin{tabular}{@{}lll@{}}
\toprule
\textbf{Brücke} & \textbf{Von → Nach} & \textbf{Verantwortung} \\
\midrule
1 & Theorie → Formalisierung & Fehr \\
2 & Formalisierung → Software & Gall \\
3 & Software → Validierung & Luger \\
4 & Validierung → Markt & FehrAdvice \\
\bottomrule
\end{tabular}
\end{table}

Diese Arbeitsteilung ist kein Zufall, sondern das Ergebnis jahrelanger Zusammenarbeit im Zürcher Ökosystem.

%------------------------------------------------------------------------------
\subsubsection{Konkrete Rollenverteilung im BCM/BEATRIX-Projekt}
%------------------------------------------------------------------------------

Die folgende Tabelle zeigt, wie die Verantwortlichkeiten für die einzelnen Komponenten des \BCM{} 2.2 verteilt sind:

\begin{table}[h]
\centering
\caption{Rollenverteilung im BCM/BEATRIX-Projekt}
\begin{tabular}{@{}lll@{}}
\toprule
\textbf{Komponente} & \textbf{Primär} & \textbf{Sekundär} \\
\midrule
META-Modul (Axiome) & Fehr & FehrAdvice \\
Operative Module (INU, KNU, IDN) & Fehr + FehrAdvice & Luger \\
ESTIMATION-Modul (LLM-Pipeline) & Gall & FehrAdvice \\
Empirische Validierung & Luger & Fehr \\
BEATRIX-Produkt & FehrAdvice & Gall \\
Akademische Publikation & Fehr + Luger & Gall \\
Go-to-Market & FehrAdvice & Luger \\
\bottomrule
\end{tabular}
\end{table}

%------------------------------------------------------------------------------
\subsubsection{Der Zeitpunkt: Warum gerade jetzt?}
%------------------------------------------------------------------------------

Die Zürich-Konstellation existiert nicht erst seit gestern. Warum entsteht das \BCM{} META-Modul gerade \textit{jetzt}?

Die Antwort liegt in der Konvergenz dreier Entwicklungen:

\paragraph{1. Theoretische Reife (1999--2020).}

Die Fehr-Tradition hat über zwei Jahrzehnte die theoretischen Bausteine geliefert: Ungleichheitsaversion, Conditional Cooperation, soziale Präferenzen, Reziprozität. Diese Konzepte sind mittlerweile etabliert, repliziert und in Lehrbüchern kanonisiert. Die \textit{theoretische Vorarbeit} ist abgeschlossen.

\paragraph{2. Praktische Validierung (2010--2024).}

FehrAdvice hat in 15 Jahren Hunderte von Projekten durchgeführt. Dabei wurde \textit{praktisches Wissen} akkumuliert: Welche Interventionen funktionieren in welchen Kontexten? Wo sind die Grenzen? Was sind typische Fehlerquellen? Dieses Erfahrungswissen ist jetzt so dicht, dass es formalisiert werden kann.

\paragraph{3. Technologische Ermöglichung (2022--heute).}

Die LLM-Revolution ab GPT-3.5/GPT-4 (2022/2023) liefert erstmals die \textit{technologische Infrastruktur}, um axiomatische Verhaltensmodelle zu operationalisieren. Was vorher ``nur Theorie'' gewesen wäre, kann jetzt implementiert werden.

Diese drei Entwicklungen konvergieren im Jahr 2024. Das Fenster ist offen -- aber es wird sich schließen, sobald ein anderer Akteur den Standard setzt.

%------------------------------------------------------------------------------
\subsubsection{Der BEATRIX-Prototyp: Proof of Concept}
%------------------------------------------------------------------------------

\epigraph{\textit{``In theory, theory and practice are the same. In practice, they are not.''}}{--- Yogi Berra (zugeschrieben)}

Ein entscheidender Unterschied zwischen der Zürich-Konstellation und potenziellen Wettbewerbern liegt darin, dass \BEATRIX{} kein Konzeptpapier ist, sondern ein \textbf{funktionierender Prototyp}.

FehrAdvice hat in den Jahren 2022--2024 eine intensive \textbf{explorative Entwicklungsphase} durchlaufen, in der das gesamte Ideenkonstrukt des \BCM{} systematisch ausexperimentiert wurde. Das Ergebnis ist ein lauffähiges System mit folgenden Komponenten:

\paragraph{Die BEATRIX-API.}

Im Kern von \BEATRIX{} steht eine \textbf{REST-API}, die verhaltensökonomische Analysen programmatisch zugänglich macht. Die API ermöglicht:

\begin{itemize}
    \item \textbf{Persona-Analyse}: Eingabe einer textuellen Persona-Beschreibung, Ausgabe der geschätzten BCM-Parameter ($\lambda$, FEPSDE-Gewichte, Bias-Suszeptibilitäten, Kooperationstyp)

    \item \textbf{Szenario-Simulation}: Vorhersage von Verhaltenswahrscheinlichkeiten für definierte Entscheidungssituationen

    \item \textbf{Interventions-Design}: Generierung und Bewertung von Interventionsvorschlägen basierend auf Persona und Zielverhalten

    \item \textbf{Batch-Processing}: Analyse ganzer Kundensegmente oder Mitarbeiterpopulationen
\end{itemize}

Die API-Architektur folgt dem Prinzip der \textit{Constrained Generation}: Das LLM operiert nicht frei, sondern wird durch die META-Axiome, empirische Base Rates und Pydantic-Schemata eingeschränkt. Dies reduziert Halluzinationen und gewährleistet Konsistenz.

\paragraph{Kontextvektoren: Die semantische Repräsentation.}

Ein zentrales technisches Konzept sind die \textbf{Kontextvektoren} -- hochdimensionale Repräsentationen, die den verhaltensrelevanten Kontext einer Situation erfassen:

\begin{equation}
\vec{c} = f_{\text{embed}}(\rho_{\text{Situation}}, \rho_{\text{Akteur}}, \rho_{\text{Umfeld}}, \rho_{\text{Zeit}})
\end{equation}

Diese Vektoren ermöglichen:

\begin{itemize}
    \item \textbf{Ähnlichkeitssuche}: Finden von analogen Situationen in der Wissensbasis (``In welchen früheren Projekten hatten wir eine ähnliche Konstellation?'')

    \item \textbf{Transfer Learning}: Übertragung von Erkenntnissen aus einem Kontext in einen anderen, unter Berücksichtigung der kontextuellen Distanz

    \item \textbf{Clustering}: Automatische Segmentierung von Personas oder Situationen basierend auf verhaltensrelevanten Merkmalen

    \item \textbf{Anomalie-Detektion}: Identifikation von Fällen, die außerhalb des trainierten Bereichs liegen und besondere Vorsicht erfordern
\end{itemize}

Die Kontextvektoren operationalisieren das ontologische Axiom MA-ONT-04 (``Kontext existiert'') in einer für LLMs verarbeitbaren Form.

\paragraph{Die explorative Phase: Lessons Learned.}

Die zweijährige Experimentierphase hat wertvolle Erkenntnisse geliefert, die in das aktuelle \BCM{} 2.2 eingeflossen sind:

\begin{enumerate}
    \item \textbf{Few-Shot Anchoring ist kritisch}: Ohne explizite Kalibrierungsbeispiele im Prompt variieren LLM-Schätzungen erheblich. Die Integration von Base Rates und Referenzfällen stabilisiert die Outputs.

    \item \textbf{Ensemble-Methoden erhöhen Reliabilität}: Einzelne LLM-Runs haben hohe Varianz. Erst durch N-fache Wiederholung und Median-Aggregation erreicht das System akzeptable Inter-Rater-Reliabilität (IRR > 0.85).

    \item \textbf{Structured Output ist unverzichtbar}: Freie Textgenerierung führt zu nicht-parseabaren Ergebnissen. Pydantic-Schemata erzwingen strukturierte, maschinenlesbare Outputs.

    \item \textbf{Bimodal Detection funktioniert}: Die Unterscheidung zwischen ``unsicher'' (unimodale Verteilung) und ``ambivalent'' (bimodale Verteilung) ist operationalisierbar und diagnostisch wertvoll.

    \item \textbf{Social Desirability muss korrigiert werden}: Persona-Beschreibungen überschätzen systematisch prosoziale Eigenschaften. Explizite Korrekturterme verbessern die Prädiktionsqualität.

    \item \textbf{Die 5-Ebenen-Hierarchie ist praktisch relevant}: Interventionen scheitern oft, weil sie die falsche Ebene adressieren. Die explizite Modellierung von META/MAKRO/MESO/MIKRO/IND verbessert das Interventionsdesign.
\end{enumerate}

\paragraph{Vom Prototyp zum Produkt.}

Der aktuelle Stand ist ein \textit{funktionierender Prototyp} -- technisch lauffähig, aber noch nicht produktionsreif. Die Weiterentwicklung zum skalierbaren Produkt erfordert:

\begin{itemize}
    \item \textbf{Robustheit}: Härtung gegen Edge Cases, adversariale Inputs und unerwartete Nutzungsmuster
    \item \textbf{Skalierbarkeit}: Architektur für hohe Nutzerzahlen und niedrige Latenz
    \item \textbf{Compliance}: DSGVO-Konformität, Audit-Trails, Erklärbarkeit
    \item \textbf{UX}: Intuitive Interfaces für Nicht-Experten
    \item \textbf{Dokumentation}: Umfassende API-Docs, Tutorials, Best Practices
\end{itemize}

Die gute Nachricht: Diese Herausforderungen sind \textit{bekannte Probleme} des Software Engineering -- keine Forschungsfragen. Mit der Expertise von Harald Gall und dem Institut für Informatik der UZH sind sie lösbar.

\paragraph{Der strategische Wert des Prototyps.}

Die Existenz eines funktionierenden Prototyps hat mehrere strategische Implikationen:

\begin{enumerate}
    \item \textbf{Proof of Feasibility}: Das Konzept ist nicht nur theoretisch plausibel, sondern praktisch umsetzbar. Skeptiker können überzeugt werden.

    \item \textbf{Iterationsbasis}: Weiterentwicklung kann auf einem funktionierenden System aufbauen, statt bei Null zu beginnen.

    \item \textbf{Early Adopter}: Ausgewählte Kunden können bereits mit dem Prototyp arbeiten und Feedback liefern.

    \item \textbf{IP-Sicherung}: Funktionierender Code ist konkreter Intellectual Property -- schwerer zu kopieren als ein Paper.

    \item \textbf{Investoren-Attraktivität}: Für potenzielle Partner oder Investoren ist ein Prototyp überzeugender als ein Konzept.
\end{enumerate}

Die Zürich-Konstellation startet nicht bei Null. Sie startet mit einem zweijährigen Vorsprung, kodifiziert in funktionierendem Code.

%------------------------------------------------------------------------------
\subsubsection{Die strategische Positionierung}
%------------------------------------------------------------------------------

\epigraph{\textit{``In the middle of difficulty lies opportunity.''}}{--- Albert Einstein (zugeschrieben)}

Die Zürich-Konstellation hat eine klare strategische Positionierung:

\begin{itemize}
    \item \textbf{Offen, aber nicht Open Source}: Das \BCM{} wird publiziert und transparent gemacht, aber die \BEATRIX-Implementierung ist kommerziell. Dies folgt dem Modell erfolgreicher Standards (vgl. HTTP ist offen, aber Cloudflare monetarisiert).

    \item \textbf{Akademisch, aber nicht akademisiert}: Publikation in Top-Journals sichert Legitimität, aber die Entwicklung folgt nicht dem langsamen akademischen Rhythmus. Preprints und Working Papers ermöglichen schnelle Iteration.

    \item \textbf{Europäisch, aber nicht provinziell}: DSGVO-Kompatibilität und europäische Werte (Transparenz, Autonomie) sind Differenzierungsmerkmale gegenüber US-Big-Tech. Aber der Anspruch ist global.

    \item \textbf{Wissenschaftlich fundiert, aber praxisorientiert}: Jedes Axiom ist literaturbasiert, aber die Auswahl folgt praktischer Relevanz, nicht akademischer Vollständigkeit.
\end{itemize}

Die Zürich-Konstellation ist nicht die einzige, die ein verhaltensökonomisches AI-Framework bauen \textit{könnte}. Aber sie ist die einzige, die alle vier erforderlichen Kompetenzen auf Weltklasse-Niveau \textit{vereint} -- und die Incentives hat, sie auch tatsächlich zu integrieren.

\subsection{Architektur des BCM}

Die Gesamtarchitektur des \BCM{} 2.2 lässt sich als Schichtenmodell darstellen:

\begin{figure}[h]
\centering
\begin{tikzpicture}[
    layer/.style={rectangle, draw, minimum width=12cm, minimum height=1.2cm, align=center},
    arrow/.style={-{Stealth}, thick}
]
    % Layers
    \node[layer, fill=blue!10] (meta) at (0,4) {\textbf{META} (115 Axiome) \\ ONT | EPI | GOV | DYN};
    \node[layer, fill=green!10] (ops) at (0,2.5) {\textbf{Operative Module} \\ INU | KNU | IDN | LAMBDA};
    \node[layer, fill=orange!10] (est) at (0,1) {\textbf{ESTIMATION} (55 Axiome) \\ LLM-basierte Parameterschätzung};
    \node[layer, fill=red!10] (app) at (0,-0.5) {\textbf{BEATRIX Application} \\ Verhaltensvorhersage \& Intervention};

    % Arrows
    \draw[arrow] (meta) -- (ops) node[midway, right] {validiert};
    \draw[arrow] (ops) -- (est) node[midway, right] {parametrisiert};
    \draw[arrow] (est) -- (app) node[midway, right] {operationalisiert};
\end{tikzpicture}
\caption{Schichtenarchitektur des BCM 2.2}
\label{fig:architecture}
\end{figure}

\subsection{Was ist das META-Modul? Eine Übersicht}

\epigraph{\textit{``Give me a place to stand, and I will move the Earth.''}}{--- Archimedes von Syrakus (ca. 287--212 v. Chr.)}

Bevor wir in die Details der einzelnen Kategorien eintauchen, ist es hilfreich, das \META-Modul als Ganzes zu verstehen: Was ist es, warum brauchen wir es, und wie funktioniert es im Zusammenspiel mit \BEATRIX?

\subsubsection{Die Kernidee: Explizite Fundamente}

Die Verhaltensökonomie seit Kahneman \& Tversky (1979) hat viele Effekte dokumentiert -- Loss Aversion, Framing, Anchoring, Present Bias -- aber \textbf{nie ein kohärentes alternatives Theoriegebäude} zum Homo Oeconomicus formuliert. Die Disziplin blieb eine Sammlung von Anomalien, ein ``Kuriositätenkabinett'' ohne architektonische Einheit.

Das \META-Modul schließt diese Lücke. Es macht \textbf{explizit, was bisher implizit} blieb -- und schafft damit die Voraussetzung für maschinelle Operationalisierung.

\begin{table}[h]
\centering
\caption{Die vier Kategorien des META-Moduls}
\begin{tabular}{@{}llll@{}}
\toprule
\textbf{Kategorie} & \textbf{Axiome} & \textbf{Kernfrage} & \textbf{Beispiele} \\
\midrule
\textbf{ONT} (Ontologie) & 26 & Was existiert? & Subjekt, Verhalten, FEPSDE, 5-Ebenen \\
\textbf{EPI} (Epistemologie) & 38 & Was können wir wissen? & Biases, System 1/2, AWX-WAX-WTX \\
\textbf{GOV} (Governance) & 25 & Was ist erlaubt? & Autonomie, 8 Interventionstypen \\
\textbf{DYN} (Dynamik) & 25 & Wie verändern sich Dinge? & Tipping Points, Habit Loops \\
\midrule
\textbf{Total} & \textbf{115} & & \\
\bottomrule
\end{tabular}
\end{table}

\subsubsection{Die historische Analogie}

Das \META-Modul steht in einer langen Tradition axiomatischer Wissenschaft:

\begin{table}[h]
\centering
\caption{Axiomatische Fundamente in der Wissenschaftsgeschichte}
\begin{tabular}{@{}llll@{}}
\toprule
\textbf{Disziplin} & \textbf{Axiomatisierung} & \textbf{Axiome} & \textbf{Konsequenz} \\
\midrule
Geometrie & Euklid (300 v. Chr.) & 5 Postulate & 465 Sätze ableitbar \\
Arithmetik & Peano (1889) & 5 Axiome & Alle natürlichen Zahlen \\
Wahrscheinlichkeit & Kolmogorov (1933) & 3 Axiome & Moderne Stochastik \\
Entscheidungstheorie & vNM (1944) & 4 Axiome & Expected Utility \\
\textbf{Verhaltensökonomie} & \textbf{BCM META (2024)} & \textbf{115 Axiome} & \textbf{BEATRIX-System} \\
\bottomrule
\end{tabular}
\end{table}

Die Lektion aus der Wissenschaftsgeschichte: \textbf{Explizite Axiome ermöglichen systematische Variation.} Nur weil Euklid das Parallelenpostulat explizit machte, konnten Lobatschewski und Riemann es ändern und nicht-euklidische Geometrien entdecken. Implizite Annahmen können nicht systematisch variiert werden -- man weiß nicht einmal, dass man sie macht.

\subsubsection{Die vier Kernfragen}

Jede der vier Kategorien adressiert eine fundamentale Frage der Verhaltensmodellierung:

\paragraph{ONT: Was existiert?}
Die Ontologie definiert die ``Möbel'' des verhaltensökonomischen Universums:
\begin{itemize}
    \item \textbf{Subjekte} ($\sigma$): Wer trifft Entscheidungen?
    \item \textbf{Verhalten} ($V$): Was können Subjekte tun?
    \item \textbf{Nutzen} ($U$): Was bewerten Subjekte?
    \item \textbf{Dimensionen} (FEPSDE): Financial, Emotional, Physical, Social, Digital, Ecological
    \item \textbf{Ebenen}: META $\rightarrow$ MAKRO $\rightarrow$ MESO $\rightarrow$ MIKRO $\rightarrow$ IND
\end{itemize}

\paragraph{EPI: Was können wir wissen?}
Die Epistemologie definiert die Grenzen menschlicher Rationalität:
\begin{itemize}
    \item \textbf{Bounded Rationality}: Rationalität ist durch Kapazität, Zeit und Information beschränkt
    \item \textbf{Zwei-Systeme-Theorie}: System 1 (schnell, automatisch) vs. System 2 (langsam, deliberativ)
    \item \textbf{Systematische Biases}: 26 formalisierte kognitive Verzerrungen
    \item \textbf{AWX-WAX-WTX}: Der Pfad von Awareness über Willingness zu Transition
    \item \textbf{Evidenzhierarchie}: Klassifikation der Robustheit von Parametern
\end{itemize}

\paragraph{GOV: Was ist erlaubt?}
Die Governance definiert ethische Leitplanken für Interventionen:
\begin{itemize}
    \item \textbf{Autonomie-Primat}: Individuelle Freiheit hat Vorrang
    \item \textbf{Nicht-Normativität}: Das BCM beschreibt, es schreibt nicht vor
    \item \textbf{Transparenz}: Verdeckte Manipulation ist zu vermeiden
    \item \textbf{Subsidiarität}: Die mildeste wirksame Intervention ist zu wählen
    \item \textbf{8 Interventionstypen}: NUDGE, BOOST, THINK, INCENTIVE, REGULATE, DESIGN, SOCIAL, IDENTITY
\end{itemize}

\paragraph{DYN: Wie verändern sich Dinge?}
Die Dynamik definiert zeitliche Prozesse und Veränderungsmechanismen:
\begin{itemize}
    \item \textbf{Pfadabhängigkeit}: Geschichte beeinflusst Gegenwart
    \item \textbf{Multiple Equilibria}: Es gibt mehrere stabile Zustände
    \item \textbf{Tipping Points}: Kleine Änderungen können große Wirkungen haben
    \item \textbf{Habit Loops}: Cue $\rightarrow$ Routine $\rightarrow$ Reward
    \item \textbf{Hyperbolic Discounting}: Zeitinkonsistente Präferenzen
\end{itemize}

\subsubsection{Die Rolle im BEATRIX-System}

Das \META-Modul erfüllt vier kritische Funktionen für \BEATRIX:

\begin{enumerate}
    \item \textbf{Constraining the LLM}: Die 115 Axiome geben dem LLM ``Leitplanken''. Es kann nicht beliebig halluzinieren, sondern muss innerhalb des axiomatischen Rahmens operieren. Wenn \BEATRIX{} einen Parameter schätzt, muss diese Schätzung mit den \META-Axiomen konsistent sein.

    \item \textbf{Base Rates}: Die EPI-Axiome liefern empirisch fundierte Default-Werte. Zum Beispiel:
    \begin{itemize}
        \item 50\% Conditional Cooperators (Fischbacher et al., 2001)
        \item Loss Aversion $\lambda \approx 2.0$ (Kahneman \& Tversky, 1979)
        \item Present Bias $\beta \approx 0.7$ (Laibson et al., 1997)
    \end{itemize}
    Das LLM startet nicht bei Null, sondern bei wissenschaftlich fundierten Ankerwerten.

    \item \textbf{Evidenzhierarchie}: Das neue Axiom MA-EPI-38 klassifiziert, wie robust welche Parameter sind:
    \begin{itemize}
        \item $\star\star\star\star$: Labor + Feld konvergent (z.B. Loss Aversion als Mechanismus)
        \item $\star\star\star$: Replizierte Labor-RCTs (z.B. Conditional Cooperators)
        \item $\star\star$: Einzelne RCTs (z.B. spezifische Nudge-Effekte)
        \item $\star$: Nur theoretische Ableitung (z.B. FEPSDE-Gewichtungen)
    \end{itemize}
    \BEATRIX{} kann Konfidenzintervalle der Evidenzqualität anpassen.

    \item \textbf{Ethische Grenzen}: Die GOV-Axiome definieren, welche Interventionen \BEATRIX{} vorschlagen darf und welche nicht. Das System kann keine manipulativen oder autonomieverletzenden Maßnahmen empfehlen, ohne die Governance-Constraints zu verletzen.
\end{enumerate}

\subsubsection{Zusammenfassung: Das META-Modul auf einen Blick}

\begin{tcolorbox}[colback=blue!5!white, colframe=blue!75!black, title=Das META-Modul in einem Satz]
Das \META-Modul ist das \textbf{axiomatische Fundament} des BCM -- 115 explizite Annahmen über Existenz (ONT), Wissen (EPI), Erlaubtes (GOV) und Veränderung (DYN), die alle operativen Module validieren und dem \BEATRIX-LLM wissenschaftlich fundierte Leitplanken geben.
\end{tcolorbox}

\subsection{Die Meta-Formel}

Die übergreifende Struktur des \META-Moduls wird durch die \textbf{Meta-Formel} ausgedrückt:

\begin{equation}
\Phi_{\text{META}} = \Omega_{\text{ONT}} \cdot \Omega_{\text{EPI}} \cdot \Omega_{\text{GOV}} \cdot \Psi(U, \omega, \rho, \sigma) \cdot \Gamma \cdot \tau \cdot \partial \cdot \mu \cdot \kappa \cdot \lambda
\label{eq:meta_formula}
\end{equation}

wobei:
\begin{itemize}
    \item $\Omega_{\text{ONT}}$: Ontologische Constraints (was existiert)
    \item $\Omega_{\text{EPI}}$: Epistemologische Constraints (was wissbar ist)
    \item $\Omega_{\text{GOV}}$: Governance-Constraints (was erlaubt ist)
    \item $\Psi$: Nutzenfunktion mit Gewichten $\omega$, Kontext $\rho$, Subjekt $\sigma$
    \item $\Gamma$: Verhaltensbeschränkungen
    \item $\tau$: Zeit
    \item $\partial$: Veränderungsoperator
    \item $\mu, \kappa, \lambda$: Modulationsparameter
\end{itemize}

%==============================================================================
\subsection{BCM/BEATRIX in der Beratungspraxis: Satisficing statt Wahrheitsfindung}
%==============================================================================

\epigraph{\textit{``The goal is not to be right, but to be useful.''}}{--- George Box (paraphrasiert)}

Bevor wir in die technischen Details der vier Axiom-Kategorien eintauchen, ist eine \textbf{kritische Reflexion} über die Rolle von BCM/BEATRIX in der Beratungspraxis notwendig. Diese Reflexion korrigiert eine zu akademische Perspektive und positioniert das System realistisch.

\subsubsection{Die Realität der Beratung: Verhandelte Rationalität}

\paragraph{Die falsche Annahme.}
Eine naive Sicht betrachtet Berater als ``Wissenschaftler'', die:
\begin{itemize}
    \item Die ``wahren'' Verhaltensparameter finden wollen
    \item Objektive Prognosen erstellen
    \item Evidenz über alles stellen
    \item Am ``richtigen'' Ergebnis gemessen werden
\end{itemize}

\paragraph{Die Realität.}
Berater sind \textbf{Dienstleister}, deren primäre Funktion nicht Wahrheitsfindung, sondern \textbf{Satisficing} des Kunden ist:
\begin{itemize}
    \item \textbf{Legitimation}: ``McKinsey sagt, wir sollen X tun'' ermöglicht Management, Entscheidungen zu rechtfertigen
    \item \textbf{Story-Telling}: Komplexe Situationen werden in kohärente Narrative überführt
    \item \textbf{Belief-Validierung}: Management hat Intuition, Berater liefert ``Evidenz''
    \item \textbf{Konsensbildung}: Verschiedene Stakeholder-Perspektiven werden auf gemeinsame Faktenbasis gebracht
    \item \textbf{Politische Munition}: Externe Analyse gibt internen Machtkämpfen Legitimation
\end{itemize}

Das Ziel ist nicht: ``Was ist wahr?''\\
Das Ziel ist: ``Was macht den Kunden entscheidungsfähig?''

\paragraph{Wie BCM/BEATRIX jede Funktion unterstützt.}

\begin{table}[h]
\centering
\caption{Die fünf Funktionen von Beratung und BCM-Unterstützung}
\small
\begin{tabular}{@{}p{2.5cm}p{4.5cm}p{5.5cm}@{}}
\toprule
\textbf{Funktion} & \textbf{Klassisch} & \textbf{Mit BCM/BEATRIX} \\
\midrule
\textbf{Legitimation} &
``Die Berater haben analysiert'' &
``Die Analyse basiert auf 115 Axiomen, kalibriert an empirischen Base Rates'' -- \textit{stärkere} Legitimation \\[0.5em]

\textbf{Story-Telling} &
Narrative aus Erfahrung und Intuition &
Narrative mit \textit{expliziter} Mechanismus-Kette: ``Loss Aversion $\rightarrow$ höhere Churn'' -- überzeugendere Story \\[0.5em]

\textbf{Belief-Validierung} &
``Unsere Analyse bestätigt...'' (vage) &
``Ihre Intuition ergibt $\lambda = 1.8$, empirisch wäre 2.0 -- Sie liegen im plausiblen Bereich'' -- \textit{kalibrierte} Validierung \\[0.5em]

\textbf{Konsensbildung} &
Verhandlung bis Kompromiss &
Verschiedene Beliefs werden \textit{explizit}: ``CFO: $\lambda = 2.5$, CMO: $\lambda = 1.5$ -- woher kommt die Differenz?'' \\[0.5em]

\textbf{Politische Munition} &
``Externe Analyse zeigt...'' &
``Die Parameter-Sensitivität zeigt, dass $\lambda$ kritisch ist -- \textit{wer hat Daten dazu}?'' -- objektiviert Diskussion \\
\bottomrule
\end{tabular}
\end{table}

\textbf{Kernerkenntnis}: BCM/BEATRIX \textit{ersetzt} keine dieser Funktionen -- es macht sie alle \textbf{robuster} und \textbf{nachvollziehbarer}.

\subsubsection{Wie Verhaltensannahmen wirklich entstehen}

\paragraph{Die ``Methodik'' der klassischen Beratung.}
Ein typischer Workshop illustriert den Prozess:

\begin{quote}
\textit{Berater}: ``Wie werden Ihre Kunden auf die Preiserhöhung reagieren?''\\
\textit{CMO}: ``Unsere treuen Kunden werden bleiben. Die sind loyal.''\\
\textit{CFO}: ``Ich rechne mit 5\% Volumenverlust.''\\
\textit{CEO}: ``Maximal 3\%. Unsere Marke ist stark.''\\[0.5em]
\textit{Berater}: ``Das deckt sich mit unserer Erfahrung aus ähnlichen Projekten. Wir würden konservativ mit 4--5\% rechnen, dann haben wir einen Puffer.''
\end{quote}

Was hier passiert:
\begin{enumerate}
    \item \textbf{Belief Anchoring}: Berater nimmt Management-Schätzungen als Anker (CEO: 3\%, CFO: 5\% $\rightarrow$ Berater: ``4--5\%'')
    \item \textbf{Experience Overlay}: Berater ergänzt eigene Intuition (``ähnliche Projekte'' = vage Referenz)
    \item \textbf{Plausibilitäts-Check}: ``Klingt das für alle vernünftig?'' -- Kopfnicken $\rightarrow$ Konsens
    \item \textbf{Safety Margin}: ``Konservativ rechnen'' = Puffer für Unsicherheit
\end{enumerate}

Das Ergebnis ist eine Zahl (4--5\%), die \textbf{niemandem gehört}, aber \textbf{alle akzeptieren} können. Das ist keine Schätzung -- das ist \textbf{Verhandlung}.

\paragraph{Visualisierung: Der Prozess der verhandelten Realität.}

\begin{figure}[h]
\centering
\begin{tikzpicture}[
    node distance=1.5cm,
    box/.style={rectangle, draw, rounded corners, minimum width=2.8cm, minimum height=0.9cm, align=center, font=\small},
    inputbox/.style={box, fill=blue!10},
    processbox/.style={box, fill=orange!20},
    outputbox/.style={box, fill=green!15},
    arrow/.style={->, thick, >=stealth},
    label/.style={font=\scriptsize\itshape, text=gray}
]

% Input boxes (left side)
\node[inputbox] (mgmt) at (0, 1.5) {Management\\Beliefs};
\node[inputbox] (berater) at (0, 0) {Berater-\\Intuition};
\node[inputbox] (erfahrung) at (0, -1.5) {``Ähnliche\\Projekte''};

% Process box (center)
\node[processbox] (verhandlung) at (4.5, 0) {\textbf{Verhandlung}\\im Workshop};

% Output boxes (right side)
\node[outputbox] (zahl) at (9, 0.75) {``4--5\% Churn''};
\node[outputbox] (story) at (9, -0.75) {PowerPoint-\\Story};

% Arrows from inputs to process
\draw[arrow] (mgmt.east) -- node[above, label] {CEO: 3\%} (verhandlung.west |- mgmt);
\draw[arrow] (berater.east) -- node[above, label] {``passt''} (verhandlung.west);
\draw[arrow] (erfahrung.east) -- node[above, label] {vage Referenz} (verhandlung.west |- erfahrung);

% Arrows from process to outputs
\draw[arrow] (verhandlung.east) -- (zahl.west);
\draw[arrow] (verhandlung.east) -- (story.west);

% Labels
\node[below=0.3cm of erfahrung, font=\scriptsize, text=gray] {Inputs (subjektiv)};
\node[below=0.3cm of verhandlung, font=\scriptsize, text=gray] {Konsensprozess};
\node[below=0.3cm of story, font=\scriptsize, text=gray] {Outputs (``objektiv'')};

% Annotation
\node[draw=red!50, dashed, rounded corners, inner sep=3pt, font=\scriptsize, text=red!70!black] at (4.5, -2.5) {Keine dokumentierte Reasoning-Chain};

\end{tikzpicture}
\caption{Der Prozess der ``verhandelten Realität'' in der klassischen Beratung. Die finale Zahl trägt den Anschein von Objektivität, ist aber Produkt eines sozialen Prozesses.}
\label{fig:verhandelte-realitaet}
\end{figure}

\textbf{Das Problem}: Diese ``verhandelte Realität'' ist nicht per se falsch -- oft funktioniert sie gut. Aber sie ist:
\begin{itemize}
    \item \textbf{Nicht nachvollziehbar}: Warum genau 4--5\%? Was waren die Annahmen?
    \item \textbf{Nicht lernfähig}: Wenn es 7\% werden, wissen wir nicht, \textit{warum} wir falsch lagen
    \item \textbf{Personenabhängig}: Ein anderer Berater hätte vielleicht 3--4\% geschätzt
    \item \textbf{Anfällig für HiPPO}: Die Meinung des CEO (3\%) dominiert, CFO (5\%) wird marginalisiert
\end{itemize}

\paragraph{``Rationalität'' als Kommunikations-Konvention.}
Die Annahme, dass ``Kunden rational auf Preissenkung reagieren werden'', ist keine wissenschaftliche Position, sondern ein \textbf{Plausibilitäts-Marker}. Sie signalisiert: ``Wir haben darüber nachgedacht und es ergibt Sinn.'' Sie bedeutet \textit{nicht}: ``Wir haben es wissenschaftlich überprüft.''

\subsubsection{Die Constraints des Beratungsgeschäfts}

Warum funktioniert Beratung so -- und muss so funktionieren?

\begin{enumerate}
    \item \textbf{Kunde muss zustimmen}: Eine Empfehlung, die der Kunde nicht akzeptiert, ist wertlos. Berater müssen innerhalb des \textit{Overton Window} des Kunden arbeiten.

    \item \textbf{Entscheidungsgremium muss überzeugt werden}: Nicht ein Entscheider, sondern ein Gremium mit verschiedenen Interessen, Hintergründen und Risikoappetiten. Die Story muss für \textit{alle} funktionieren.

    \item \textbf{Zeitdruck}: Board Meeting in 6 Wochen $\rightarrow$ Empfehlung muss stehen. ``Good enough'' schlägt ``perfekt aber zu spät''.

    \item \textbf{Folgeaufträge}: Berater will wiederkommen. Kunde darf nicht ``dumm'' aussehen. Beziehung wichtiger als ``Wahrheit''.

    \item \textbf{Präsentierbarkeit}: ``Das muss auf 20 Slides passen.'' Nuancen gehen verloren, Unsicherheit wird weggelassen.
\end{enumerate}

\subsubsection{Wann funktioniert die klassische Methodik?}

\paragraph{Gut funktioniert sie bei:}
\begin{itemize}
    \item Incrementalen Entscheidungen (kleine Anpassungen, Fehler korrigierbar)
    \item Bekanntem Terrain (Management hat genuine Erfahrung)
    \item Schnellem Feedback (Fehler werden schnell sichtbar)
    \item Homogenen Stakeholdern (ähnliche mentale Modelle)
    \item Stabiler Umwelt (Vergangenheit ist guter Prädiktor)
\end{itemize}

\paragraph{Schlecht funktioniert sie bei:}
\begin{itemize}
    \item Großen, irreversiblen Entscheidungen (Merger, Markteintritte)
    \item Neuem Terrain (wo Intuition nicht kalibriert ist)
    \item Langsamem Feedback (Strategien, deren Erfolg erst in Jahren sichtbar wird)
    \item Heterogenen Stakeholdern (Kunden, Kulturen, die Management nicht kennt)
    \item Disruption (Vergangenheit kein Guide für Zukunft)
    \item \textbf{Verhaltensintensiven Interventionen} (Preisgestaltung, Nudging, Change)
\end{itemize}

\subsubsection{Das HiPPO-Problem: Wenn Status entscheidet}

\paragraph{HiPPO = Highest Paid Person's Opinion.}
In der klassischen Beratung hat das HiPPO-Phänomen erhebliche Auswirkungen auf die Qualität von Verhaltensannahmen:

\begin{quote}
\textit{Workshop-Situation}: CEO schätzt Churn auf 3\%, CFO auf 5\%, Marketing-Chef auf 6\%.\\
\textit{Ergebnis}: ``Wir nehmen konservativ 3.5--4\% an.'' \\
\textit{Realität}: Die CEO-Schätzung (3\%) wirkt als Anker, andere Meinungen werden marginalisiert.
\end{quote}

\paragraph{Warum HiPPO problematisch ist.}
\begin{itemize}
    \item \textbf{Keine Expertise-Korrelation}: Gehalt und hierarchische Position korrelieren nicht mit der Qualität von Verhaltensschätzungen
    \item \textbf{Optimismus-Bias}: Führungskräfte sind oft systematisch optimistischer (``unsere Kunden sind loyal'', ``unsere Marke ist stark'')
    \item \textbf{Confirmation Bias}: Berater bestätigen tendenziell die Meinung dessen, der die Rechnung bezahlt
    \item \textbf{Groupthink}: Widerspruch zum CEO ist karriereschädlich -- also schweigt der CFO
    \item \textbf{Fehlende Accountability}: Wenn es schiefgeht, war es ``unvorhersehbar''
\end{itemize}

\paragraph{Wie BCM/BEATRIX HiPPO demokratisiert.}

\begin{table}[h]
\centering
\caption{HiPPO-Dynamik: Klassisch vs. BCM/BEATRIX}
\begin{tabular}{@{}p{6cm}p{6.5cm}@{}}
\toprule
\textbf{Klassisches Setting} & \textbf{Mit BCM/BEATRIX} \\
\midrule
CEO: ``3\%'' $\rightarrow$ Anker & CEO: ``3\%'' $\rightarrow$ ``Das entspricht $\lambda = 1.5$'' \\[0.3em]
CFO: ``5\%'' $\rightarrow$ ``zu pessimistisch'' & CFO: ``5\%'' $\rightarrow$ ``Das entspricht $\lambda = 2.5$'' \\[0.3em]
Berater: ``4\% ist ein guter Kompromiss'' & Berater: ``Die Differenz kommt von unterschiedlichen $\lambda$-Annahmen. Welche Evidenz haben wir?'' \\[0.3em]
Keine Begründung erforderlich & Begründung wird \textit{explizit} eingefordert \\[0.3em]
Status entscheidet & \textbf{Evidenz und Logik} entscheiden \\
\bottomrule
\end{tabular}
\end{table}

\paragraph{Der Demokratisierungseffekt.}
BCM/BEATRIX verschiebt die Diskussion von ``Wer hat Recht?'' zu ``Welche Annahmen sind plausibel?'':

\begin{enumerate}
    \item \textbf{Explizite Parameter}: Jede Meinung wird in BCM-Parameter übersetzt -- dadurch sichtbar und diskutierbar
    \item \textbf{Base Rates als Referenz}: ``Der empirische Durchschnitt für Loss Aversion ist $\lambda = 2.0$. CEO liegt darunter, CFO darüber -- beide müssen begründen warum.''
    \item \textbf{Sensitivitätsanalyse}: ``Mit CEO-Annahmen: 3.2\% Churn. Mit CFO-Annahmen: 5.8\%. Die Differenz ist materiell -- wir brauchen bessere Evidenz.''
    \item \textbf{Dokumentation}: Alle Annahmen werden festgehalten -- späteres ``ich hatte immer Recht'' wird überprüfbar
\end{enumerate}

\textbf{Kernerkenntnis}: BCM/BEATRIX kann HiPPO nicht eliminieren (hierarchische Dynamiken bleiben), aber es kann die \textbf{Kosten} von HiPPO-Entscheidungen \textbf{sichtbar} machen und eine \textbf{faktenbasierte Alternative} bieten.

\subsubsection{Entscheidungsmatrix: Wann welche Methodologie?}

Die folgende Matrix hilft bei der Entscheidung, wann klassische Beratungsmethodik ausreicht und wann BCM/BEATRIX einen kritischen Mehrwert bietet:

\begin{table}[h]
\centering
\caption{Entscheidungsmatrix: Klassische Methodik vs. BCM/BEATRIX}
\small
\begin{tabular}{@{}p{3.5cm}ccp{5cm}@{}}
\toprule
\textbf{Projektmerkmal} & \textbf{Klassisch} & \textbf{BCM} & \textbf{Rationale} \\
\midrule
\multicolumn{4}{l}{\textit{Entscheidungstyp}} \\
\quad Inkrementell & \checkmark & $\circ$ & Fehler schnell korrigierbar \\
\quad Irreversibel & $\circ$ & \checkmark & Hohe Kosten bei Fehleinschätzung \\[0.5em]

\multicolumn{4}{l}{\textit{Verhaltensrelevanz}} \\
\quad Operational & \checkmark & $\circ$ & Verhalten ist nicht kritischer Faktor \\
\quad Verhaltensintensiv & $\circ$ & \checkmark & Verhalten bestimmt Erfolg/Misserfolg \\[0.5em]

\multicolumn{4}{l}{\textit{Erfahrung des Managements}} \\
\quad Bekanntes Terrain & \checkmark & $\circ$ & Intuition ist gut kalibriert \\
\quad Neues Terrain & $\circ$ & \checkmark & Intuition nicht zuverlässig \\[0.5em]

\multicolumn{4}{l}{\textit{Feedback-Geschwindigkeit}} \\
\quad Schnelles Feedback & \checkmark & $\circ$ & Lernen aus Fehlern möglich \\
\quad Langsames Feedback & $\circ$ & \checkmark & Fehler erst nach Jahren sichtbar \\[0.5em]

\multicolumn{4}{l}{\textit{Stakeholder-Heterogenität}} \\
\quad Homogenes Team & \checkmark & $\circ$ & Konsens leicht erreichbar \\
\quad Heterogene Stakeholder & $\circ$ & \checkmark & BCM als gemeinsame Sprache \\[0.5em]

\multicolumn{4}{l}{\textit{Dokumentationsbedarf}} \\
\quad Intern/informell & \checkmark & $\circ$ & Keine Audit-Anforderungen \\
\quad Extern/reguliert & $\circ$ & \checkmark & Nachvollziehbarkeit erforderlich \\[0.5em]

\multicolumn{4}{l}{\textit{Projektbudget}} \\
\quad Klein ($<$100k) & \checkmark & $\circ$ & BCM-Setup lohnt nicht \\
\quad Groß ($>$500k) & $\circ$ & \checkmark & Fehlerkosten rechtfertigen Investment \\
\bottomrule
\end{tabular}
\\[0.5em]
\raggedright\footnotesize
\checkmark = Bevorzugt \quad $\circ$ = Weniger geeignet
\end{table}

\paragraph{Die praktische Entscheidungsregel.}

\begin{tcolorbox}[colback=yellow!5!white, colframe=yellow!75!black, title=Wann BCM/BEATRIX einsetzen?]
\textbf{Grundregel}: Setze BCM/BEATRIX ein, wenn mindestens \textbf{zwei} der folgenden Bedingungen erfüllt sind:
\begin{enumerate}
    \item Die Entscheidung ist \textbf{materiell} ($>$500k CHF Auswirkung)
    \item Der Erfolg hängt \textbf{kritisch} von Verhaltensannahmen ab
    \item Das Management betritt \textbf{neues Terrain} (neue Märkte, Kundengruppen, Produkte)
    \item Es gibt \textbf{kein schnelles Feedback} (Strategie, langfristige Investments)
    \item \textbf{Regulatorische oder Governance-Anforderungen} verlangen Dokumentation
    \item Es gibt \textbf{signifikante Meinungsverschiedenheiten} im Team (HiPPO-Risiko)
\end{enumerate}
\textbf{Korollar}: Für schnelle, kleine, inkrementelle Entscheidungen in bekanntem Terrain ist klassische Methodik oft ausreichend und effizienter.
\end{tcolorbox}

\subsubsection{Die korrigierte Perspektive auf BCM/BEATRIX}

\paragraph{Was BCM/BEATRIX NICHT ist.}
\begin{itemize}
    \item[$\times$] ``Wir ersetzen eure dumme Intuition durch Wissenschaft'' (arrogant und falsch)
    \item[$\times$] Eine Methode, um ``die Wahrheit'' zu finden
    \item[$\times$] Ersatz für lokales Wissen und Erfahrung
    \item[$\times$] Garantie für perfekte Prognosen
\end{itemize}

\paragraph{Was BCM/BEATRIX IST.}
Ein System, das kollektive Intuition \textbf{explizit}, \textbf{strukturiert} und \textbf{lernfähig} macht:

\begin{table}[h]
\centering
\caption{Klassische Beratung vs. BCM/BEATRIX-gestützte Beratung}
\begin{tabular}{@{}p{6.5cm}p{6.5cm}@{}}
\toprule
\textbf{Klassische Beratung} & \textbf{Mit BCM/BEATRIX} \\
\midrule
``4--5\% Churn'' & ``4.2\% Churn [2.8\%, 5.9\%]'' \\[0.3em]
Woher die Zahl? \textit{Verhandlung} & Woher? \textit{Dokumentiert:} Base Rate 3\%, $\lambda = 2.3$, Loss Frame +0.5\% \\[0.3em]
Wie sicher? \textit{Keine Ahnung} & [2.8\%, 5.9\%] = ``wahrscheinlich 3--6\%'' \\[0.3em]
Was sind Treiber? \textit{Implizit} & Sensitivität: $\lambda$ erklärt 60\% der Unsicherheit \\[0.3em]
Was lernen wir bei Fehler? \textit{Nichts systematisch} & Outcome 6.5\% $\rightarrow$ $\lambda$ war unterschätzt $\rightarrow$ Update \\
\bottomrule
\end{tabular}
\end{table}

\paragraph{Die Integration von Intuition und Modell.}
BCM/BEATRIX ist ein \textbf{Verstärker}, kein Ersatz:

\begin{quote}
\textit{Berater}: ``Ihre Intuition, dass Kunden loyal sind, ergibt $\lambda = 1.8$ -- unter dem Durchschnitt von 2.0. Aber: Preiserhöhungen werden als Verlust wahrgenommen, was Loss Aversion aktiviert. Das erhöht $\lambda$ effektiv auf 2.2. Deshalb kommen wir auf 4.2\%, nicht 3\%. Sind Sie einverstanden mit dieser Logik?''
\end{quote}

Was hier passiert:
\begin{itemize}
    \item Intuition wird \textbf{herausgefordert}, aber \textbf{respektiert}
    \item Mechanismus wird \textbf{explizit} gemacht
    \item Management kann \textbf{nachvollziehen} und \textbf{korrigieren}
    \item Beide Wissensquellen werden \textbf{integriert}
\end{itemize}

\begin{axiom}[Intuitions-Integration: MA-ORG-02]
\begin{enumerate}
    \item Management-Intuition und Berater-Erfahrung sind \textbf{wertvolle Informationsquellen}, die lokales Wissen enthalten, das in empirischen Base Rates nicht abgebildet ist.

    \item BCM/BEATRIX \textbf{ersetzt} Intuition nicht, sondern:
    \begin{itemize}
        \item \textbf{Expliziert} sie (macht Annahmen sichtbar)
        \item \textbf{Kalibriert} sie (gegen empirische Base Rates)
        \item \textbf{Strukturiert} sie (in BCM-Dimensionen)
        \item \textbf{Dokumentiert} sie (für Nachvollziehbarkeit und Lernen)
    \end{itemize}

    \item Konflikte zwischen Intuition und Empirie werden \textbf{transparent} gemacht und diskutiert, nicht einseitig aufgelöst.

    \item Der Mehrwert liegt nicht in ``besseren Zahlen'', sondern in:
    \begin{itemize}
        \item Nachvollziehbarkeit der Reasoning-Kette
        \item Quantifizierung der Unsicherheit
        \item Strukturiertem Lernen aus Outcomes
        \item Überbrückung verschiedener mentaler Modelle
    \end{itemize}

    \item ``Satisficing'' des Kunden bleibt ein \textbf{legitimes Ziel}. BCM/BEATRIX macht diesen Prozess \textbf{robuster}, nicht ``wissenschaftlicher''.
\end{enumerate}
\end{axiom}

\subsubsection{Das vollständige Value Proposition}

Die Wertschöpfung von BCM/BEATRIX erstreckt sich über mehrere Dimensionen:

\paragraph{1. Geschwindigkeit und Effizienz.}

\begin{table}[h]
\centering
\caption{Zeitersparnis durch BCM/BEATRIX}
\begin{tabular}{@{}lcc@{}}
\toprule
\textbf{Aktivität} & \textbf{Klassisch} & \textbf{Mit BEATRIX} \\
\midrule
Verhaltensparameter-Schätzung & 2--4 Wochen (Primärforschung) & 2--4 Stunden \\
Szenario-Modellierung & 1--2 Wochen & Echtzeit \\
Sensitivitätsanalyse & Oft gar nicht & Automatisch \\
Dokumentation der Annahmen & Nachträglich, unvollständig & Integriert \\
Cross-Check mit Literatur & Manuell, sporadisch & Systematisch \\
\bottomrule
\end{tabular}
\end{table}

\textbf{Implikation}: Projekte können schneller starten, mehr Iterationen durchlaufen, und Kunden bekommen früher Ergebnisse. Zeit ist Geld -- und Geschwindigkeit ist Wettbewerbsvorteil.

\paragraph{2. Bessere Entscheidungen.}

BCM/BEATRIX verbessert die Entscheidungsqualität durch:

\begin{itemize}
    \item \textbf{Reduktion von Blind Spots}: Systematische Berücksichtigung von 38 dokumentierten Biases verhindert, dass wichtige Verhaltenseffekte übersehen werden
    \item \textbf{Explizite Unsicherheit}: Konfidenzintervalle statt Punktschätzungen ermöglichen risikoadjustierte Entscheidungen
    \item \textbf{Kontrafaktische Analyse}: ``Was wäre wenn $\lambda$ höher ist?'' kann in Sekunden beantwortet werden
    \item \textbf{Pre-Mortem-Fähigkeit}: Identifikation von Implementierungsbarrieren \textit{vor} dem Scheitern
    \item \textbf{Konsistenz}: Dieselbe Situation wird nicht je nach Berater unterschiedlich eingeschätzt
\end{itemize}

\paragraph{3. Höhere Interventionswirksamkeit.}

\begin{table}[h]
\centering
\caption{Hebel für Interventionserfolg}
\begin{tabular}{@{}ll@{}}
\toprule
\textbf{Problem} & \textbf{BCM-Lösung} \\
\midrule
Falsche Diagnose (AWX vs. WAX vs. WTX) & Systematische Engpass-Identifikation \\
Falscher Interventionstyp & 8-Typen-Taxonomie mit Matching \\
Überschätzung von Nudge-Effekten & Lab-Field-Gap-Korrektur (MA-EPI-38) \\
Ignorieren von Identitätsdynamik & IDN-Modul-Integration \\
Fehlende Nachhaltigkeit & DYN-Axiome für Habit Formation \\
\bottomrule
\end{tabular}
\end{table}

\textbf{Quantifizierung}: Wenn BCM/BEATRIX die Interventionswirksamkeit um 20\% steigert (konservative Schätzung basierend auf besserer Diagnose), und ein typisches Projekt 500k CHF Interventionsbudget hat, sind das 100k CHF zusätzlicher Wert -- pro Projekt.

\paragraph{4. Risikoreduktion.}

\begin{itemize}
    \item \textbf{Vermeidung teurer Fehlschläge}: Verhaltensblinde Strategien scheitern oft spektakulär (z.B. Produktlaunches, die Kundenverhalten falsch einschätzen)
    \item \textbf{Früherkennung}: BEATRIX kann warnen, wenn Annahmen unplausibel sind (Plausibility Check)
    \item \textbf{Dokumentierte Reasoning-Chain}: Bei Problemen kann nachvollzogen werden, \textit{warum} eine Entscheidung getroffen wurde
    \item \textbf{Regulatory Compliance}: Für regulierte Branchen (Banken, Versicherungen) ist dokumentierte Methodik zunehmend wichtig
\end{itemize}

\paragraph{5. Skalierbarkeit und Konsistenz.}

\begin{itemize}
    \item \textbf{Über Projekte hinweg}: Dieselbe Methodik für Pricing, Change Management, Produktdesign, Policy
    \item \textbf{Über Teams hinweg}: Junior-Berater können schneller auf Senior-Niveau arbeiten
    \item \textbf{Über Klienten hinweg}: Erkenntnisse aus Projekt A verbessern Projekt B
    \item \textbf{Über Zeit hinweg}: Akkumulierendes Wissen statt wiederholtes Neuerfinden
\end{itemize}

\paragraph{6. Wissensmanagement und Lernen.}

\begin{table}[h]
\centering
\caption{Lernschleifen im BCM/BEATRIX-System}
\begin{tabular}{@{}lll@{}}
\toprule
\textbf{Ebene} & \textbf{Was wird gelernt?} & \textbf{Wie?} \\
\midrule
Projekt & War unsere Prognose korrekt? & Outcome-Tracking \\
Kunde & Welche Parameter sind kundenspezifisch? & Kalibrierung über Zeit \\
Industrie & Gibt es branchenspezifische Patterns? & Cross-Projekt-Analyse \\
Methodik & Welche Base Rates stimmen? & Meta-Analyse der Predictions \\
\bottomrule
\end{tabular}
\end{table}

\textbf{Der Flywheel-Effekt}: Je mehr Projekte mit BCM/BEATRIX durchgeführt werden, desto besser werden die Base Rates, desto präziser die Prognosen, desto höher der Wert für Kunden.

\paragraph{7. Talent und Training.}

\begin{itemize}
    \item \textbf{Schnelleres Onboarding}: Neue Berater lernen BCM-Denken systematisch statt durch Osmose
    \item \textbf{Konsistente Qualität}: Weniger Abhängigkeit von individueller Erfahrung einzelner Partner
    \item \textbf{Codifiziertes Expertenwissen}: Das Wissen von Fehr, Fischbacher et al. ist im System, nicht nur in deren Köpfen
    \item \textbf{Retention}: Wissen bleibt in der Firma, auch wenn Berater gehen
\end{itemize}

\paragraph{8. Differenzierung und Positionierung.}

\begin{itemize}
    \item \textbf{Unique Selling Proposition}: Kein anderer Berater hat ein vergleichbares System
    \item \textbf{Wissenschaftliche Glaubwürdigkeit}: Zürich-Konstellation (ETH, UZH, FehrAdvice) ist einzigartig
    \item \textbf{Defensible Methodology}: Schwer zu kopieren, weil auf Jahrzehnten Forschung aufbauend
    \item \textbf{Premium-Positionierung}: Rechtfertigt höhere Tagessätze
\end{itemize}

\paragraph{9. Client Experience.}

\begin{itemize}
    \item \textbf{Transparenz}: Kunden verstehen, \textit{warum} Empfehlungen gemacht werden
    \item \textbf{Ownership}: Kunden können Parameter anpassen und Konsequenzen sehen
    \item \textbf{Vertrauen}: Wissenschaftliche Fundierung reduziert ``Bauchgefühl''-Skepsis
    \item \textbf{Kontinuität}: Bei Folgeprojekten kann auf vorherige Analysen aufgebaut werden
\end{itemize}

\subsubsection{Quantifizierung des Gesamtwertes}

\begin{table}[h]
\centering
\caption{Illustrative Wertschöpfung pro Projekt (konservative Schätzung)}
\begin{tabular}{@{}lr@{}}
\toprule
\textbf{Werttreiber} & \textbf{CHF} \\
\midrule
Zeitersparnis (40h $\times$ 300 CHF/h) & 12'000 \\
Höhere Interventionswirksamkeit (+10\% auf 500k Budget) & 50'000 \\
Vermiedene Fehlentscheidungen (10\% Wahrscheinlichkeit $\times$ 200k Schaden) & 20'000 \\
Bessere Client Experience (höhere Retention, Upselling) & 15'000 \\
\midrule
\textbf{Geschätzter Mehrwert pro Projekt} & \textbf{97'000} \\
\bottomrule
\end{tabular}
\end{table}

\textit{Anmerkung}: Diese Zahlen sind illustrativ und projektabhängig. Der tatsächliche Wert variiert je nach Projektgröße, Komplexität und Anwendungsbereich.

\subsubsection{Was BCM/BEATRIX NICHT verspricht}

Zur Vollständigkeit -- und um unrealistische Erwartungen zu vermeiden:

\begin{itemize}
    \item[$\times$] \textbf{``Die Wahrheit''}: BCM liefert kalibrierte Schätzungen, keine Gewissheit
    \item[$\times$] \textbf{Perfekte Prognosen}: Menschliches Verhalten bleibt stochastisch
    \item[$\times$] \textbf{Ersatz für Domänenwissen}: BCM braucht lokalen Input, um zu funktionieren
    \item[$\times$] \textbf{Automatisierung von Beratung}: BEATRIX unterstützt Berater, ersetzt sie nicht
    \item[$\times$] \textbf{Eliminierung von Unsicherheit}: Es macht Unsicherheit \textit{sichtbar}, nicht \textit{weg}
\end{itemize}

\begin{tcolorbox}[colback=blue!5!white, colframe=blue!75!black, title=Das BCM/BEATRIX Value Proposition in einem Satz]
BCM/BEATRIX macht verhaltensökonomische Beratung \textbf{schneller} (Stunden statt Wochen), \textbf{besser} (systematisch statt ad-hoc), \textbf{skalierbarer} (über Teams und Projekte), \textbf{lernfähiger} (akkumulierendes Wissen), und \textbf{nachvollziehbarer} (dokumentierte Reasoning-Chain) -- während es die unverzichtbare Rolle von Berater-Expertise und Kunden-Intuition \textbf{verstärkt statt ersetzt}.
\end{tcolorbox}

%==============================================================================


% Kapitel 2: Ontologie (ONT)
\section{Kategorie 1: Ontologie (ONT)}
%==============================================================================

\epigraph{\textit{``The question `What is there?' can be answered, moreover, in a word -- `Everything' -- and everyone will accept this answer as true.''}}{--- Willard Van Orman Quine, \textit{On What There Is} (1948)}

Die Ontologie -- von griechisch \textit{on} (seiend) und \textit{logos} (Lehre) -- ist die philosophische Disziplin, die fragt: \textit{Was existiert?} Welche Entitäten, Eigenschaften und Beziehungen sind fundamental? Was ist die grundlegende Möblierung des Universums?

Für das \BCM{} ist die Ontologie-Kategorie (26 Axiome: MA-ONT-01 bis MA-ONT-26) konstitutiv: Sie definiert die ``Möbel'' des verhaltensökonomischen Universums -- die Bausteine, aus denen alle weiteren Modellierungen zusammengesetzt werden.

%------------------------------------------------------------------------------
\subsection{Philosophischer Hintergrund: Von Aristoteles zu Quine}
%------------------------------------------------------------------------------

\epigraph{\textit{``Being is said in many ways.''}}{--- Aristoteles, \textit{Metaphysik} IV.2 (ca. 350 v. Chr.)}

Die ontologische Tradition in der westlichen Philosophie beginnt mit Aristoteles' \textit{Kategorienlehre}. Aristoteles unterschied zehn fundamentale Kategorien des Seins: Substanz, Quantität, Qualität, Relation, Ort, Zeit, Lage, Haben, Wirken und Leiden. Diese Kategorien sollten erschöpfend und gegenseitig exklusiv sein -- jedes Seiende fällt unter genau eine primäre Kategorie.

Die moderne Ontologie wurde durch Quines Aufsatz ``On What There Is'' (1948) revolutioniert. Quine formulierte das berühmte Kriterium:

\begin{quote}
\textit{``To be is to be the value of a bound variable.''}
\end{quote}

Das bedeutet: Wir verpflichten uns ontologisch auf genau die Entitäten, über die wir in unseren besten Theorien quantifizieren. Das \BCM{} META-Modul folgt diesem quineanischen Ansatz: Die ONT-Axiome explizieren, auf welche Entitäten sich das \BCM{} ontologisch verpflichtet.

%------------------------------------------------------------------------------
\subsection{Fundamentale Entitäten: Das ontologische Fundament}
%------------------------------------------------------------------------------

Die ersten sechs ONT-Axiome etablieren die grundlegenden Existenzaussagen des \BCM:

\begin{axiom}[Existenz des Subjekts: MA-ONT-01]
\begin{equation}
\exists \sigma: \sigma \in \Sigma, |\Sigma| \geq 1
\end{equation}
Es existiert mindestens ein Entscheidungssubjekt $\sigma$.
\end{axiom}

\paragraph{Philosophische Begründung.}
Dieses Axiom macht eine starke Annahme: Es gibt distinkte, identifizierbare Entscheidungsträger. Für das \BCM{} ist das Subjekt fundamental: Ohne Subjekte gibt es niemanden, der Entscheidungen trifft, Nutzen erfährt oder Verhalten zeigt.

\paragraph{BCM-Relevanz.}
Das Subjekt $\sigma$ ist der Anker aller weiteren Modellierung: Das INU-Modul modelliert den individuellen Nutzen \textit{dieses} Subjekts, das KNU-Modul modelliert, wie \textit{dieses} Subjekt den Nutzen anderer berücksichtigt, das IDN-Modul modelliert die Identität(en) \textit{dieses} Subjekts.

\begin{axiom}[Existenz von Verhalten: MA-ONT-02]
\begin{equation}
\exists V: V \in \mathcal{V}, |\mathcal{V}| \geq 1
\end{equation}
Es gibt beobachtbare Verhaltensweisen $V$.
\end{axiom}

\paragraph{Philosophische Begründung.}
Dieses Axiom trennt das \BCM{} von rein mentalistischen Ansätzen. Verhalten ist \textit{beobachtbar} -- nicht nur subjektiv erlebt. Die Menge $\mathcal{V}$ enthält manifeste Handlungen, Unterlassungen, verbale Äußerungen, physiologische Reaktionen und digitale Interaktionen.

\begin{axiom}[Existenz von Nutzen: MA-ONT-03]
\begin{equation}
U: \mathcal{V} \times \Sigma \rightarrow \mathbb{R}
\end{equation}
Subjekte haben Nutzenfunktionen, die Verhalten bewerten.
\end{axiom}

\paragraph{Philosophische Begründung.}
Das \BCM{} folgt der von-Neumann-Morgenstern-Tradition: Nutzen ist eine reellwertige Funktion, die Verhalten relativ zu einem Subjekt bewertet. Die Notation $U: \mathcal{V} \times \Sigma \rightarrow \mathbb{R}$ macht drei Dinge explizit: Nutzen ist \textit{verhaltensrelativ}, \textit{subjektrelativ} und \textit{reellwertig}.

\begin{axiom}[Existenz von Kontext: MA-ONT-04]
\begin{equation}
\exists \rho: \rho \in \mathcal{P}, |\mathcal{P}| \geq 1
\end{equation}
Verhalten findet in Kontexten statt, die den Nutzen modulieren.
\end{axiom}

\paragraph{Philosophische Begründung.}
Dieses Axiom kodifiziert eine der zentralen Einsichten der Verhaltensökonomie: \textit{Context matters}. Der Homo Oeconomicus war kontextblind -- seine Präferenzen waren stabil und situationsunabhängig. Die empirische Forschung hat diese Annahme widerlegt.

\begin{axiom}[Existenz von Zeit: MA-ONT-05]
\begin{equation}
\exists \tau: \tau \in T, T \subseteq \mathbb{R}^+
\end{equation}
Entscheidungen und Verhalten finden in der Zeit statt.
\end{axiom}

\begin{axiom}[Existenz von Constraints: MA-ONT-06]
\begin{equation}
\exists \Gamma: \Gamma(\sigma, \rho) \rightarrow \mathcal{V}' \subseteq \mathcal{V}
\end{equation}
Nicht alle Verhaltensweisen sind für alle Subjekte in allen Kontexten verfügbar.
\end{axiom}

%------------------------------------------------------------------------------
\subsection{Strukturelle Axiome: Die Architektur des Nutzens}
%------------------------------------------------------------------------------

\begin{axiom}[Drei-Komponenten-Nutzen: MA-ONT-07]
\begin{equation}
U = U_{\text{INU}} \oplus U_{\text{KNU}} \oplus U_{\text{IDN}}
\end{equation}
Der Gesamtnutzen setzt sich aus individuellem Nutzen (INU), kollektivem Nutzen (KNU) und Identitätsnutzen (IDN) zusammen.
\end{axiom}

\paragraph{Theoretischer Hintergrund.}
Diese Dreiteilung synthetisiert drei Forschungstraditionen: (1) \textbf{Individueller Nutzen}: Die klassische ökonomische Tradition von Bentham bis Samuelson; (2) \textbf{Kollektiver Nutzen}: Die Forschung zu sozialen Präferenzen (Fehr \& Schmidt, 1999); (3) \textbf{Identitätsnutzen}: Die Identitätsökonomie (Akerlof \& Kranton, 2000).

\begin{axiom}[FEPSDE-Dimensionen: MA-ONT-08]
\begin{equation}
U_{\text{INU}} = \sum_{d \in \text{FEPSDE}} \omega_d \cdot u_d, \quad \sum_{d \in \text{FEPSDE}} \omega_d = 1
\end{equation}
Der individuelle Nutzen hat sechs Dimensionen mit normierter Gewichtung.
\end{axiom}

\begin{table}[h]
\centering
\caption{Die FEPSDE-Nutzen-Dimensionen}
\begin{tabular}{@{}lp{5cm}p{5cm}@{}}
\toprule
\textbf{Dimension} & \textbf{Beschreibung} & \textbf{Beispiele} \\
\midrule
\textbf{F}inancial & Monetärer/materieller Nutzen & Einkommen, Vermögen, Ersparnisse \\
\textbf{E}motional & Affektiver Nutzen & Freude, Zufriedenheit, Stolz \\
\textbf{P}hysical & Körperlicher Nutzen & Gesundheit, Energie, Schmerzfreiheit \\
\textbf{S}ocial & Sozialer Nutzen & Zugehörigkeit, Status, Anerkennung \\
\textbf{D}igital & Digitaler Nutzen & Convenience, Vernetzung, Datenschutz \\
\textbf{E}cological & Ökologischer Nutzen & Nachhaltigkeit, Umweltschutz \\
\bottomrule
\end{tabular}
\end{table}

%------------------------------------------------------------------------------
\subsection{Soziale Ontologie: Gruppen, Identitäten, Normen}
%------------------------------------------------------------------------------

\epigraph{\textit{``No man is an island, entire of itself.''}}{--- John Donne, \textit{Meditation XVII} (1624)}

\begin{axiom}[Multiple Identitäten: MA-ONT-09]
\begin{equation}
\forall \sigma: I(\sigma) = \{i_1, i_2, \ldots, i_n\}, \quad n \geq 1
\end{equation}
Jedes Subjekt hat mindestens eine, typischerweise mehrere Identitäten.
\end{axiom}

\begin{axiom}[Existenz von Gruppen: MA-ONT-10]
\begin{equation}
\exists G: G \subseteq \Sigma, |G| \geq 2
\end{equation}
Es existieren Gruppen als Aggregate von Subjekten.
\end{axiom}

\begin{axiom}[Existenz von Normen: MA-ONT-11]
\begin{equation}
\exists N: N \in \mathcal{N}, \mathcal{N} = \{(V, G, \text{Sanktion})\}
\end{equation}
Normen sind Verhaltenserwartungen mit Sanktionen bei Abweichung.
\end{axiom}

%------------------------------------------------------------------------------
\subsection{Die 5-Ebenen-Hierarchie (MA-ONT-26)}
%------------------------------------------------------------------------------

\begin{axiom}[BCM-Ebenen-Hierarchie: MA-ONT-26]
\begin{equation}
E = \{\text{META}, \text{MAKRO}, \text{MESO}, \text{MIKRO}, \text{IND}\}
\end{equation}
\begin{equation}
\forall e_i, e_j \in E: i < j \Rightarrow e_i \text{ kontextualisiert } e_j
\end{equation}
\end{axiom}

\begin{table}[h]
\centering
\caption{Die 5 BCM-Ebenen im Detail}
\label{tab:ebenen_detail}
\begin{tabular}{@{}lp{3.5cm}p{3.5cm}p{4cm}@{}}
\toprule
\textbf{Ebene} & \textbf{Beschreibung} & \textbf{Zeitskala} & \textbf{Beispiele} \\
\midrule
META & Kulturelle Tiefenstruktur & Generationen & Individualismus/Kollektivismus, Machtdistanz (Hofstede) \\
MAKRO & Institutionelle Ordnung & Dekaden & Verfassung, Rechtssystem, Marktwirtschaft \\
MESO & Organisationale Regeln & Jahre & Unternehmenskultur, Anreizsysteme \\
MIKRO & Gruppenkontext & Monate & Teamdynamik, Familienstrukturen \\
IND & Personenmerkmale & Stabil/variabel & Traits, Präferenzen, Identitäten \\
\bottomrule
\end{tabular}
\end{table}

\paragraph{Das Prinzip der Kontextualisierung.}
Das Axiom postuliert: \textit{Höhere Ebenen kontextualisieren niedrigere Ebenen}. META setzt Defaults für alle tieferen Ebenen; MAKRO kann META-Defaults modulieren, aber nicht vollständig überschreiben; usw.

\paragraph{BEATRIX-Relevanz.}
Das 5-Ebenen-Modell ist zentral für die \BEATRIX-Parameterschätzung: (1) \textbf{Default-Kaskade}: Fehlende Information auf niedriger Ebene wird durch höhere Ebenen kompensiert; (2) \textbf{Plausibilitätsprüfung}: Schätzungen müssen mit höheren Ebenen konsistent sein; (3) \textbf{Interventionsdesign}: Die richtige Ebene für Interventionen muss identifiziert werden.

%==============================================================================


% Kapitel 3: Epistemologie (EPI)
\section{Kategorie 2: Epistemologie (EPI)}
%==============================================================================

\epigraph{\textit{``Ich weiß, dass ich nichts weiß.''}}{--- Sokrates, nach Platon, \textit{Apologie}}

Die Epistemologie-Kategorie (38 Axiome: MA-EPI-01 bis MA-EPI-38) definiert, \textbf{was wir wissen können}, welche Grenzen Wissen hat, und welche systematischen Verzerrungen (Biases) existieren.

\subsection{Philosophischer Hintergrund: Von Platon zu Kahneman}

\epigraph{\textit{``The mind is not a vessel to be filled, but a fire to be kindled.''}}{--- Plutarch, \textit{On Listening}}

Die Frage nach den Grenzen menschlichen Wissens durchzieht die gesamte Philosophiegeschichte. Bereits Platons Höhlengleichnis illustriert, dass Menschen systematisch einer verzerrten Wahrnehmung der Realität unterliegen -- eine Einsicht, die in der modernen Kognitionspsychologie ihre empirische Bestätigung findet.

\paragraph{Die rationalistische Tradition.}
Descartes' \textit{cogito ergo sum} etablierte den methodischen Zweifel als epistemologisches Werkzeug. Interessanterweise führte der Rationalismus zu einer Überschätzung menschlicher Rationalität, die erst durch die Verhaltensökonomie korrigiert wurde. Leibniz' Konzept der \textit{petites perceptions} -- unbewusste Wahrnehmungen, die unser Denken beeinflussen -- nimmt jedoch bereits System-1-Prozesse vorweg.

\paragraph{Der empiristische Einwand.}
David Humes radikaler Empirismus stellte die Frage, ob \textit{irgendwelches} Wissen über kausale Zusammenhänge möglich sei. Seine Analyse der Induktion als psychologischer Gewohnheit (habit) statt logischer Notwendigkeit ist ein direkter Vorläufer der Heuristiken-Forschung. Wenn wir erwarten, dass die Sonne morgen aufgeht, tun wir dies nicht aus logischen Gründen, sondern aufgrund der \textit{Verfügbarkeitsheuristik} vergangener Sonnenaufgänge.

\paragraph{Kants Synthese und die Grenzen der Vernunft.}
Immanuel Kants \textit{Kritik der reinen Vernunft} etablierte, dass menschliche Erkenntnis durch die Struktur unseres Geistes geformt wird. Seine Kategorien der Anschauung (Raum, Zeit) und des Verstandes (Kausalität, Substanz) sind \textit{a priori} -- sie strukturieren unsere Erfahrung, bevor wir sie machen. Dies ist die philosophische Grundlage für die BCM-Annahme, dass Menschen die Realität nicht ``objektiv'' wahrnehmen, sondern durch kognitive Schemata.

\paragraph{Der pragmatische Turn.}
William James und John Dewey betonten, dass Wissen nicht Abbild der Realität ist, sondern Werkzeug zur Bewältigung von Problemen. Eine Überzeugung ist ``wahr'', wenn sie \textit{funktioniert}. Diese Perspektive legitimiert Heuristiken: Sie mögen logisch fehlerhaft sein, aber sie funktionieren ``gut genug'' in den meisten Situationen -- ein Kerngedanke von Gerd Gigerenzers ``ökologischer Rationalität''.

\paragraph{Von der Philosophie zur Psychologie.}
Der Übergang zur empirischen Kognitionspsychologie erfolgte durch:
\begin{itemize}
    \item \textbf{Herbert Simon (1955)}: Bounded Rationality als Alternative zum homo oeconomicus
    \item \textbf{Amos Tversky \& Daniel Kahneman (1974)}: Systematische Dokumentation von Heuristiken und Biases
    \item \textbf{Gerd Gigerenzer (1990er)}: Rehabilitation von Heuristiken als ``adaptiver Werkzeugkasten''
    \item \textbf{Daniel Kahneman (2011)}: Synthese zur Zwei-Systeme-Theorie in \textit{Thinking, Fast and Slow}
\end{itemize}

Das \META-Modul integriert diese Traditionen: Es akzeptiert die \textit{Grenzen} menschlicher Rationalität (gegen den naiven Rationalismus), die \textit{Systematik} kognitiver Verzerrungen (gegen den naiven Empirismus), und die \textit{Funktionalität} von Heuristiken (mit dem Pragmatismus).

\subsection{Grundlegende epistemische Axiome}

\begin{axiom}[Unvollständiges Wissen: MA-EPI-01]
\begin{equation}
K(U) \subset U
\end{equation}
Wissen über den Nutzen $K(U)$ ist immer eine echte Teilmenge des tatsächlichen Nutzens $U$. Vollständiges Wissen ist unerreichbar.

\paragraph{Philosophische Begründung.}
Dies ist die epistemologische Kernaussage, die von Sokrates bis zu den Grenzen des Laplaceschen Dämons reicht. In der Ökonomie manifestiert sie sich als ``Informationsasymmetrie'' (Akerlof, 1970), in der Entscheidungstheorie als ``Unsicherheit vs. Risiko'' (Knight, 1921).

\paragraph{BCM-Relevanz.}
Jede Parameterschätzung im BCM muss mit Unsicherheit operieren. Die BEATRIX-Pipeline verwendet daher Ensemble-Methoden und berichtet Konfidenzintervalle, nicht Punktschätzungen.
\end{axiom}

\begin{axiom}[Bounded Rationality: MA-EPI-02]
\begin{equation}
\text{Rationalität}(\sigma) \leq R_{\max}(\text{Kapazität}, \text{Zeit}, \text{Information})
\end{equation}
Rationalität ist durch kognitive Kapazität, verfügbare Zeit und Information beschränkt (Simon, 1955).

\paragraph{Philosophische Begründung.}
Herbert Simon formulierte Bounded Rationality als Antwort auf das unrealistische Modell des homo oeconomicus. Menschen \textit{satisficen} (``gut genug'') statt zu optimieren. Dies ist keine Schwäche, sondern eine adaptive Strategie angesichts begrenzter Ressourcen.

\paragraph{Empirische Evidenz.}
Miller (1956) zeigte, dass das Arbeitsgedächtnis auf $7 \pm 2$ Chunks begrenzt ist. Cowan (2001) reduzierte diese Schätzung auf 4 Chunks. Diese fundamentale Kapazitätsbeschränkung erklärt, warum Menschen Heuristiken verwenden \textit{müssen}.

\paragraph{BCM-Relevanz.}
Das BCM modelliert keine perfekt rationalen Agenten, sondern ``boundedly rational'' Entscheider. Die Komplexität von Interventionen muss der kognitiven Kapazität der Zielgruppe angepasst sein.
\end{axiom}

\begin{axiom}[Zwei-Systeme-Theorie: MA-EPI-05]
Kognition operiert in zwei Modi:
\begin{itemize}
    \item \textbf{System 1}: Schnell, automatisch, heuristisch, mühelos, emotional
    \item \textbf{System 2}: Langsam, deliberativ, analytisch, anstrengend, rational
\end{itemize}
Die meisten Alltagsentscheidungen werden von System 1 getroffen.

\paragraph{Philosophische Wurzeln.}
Die Unterscheidung zwischen automatischem und deliberativem Denken findet sich bereits bei:
\begin{itemize}
    \item Aristoteles: \textit{hexis} (Gewohnheit) vs. \textit{phronesis} (praktische Klugheit)
    \item Hume: Passion vs. Reason (``Reason is, and ought only to be the slave of the passions'')
    \item Freud: Es vs. Ich (mit anderem Fokus, aber ähnlicher Struktur)
    \item William James: Habit vs. Will
\end{itemize}

\paragraph{Empirische Operationalisierung.}
Kahneman \& Frederick (2002) operationalisieren die Unterscheidung durch:
\begin{itemize}
    \item \textbf{Cognitive Reflection Test (CRT)}: ``Ein Schläger und ein Ball kosten zusammen 1.10€. Der Schläger kostet 1€ mehr als der Ball. Was kostet der Ball?'' (System-1-Antwort: 0.10€, korrekt: 0.05€)
    \item \textbf{Response Time}: System-1-Antworten kommen schneller
    \item \textbf{Pupillendilation}: System-2-Aktivierung korreliert mit Pupillenerweiterung
\end{itemize}

\paragraph{BCM-Relevanz.}
Die Interventionstypen NUDGE und DESIGN zielen auf System 1, während BOOST und THINK System 2 aktivieren. Die Wahl des Interventionstyps muss berücksichtigen, welches System die Zielverhalten typischerweise steuert.
\end{axiom}

\subsection{Die AWX-WAX-WTX-Sequenz: Ein epistemischer Pfad}

Das BCM modelliert die Transformation von Information zu Verhalten als dreistufige Sequenz:

\begin{axiom}[AWX -- Awareness: MA-EPI-07]
\begin{equation}
\text{AWX}(\sigma, I) = \mathbb{P}(\sigma \text{ kennt Information } I)
\end{equation}
Awareness ist die Wahrscheinlichkeit, dass ein Agent $\sigma$ eine relevante Information $I$ überhaupt wahrnimmt.

\paragraph{Kognitive Mechanismen.}
Awareness wird beeinflusst durch:
\begin{itemize}
    \item \textbf{Aufmerksamkeit}: Selektive Wahrnehmung filtert >99\% der Umgebungsinformation
    \item \textbf{Salienz}: Hervorstechende Stimuli werden eher wahrgenommen
    \item \textbf{Priming}: Vorherige Aktivierung erhöht Wahrnehmungswahrscheinlichkeit
    \item \textbf{Inattentional Blindness}: Gorilla-Experiment (Simons \& Chabris, 1999)
\end{itemize}

\paragraph{Interventions-Implikation.}
Bevor komplexere Interventionen greifen können, muss \textit{überhaupt} Awareness geschaffen werden. Dies ist oft der Engpass.
\end{axiom}

\begin{axiom}[WAX -- Willingness: MA-EPI-08]
\begin{equation}
\text{WAX}(\sigma, V | I) = \mathbb{P}(\sigma \text{ ist bereit zu } V | \sigma \text{ kennt } I)
\end{equation}
Willingness ist die bedingte Wahrscheinlichkeit, dass ein Agent zur Verhaltensänderung bereit ist, gegeben er hat die relevante Information.

\paragraph{Motivationale Determinanten.}
Willingness wird bestimmt durch:
\begin{itemize}
    \item Wahrgenommene Kosten-Nutzen-Relation
    \item Selbstwirksamkeitserwartung (Bandura, 1977)
    \item Identitätskongruenz (passt die Handlung zu meinem Selbstbild?)
    \item Soziale Normen (was tun andere?)
    \item Emotionale Valenz der Information
\end{itemize}

\paragraph{Der Knowledge-Attitude Gap.}
Empirisch ist $\text{WAX} \ll \text{AWX}$ häufig: Menschen \textit{wissen}, dass Rauchen schädlich ist, sind aber nicht \textit{bereit}, aufzuhören.
\end{axiom}

\begin{axiom}[WTX -- Willingness to Transit: MA-EPI-09]
\begin{equation}
\text{WTX}(\sigma, V | \text{WAX}) = \mathbb{P}(\sigma \text{ handelt } V | \sigma \text{ ist bereit})
\end{equation}
WTX ist die Wahrscheinlichkeit, dass aus Bereitschaft tatsächlich Handlung wird.

\paragraph{Der Intention-Action Gap.}
Der Übergang von Intention zu Handlung scheitert oft an:
\begin{itemize}
    \item \textbf{Procrastination}: Aufschieben trotz Intention (Present Bias)
    \item \textbf{Friction}: Kleine Hindernisse blockieren große Intentionen
    \item \textbf{Forgetting}: Die Intention wird ``vergessen'' im Moment der Entscheidung
    \item \textbf{Self-Control Failure}: Kurzfristige Impulse überwinden langfristige Intentionen
\end{itemize}

\paragraph{Interventions-Implikation.}
Interventionen, die nur AWX oder WAX adressieren, scheitern oft am WTX-Bottleneck. Implementation Intentions (``Wenn Situation X, dann Handlung Y'') können den Intention-Action Gap schließen.
\end{axiom}

\begin{axiom}[Die vollständige Sequenz: MA-EPI-10]
\begin{equation}
\mathbb{P}(V) = \text{AWX} \times \text{WAX} \times \text{WTX}
\end{equation}
Die Wahrscheinlichkeit eines Verhaltens ist das Produkt aller drei Faktoren.

\paragraph{Diagnostische Implikation.}
Wenn ein Zielverhalten nicht auftritt, muss diagnostiziert werden, \textit{wo} die Sequenz bricht:
\begin{itemize}
    \item AWX-Problem: ``Die wissen gar nicht, dass...'' $\rightarrow$ Information, Salienz
    \item WAX-Problem: ``Die wollen einfach nicht...'' $\rightarrow$ Motivation, Framing, Identität
    \item WTX-Problem: ``Die nehmen es sich vor, aber...'' $\rightarrow$ Friction, Defaults, Commitment
\end{itemize}
\end{axiom}

\subsection{Systematische Biases (MA-EPI-12 bis MA-EPI-37)}

Das \META-Modul formalisiert 26 kognitive Biases, die für das \BCM{} und \BEATRIX{} relevant sind:

\subsubsection{Klassische Heuristiken und Biases}

\begin{axiom}[Verfügbarkeitsheuristik: MA-EPI-12]
\begin{equation}
\hat{P}(E) = f(\text{Ease of Recall}(E))
\end{equation}
Die geschätzte Wahrscheinlichkeit eines Ereignisses korreliert mit der Leichtigkeit, Beispiele abzurufen.
\end{axiom}

\begin{axiom}[Framing-Effekte: MA-EPI-13]
\begin{equation}
U(V|\text{Frame}_A) \neq U(V|\text{Frame}_B)
\end{equation}
Identische Optionen werden unterschiedlich bewertet, abhängig von ihrer Darstellung.
\end{axiom}

\begin{axiom}[Anchoring: MA-EPI-14]
\begin{equation}
\mathbb{E}[X|\text{Anchor}] = f(\text{Anchor})
\end{equation}
Schätzungen werden systematisch von einem Ankerwert beeinflusst.
\end{axiom}

\begin{axiom}[Loss Aversion: MA-EPI-15]
\begin{equation}
|U(\text{Loss})| > |U(\text{Gain})| \quad \text{für } |\text{Loss}| = |\text{Gain}|
\end{equation}
Verluste wiegen subjektiv schwerer als betragsmäßig gleiche Gewinne. Typischerweise gilt $\lambda \approx 2$ (Kahneman \& Tversky, 1979).
\end{axiom}

\subsubsection{BCM-kritische Biases (Version 2.2)}

Die folgenden Biases wurden in Version 2.2 des \META-Moduls hinzugefügt aufgrund ihrer besonderen Relevanz für \BEATRIX:

\begin{axiom}[Base Rate Neglect: MA-EPI-26]
\begin{equation}
\hat{P}(H|E) \approx P(E|H), \quad \text{ignoriert } P(H)
\end{equation}
Prior-Wahrscheinlichkeiten (Base Rates) werden zugunsten von Einzelfall-Information ignoriert.

\textbf{BCM-Relevanz}: Bei der LLM-basierten Typen-Klassifikation muss explizit auf Base Rates (z.B. 50\% Conditional Cooperators, 30\% Self-Interested, 10\% Altruisten nach Fischbacher et al., 2001) geachtet werden.
\end{axiom}

\begin{axiom}[Self-Serving Bias: MA-EPI-31]
\begin{equation}
P(\text{Internal}|\text{Success}) > P(\text{Internal}|\text{Failure})
\end{equation}
Erfolge werden internal attribuiert (``Ich bin gut''), Misserfolge external (``Die Umstände waren schuld'').

\textbf{BCM-Relevanz}: Dies ist der Kern-Mechanismus für \textbf{Motivated Inference} im IDN-Modul. Agenten verzerren ihre Wahrnehmung der Norm $\phi$, um das Selbstbild zu schützen.
\end{axiom}

\begin{axiom}[Social Desirability Bias: MA-EPI-32]
\begin{equation}
\text{Report}(V) = V_{\text{true}} + \varepsilon_{SD}; \quad \varepsilon_{SD} = f(\text{Visibility}, \text{Norm})
\end{equation}
Selbstberichte sind systematisch in Richtung sozialer Erwünschtheit verzerrt.

\textbf{BCM-Relevanz}: \textbf{Kritisch für BEATRIX}: Persona-Beschreibungen, die das LLM analysiert, enthalten systematische Verzerrungen. Das LLM muss zwischen ``stated values'' und ``revealed behavior'' unterscheiden.
\end{axiom}

\begin{axiom}[Moral Licensing: MA-EPI-34]
\begin{equation}
P(\text{BadAction}|\text{PriorGoodAction}) > P(\text{BadAction}|\text{NoPriorAction})
\end{equation}
Eine vorherige gute Tat ``lizenziert'' nachfolgende schlechte Taten.

\textbf{BCM-Relevanz}: Erklärt Identitäts-Dynamik im IDN-Modul: Selbstbild-Kredit wird aufgebaut und kann ``ausgegeben'' werden.
\end{axiom}

\begin{axiom}[Scope Insensitivity: MA-EPI-36]
\begin{equation}
\frac{\partial U}{\partial N} \rightarrow 0 \quad \text{für } N \rightarrow \infty; \quad \text{WTP}(N) \approx \text{WTP}(1)
\end{equation}
Bewertung skaliert nicht proportional mit der Anzahl Betroffener.

\textbf{BCM-Relevanz}: Kollektiver Nutzen (KNU) wird systematisch unterschätzt. Dies rechtfertigt Interventionen, die KNU salient machen.
\end{axiom}

\subsection{Evidenzhierarchie und die Labor-Feld-Debatte}

\epigraph{\textit{``We argue that many recent objections against lab experiments are misguided and that even more lab experiments should be conducted.''}}{--- Falk \& Heckman, \textit{Science} (2009)}

Eine fundamentale epistemische Frage für das \BCM{} lautet: \textbf{Woher stammen die Parameter, und wie robust sind sie?} Diese Frage berührt eine zentrale methodologische Debatte der Verhaltensökonomie.

\subsubsection{Die methodologische Kontroverse}

\paragraph{Die Kritik an Laborexperimenten.}
Levitt \& List (2007) formulierten einflussreiche Einwände gegen die externe Validität von Laborexperimenten:
\begin{itemize}
    \item \textbf{Selektion}: Laborprobanden sind typischerweise WEIRD (Western, Educated, Industrialized, Rich, Democratic) -- insbesondere Studierende
    \item \textbf{Stakes}: Laborauszahlungen sind oft trivial im Vergleich zu realen Entscheidungen
    \item \textbf{Scrutiny}: Probanden wissen, dass sie beobachtet werden (Hawthorne-Effekt)
    \item \textbf{Kontext}: Das Labor abstrahiert von den sozialen und institutionellen Faktoren realer Entscheidungen
\end{itemize}

\paragraph{Die Verteidigung des Labors.}
Falk \& Heckman (2009) verteidigten Laborexperimente in \textit{Science}:
\begin{itemize}
    \item \textbf{Kontrolle}: Nur im Labor können Schlüsselaspekte des Kontexts systematisch variiert werden
    \item \textbf{Kausalität}: Randomisierung ermöglicht kausale Inferenz, die aus Beobachtungsdaten nicht möglich ist
    \item \textbf{Mechanismus-Identifikation}: Ohne Labor wüssten wir nicht, \textit{dass} Loss Aversion existiert oder \textit{wie} sie funktioniert
    \item \textbf{Replikation}: Laborergebnisse sind replizierbar unter kontrollierten Bedingungen
\end{itemize}

Ihre Kernthese: ``The objection that experiments are artificial misses the point. Any data set can be used to get a specific estimate only to the extent that the variation in the data can be controlled for extraneous factors.''

\paragraph{Die moderne Synthese: Triangulation.}
Die zeitgenössische Position -- reflektiert im Nobelpreis 2021 für Card, Angrist und Imbens -- ist:

\begin{center}
\textbf{Labor + Feld-RCT + Natürliche Experimente + Strukturelle Modelle}\\[0.3em]
$\downarrow$\\[0.3em]
\textbf{TRIANGULATION}\\[0.3em]
$\downarrow$\\[0.3em]
\textbf{Robustes Kausalwissen}
\end{center}

Konvergente Befunde über verschiedene Methoden hinweg erhöhen das Vertrauen in die Robustheit eines Effekts dramatisch.

\subsubsection{Implikationen für BCM-Parameter}

Die BCM-Base-Rates stammen primär aus Laborexperimenten. Dies hat Implikationen:

\begin{table}[h]
\centering
\caption{Validierungsstatus zentraler BCM-Parameter}
\begin{tabular}{@{}llll@{}}
\toprule
\textbf{Parameter} & \textbf{Labor-Evidenz} & \textbf{Feld-Evidenz} & \textbf{Status} \\
\midrule
$\lambda \approx 2.25$ (Loss Aversion) & Kahneman \& Tversky (1979) & Variiert nach Kontext & Mechanismus robust \\
$\beta \approx 0.7$ (Present Bias) & Laibson et al. (1997) & Kontextabhängig & Mechanismus robust \\
50\% Conditional Cooperators & Fischbacher et al. (2001) & Kulturelle Variation & Robust in WEIRD \\
Nudge-Effekte +15\% & Diverse Meta-Analysen & DellaVigna \& Linos: +2\% & \textbf{Lab-Field Gap!} \\
\bottomrule
\end{tabular}
\end{table}

\paragraph{Der Lab-Field Gap.}
DellaVigna \& Linos (2022) dokumentierten systematisch, dass Nudge-Effekte in Feld-Implementierungen oft nur 10--20\% der Labor-Effektstärken erreichen. Mögliche Gründe:
\begin{itemize}
    \item \textbf{Compliance}: Im Labor folgen 100\% den Instruktionen; im Feld deutlich weniger
    \item \textbf{Konkurrenz}: Im Feld konkurrieren Nudges mit vielen anderen Einflüssen
    \item \textbf{Heterogenität}: Labor-Effekte sind Durchschnitte; Feld-Populationen sind heterogener
    \item \textbf{Implementation}: Perfekte Labor-Implementierung vs. reale Umsetzungsprobleme
\end{itemize}

\begin{axiom}[Evidenzhierarchie für BCM-Parameter: MA-EPI-38]
BCM-Parameter sind nach Evidenzqualität zu klassifizieren:

\begin{enumerate}
    \item \textbf{Höchste Evidenz} ($\star\star\star\star$): Konvergenz von Labor-RCT + Feld-RCT + Natürlichem Experiment
    \begin{itemize}
        \item Beispiel: Loss Aversion als Mechanismus
        \item Interpretation: Robuster Parameter, direkt verwendbar
    \end{itemize}

    \item \textbf{Hohe Evidenz} ($\star\star\star$): Replizierte Labor-RCTs mit theoretischer Fundierung
    \begin{itemize}
        \item Beispiel: 50\% Conditional Cooperators (Fischbacher et al.)
        \item Interpretation: Mechanismus validiert, Effektstärke kontextabhängig
    \end{itemize}

    \item \textbf{Mittlere Evidenz} ($\star\star$): Einzelne Labor-RCTs oder nicht-replizierte Feldstudien
    \begin{itemize}
        \item Beispiel: Spezifische Nudge-Effektstärken
        \item Interpretation: Vorläufiger Parameter, Validierung nötig
    \end{itemize}

    \item \textbf{Niedrige Evidenz} ($\star$): Beobachtungsstudien oder theoretische Ableitung
    \begin{itemize}
        \item Beispiel: FEPSDE-Gewichtungen
        \item Interpretation: Arbeitshypothese, empirische Validierung erforderlich
    \end{itemize}
\end{enumerate}

\paragraph{BEATRIX-Implikation.}
Das System sollte bei Parameterschätzungen die Evidenzstufe kommunizieren und bei niedrigerer Evidenz breitere Konfidenzintervalle verwenden.
\end{axiom}

\subsubsection{Die Rolle von Armin Falk}

Die Entwicklung von Armin Falk selbst illustriert die methodologische Evolution:

\begin{table}[h]
\centering
\caption{Armin Falks methodologische Entwicklung}
\begin{tabular}{@{}llll@{}}
\toprule
\textbf{Jahr} & \textbf{Studie} & \textbf{Methode} & \textbf{Beitrag} \\
\midrule
2001 & Fehr, Fischbacher, Falk & Labor-Experiment & Punishment in Public Goods \\
2009 & Falk \& Heckman & Meta-Artikel & Verteidigung des Labors \\
2018 & Global Preference Survey & Feld-Survey (76 Länder) & Kulturelle Variation dokumentiert \\
2023 & Diverse Policy-Evaluationen & Natürliche Experimente & Reale Policy-Effekte \\
\bottomrule
\end{tabular}
\end{table}

Die Einsicht: Das Labor ist \textbf{notwendig aber nicht hinreichend}. Es identifiziert Mechanismen; Feld-Validierung quantifiziert reale Effektstärken.

\paragraph{Implikation für das BCM.}
Das \BCM{} verwendet Labor-Base-Rates als \textit{Default}-Werte, sollte aber:
\begin{enumerate}
    \item Transparenz über die Evidenzquelle herstellen
    \item Kontextuelle Anpassung ermöglichen (z.B. kulturelle Kalibrierung)
    \item Konfidenzintervalle der Evidenzqualität anpassen
    \item Feld-Validierung aktiv fördern und integrieren
\end{enumerate}

\subsection{Vollständige Liste der EPI-Axiome}

\begin{longtable}{@{}lll@{}}
\caption{Alle 38 Epistemologie-Axiome} \\
\toprule
\textbf{ID} & \textbf{Statement} & \textbf{Referenz} \\
\midrule
\endfirsthead
\toprule
\textbf{ID} & \textbf{Statement} & \textbf{Referenz} \\
\midrule
\endhead
\midrule
\multicolumn{3}{r}{\textit{Fortsetzung auf nächster Seite}} \\
\endfoot
\bottomrule
\endlastfoot
MA-EPI-01 & Unvollständiges Wissen & -- \\
MA-EPI-02 & Bounded Rationality & Simon (1955) \\
MA-EPI-03 & Irreduzible Unsicherheit & -- \\
MA-EPI-04 & Adaptive Heuristiken & Gigerenzer \\
MA-EPI-05 & Zwei-Systeme-Theorie & Kahneman (2011) \\
MA-EPI-06 & Systematische Biases & Kahneman \& Tversky \\
MA-EPI-07 & AWX: Awareness & -- \\
MA-EPI-08 & WAX: Willingness & -- \\
MA-EPI-09 & WTX: Transition & -- \\
MA-EPI-10 & Sequenz AWX → WAX → WTX → V & -- \\
MA-EPI-11 & Intention-Action Gap & -- \\
MA-EPI-12 & Verfügbarkeitsheuristik & Tversky \& Kahneman \\
MA-EPI-13 & Framing-Effekte & Tversky \& Kahneman (1981) \\
MA-EPI-14 & Anchoring & -- \\
MA-EPI-15 & Loss Aversion & Kahneman \& Tversky (1979) \\
MA-EPI-16 & Present Bias & -- \\
MA-EPI-17 & Status Quo Bias & -- \\
MA-EPI-18 & Overconfidence & -- \\
MA-EPI-19 & Lernen ist möglich & -- \\
MA-EPI-20 & Feedback verbessert Schätzungen & -- \\
MA-EPI-21 & Gewohnheiten reduzieren Kognition & -- \\
MA-EPI-22 & Debiasing ist schwierig & -- \\
MA-EPI-23 & Exponentielle Sättigung & -- \\
MA-EPI-24 & Belief-Update folgt Bayes & -- \\
MA-EPI-25 & Confirmation Bias & -- \\
\midrule
\multicolumn{3}{l}{\textit{Neu in Version 2.2:}} \\
\midrule
MA-EPI-26 & Base Rate Neglect & Kahneman \& Tversky (1973) \\
MA-EPI-27 & Representativeness Heuristic & Kahneman \& Tversky (1972) \\
MA-EPI-28 & Sunk Cost Fallacy & Arkes \& Blumer (1985) \\
MA-EPI-29 & Endowment Effect & Thaler (1980) \\
MA-EPI-30 & Optimism Bias & Weinstein (1980) \\
MA-EPI-31 & Self-Serving Bias & Miller \& Ross (1975) \\
MA-EPI-32 & Social Desirability Bias & Crowne \& Marlowe (1960) \\
MA-EPI-33 & Omission Bias & Spranca et al. (1991) \\
MA-EPI-34 & Moral Licensing & Merritt et al. (2010) \\
MA-EPI-35 & Identifiable Victim Effect & Small \& Loewenstein (2003) \\
MA-EPI-36 & Scope Insensitivity & Kahneman (1986) \\
MA-EPI-37 & Affect Heuristic & Slovic et al. (2002) \\
MA-EPI-38 & Evidenzhierarchie für BCM-Parameter & Falk \& Heckman (2009) \\
\end{longtable}

%==============================================================================


% Kapitel 4: Governance (GOV)
\section{Kategorie 3: Governance (GOV)}
%==============================================================================

\epigraph{\textit{``The only freedom which deserves the name is that of pursuing our own good in our own way.''}}{--- John Stuart Mill, \textit{On Liberty} (1859)}

Die Governance-Kategorie (25 Axiome: MA-GOV-01 bis MA-GOV-25) definiert, \textbf{was erlaubt ist} -- sowohl im Sinne von Verhaltens-Constraints als auch im Sinne ethischer Prinzipien für Interventionen.

\subsection{Philosophischer Hintergrund: Von Mill zu Thaler}

\epigraph{\textit{``Libertarian paternalism is not an oxymoron.''}}{--- Thaler \& Sunstein (2003)}

Die Frage, unter welchen Bedingungen Verhaltensinterventionen ethisch legitim sind, berührt fundamentale Fragen der politischen Philosophie. Das Governance-Modul des \META{} navigiert zwischen zwei Extrempositionen: dem radikalen Libertarismus (``keine Intervention ist legitim'') und dem harten Paternalismus (``der Staat weiß es besser'').

\paragraph{Mills Harm Principle.}
John Stuart Mills \textit{On Liberty} (1859) etablierte das ``Harm Principle'' als liberales Grundprinzip: Die einzige Rechtfertigung für Einschränkung individueller Freiheit ist die Verhinderung von Schaden an anderen. Selbstschädigendes Verhalten darf nicht reguliert werden. Diese Position ist die historische Grundlage für die Ablehnung von Paternalismus.

\paragraph{Die paternalistische Herausforderung.}
Gegen Mill wurde argumentiert, dass Menschen systematisch gegen ihre eigenen Interessen handeln -- nicht aus freier Wahl, sondern aus kognitiven Limitationen. Gerald Dworkin (1972) unterschied:
\begin{itemize}
    \item \textbf{Harter Paternalismus}: Eingriff gegen den expliziten Willen der Person (z.B. Gurtpflicht)
    \item \textbf{Weicher Paternalismus}: Eingriff nur bei nicht-autonomen Entscheidungen (z.B. bei Manipulation, Zwang, oder fehlender Information)
\end{itemize}

\paragraph{Libertärer Paternalismus.}
Thaler \& Sunstein (2003, 2008) versöhnten die Positionen mit dem Konzept des ``Libertarian Paternalism'':
\begin{itemize}
    \item \textbf{Paternalistisch}: Interventionen zielen darauf, Menschen zu helfen, bessere Entscheidungen zu treffen (nach ihren eigenen Maßstäben)
    \item \textbf{Libertär}: Die Wahlfreiheit bleibt erhalten; es werden keine Optionen entfernt
\end{itemize}

Das Schlüsselargument: Es gibt \textit{keine neutrale} Choice Architecture. Jede Präsentation von Optionen beeinflusst die Wahl. Wenn Beeinflussung unvermeidbar ist, sollte sie wenigstens in eine ``gute'' Richtung gehen.

\paragraph{Die Kritik an Nudging.}
Bovens (2009) und andere haben Nudging kritisiert:
\begin{itemize}
    \item \textbf{Manipulation}: Nudges umgehen bewusste Deliberation
    \item \textbf{Transparenz}: Viele Nudges funktionieren nur, wenn sie nicht bekannt sind
    \item \textbf{Wer definiert ``gut''?}: Das Konzept setzt voraus, dass jemand weiß, was für andere gut ist
\end{itemize}

\paragraph{Die BCM-Position.}
Das \BCM{} ist \textbf{deskriptiv, nicht normativ}. Es modelliert Verhalten, es schreibt keine Ziele vor. Die Governance-Axiome definieren \textit{Constraints} für ethisches Interventionsdesign, aber die \textit{Ziele} werden extern vorgegeben. Diese Trennung von ``Ist'' und ``Soll'' folgt dem Humeschen Guillotine-Prinzip.

\subsection{Fundamentale Governance-Prinzipien}

\begin{axiom}[Autonomie-Primat: MA-GOV-03]
Individuelle Autonomie hat grundsätzlich Vorrang. Einschränkungen bedürfen der Rechtfertigung.

\paragraph{Philosophische Begründung.}
Autonomie -- die Fähigkeit zur Selbstgesetzgebung -- ist ein Kernwert der Aufklärung. Kant formulierte: ``Handle nur nach derjenigen Maxime, durch die du zugleich wollen kannst, dass sie ein allgemeines Gesetz werde.'' Autonome Entscheidungen verdienen Respekt, auch wenn sie ``schlecht'' erscheinen.

\paragraph{Operationalisierung.}
Im BCM bedeutet dies: Jede Intervention, die Autonomie einschränkt, muss \textit{explizit} gerechtfertigt werden. Die Beweislast liegt beim Intervenierenden.
\end{axiom}

\begin{axiom}[Nicht-Normativität: MA-GOV-04]
Das \BCM{} \textbf{beschreibt} Verhalten, es \textbf{schreibt nicht vor}. Die Wahl der Ziele liegt beim Auftraggeber, nicht beim Modell.

\paragraph{Philosophische Begründung.}
Dies folgt Humes Guillotine: Aus deskriptiven Aussagen (``so ist es'') folgen keine normativen (``so soll es sein''). Das BCM kann analysieren, \textit{wie} ein Ziel erreicht werden kann, aber nicht, \textit{ob} das Ziel gut ist.

\paragraph{Praktische Implikation.}
\BEATRIX{} kann sowohl für pro-soziale als auch für kommerzielle Ziele eingesetzt werden. Die ethische Verantwortung liegt beim Auftraggeber, nicht beim Modell.
\end{axiom}

\begin{axiom}[Transparenz-Gebot: MA-GOV-05]
Interventionen sollen transparent sein. Verdeckte Manipulation ist zu vermeiden.

\paragraph{Philosophische Begründung.}
Kants Publizitätsprinzip: ``Alle auf das Recht anderer Menschen bezogenen Handlungen, deren Maxime sich nicht mit der Publizität verträgt, sind unrecht.'' Eine Intervention, die nur funktioniert, wenn sie geheim bleibt, ist ethisch problematisch.

\paragraph{Empirisches Gegenargument.}
Einige Nudges funktionieren auch nach Offenlegung (Loewenstein et al., 2015). Aber nicht alle -- und hier entsteht ein ethisches Dilemma.
\end{axiom}

\begin{axiom}[Subsidiarität: MA-GOV-06]
Die \textbf{mildeste wirksame Intervention} ist zu wählen (Principle of Minimal Intervention).

\paragraph{Philosophische Begründung.}
Dies ist die Anwendung des Proportionalitätsprinzips auf Verhaltensinterventionen. Eine stärkere Intervention (z.B. Verbot) ist nur gerechtfertigt, wenn mildere Alternativen (z.B. Information) versagen.

\paragraph{Praktische Anwendung.}
Die 8 Interventionstypen des BCM sind nach ``Milde'' geordnet. Interventionsdesign sollte diese Hierarchie respektieren.
\end{axiom}

\subsection{Interventions-Taxonomie}

\begin{axiom}[8 Interventionstypen: MA-GOV-08]
Das \BCM{} unterscheidet acht fundamentale Interventionstypen:
\begin{enumerate}
    \item \textbf{NUDGE}: Änderung der Choice Architecture
    \item \textbf{BOOST}: Stärkung kognitiver Kompetenzen
    \item \textbf{THINK}: Aktivierung von System 2
    \item \textbf{INCENTIVE}: Materielle Anreize
    \item \textbf{REGULATE}: Ge- und Verbote
    \item \textbf{DESIGN}: Änderung des physischen Umfelds
    \item \textbf{SOCIAL}: Soziale Normen und Peer Effects
    \item \textbf{IDENTITY}: Aktivierung oder Änderung von Identitäten
\end{enumerate}
\end{axiom}

\begin{axiom}[Freiheits-Hierarchie: MA-GOV-09]
Die 8 Interventionstypen sind nach Freiheitsgrad geordnet:
\begin{equation}
\text{NUDGE} > \text{BOOST} > \text{THINK} > \text{INCENTIVE} > \text{SOCIAL} > \text{IDENTITY} > \text{DESIGN} > \text{REGULATE}
\end{equation}
wobei $>$ bedeutet ``mehr Wahlfreiheit erhaltend als''.
\end{axiom}

\begin{axiom}[Eskalations-Prinzip: MA-GOV-10]
Interventionen sollen von mild nach stark eskalieren:
\begin{equation}
\text{NUDGE} \rightarrow \text{BOOST} \rightarrow \text{INCENTIVE} \rightarrow \text{REGULATE}
\end{equation}
Erst wenn mildere Interventionen versagen, sind stärkere gerechtfertigt.
\end{axiom}

\subsection{Ethische Constraints}

\begin{axiom}[Distributive Gerechtigkeit: MA-GOV-12]
Kosten und Nutzen einer Intervention sind fair zu verteilen. Interventionen dürfen keine systematische Benachteiligung erzeugen.
\end{axiom}

\begin{axiom}[Reversibilität: MA-GOV-13]
Reversible Interventionen sind irreversiblen vorzuziehen.
\end{axiom}

\begin{axiom}[Crowding-Out vermeiden: MA-GOV-15]
\begin{equation}
\frac{\partial \text{Intrinsic}}{\partial \text{Extrinsic}} < 0 \quad \text{ist möglich}
\end{equation}
Extrinsische Anreize können intrinsische Motivation zerstören. Dies ist bei Interventionsdesign zu berücksichtigen.
\end{axiom}

%==============================================================================


% Kapitel 5: Dynamik (DYN)
\section{Kategorie 4: Dynamik (DYN)}
%==============================================================================

\epigraph{\textit{``Panta rhei'' -- Alles fließt.}}{--- Heraklit von Ephesus (ca. 500 v. Chr.)}

Die Dynamik-Kategorie (25 Axiome: MA-DYN-01 bis MA-DYN-25) definiert, \textbf{wie sich Dinge verändern} -- zeitliche Prozesse, Gleichgewichte, Pfadabhängigkeiten und Evolution.

\subsection{Philosophischer Hintergrund: Von Heraklit zur Komplexitätstheorie}

\epigraph{\textit{``Time present and time past\\
Are both perhaps present in time future,\\
And time future contained in time past.''}}{--- T.S. Eliot, \textit{Four Quartets}}

Die Frage nach der Natur von Veränderung und Zeit gehört zu den ältesten philosophischen Problemen. Das Dynamik-Modul des \META{} integriert Erkenntnisse aus Philosophie, Physik und Komplexitätswissenschaft.

\paragraph{Die heraklitische Tradition.}
Heraklit von Ephesus formulierte die radikale Position, dass \textit{alles} im Fluss ist: ``Man kann nicht zweimal in denselben Fluss steigen.'' Verhalten ist kein statischer Zustand, sondern ein kontinuierlicher Prozess. Diese Einsicht steht gegen die Tendenz der neoklassischen Ökonomie, Gleichgewichtszustände als ``normal'' zu betrachten.

\paragraph{Der parmenidische Einwand.}
Parmenides argumentierte dagegen, dass wahre Veränderung unmöglich sei -- was ist, kann nicht aufhören zu sein. Diese Position mag absurd erscheinen, enthält aber einen Kern: Viele ``Veränderungen'' sind nur scheinbar, sie sind Rekonfigurationen stabiler Elemente. Im BCM: Typen ($\tau$) ändern sich selten, aber ihre Manifestation in Verhalten variiert situativ.

\paragraph{Aristoteles' Potentialität.}
Aristoteles versöhnte die Positionen mit dem Konzept von \textit{dynamis} (Potentialität) und \textit{energeia} (Aktualität). Veränderung ist die Aktualisierung von Potential. Im BCM: Eine Person \textit{hat} das Potential zur Verhaltensänderung (AWX, WAX), aber die \textit{Aktualisierung} (WTX) erfordert den richtigen Kontext.

\paragraph{Newtons Mechanik und die Idee der Vorhersagbarkeit.}
Die newtonsche Physik suggerierte ein deterministisches Universum: Gegeben Anfangsbedingungen, ist die Zukunft vollständig bestimmt. Diese Idee beeinflusste die klassische Ökonomie mit ihren Gleichgewichtsmodellen. Laplace' Dämon könnte -- mit perfektem Wissen -- alle Zukunft vorhersagen.

\paragraph{Poincaré, Lorenz und das Chaos.}
Henri Poincaré (1890er) und Edward Lorenz (1963) zeigten, dass selbst deterministische Systeme \textit{praktisch} unvorhersagbar sein können -- der berühmte ``Schmetterlingseffekt''. Kleine Unterschiede in Anfangsbedingungen führen zu radikal unterschiedlichen Entwicklungen. Dies erklärt, warum Verhaltensvorhersagen fundamental unsicher sind.

\paragraph{Komplexitätstheorie und Emergenz.}
Die Santa Fe School (Holland, Arthur, Kauffman) etablierte, dass komplexe Systeme Eigenschaften zeigen, die nicht aus ihren Teilen vorhersagbar sind. Soziale Normen \textit{emergieren} aus individuellen Interaktionen. Tipping Points entstehen, wenn lokale Änderungen globale Kaskaden auslösen.

\paragraph{Zeitliche Inkonsistenz und der gespaltene Wille.}
Die philosophische Frage nach der personalen Identität über Zeit (Locke, Parfit) verbindet sich mit dem ökonomischen Problem der zeitlichen Inkonsistenz. Bin ``Ich heute'' derselbe wie ``Ich morgen''? Wenn nicht -- warum sollte Ich-heute für Ich-morgen sparen? Hyperbolic Discounting und Present Bias sind nicht ``irrational'', sondern reflektieren die genuine Pluralität des Selbst über Zeit.

\subsection{Temporale Grundstruktur}

\begin{axiom}[Diskrete Events: MA-DYN-02]
Entscheidungen sind diskrete Events in kontinuierlicher Zeit:
\begin{equation}
V(t) = \sum_i V_i \cdot \delta(t - t_i)
\end{equation}

\paragraph{Philosophische Bedeutung.}
Dies formalisiert die Intuition, dass Verhalten aus diskreten \textit{Handlungen} besteht, nicht aus kontinuierlichen Zuständen. Eine Entscheidung ist ein ``Sprung'' -- mathematisch eine Dirac-Delta-Funktion.

\paragraph{Praktische Implikation.}
Interventionen müssen auf diese Entscheidungspunkte zielen. Der ``Moment der Wahrheit'' ist der Zeitpunkt $t_i$, an dem die Entscheidung fällt.
\end{axiom}

\begin{axiom}[Pfadabhängigkeit: MA-DYN-03]
\begin{equation}
U(V, t) = f(V, \{V_\tau : \tau < t\})
\end{equation}
Der Nutzen einer Handlung hängt von der Geschichte ab, nicht nur vom aktuellen Zustand.

\paragraph{Philosophische Begründung.}
Dies ist die formale Ablehnung des ``markovianischen'' Menschenbilds. Menschen sind nicht gedächtnislos. Vergangene Entscheidungen beeinflussen zukünftige -- durch Gewohnheit, Identität, Sunk-Cost-Effekte und Pfadabhängigkeit.

\paragraph{Empirische Evidenz.}
Arthur (1989) zeigte, dass technologische ``Lock-ins'' (QWERTY-Tastatur, VHS vs. Betamax) durch Pfadabhängigkeit entstehen. Dasselbe gilt für Verhaltens-Lock-ins.
\end{axiom}

\subsection{Gleichgewichte und Stabilität}

\begin{axiom}[Multiple Equilibria: MA-DYN-04]
Es kann mehrere stabile Zustände geben. Das System kann in verschiedenen Gleichgewichten ``gefangen'' sein.

\paragraph{Philosophische Bedeutung.}
Dies widerspricht der Idee eines \textit{einzigen} optimalen Zustands. Gesellschaften können in ``guten'' oder ``schlechten'' Gleichgewichten stecken (z.B. hohe vs. niedrige Kooperation). Der aktuelle Zustand ist nicht notwendig der ``beste''.

\paragraph{Interventions-Implikation.}
Es genügt nicht, ein ``besseres'' Gleichgewicht zu kennen. Man muss den \textit{Übergang} orchestrieren -- und dieser kann eine kritische Masse erfordern.
\end{axiom}

\begin{axiom}[Tipping Points: MA-DYN-06]
Kritische Schwellen existieren, an denen kleine Änderungen große Zustandswechsel auslösen können.

\paragraph{Mathematische Grundlage.}
Tipping Points sind Bifurkationspunkte in dynamischen Systemen. Am kritischen Punkt divergiert die ``Sensitivität'' gegen unendlich -- minimale Änderungen haben maximale Wirkung.

\paragraph{Empirische Beispiele.}
\begin{itemize}
    \item \textbf{Soziale Normen}: Rauchen in Restaurants war in den 1980ern normal, kippte in den 2000ern
    \item \textbf{Technologieadoption}: Elektroautos brauchten Jahre der Stagnation, dann exponentielles Wachstum
    \item \textbf{Sprache}: Neue Begriffe (``Gender'', ``Climate Change'') kippen von Randphänomenen zum Mainstream
\end{itemize}
\end{axiom}

\subsection{Gewohnheitsformation}

\begin{axiom}[Habit Loop: MA-DYN-08]
Gewohnheiten folgen dem Zyklus:
\begin{equation}
\text{Cue} \rightarrow \text{Routine} \rightarrow \text{Reward}
\end{equation}
Wiederholung stärkt die Assoziation zwischen Cue und Routine.

\paragraph{Neurowissenschaftliche Grundlage.}
Der Habit Loop ist im Basalganglien verankert (Yin \& Knowlton, 2006). Bei Gewohnheitsformation verschiebt sich die Aktivität von präfrontalem Kortex (deliberativ) zu Striatum (automatisch). Gewohnheiten werden ``ausgelagert'' -- sie erfordern keine bewusste Aufmerksamkeit mehr.

\paragraph{Interventions-Implikation.}
Gewohnheiten zu ändern erfordert, den \textit{Cue} zu identifizieren und entweder zu entfernen (DESIGN) oder mit einer neuen Routine zu verknüpfen (``habit substitution'').
\end{axiom}

\subsection{Soziale Dynamik}

\begin{axiom}[Soziale Diffusion: MA-DYN-10]
\begin{equation}
\frac{\partial P(V)}{\partial t} = \beta \cdot P(V) \cdot (1 - P(V)) \cdot f(\text{Network})
\end{equation}
Verhalten breitet sich sozial aus, analog zu epidemiologischen Modellen.

\paragraph{Mathematische Struktur.}
Dies ist eine logistische Wachstumsgleichung. Der Term $P(V) \cdot (1-P(V))$ bedeutet: Diffusion ist maximal, wenn 50\% adoptiert haben. Am Anfang ist zu wenig ``Ansteckungspotential'', am Ende zu wenig ``Empfängliche''.

\paragraph{Netzwerk-Effekte.}
Die Funktion $f(\text{Network})$ modelliert, dass Diffusion netzwerkabhängig ist. In dichten, homogenen Netzwerken ist $\beta$ hoch; in fragmentierten Netzwerken kann Diffusion stagnieren.
\end{axiom}

\begin{axiom}[Kritische Masse: MA-DYN-11]
Ab einer Schwelle $P^* \approx 25\%$ (Centola, 2018) wird Verhalten selbstverstärkend.

\paragraph{Empirische Evidenz.}
Centola (2018) zeigte in kontrollierten Experimenten, dass soziale Konventionen kippen, wenn ca. 25\% eine Alternative praktizieren. Dies ist robust über verschiedene Kontexte.

\paragraph{Praktische Implikation.}
Um eine Norm zu ändern, muss nicht die gesamte Population überzeugt werden -- nur eine kritische Minderheit. Dies macht Verhaltensänderung ``effizienter'' als es erscheint.
\end{axiom}

\begin{axiom}[Norm-Kaskaden: MA-DYN-12]
Normen können rapid kippen, wenn die kritische Masse erreicht wird.

\paragraph{Historische Beispiele.}
\begin{itemize}
    \item \textbf{Fall der Berliner Mauer} (1989): Jahrzehntelange Stabilität, dann Kollaps innerhalb von Wochen
    \item \textbf{MeToo-Bewegung} (2017): Jahrzehntelange Toleranz für Belästigung, dann rapides Kippen
    \item \textbf{Gay Marriage}: Von Mehrheitsablehnung zu Mehrheitsakzeptanz innerhalb einer Generation
\end{itemize}

\paragraph{Theoretische Erklärung.}
Kuran (1989) sprach von ``preference falsification'': Menschen verbergen ihre wahren Präferenzen, wenn sie glauben, in der Minderheit zu sein. Kaskaden entstehen, wenn ``pluralistische Ignoranz'' aufgedeckt wird.
\end{axiom}

\subsection{Temporale Biases}

\begin{axiom}[Hyperbolisches Discounting: MA-DYN-14]
\begin{equation}
D(t) = \frac{1}{1 + k \cdot t} \quad \text{statt} \quad D(t) = e^{-\delta t}
\end{equation}
Die tatsächliche Diskontierung ist hyperbolisch, nicht exponentiell. Dies führt zu zeitinkonsistenten Präferenzen.

\paragraph{Philosophische Implikation.}
Hyperbolisches Discounting bedeutet, dass Präferenzen zeitinkonsistent sind: Heute präferiere ich 100€ heute über 110€ morgen, aber gleichzeitig 110€ in 31 Tagen über 100€ in 30 Tagen. Dies ist ``irrational'' im Sinne der Erwartungsnutzentheorie, aber universell.

\paragraph{Evolutionäre Erklärung.}
In der evolutionären Umgebung war die Zukunft radikal unsicher. Immediate Belohnung hatte Überlebenswert. Die ``Irrationalität'' ist ein evolutionäres Erbe.

\paragraph{Interventions-Implikation.}
Commitment Devices (Odysseus am Mast) helfen, die ``irrationale'' zukünftige Person zu binden. Pre-Commitment ist eine rationale Strategie für ein irrationales Selbst.
\end{axiom}

\begin{axiom}[Hedonic Treadmill: MA-DYN-20]
\begin{equation}
\lim_{t \rightarrow \infty} U(t|\text{Change}) = U_{\text{baseline}}
\end{equation}
Glück kehrt nach positiven oder negativen Events zur Baseline zurück (Kahneman, 1999).

\paragraph{Philosophische Bedeutung.}
Dies ist eine empirische Bestätigung der stoischen und buddhistischen Intuition, dass externe Güter kein dauerhaftes Glück bringen. Hedonic Adaptation bedeutet: Lottogewinner sind langfristig nicht glücklicher, Querschnittsgelähmte nicht unglücklicher.

\paragraph{Interventions-Implikation.}
Interventionen, die auf dauerhaftes ``Glück'' durch Konsum zielen, sind zum Scheitern verurteilt. Nachhaltige Wohlbefindenssteigerung erfordert andere Strategien (Flow, Meaning, Relationships).
\end{axiom}

%==============================================================================


% Kapitel 6: Integration mit operativen Modulen
\section{Integration mit operativen Modulen}
%==============================================================================

\subsection{Validierungs-Beziehungen}

Jedes \META-Axiom spezifiziert im Feld \texttt{validates}, welche operativen Module es fundiert:

\begin{table}[h]
\centering
\caption{Validierungs-Matrix: META → Operative Module}
\begin{tabular}{@{}lcccc@{}}
\toprule
\textbf{META-Kategorie} & \textbf{INU} & \textbf{KNU} & \textbf{IDN} & \textbf{LAMBDA} \\
\midrule
ONT (26 Axiome) & \checkmark & \checkmark & \checkmark & \checkmark \\
EPI (38 Axiome) & \checkmark & \checkmark & \checkmark & -- \\
GOV (25 Axiome) & -- & -- & -- & -- \\
DYN (25 Axiome) & \checkmark & \checkmark & \checkmark & \checkmark \\
\bottomrule
\end{tabular}
\end{table}

\subsection{Die BCM 2.2 Kernformel}

Die Integration aller Module kulminiert in der BCM 2.2 Kernformel für den effektiven Nutzen:

\begin{equation}
U_{\text{eff}} = (1-\lambda) \cdot U_{\text{INU}} + \lambda \cdot U_{\text{KNU}} + U_{\text{IDN}}(a, \sigma, \phi, \tilde{\mu}; F, s, m)
\label{eq:bcm_core}
\end{equation}

wobei der Identitätsnutzen $U_{\text{IDN}}$ nach dem BCM 2.2 Vollmodell berechnet wird:

\begin{equation}
U_{\text{IDN}} = U_{\text{mat}} - \mu_{\text{ID}}\left[L_{\text{aff}} + \omega_1 \kappa (a - \phi)^2 + \omega_3 \psi \cdot \mathbb{E}[(\phi - \eta)^2]\right] - C_{\text{MI}}(\tilde{\mu}) - C_{\text{info}}(\sigma)
\end{equation}

\subsection{META-Axiome als Constraints}

Die \META-Axiome fungieren als \textbf{Constraints} für die operativen Module:

\begin{itemize}
    \item \textbf{ONT-Constraints}: Definieren, welche Entitäten (Subjekte, Gruppen, Identitäten) modelliert werden können
    \item \textbf{EPI-Constraints}: Definieren, welche kognitiven Verzerrungen berücksichtigt werden müssen
    \item \textbf{GOV-Constraints}: Definieren, welche Interventionen ethisch zulässig sind
    \item \textbf{DYN-Constraints}: Definieren, welche zeitlichen Prozesse modelliert werden müssen
\end{itemize}

%==============================================================================


% Kapitel 7: META im BEATRIX-System
\section{META im BEATRIX-System}
%==============================================================================

\subsection{Rolle bei der LLM-basierten Estimation}

Das \BEATRIX-System verwendet ein LLM zur Parameterschätzung. Das \META-Modul informiert diesen Prozess auf mehreren Ebenen:

\begin{enumerate}
    \item \textbf{Base Rate Awareness (MA-EPI-26)}: Das LLM muss empirische Base Rates für Typen-Verteilungen berücksichtigen
    \item \textbf{Social Desirability Korrektur (MA-EPI-32)}: Persona-Beschreibungen müssen auf systematische Verzerrungen geprüft werden
    \item \textbf{Anchoring-Prävention (MA-EPI-14)}: Few-Shot Anchoring im Prompt verhindert willkürliche Schätzungen
\end{enumerate}

\subsection{Constrained Generation Pipeline}

Die \BEATRIX{} Estimation Pipeline (ESTIMATION-Modul) implementiert die epistemologischen Constraints:

\begin{enumerate}
    \item \textbf{PreFlight}: Prüfung auf plausibles Szenario (verhindert Halluzination)
    \item \textbf{Few-Shot Anchoring}: Kalibrierungsskalen im Prompt
    \item \textbf{Pydantic Schema}: Strukturierte Outputs
    \item \textbf{Ensemble Voting}: N parallele Runs, Median-Aggregation
    \item \textbf{IRR-Check}: Inter-Rater Reliability > 0.85
\end{enumerate}

\subsection{Ambiguity Detection}

Die Unterscheidung zwischen \textbf{Unsicherheit} (unimodale Verteilung) und \textbf{Wertekonflikt} (bimodale Verteilung) basiert auf:

\begin{equation}
BC = \frac{s^2 + 1}{k + 3 \cdot \frac{(n-1)^2}{(n-2)(n-3)}}
\end{equation}

wobei $s$ = Schiefe, $k$ = Kurtosis. Bei $BC > 0.555$ liegt Polarisierung vor.

%==============================================================================


% Kapitel 8: Von Handwerk zu Engineering - Die BEATRIX-Transformation
\section{Von Handwerk zu Engineering: Die BEATRIX-Transformation}
%==============================================================================

\epigraph{\textit{``Software engineering is what happens to programming when you add time and other programmers.''}}{--- Russ Cox}

Die bisherigen Abschnitte haben die axiomatischen Grundlagen des BCM und die technischen Komponenten von BEATRIX beschrieben. Dieser Abschnitt adressiert eine fundamentale Frage: \textbf{Warum war eine systematische, skalierbare Verhaltensparameter-Schätzung vor LLMs nicht möglich?}

Die Antwort liegt in der Struktur des klassischen BCM-Prozesses, der als \textit{fragmentiertes Handwerk} charakterisiert werden kann -- im Gegensatz zum \textit{integrierten Engineering-Prozess}, den BEATRIX ermöglicht.

\subsection{Der Pre-LLM BCM-Prozess: Fragmentiertes Handwerk}

\subsubsection{Die fragmentierte Wertschöpfungskette}

Der klassische Prozess der Verhaltensparameter-Schätzung besteht aus vier Schritten mit systematischen Bruchstellen:

\begin{figure}[h]
\centering
\begin{tikzpicture}[
    node distance=0.8cm,
    step/.style={rectangle, draw, rounded corners, minimum width=2.2cm, minimum height=1.2cm, align=center, font=\small},
    break/.style={draw=red!70, dashed, thick},
    arrow/.style={->, thick, >=stealth}
]

% Steps
\node[step, fill=blue!10] (s1) at (0, 0) {Kontext\\verstehen};
\node[step, fill=blue!10] (s2) at (3.5, 0) {Verhaltens-\\annahmen\\bilden};
\node[step, fill=blue!10] (s3) at (7, 0) {Parameter-\\schätzung};
\node[step, fill=blue!10] (s4) at (10.5, 0) {BCM-\\Modellierung};

% Arrows with breaks
\draw[arrow] (s1) -- node[above, font=\scriptsize, red!70] {Bruch 1} (s2);
\draw[arrow] (s2) -- node[above, font=\scriptsize, red!70] {Bruch 2} (s3);
\draw[arrow] (s3) -- node[above, font=\scriptsize, red!70] {Bruch 3} (s4);

% Outputs below
\node[below=0.5cm of s1, font=\scriptsize, text=gray] {Dokument};
\node[below=0.5cm of s2, font=\scriptsize, text=gray] {Kopf des\\Beraters};
\node[below=0.5cm of s3, font=\scriptsize, text=gray] {Notiz\\(vielleicht)};
\node[below=0.5cm of s4, font=\scriptsize, text=gray] {Excel/\\Python};

% Missing feedback
\node[below=1.8cm of s2, font=\scriptsize, red!70] {Bruch 4: Kein Feedback-Loop};
\draw[break, <->] (0, -1.8) -- (10.5, -1.8);

\end{tikzpicture}
\caption{Die fragmentierte Wertschöpfungskette des Pre-LLM BCM-Prozesses mit vier systematischen Bruchstellen.}
\label{fig:fragmented-chain}
\end{figure}

\paragraph{Die Handwerk-Probleme im Überblick.}
\begin{itemize}
    \item \textbf{Nicht reproduzierbar}: Jeder Berater macht es anders
    \item \textbf{Nicht skalierbar}: Linear mit Beraterkapazität
    \item \textbf{Qualitätsvariation}: Abhängig von Tagesform, Erfahrung, Zeitdruck
    \item \textbf{Wissenserosion}: Wissen bleibt in Köpfen, nicht im System
    \item \textbf{Keine Lernschleife}: Fehler werden nicht systematisch erfasst
\end{itemize}

\subsubsection{Bruch 1: Kontext → Verhaltensannahmen}

\paragraph{Das Problem: Informationsverlust.}
Ein typischer 2-Stunden-Workshop mit Stakeholdern produziert 50+ verhaltensrelevante Aussagen:

\begin{quote}
\textit{``Unsere Kunden sind sehr preisbewusst''} \\
\textit{``Die älteren Kunden sind loyaler''} \\
\textit{``Im B2B-Segment entscheiden Einkäufer rational''} \\
\textit{``Qualität ist wichtiger als Preis''} (Widerspruch zu Aussage 1!) \\
... und 46 weitere Aussagen ...
\end{quote}

\paragraph{Was der Berater mitnimmt:}
\begin{itemize}
    \item Die 3--4 Aussagen, die er sich gemerkt hat
    \item Gefiltert durch seine eigenen kognitiven Biases
    \item Ohne Gewichtung nach Evidenzstärke
    \item Ohne Dokumentation der Widersprüche
\end{itemize}

\textbf{Geschätzter Informationsverlust}: $\approx 80\%$

\subsubsection{Bruch 2: Verhaltensannahmen → Parameter}

\paragraph{Das Problem: Willkürliche Übersetzung.}

Der Berater denkt: \textit{``Die Kunden sind preisbewusst.''}

Was bedeutet das für BCM-Parameter?

\begin{table}[h]
\centering
\caption{Willkürliche Parameter-Übersetzung im klassischen Prozess}
\small
\begin{tabular}{@{}lll@{}}
\toprule
\textbf{Parameter} & \textbf{Berater-Überlegung} & \textbf{Basis} \\
\midrule
$\lambda$ (Loss Aversion) & ``Höher als normal? Vielleicht 2.5?'' & Intuition \\
Reference Point & ``Was ist der Referenzpreis? Keine Ahnung.'' & Ignoriert \\
$w_F$ (Financial) & ``Financial wichtig... 0.4?'' & Aus der Luft \\
Elastizität & ``Keine Daten, ich nehme $-1.5$'' & ``Branchenstandard'' \\
\bottomrule
\end{tabular}
\end{table}

\textbf{Kernprobleme}:
\begin{itemize}
    \item Keine systematische Übersetzungslogik
    \item Keine Dokumentation der Reasoning-Kette
    \item Keine Base Rates als empirischer Anker
\end{itemize}

\subsubsection{Bruch 3: Parameter → Modell}

\paragraph{Das Problem: Händische Übertragung.}

Der Berater hat (irgendwo) notiert: \textit{``$\lambda \approx 2.5$, preissensitiv, ältere loyaler''}

Der Analyst muss ins Excel/Python-Modell übertragen:

\begin{verbatim}
lambda = 2.5   # TODO: check with PM
beta = 0.7     # default
alpha = ???    # was ist das?
\end{verbatim}

\textbf{Kernprobleme}:
\begin{itemize}
    \item Übertragungsfehler (Tippfehler, Verwechslungen)
    \item Unvollständige Spezifikation (fehlende Parameter)
    \item Keine Validierung (inkonsistente Werte möglich)
    \item Kontextverlust (``warum 2.5?'' nicht dokumentiert)
\end{itemize}

\subsubsection{Bruch 4: Kein Feedback-Loop}

\paragraph{Das Problem: Kein organisationales Lernen.}

\begin{quote}
\textbf{Modell-Output} (Januar): ``Bei 5\% Preiserhöhung $\rightarrow$ 3\% Churn'' \\
\textbf{Tatsächlicher Outcome} (Juli): 7\% Churn
\end{quote}

Was passiert typischerweise?

\begin{table}[h]
\centering
\caption{Typische Reaktionen auf Prognose-Fehler}
\begin{tabular}{@{}ll@{}}
\toprule
\textbf{Option} & \textbf{Konsequenz} \\
\midrule
A: Niemand vergleicht & ``Das Projekt ist abgeschlossen'' \\
B: Jemand merkt es & ``Das lag an externen Faktoren'' (keine Analyse) \\
C: Erkenntnis geht verloren & Berater nicht mehr im Projekt, Excel weg \\
D: Nächstes Projekt & ``Wir machen wieder einen Workshop...'' \\
\bottomrule
\end{tabular}
\end{table}

\textbf{Ergebnis}: \textbf{Kein organisationales Lernen}. Jedes Projekt startet bei Null.

\subsection{Die Analogie: Software 1980er vs. 2020er}

Die Transformation von BCM-Handwerk zu BEATRIX-Engineering entspricht der Evolution der Softwareentwicklung:

\begin{table}[h]
\centering
\caption{Analogie: Softwareentwicklung und BCM-Prozess}
\begin{tabular}{@{}p{5.5cm}p{5.5cm}@{}}
\toprule
\textbf{Software 1980er (Handwerk)} & \textbf{Software 2020er (Engineering)} \\
\midrule
Code auf Papier & Git Repository \\
Mündliche Übergabe & CI/CD Pipeline \\
``Funktioniert auf meinem Rechner'' & Docker Container \\
Keine Tests & Unit Tests, Integration Tests \\
Keine Dokumentation & Auto-generierte Docs \\
Wissen in Köpfen & Knowledge Base, Wiki \\
\bottomrule
\end{tabular}
\end{table}

\textbf{BCM heute = Software 1980er}:
\begin{itemize}
    \item Parameter in Köpfen / Notizen
    \item Mündliche Übergabe zwischen Berater und Analyst
    \item ``Funktioniert für diesen Kunden''
    \item Keine systematische Dokumentation
    \item Wissen geht mit Mitarbeitern
\end{itemize}

\subsection{Die Blockade: Warum war End-to-End nicht möglich?}

\subsubsection{Das Schwarze Loch: Natürliche Sprache → Parameter}

Der kritische Engpass lag in der Übersetzung von unstrukturiertem Text zu strukturierten Parametern:

\begin{figure}[h]
\centering
\begin{tikzpicture}[
    node distance=1.5cm
]

% Input
\node[rectangle, draw, rounded corners, fill=blue!10, minimum width=10cm, minimum height=1.5cm, align=center] (input) at (0, 2) {
    \textbf{INPUT} (natürliche Sprache):\\
    ``Anna ist 45, Lehrerin, umweltbewusst, achtet auf Preise, hat zwei Kinder...''
};

% Black hole
\node[rectangle, draw=red!70, dashed, thick, rounded corners, minimum width=4cm, minimum height=1.5cm, align=center, fill=gray!10] (hole) at (0, 0) {
    \textbf{SCHWARZES LOCH}\\
    \small Wie kommt man von\\
    TEXT zu ZAHLEN?
};

% Output
\node[rectangle, draw, rounded corners, fill=green!10, minimum width=10cm, minimum height=1.5cm, align=center] (output) at (0, -2) {
    \textbf{OUTPUT} (Parameter):\\
    $\lambda = ???$, $\beta = ???$, $w_F = ???$, $w_{\ddot{O}} = ???$, ...
};

% Arrows
\draw[->, thick] (input) -- (hole);
\draw[->, thick] (hole) -- (output);

\end{tikzpicture}
\caption{Das ``Schwarze Loch'' der Parameter-Schätzung: Der Übergang von natürlicher Sprache zu quantifizierten Parametern war der kritische Engpass.}
\label{fig:black-hole}
\end{figure}

\subsubsection{Pre-LLM Optionen -- und warum sie scheiterten}

\paragraph{Option 1: Mensch (Berater).}
\begin{itemize}
    \item[$-$] Langsam: 30--60 Minuten pro Persona
    \item[$-$] Teuer: Beraterhonorar (\EUR{50--200} pro Schätzung)
    \item[$-$] Inkonsistent: Verschiedene Berater $\rightarrow$ verschiedene Ergebnisse
    \item[$-$] Nicht dokumentiert: Reasoning im Kopf
    \item[$-$] Nicht skalierbar
\end{itemize}

\paragraph{Option 2: Regelbasiertes System.}
\begin{verbatim}
WENN Alter > 60 DANN lambda += 0.1
WENN "preisbewusst" IN text DANN w_F += 0.1
...
\end{verbatim}

\begin{itemize}
    \item[$-$] Kombinatorische Explosion: Tausende von Regeln nötig
    \item[$-$] Kontext wird ignoriert: ``preisbewusst'' bedeutet in verschiedenen Kontexten Verschiedenes
    \item[$-$] Neue Fälle nicht abgedeckt
    \item[$-$] Wurde versucht, hat nie skaliert
\end{itemize}

\paragraph{Option 3: Machine Learning (klassisch).}
Theoretisch möglich, aber:
\begin{itemize}
    \item[$-$] Bräuchte Labeled Training Data: Tausende Personas mit ``echten'' Parametern
    \item[$-$] \textbf{Diese Daten existieren nicht!}
    \item[$-$] Kein Transfer auf neue Domänen
    \item[$-$] Praktisch nicht umsetzbar
\end{itemize}

\textbf{Ergebnis}: Dieser Schritt blieb \textbf{manuelles Handwerk}.

\subsection{LLM als fehlende Brücke}

Large Language Models lösen das ``Schwarze Loch''-Problem durch eine einzigartige Kombination von Fähigkeiten:

\begin{figure}[h]
\centering
\begin{tikzpicture}[
    node distance=1.2cm
]

% Input
\node[rectangle, draw, rounded corners, fill=blue!10, minimum width=10cm, minimum height=1.2cm, align=center] (input) at (0, 3) {
    \textbf{INPUT}: ``Anna ist 45, Lehrerin, umweltbewusst...''
};

% LLM box
\node[rectangle, draw=green!70!black, thick, rounded corners, fill=green!10, minimum width=10cm, minimum height=3cm, align=left] (llm) at (0, 0) {
    \textbf{LLM (mit BCM-Kontext)} kann:\\[0.3em]
    $\bullet$ Natürliche Sprache verstehen\\
    $\bullet$ Implizite Informationen inferieren\\
    $\bullet$ BCM-Ontologie anwenden\\
    $\bullet$ Base Rates als Anker nutzen\\
    $\bullet$ Reasoning dokumentieren\\
    $\bullet$ Unsicherheit quantifizieren
};

% Output
\node[rectangle, draw, rounded corners, fill=yellow!20, minimum width=10cm, minimum height=2cm, align=left, font=\small\ttfamily] (output) at (0, -3) {
    \{\\
    \quad "lambda": 2.2,\\
    \quad "lambda\_ci": [1.8, 2.6],\\
    \quad "reasoning": "Preisbewusst → erhöhte Loss Aversion...",\\
    \quad "confidence": "medium"\\
    \}
};

% Arrows
\draw[->, thick] (input) -- (llm);
\draw[->, thick] (llm) -- (output);

\end{tikzpicture}
\caption{LLM als Brücke zwischen natürlicher Sprache und strukturierten BCM-Parametern.}
\label{fig:llm-bridge}
\end{figure}

\paragraph{Die kritischen Eigenschaften.}
\begin{enumerate}
    \item \textbf{Schnell}: Sekunden statt Stunden
    \item \textbf{Konsistent}: Gleicher Input $\rightarrow$ gleiches Output (bei fixierter Temperatur)
    \item \textbf{Dokumentiert}: Reasoning-Chain ist Teil des Outputs
    \item \textbf{Skalierbar}: 1000 Personas parallel möglich
    \item \textbf{Integrierbar}: API-basiert, kein Medienbruch
\end{enumerate}

\subsection{Die BEATRIX End-to-End Pipeline}

BEATRIX implementiert einen durchgängigen 6-Phasen-Prozess:

\subsubsection{Phase 1: Input Layer}

\begin{itemize}
    \item Akzeptiert verschiedene Formate: Workshop-Protokolle (.docx), CRM-Daten (.csv), Interview-Transkripte (.txt)
    \item BEATRIX Parser strukturiert unstrukturierten Input
    \item Keine manuelle Vorverarbeitung nötig
\end{itemize}

\subsubsection{Phase 2: Extraction Layer}

LLM extrahiert aus unstrukturiertem Input:
\begin{itemize}
    \item Demografische Merkmale
    \item Verhaltensrelevante Aussagen
    \item Kontextinformationen
    \item Widersprüche und Unsicherheiten
    \item Fehlende Informationen (Gaps)
\end{itemize}

\textbf{Output}: Strukturiertes Persona-Objekt

\subsubsection{Phase 3: Estimation Layer}

Die 4-Stufen-Architektur (vgl. Abschnitt~\ref{sec:4-stufen}) wird angewendet:

\begin{table}[h]
\centering
\caption{Parameter-Schätzung am Beispiel $\lambda$}
\begin{tabular}{@{}llr@{}}
\toprule
\textbf{Stufe} & \textbf{Quelle} & \textbf{$\lambda$-Beitrag} \\
\midrule
1: Base Rate & Gächter et al. & 2.00 \\
2: Demografie & Alter 45 & +0.10 \\
3a: Kontext (empirisch) & Low Stakes & $-0.10$ \\
3b: Kontext (LLM) & Umweltentscheidung & +0.15 \\
4: Persona (LLM) & ``preisbewusst'' & +0.05 \\
\midrule
\textbf{Final} & & \textbf{2.20} [1.7, 2.7] \\
\bottomrule
\end{tabular}
\\[0.3em]
\footnotesize Anteil empirisch: 81\%, Anteil LLM: 19\%
\end{table}

\textit{Wiederholen für alle Parameter: $\beta$, $\alpha$, $\beta_{FS}$, $w_F$, $w_E$, $w_P$, $w_S$, $w_D$, $w_{\ddot{O}}$, ...}

\subsubsection{Phase 4: Modeling Layer}

Parameter fließen \textbf{automatisch} in BCM-Modell:
\begin{itemize}
    \item Keine manuelle Übertragung
    \item Validierung gegen Constraints (z.B. $\beta_{FS} \leq \alpha$)
    \item Sensitivitätsanalyse automatisch
    \item Monte-Carlo für Unsicherheitspropagation
\end{itemize}

\subsubsection{Phase 5: Output \& Validation Layer}

Drei Output-Typen:
\begin{enumerate}
    \item \textbf{Prognose-Report}: ``Bei 5\% Preiserhöhung $\rightarrow$ 4\% Churn [2\%, 6\%]''
    \item \textbf{Audit-Trail}: Vollständige Dokumentation aller Schritte
    \item \textbf{Interventions-Empfehlung}: ``Framing auf Verlustvermeidung optimieren''
\end{enumerate}

\subsubsection{Phase 6: Feedback Layer}

Nach Implementierung:
\begin{itemize}
    \item Tatsächlicher Outcome wird erfasst
    \item Vergleich mit Prognose
    \item Optional: Bayesian Update der Base Rates
    \item Organisationales Lernen wird ermöglicht
\end{itemize}

\begin{quote}
\textit{``Prognose war 4\% Churn, tatsächlich 6\%''} \\
$\rightarrow$ \textit{``$\lambda$ war unterschätzt, Update: $\lambda_{prior} \mathrel{+}= 0.15$''}
\end{quote}

\subsection{Der Kontrast: Handwerk vs. Engineering}

\begin{table}[h]
\centering
\caption{Vergleich: Pre-LLM Handwerk vs. BEATRIX End-to-End}
\small
\begin{tabular}{@{}lcc@{}}
\toprule
\textbf{Dimension} & \textbf{Pre-LLM (Handwerk)} & \textbf{BEATRIX (E2E)} \\
\midrule
Medienbrüche & 3--4 pro Projekt & 0 \\
Dokumentation & Fragmentiert, oft fehlend & Automatisch, vollständig \\
Reproduzierbarkeit & Niedrig (personenabhängig) & Hoch (deterministisch) \\
Zeit pro Persona & 30--60 Minuten & 10--30 Sekunden \\
Skalierbarkeit & Linear mit Kapazität & Quasi-unbegrenzt \\
Konsistenz & Variiert & Garantiert \\
Wissenstransfer & Mündlich, informell & Systemisch, persistent \\
Lernfähigkeit & Ad-hoc, unsystematisch & Strukturiert, messbar \\
Audit-Fähigkeit & Schwierig bis unmöglich & Built-in \\
Kosten pro Schätzung & \EUR{50--200} & \EUR{0.01--0.10} \\
\bottomrule
\end{tabular}
\end{table}

\subsection{Das Axiom: MA-INF-08 End-to-End Integration}

Die Transformation von Handwerk zu Engineering wird im folgenden Axiom formalisiert:

\begin{axiom}[End-to-End Integration: MA-INF-08]
\begin{enumerate}
    \item BEATRIX implementiert einen \textbf{durchgängigen Prozess} von unstrukturiertem Input (Text, Daten) zu quantifizierten Verhaltensparametern \textbf{ohne manuelle Zwischenschritte}.

    \item Jeder Prozessschritt ist \textbf{automatisch dokumentiert} und nachvollziehbar (Audit Trail).

    \item Die Pipeline garantiert \textbf{Konsistenz}: Gleicher Input führt zu gleichem Output (bei fixiertem System-State).

    \item \textbf{Validierung} ist integriert:
    \begin{itemize}
        \item Constraint-Checks (z.B. $\beta_{FS} \leq \alpha$)
        \item Plausibilitätsprüfungen gegen Base Rates
        \item Automatische Sensitivitätsanalysen
    \end{itemize}

    \item \textbf{Feedback} ist strukturiert: Outcomes werden erfasst und zur Kalibrierung genutzt.

    \item Das System \textbf{ersetzt fragmentiertes Handwerk} durch einen \textbf{ingenieurmäßigen Prozess}.
\end{enumerate}
\end{axiom}

\begin{tcolorbox}[colback=green!5!white, colframe=green!75!black, title=Die Transformation in einem Satz]
BEATRIX transformiert BCM-Parameter-Schätzung von \textbf{fragmentiertem Handwerk} (personenabhängig, undokumentiert, nicht skalierbar) zu \textbf{integriertem Engineering} (systematisch, nachvollziehbar, skalierbar) -- ermöglicht durch LLMs als Brücke zwischen natürlicher Sprache und strukturierten Parametern.
\end{tcolorbox}

\subsection{Die Drei-Phasen-Architektur: Technologische Realisierung}

Die in Abschnitt 9.5 beschriebene \textbf{6-Phasen-Pipeline} (Input $\rightarrow$ Extraction $\rightarrow$ Estimation $\rightarrow$ Modeling $\rightarrow$ Output $\rightarrow$ Feedback) beschreibt die \textit{funktionale} Perspektive -- was passiert in jedem Schritt.
Komplementär dazu beschreibt die \textbf{Drei-Phasen-Architektur} die \textit{technologische} Perspektive -- welche Technologien werden wie orchestriert.

\begin{tcolorbox}[colback=blue!5!white, colframe=blue!75!black, title=Zwei komplementäre Perspektiven]
\textbf{6 Phasen (funktional):} Was passiert? -- Input, Extraction, Estimation, Modeling, Output, Feedback

\textbf{3 Phasen (technologisch):} Womit? -- LLM (Parsing) $\rightarrow$ Code Interpreter (Berechnung) $\rightarrow$ LLM (Interpretation)
\end{tcolorbox}

\subsubsection{Das Pre-LLM Problem revisited}

Das BCM konnte immer \textbf{berechnet} werden -- gegeben die richtigen Parameter, liefert die Mathematik deterministische Ergebnisse.
Das Problem lag an zwei Stellen:

\begin{enumerate}
    \item \textbf{Input-Problem}: Wie kommt man von einer Persona-Beschreibung (``Anna ist 45, Lehrerin, umweltbewusst...'') zu den Modellparametern ($\lambda = 2.2$, $\alpha = 0.88$, ...)?

    \item \textbf{Output-Problem}: Wie erklärt man einem Nicht-Experten, was ``$U_{\text{eff}} = -0.34$ mit Sensitivität $\partial U / \partial \lambda = 0.6$'' bedeutet?
\end{enumerate}

Vor LLMs erforderten beide Schritte \textbf{menschliche Experten} -- Verhaltensökonomen, die Parameter schätzen und Ergebnisse interpretieren konnten.
Dies machte BCM-Analysen langsam, teuer und nicht skalierbar.

\subsubsection{Die BEATRIX-Lösung: Drei Phasen}

BEATRIX löst dieses Problem durch eine orchestrierte Drei-Phasen-Architektur:

\begin{figure}[h]
\centering
\begin{tikzpicture}[
    node distance=1.5cm,
    phase/.style={rectangle, draw, rounded corners, minimum width=3.5cm, minimum height=2cm, align=center},
    arrow/.style={->, thick, >=stealth}
]

% Phases
\node[phase, fill=blue!15] (p1) at (0, 0) {\textbf{Phase 1}\\LLM\\[0.2em]\footnotesize Input-Parsing};
\node[phase, fill=green!15] (p2) at (5, 0) {\textbf{Phase 2}\\Code Interpreter\\[0.2em]\footnotesize BCM-Berechnung};
\node[phase, fill=orange!15] (p3) at (10, 0) {\textbf{Phase 3}\\LLM\\[0.2em]\footnotesize Interpretation};

% Input/Output
\node[above=0.8cm of p1, align=center] (input) {\footnotesize ``Anna ist 45,\\Lehrerin, ...''};
\node[above=0.8cm of p2, align=center] (params) {\footnotesize $\lambda=2.2$, $\alpha=0.88$\\$w_F=0.35$, ...};
\node[above=0.8cm of p3, align=center] (results) {\footnotesize Churn: 4.2\%\\CI: [2.8\%, 5.9\%]};
\node[below=0.8cm of p3, align=center] (output) {\footnotesize ``Bei 5\% Preiserhöhung\\erwarten wir...''};

% Arrows
\draw[arrow] (input) -- (p1);
\draw[arrow] (p1) -- node[above, font=\footnotesize] {JSON} (p2);
\draw[arrow] (p2) -- node[above, font=\footnotesize] {Dict} (p3);
\draw[arrow] (p3) -- (output);

% Labels below
\node[below=0.5cm of p1, font=\footnotesize\itshape, text=gray] {Probabilistisch};
\node[below=0.5cm of p2, font=\footnotesize\itshape, text=gray] {Deterministisch};
\node[below=0.5cm of p3, font=\footnotesize\itshape, text=gray] {Probabilistisch};

\end{tikzpicture}
\caption{Die Drei-Phasen-Architektur von BEATRIX. LLMs übernehmen die ``Übersetzung'' zwischen natürlicher Sprache und Modellparametern; die eigentliche Berechnung erfolgt deterministisch.}
\label{fig:three-phases}
\end{figure}

\begin{tcolorbox}[colback=blue!5!white, colframe=blue!75!black, title=Kernprinzip]
Das LLM ist \textbf{kein Taschenrechner} -- es führt keine BCM-Berechnungen durch.\\
Das LLM ist ein \textbf{Übersetzer}:
\begin{itemize}
    \item Phase 1: Mensch $\rightarrow$ Modell (natürliche Sprache $\rightarrow$ Parameter)
    \item Phase 3: Modell $\rightarrow$ Mensch (Ergebnisse $\rightarrow$ Narrative)
\end{itemize}
Die \textbf{Mathematik bleibt deterministisch} im Code Interpreter (Phase 2).
\end{tcolorbox}

\subsubsection{Historische Entwicklung: Seit wann ist das möglich?}

Die Drei-Phasen-Architektur wurde erst durch eine Konvergenz mehrerer technologischer Entwicklungen möglich:

\begin{table}[h]
\centering
\caption{Meilensteine der LLM-Entwicklung für BEATRIX}
\small
\begin{tabular}{@{}llp{7cm}@{}}
\toprule
\textbf{Datum} & \textbf{Entwicklung} & \textbf{Bedeutung für BEATRIX} \\
\midrule
Jun 2017 & Transformer (Vaswani et al.) & Grundarchitektur für alle modernen LLMs \\
Jun 2020 & GPT-3 (175B Parameter) & In-Context Learning entdeckt \\
\textbf{Mär 2022} & \textbf{InstructGPT / ChatGPT} & \textbf{Instruction Following wird zuverlässig} \\
\textbf{Jun 2023} & \textbf{Function Calling (OpenAI)} & \textbf{Strukturierte JSON-Outputs werden zuverlässig} \\
\textbf{Jul 2023} & \textbf{Code Interpreter (ChatGPT)} & \textbf{Python-Ausführung in Sandbox} \\
\textbf{Aug 2024} & \textbf{Structured Outputs (OpenAI)} & \textbf{JSON Schema Enforcement -- garantierte Struktur} \\
\bottomrule
\end{tabular}
\end{table}

\begin{tcolorbox}[colback=yellow!5!white, colframe=yellow!75!black, title=Zeitliche Einordnung]
\textbf{Frühester Zeitpunkt} (Proof of Concept): \textbf{Juni 2023}\\
Mit Function Calling und Code Interpreter waren alle drei Phasen grundsätzlich umsetzbar.

\textbf{Production-Ready Zeitpunkt}: \textbf{August 2024}\\
Mit Structured Outputs werden Parameter-Extraktionen \textit{garantiert} strukturiert.

\textbf{Die BEATRIX Drei-Phasen-Architektur ist erst seit $\approx$18 Monaten praktisch umsetzbar.}
\end{tcolorbox}

\subsubsection{Wissenschaftliche Grundlagen der drei Phasen}

Die Drei-Phasen-Architektur stützt sich auf Erkenntnisse aus zehn Forschungsbereichen:

\paragraph{Phase 1: Natural Language $\rightarrow$ Parameter}

\begin{enumerate}
    \item \textbf{Information Extraction (IE)}: Named Entity Recognition, Relation Extraction, Slot Filling -- BEATRIX nutzt LLMs für erweiterte Slot-Filling (BCM-Parameter als ``Slots'')

    \item \textbf{Semantic Parsing}: Transformation natürlicher Sprache in formale Repräsentationen (``Anna ist preisbewusst'' $\rightarrow$ \texttt{w\_financial: 0.4})

    \item \textbf{In-Context Learning}: Brown et al. (2020) zeigten, dass LLMs neue Aufgaben ohne Finetuning erlernen können durch Few-Shot-Beispiele im Prompt

    \item \textbf{Instruction Tuning}: RLHF transformiert Base-Modelle in Instruction-Following-Assistenten (InstructGPT, Constitutional AI)
\end{enumerate}

\paragraph{Phase 2: Parameter $\rightarrow$ Berechnung}

\begin{enumerate}
    \item \textbf{Behavioral Economics Models}: Prospect Theory, Reference-Dependent Utility, Social Preferences, Identity Utility

    \item \textbf{Uncertainty Quantification}: Monte Carlo Simulation, Sensitivity Analysis, Confidence Intervals

    \item \textbf{Tool-Augmented LLMs}: Toolformer (Schick et al., 2023), PAL (Gao et al., 2023) -- LLMs erkennen, \textit{wann} Tools nötig sind
\end{enumerate}

\paragraph{Phase 3: Ergebnis $\rightarrow$ Narrative}

\begin{enumerate}
    \item \textbf{Natural Language Generation (NLG)}: Data-to-Text Generation transformiert strukturierte Daten in natürliche Sprache

    \item \textbf{Explainable AI (XAI)}: Contrastive Explanations (``Warum 4\% und nicht 2\%?''), Counterfactual Explanations, Feature Attribution

    \item \textbf{Grounded Generation}: Outputs müssen auf tatsächlichen BCM-Ergebnissen basieren, nicht halluziniert sein
\end{enumerate}

\subsubsection{Technische Implementierungsvarianten}

BEATRIX unterstützt drei Implementierungsvarianten mit steigender Komplexität:

\begin{table}[h]
\centering
\caption{Vergleich der drei Implementierungsvarianten}
\small
\begin{tabular}{@{}llll@{}}
\toprule
\textbf{Aspekt} & \textbf{A: Monolithisch} & \textbf{B: API-Orchestriert} & \textbf{C: Agentic} \\
\midrule
Architektur & Alles in einer Session & Explizite Phasentrennung & Autonomer Agent \\
BCM-Code & Jedes Mal generiert & Statisch validiert & Statisch mit Iteration \\
Kontrolle & Niedrig & Hoch & Mittel \\
Reproduzierbarkeit & Niedrig & Hoch & Mittel \\
Flexibilität & Hoch & Mittel & Sehr hoch \\
Kosten/Anfrage & Niedrig & Mittel & Hoch \\
\midrule
\textbf{Eignung} & Demos, Exploration & \textbf{Production} & Komplexe Szenarien \\
\bottomrule
\end{tabular}
\end{table}

\paragraph{Variante B: API-Orchestriert (Production)}

Die empfohlene Production-Variante trennt die drei Phasen explizit:

\begin{lstlisting}[caption={BEATRIX Orchestrator: Integration der drei Phasen}]
def beatrix_analyze(persona: str, scenario: Dict) -> Dict:
    """Hauptfunktion: Orchestriert alle drei Phasen"""

    # Phase 1: Extract (LLM)
    params = extract_parameters(persona)  # -> BCMParameters

    # Validierung zwischen Phasen
    assert 1.0 <= params.lambda_loss_aversion <= 3.0

    # Phase 2: Calculate (Deterministisch)
    results = calculate_bcm(params, scenario)  # -> Dict

    # Phase 3: Interpret (LLM)
    interpretation = interpret_results(params, results, persona)

    return {
        "parameters": params.model_dump(),
        "results": results,
        "interpretation": interpretation,
        "audit_trail": {
            "phase1_model": "claude-sonnet-4-20250514",
            "phase2_deterministic": True,
            "timestamp": datetime.now().isoformat()
        }
    }
\end{lstlisting}

\begin{tcolorbox}[colback=green!5!white, colframe=green!75!black, title=Schlussfolgerung]
Die Drei-Phasen-Architektur von BEATRIX transformiert das BCM von einem Experten-Tool zu einem zugänglichen System:

\textbf{Pre-BEATRIX}: Experte $\rightarrow$ Parameter $\rightarrow$ Excel $\rightarrow$ Experte $\rightarrow$ Bericht

\textbf{Mit BEATRIX}: User $\rightarrow$ LLM $\rightarrow$ BCM $\rightarrow$ LLM $\rightarrow$ Verständlicher Output

Der Schlüssel ist die \textbf{kluge Arbeitsteilung}: LLMs für Sprache, Code für Mathematik.
\end{tcolorbox}

%==============================================================================


% Kapitel 9: Zusammenfassung und Ausblick
\section{Zusammenfassung und Ausblick}
%==============================================================================

\subsection{Kernbeiträge des META-Moduls}

Das \META-Modul leistet folgende Beiträge zum \BCM:

\begin{enumerate}
    \item \textbf{Axiomatische Fundierung}: 115 explizite Annahmen machen das Modell transparent und kritisierbar
    \item \textbf{Interdisziplinäre Verankerung}: Die vier Kategorien (ONT, EPI, GOV, DYN) verbinden Verhaltensökonomie mit klassischen philosophischen Disziplinen
    \item \textbf{Bias-Katalog}: 37 formalisierte kognitive Biases ermöglichen systematische Berücksichtigung menschlicher Irrationalität
    \item \textbf{Ethische Leitplanken}: Die GOV-Kategorie definiert klare Prinzipien für Interventionsdesign
    \item \textbf{LLM-Guidance}: Die EPI-Axiome informieren die BEATRIX-Parameterschätzung
\end{enumerate}

\subsection{Versionierung}

\begin{table}[h]
\centering
\caption{META-Modul Versionshistorie}
\begin{tabular}{@{}llll@{}}
\toprule
\textbf{Version} & \textbf{Datum} & \textbf{Axiome} & \textbf{Änderungen} \\
\midrule
1.0 & 2023 & 100 & Initiale Version \\
2.0 & 2024 & 100 & BCM 2.2 Integration \\
2.1 & Dez 2024 & 101 & MA-ONT-26: 5-Ebenen-Hierarchie \\
2.2 & Dez 2024 & 114 & +12 Biases (EPI-26 bis EPI-37), +Evidenzhierarchie (EPI-38) \\
2.3 & Dez 2024 & 115 & +End-to-End Integration (INF-08) \\
\bottomrule
\end{tabular}
\end{table}

\subsection{Zukünftige Erweiterungen}

Geplante Erweiterungen für \META{} 3.0:
\begin{itemize}
    \item Integration von Affekt-Dynamik-Axiomen
    \item Formalisierung von Sozialkapital-Konzepten
    \item Erweiterung der Governance-Kategorie um KI-spezifische Ethik
    \item Cross-kulturelle Validierung der Bias-Universalität
\end{itemize}

%==============================================================================
% References
%==============================================================================

\section*{Literaturverzeichnis}

\subsection*{Grundlagenwerke Verhaltensökonomie}

\begin{itemize}[leftmargin=*]
    \item Kahneman, D. (2011). \textit{Thinking, Fast and Slow}. Farrar, Straus and Giroux.
    \item Kahneman, D., \& Tversky, A. (1979). Prospect theory: An analysis of decision under risk. \textit{Econometrica}, 47(2), 263-291.
    \item Thaler, R. H. (2015). \textit{Misbehaving: The Making of Behavioral Economics}. W. W. Norton.
    \item Thaler, R. H., \& Sunstein, C. R. (2008). \textit{Nudge: Improving Decisions About Health, Wealth, and Happiness}. Yale University Press.
    \item Ariely, D. (2008). \textit{Predictably Irrational: The Hidden Forces That Shape Our Decisions}. HarperCollins.
    \item Camerer, C. F., Loewenstein, G., \& Rabin, M. (Eds.). (2004). \textit{Advances in Behavioral Economics}. Princeton University Press.
    \item Gilovich, T., Griffin, D., \& Kahneman, D. (Eds.). (2002). \textit{Heuristics and Biases: The Psychology of Intuitive Judgment}. Cambridge University Press.
\end{itemize}

\subsection*{Bounded Rationality \& Heuristiken (MA-EPI-01 bis MA-EPI-06)}

\begin{itemize}[leftmargin=*]
    \item Simon, H. A. (1955). A behavioral model of rational choice. \textit{Quarterly Journal of Economics}, 69(1), 99-118.
    \item Simon, H. A. (1957). \textit{Models of Man: Social and Rational}. Wiley.
    \item Gigerenzer, G. (2007). \textit{Gut Feelings: The Intelligence of the Unconscious}. Viking.
    \item Gigerenzer, G., \& Selten, R. (Eds.). (2001). \textit{Bounded Rationality: The Adaptive Toolbox}. MIT Press.
    \item Gigerenzer, G., Todd, P. M., \& ABC Research Group. (1999). \textit{Simple Heuristics That Make Us Smart}. Oxford University Press.
    \item Tversky, A., \& Kahneman, D. (1974). Judgment under uncertainty: Heuristics and biases. \textit{Science}, 185(4157), 1124-1131.
    \item Stanovich, K. E., \& West, R. F. (2000). Individual differences in reasoning: Implications for the rationality debate. \textit{Behavioral and Brain Sciences}, 23(5), 645-665.
    \item Evans, J. S. B. T. (2008). Dual-processing accounts of reasoning, judgment, and social cognition. \textit{Annual Review of Psychology}, 59, 255-278.
\end{itemize}

\subsection*{Kognitive Biases -- Klassische Heuristiken (MA-EPI-12 bis MA-EPI-18)}

\begin{itemize}[leftmargin=*]
    \item Tversky, A., \& Kahneman, D. (1973). Availability: A heuristic for judging frequency and probability. \textit{Cognitive Psychology}, 5(2), 207-232.
    \item Tversky, A., \& Kahneman, D. (1981). The framing of decisions and the psychology of choice. \textit{Science}, 211(4481), 453-458.
    \item Tversky, A., \& Kahneman, D. (1992). Advances in prospect theory: Cumulative representation of uncertainty. \textit{Journal of Risk and Uncertainty}, 5(4), 297-323.
    \item Kahneman, D., Knetsch, J. L., \& Thaler, R. H. (1991). Anomalies: The endowment effect, loss aversion, and status quo bias. \textit{Journal of Economic Perspectives}, 5(1), 193-206.
    \item Samuelson, W., \& Zeckhauser, R. (1988). Status quo bias in decision making. \textit{Journal of Risk and Uncertainty}, 1(1), 7-59.
    \item Strack, F., \& Mussweiler, T. (1997). Explaining the enigmatic anchoring effect: Mechanisms of selective accessibility. \textit{Journal of Personality and Social Psychology}, 73(3), 437-446.
    \item Epley, N., \& Gilovich, T. (2006). The anchoring-and-adjustment heuristic: Why the adjustments are insufficient. \textit{Psychological Science}, 17(4), 311-318.
    \item Moore, D. A., \& Healy, P. J. (2008). The trouble with overconfidence. \textit{Psychological Review}, 115(2), 502-517.
    \item Fischhoff, B., Slovic, P., \& Lichtenstein, S. (1977). Knowing with certainty: The appropriateness of extreme confidence. \textit{Journal of Experimental Psychology: Human Perception and Performance}, 3(4), 552-564.
\end{itemize}

\subsection*{Kognitive Biases -- Erweiterte Liste (MA-EPI-26 bis MA-EPI-37)}

\begin{itemize}[leftmargin=*]
    \item Kahneman, D., \& Tversky, A. (1973). On the psychology of prediction. \textit{Psychological Review}, 80(4), 237-251. [Base Rate Neglect]
    \item Tversky, A., \& Kahneman, D. (1972). Subjective probability: A judgment of representativeness. \textit{Cognitive Psychology}, 3(3), 430-454. [Representativeness]
    \item Arkes, H. R., \& Blumer, C. (1985). The psychology of sunk cost. \textit{Organizational Behavior and Human Decision Processes}, 35(1), 124-140. [Sunk Cost]
    \item Thaler, R. (1980). Toward a positive theory of consumer choice. \textit{Journal of Economic Behavior \& Organization}, 1(1), 39-60. [Endowment Effect]
    \item Kahneman, D., Knetsch, J. L., \& Thaler, R. H. (1990). Experimental tests of the endowment effect and the Coase theorem. \textit{Journal of Political Economy}, 98(6), 1325-1348.
    \item Weinstein, N. D. (1980). Unrealistic optimism about future life events. \textit{Journal of Personality and Social Psychology}, 39(5), 806-820. [Optimism Bias]
    \item Sharot, T. (2011). The optimism bias. \textit{Current Biology}, 21(23), R941-R945.
    \item Miller, D. T., \& Ross, M. (1975). Self-serving biases in the attribution of causality. \textit{Psychological Bulletin}, 82(2), 213-225. [Self-Serving Bias]
    \item Crowne, D. P., \& Marlowe, D. (1960). A new scale of social desirability. \textit{Journal of Consulting Psychology}, 24(4), 349-354. [Social Desirability]
    \item Paulhus, D. L. (1984). Two-component models of socially desirable responding. \textit{Journal of Personality and Social Psychology}, 46(3), 598-609.
    \item Spranca, M., Minsk, E., \& Baron, J. (1991). Omission and commission in judgment and choice. \textit{Journal of Experimental Social Psychology}, 27(1), 76-105. [Omission Bias]
    \item Ritov, I., \& Baron, J. (1990). Reluctance to vaccinate: Omission bias and ambiguity. \textit{Journal of Behavioral Decision Making}, 3(4), 263-277.
    \item Merritt, A. C., Effron, D. A., \& Monin, B. (2010). Moral self-licensing: When being good frees us to be bad. \textit{Social and Personality Psychology Compass}, 4(5), 344-357. [Moral Licensing]
    \item Sachdeva, S., Iliev, R., \& Medin, D. L. (2009). Sinning saints and saintly sinners: The paradox of moral self-regulation. \textit{Psychological Science}, 20(4), 523-528.
    \item Small, D. A., \& Loewenstein, G. (2003). Helping a victim or helping the victim: Altruism and identifiability. \textit{Journal of Risk and Uncertainty}, 26(1), 5-16. [Identifiable Victim]
    \item Jenni, K., \& Loewenstein, G. (1997). Explaining the identifiable victim effect. \textit{Journal of Risk and Uncertainty}, 14(3), 235-257.
    \item Kahneman, D. (1986). Comments on the contingent valuation method. In R. G. Cummings et al. (Eds.), \textit{Valuing Environmental Goods}. Rowman \& Allanheld. [Scope Insensitivity]
    \item Desvousges, W. H., et al. (1993). Measuring nonuse damages using contingent valuation. \textit{Journal of Environmental Economics and Management}, 23(3), 177-194.
    \item Slovic, P., Finucane, M., Peters, E., \& MacGregor, D. G. (2002). The affect heuristic. In T. Gilovich et al. (Eds.), \textit{Heuristics and Biases} (pp. 397-420). Cambridge University Press. [Affect Heuristic]
    \item Slovic, P., Finucane, M. L., Peters, E., \& MacGregor, D. G. (2007). The affect heuristic. \textit{European Journal of Operational Research}, 177(3), 1333-1352.
    \item Nickerson, R. S. (1998). Confirmation bias: A ubiquitous phenomenon in many guises. \textit{Review of General Psychology}, 2(2), 175-220. [Confirmation Bias]
\end{itemize}

\subsection*{Identitätsökonomie \& Soziale Präferenzen (MA-ONT-09 bis MA-ONT-11)}

\begin{itemize}[leftmargin=*]
    \item Akerlof, G. A., \& Kranton, R. E. (2000). Economics and identity. \textit{Quarterly Journal of Economics}, 115(3), 715-753.
    \item Akerlof, G. A., \& Kranton, R. E. (2010). \textit{Identity Economics: How Our Identities Shape Our Work, Wages, and Well-Being}. Princeton University Press.
    \item Bénabou, R., \& Tirole, J. (2011). Identity, morals, and taboos: Beliefs as assets. \textit{Quarterly Journal of Economics}, 126(2), 805-855.
    \item Fehr, E., \& Schmidt, K. M. (1999). A theory of fairness, competition, and cooperation. \textit{Quarterly Journal of Economics}, 114(3), 817-868.
    \item Fehr, E., \& Fischbacher, U. (2002). Why social preferences matter -- the impact of non-selfish motives on competition, cooperation and incentives. \textit{Economic Journal}, 112(478), C1-C33.
    \item Fischbacher, U., Gächter, S., \& Fehr, E. (2001). Are people conditionally cooperative? Evidence from a public goods experiment. \textit{Economics Letters}, 71(3), 397-404.
    \item Bolton, G. E., \& Ockenfels, A. (2000). ERC: A theory of equity, reciprocity, and competition. \textit{American Economic Review}, 90(1), 166-193.
    \item Charness, G., \& Rabin, M. (2002). Understanding social preferences with simple tests. \textit{Quarterly Journal of Economics}, 117(3), 817-869.
    \item Tajfel, H. (1978). \textit{Differentiation Between Social Groups}. Academic Press.
    \item Tajfel, H., \& Turner, J. C. (1979). An integrative theory of intergroup conflict. In W. G. Austin \& S. Worchel (Eds.), \textit{The Social Psychology of Intergroup Relations} (pp. 33-47). Brooks/Cole.
\end{itemize}

\subsection*{Soziale Normen \& Kooperation (MA-ONT-11, MA-DYN-10 bis MA-DYN-12)}

\begin{itemize}[leftmargin=*]
    \item Bicchieri, C. (2006). \textit{The Grammar of Society: The Nature and Dynamics of Social Norms}. Cambridge University Press.
    \item Bicchieri, C. (2017). \textit{Norms in the Wild: How to Diagnose, Measure, and Change Social Norms}. Oxford University Press.
    \item Cialdini, R. B., Reno, R. R., \& Kallgren, C. A. (1990). A focus theory of normative conduct: Recycling the concept of norms. \textit{Journal of Personality and Social Psychology}, 58(6), 1015-1026.
    \item Cialdini, R. B., \& Goldstein, N. J. (2004). Social influence: Compliance and conformity. \textit{Annual Review of Psychology}, 55, 591-621.
    \item Centola, D. (2018). \textit{How Behavior Spreads: The Science of Complex Contagions}. Princeton University Press.
    \item Centola, D., \& Macy, M. (2007). Complex contagions and the weakness of long ties. \textit{American Journal of Sociology}, 113(3), 702-734.
    \item Granovetter, M. (1978). Threshold models of collective behavior. \textit{American Journal of Sociology}, 83(6), 1420-1443.
    \item Young, H. P. (2015). The evolution of social norms. \textit{Annual Review of Economics}, 7, 359-387.
    \item Ostrom, E. (1990). \textit{Governing the Commons: The Evolution of Institutions for Collective Action}. Cambridge University Press.
\end{itemize}

\subsection*{Temporale Entscheidungen \& Discounting (MA-DYN-13 bis MA-DYN-16)}

\begin{itemize}[leftmargin=*]
    \item Frederick, S., Loewenstein, G., \& O'Donoghue, T. (2002). Time discounting and time preference: A critical review. \textit{Journal of Economic Literature}, 40(2), 351-401.
    \item Laibson, D. (1997). Golden eggs and hyperbolic discounting. \textit{Quarterly Journal of Economics}, 112(2), 443-478.
    \item O'Donoghue, T., \& Rabin, M. (1999). Doing it now or later. \textit{American Economic Review}, 89(1), 103-124.
    \item Strotz, R. H. (1955). Myopia and inconsistency in dynamic utility maximization. \textit{Review of Economic Studies}, 23(3), 165-180.
    \item Loewenstein, G. (1996). Out of control: Visceral influences on behavior. \textit{Organizational Behavior and Human Decision Processes}, 65(3), 272-292.
    \item Loewenstein, G., \& Prelec, D. (1992). Anomalies in intertemporal choice: Evidence and an interpretation. \textit{Quarterly Journal of Economics}, 107(2), 573-597.
    \item Loewenstein, G., O'Donoghue, T., \& Rabin, M. (2003). Projection bias in predicting future utility. \textit{Quarterly Journal of Economics}, 118(4), 1209-1248.
    \item Van Boven, L., \& Loewenstein, G. (2003). Social projection of transient drive states. \textit{Personality and Social Psychology Bulletin}, 29(9), 1159-1168.
\end{itemize}

\subsection*{Gewohnheiten \& Selbstkontrolle (MA-DYN-07 bis MA-DYN-09, MA-DYN-17 bis MA-DYN-18)}

\begin{itemize}[leftmargin=*]
    \item Duhigg, C. (2012). \textit{The Power of Habit: Why We Do What We Do in Life and Business}. Random House.
    \item Wood, W., \& Rünger, D. (2016). Psychology of habit. \textit{Annual Review of Psychology}, 67, 289-314.
    \item Neal, D. T., Wood, W., \& Quinn, J. M. (2006). Habits -- A repeat performance. \textit{Current Directions in Psychological Science}, 15(4), 198-202.
    \item Baumeister, R. F., Bratslavsky, E., Muraven, M., \& Tice, D. M. (1998). Ego depletion: Is the active self a limited resource? \textit{Journal of Personality and Social Psychology}, 74(5), 1252-1265.
    \item Baumeister, R. F., Vohs, K. D., \& Tice, D. M. (2007). The strength model of self-control. \textit{Current Directions in Psychological Science}, 16(6), 351-355.
    \item Gollwitzer, P. M. (1999). Implementation intentions: Strong effects of simple plans. \textit{American Psychologist}, 54(7), 493-503.
    \item Gollwitzer, P. M., \& Sheeran, P. (2006). Implementation intentions and goal achievement: A meta-analysis. \textit{Advances in Experimental Social Psychology}, 38, 69-119.
    \item Bryan, G., Karlan, D., \& Nelson, S. (2010). Commitment devices. \textit{Annual Review of Economics}, 2(1), 671-698.
\end{itemize}

\subsection*{Hedonische Adaptation \& Wohlbefinden (MA-DYN-19 bis MA-DYN-21)}

\begin{itemize}[leftmargin=*]
    \item Brickman, P., \& Campbell, D. T. (1971). Hedonic relativism and planning the good society. In M. H. Appley (Ed.), \textit{Adaptation-Level Theory} (pp. 287-305). Academic Press.
    \item Brickman, P., Coates, D., \& Janoff-Bulman, R. (1978). Lottery winners and accident victims: Is happiness relative? \textit{Journal of Personality and Social Psychology}, 36(8), 917-927.
    \item Diener, E., Lucas, R. E., \& Scollon, C. N. (2006). Beyond the hedonic treadmill: Revising the adaptation theory of well-being. \textit{American Psychologist}, 61(4), 305-314.
    \item Kahneman, D., Wakker, P. P., \& Sarin, R. (1997). Back to Bentham? Explorations of experienced utility. \textit{Quarterly Journal of Economics}, 112(2), 375-406.
    \item Kahneman, D., Fredrickson, B. L., Schreiber, C. A., \& Redelmeier, D. A. (1993). When more pain is preferred to less: Adding a better end. \textit{Psychological Science}, 4(6), 401-405. [Peak-End Rule]
    \item Redelmeier, D. A., \& Kahneman, D. (1996). Patients' memories of painful medical treatments: Real-time and retrospective evaluations. \textit{Pain}, 66(1), 3-8.
\end{itemize}

\subsection*{Nudging \& Interventionsdesign (MA-GOV-07 bis MA-GOV-10)}

\begin{itemize}[leftmargin=*]
    \item Thaler, R. H., \& Sunstein, C. R. (2008). \textit{Nudge: Improving Decisions About Health, Wealth, and Happiness}. Yale University Press.
    \item Sunstein, C. R. (2014). \textit{Why Nudge? The Politics of Libertarian Paternalism}. Yale University Press.
    \item Johnson, E. J., \& Goldstein, D. (2003). Do defaults save lives? \textit{Science}, 302(5649), 1338-1339.
    \item Madrian, B. C., \& Shea, D. F. (2001). The power of suggestion: Inertia in 401(k) participation and savings behavior. \textit{Quarterly Journal of Economics}, 116(4), 1149-1187.
    \item Benartzi, S., \& Thaler, R. H. (2013). Behavioral economics and the retirement savings crisis. \textit{Science}, 339(6124), 1152-1153.
    \item Hertwig, R., \& Grüne-Yanoff, T. (2017). Nudging and boosting: Steering or empowering good decisions. \textit{Perspectives on Psychological Science}, 12(6), 973-986.
    \item Grüne-Yanoff, T., \& Hertwig, R. (2016). Nudge versus boost: How coherent are policy and theory? \textit{Minds and Machines}, 26(1-2), 149-183.
    \item Hallsworth, M., \& Kirkman, E. (2020). \textit{Behavioral Insights}. MIT Press.
\end{itemize}

\subsection*{Ethik \& Governance von Interventionen (MA-GOV-03 bis MA-GOV-06, MA-GOV-11 bis MA-GOV-18)}

\begin{itemize}[leftmargin=*]
    \item Sunstein, C. R. (2016). \textit{The Ethics of Influence: Government in the Age of Behavioral Science}. Cambridge University Press.
    \item Hausman, D. M., \& Welch, B. (2010). Debate: To nudge or not to nudge. \textit{Journal of Political Philosophy}, 18(1), 123-136.
    \item Hansen, P. G., \& Jespersen, A. M. (2013). Nudge and the manipulation of choice. \textit{European Journal of Risk Regulation}, 4(1), 3-28.
    \item Bovens, L. (2009). The ethics of nudge. In T. Grüne-Yanoff \& S. O. Hansson (Eds.), \textit{Preference Change} (pp. 207-219). Springer.
    \item Frey, B. S., \& Jegen, R. (2001). Motivation crowding theory. \textit{Journal of Economic Surveys}, 15(5), 589-611. [Crowding Out]
    \item Deci, E. L., Koestner, R., \& Ryan, R. M. (1999). A meta-analytic review of experiments examining the effects of extrinsic rewards on intrinsic motivation. \textit{Psychological Bulletin}, 125(6), 627-668.
    \item Gneezy, U., Meier, S., \& Rey-Biel, P. (2011). When and why incentives (don't) work to modify behavior. \textit{Journal of Economic Perspectives}, 25(4), 191-210.
    \item Bowles, S., \& Polania-Reyes, S. (2012). Economic incentives and social preferences: Substitutes or complements? \textit{Journal of Economic Literature}, 50(2), 368-425.
\end{itemize}

\subsection*{Dynamische Systeme \& Komplexität (MA-DYN-04 bis MA-DYN-06)}

\begin{itemize}[leftmargin=*]
    \item Arthur, W. B. (1994). \textit{Increasing Returns and Path Dependence in the Economy}. University of Michigan Press.
    \item David, P. A. (1985). Clio and the economics of QWERTY. \textit{American Economic Review}, 75(2), 332-337.
    \item Gladwell, M. (2000). \textit{The Tipping Point: How Little Things Can Make a Big Difference}. Little, Brown.
    \item Scheffer, M. (2009). \textit{Critical Transitions in Nature and Society}. Princeton University Press.
    \item Page, S. E. (2006). Path dependence. \textit{Quarterly Journal of Political Science}, 1(1), 87-115.
\end{itemize}

\subsection*{Methodologie \& Evaluation (MA-GOV-20 bis MA-GOV-24)}

\begin{itemize}[leftmargin=*]
    \item Angrist, J. D., \& Pischke, J. S. (2009). \textit{Mostly Harmless Econometrics: An Empiricist's Companion}. Princeton University Press.
    \item Duflo, E., Glennerster, R., \& Kremer, M. (2007). Using randomization in development economics research: A toolkit. \textit{Handbook of Development Economics}, 4, 3895-3962.
    \item List, J. A. (2011). Why economists should conduct field experiments and 14 tips for pulling one off. \textit{Journal of Economic Perspectives}, 25(3), 3-16.
    \item DellaVigna, S., \& Linos, E. (2022). RCTs to scale: Comprehensive evidence from two nudge units. \textit{Econometrica}, 90(1), 81-116.
    \item Banerjee, A. V., \& Duflo, E. (2009). The experimental approach to development economics. \textit{Annual Review of Economics}, 1, 151-178.
\end{itemize}

\subsection*{Kulturelle \& Institutionelle Ebenen (MA-ONT-26)}

\begin{itemize}[leftmargin=*]
    \item Hofstede, G. (2001). \textit{Culture's Consequences: Comparing Values, Behaviors, Institutions and Organizations Across Nations} (2nd ed.). Sage.
    \item Henrich, J., Heine, S. J., \& Norenzayan, A. (2010). The weirdest people in the world? \textit{Behavioral and Brain Sciences}, 33(2-3), 61-83.
    \item North, D. C. (1990). \textit{Institutions, Institutional Change and Economic Performance}. Cambridge University Press.
    \item Acemoglu, D., \& Robinson, J. A. (2012). \textit{Why Nations Fail: The Origins of Power, Prosperity, and Poverty}. Crown Business.
    \item Alesina, A., \& Giuliano, P. (2015). Culture and institutions. \textit{Journal of Economic Literature}, 53(4), 898-944.
    \item Gelfand, M. J. (2018). \textit{Rule Makers, Rule Breakers: How Tight and Loose Cultures Wire Our World}. Scribner.
\end{itemize}

%==============================================================================
% Appendix
%==============================================================================

\appendix



% Kapitel 10: Vollständige Axiom-Liste
\section{Vollständige Axiom-Liste nach ID}

\subsection{Ontologie (ONT): 26 Axiome}
\begin{enumerate}[label=MA-ONT-\arabic*:, leftmargin=3cm]
    \item $\exists \sigma$: Subjekt existiert
    \item $\exists V$: Verhalten existiert
    \item $\exists U$: Nutzen existiert
    \item $\exists \rho$: Kontext existiert
    \item $\exists \tau$: Zeit existiert
    \item $\exists \Gamma$: Constraints existieren
    \item $U = U_{\text{INU}} \oplus U_{\text{KNU}} \oplus U_{\text{IDN}}$
    \item FEPSDE: 6 Nutzen-Dimensionen
    \item $\exists I$: Identitäten sind multipel
    \item $\exists G$: Gruppen existieren
    \item $\exists N$: Normen existieren
    \item $\exists \kappa$: Komplementarität existiert
    \item $\Psi$: Kohärenz-Operator
    \item $\lambda$: Kooperations-Regulator
    \item $\lambda \in [0,1]$: Hybrid-Kontinuum
    \item $\omega$: Gewichtungsfunktion
    \item $\pi$: Identitäts-Aktivierung
    \item $A$: Alignment-Funktion
    \item $\partial U / \partial \tau$: Nutzen ist zeitabhängig
    \item Habituation existiert
    \item Feedback-Schleifen existieren
    \item Spillover existiert
    \item Identitäten evolvieren
    \item Normen evolvieren
    \item Meta-Evolution existiert
    \item \textbf{BCM operiert auf 5 hierarchischen Ebenen}
\end{enumerate}

\subsection{Epistemologie (EPI): 37 Axiome}
\begin{enumerate}[label=MA-EPI-\arabic*:, leftmargin=3cm]
    \item Unvollständiges Wissen
    \item Bounded Rationality
    \item Unsicherheit ist irreduzibel
    \item Heuristiken sind adaptiv
    \item Zwei Systeme (System 1 / System 2)
    \item Biases sind systematisch
    \item AWX: Awareness
    \item WAX: Willingness
    \item WTX: Transition
    \item Sequenz: AWX → WAX → WTX → V
    \item Intention-Action Gap
    \item Verfügbarkeitsheuristik
    \item Framing-Effekte
    \item Anchoring
    \item Loss Aversion
    \item Present Bias
    \item Status Quo Bias
    \item Overconfidence
    \item Lernen ist möglich
    \item Feedback verbessert Schätzungen
    \item Gewohnheiten reduzieren Kognition
    \item Debiasing ist schwierig
    \item Exponentielle Sättigung
    \item Belief-Update folgt Bayes
    \item Confirmation Bias
    \item \textbf{Base Rate Neglect}
    \item \textbf{Representativeness Heuristic}
    \item \textbf{Sunk Cost Fallacy}
    \item \textbf{Endowment Effect}
    \item \textbf{Optimism Bias}
    \item \textbf{Self-Serving Bias}
    \item \textbf{Social Desirability Bias}
    \item \textbf{Omission Bias}
    \item \textbf{Moral Licensing}
    \item \textbf{Identifiable Victim Effect}
    \item \textbf{Scope Insensitivity}
    \item \textbf{Affect Heuristic}
\end{enumerate}

\subsection{Governance (GOV): 25 Axiome}
\begin{enumerate}[label=MA-GOV-\arabic*:, leftmargin=3cm]
    \item Verhalten unterliegt Constraints: $\Gamma: V \rightarrow \{0, 1\}$
    \item Constraints sind kontextabhängig: $\Gamma = f(\rho)$
    \item Autonomie-Primat: Individuelle Autonomie hat Vorrang
    \item Nicht-Normativität: BCM beschreibt, schreibt nicht vor
    \item Transparenz-Gebot: Interventionen sollen transparent sein
    \item Subsidiarität: Mildeste wirksame Intervention wählen
    \item 8 Interventionstypen: NUDGE, BOOST, THINK, INCENTIVE, REGULATE, DESIGN, SOCIAL, IDENTITY
    \item Typenwahl kontextabhängig: $I^* = \arg\max_{I} E[Effect(I)|\rho, \sigma]$
    \item Freiheits-Hierarchie: Ordnung nach Freiheitsgrad
    \item Eskalations-Prinzip: Von mild nach stark eskalieren
    \item Distributive Gerechtigkeit: $\forall \sigma: \text{Burden}(\sigma) \propto \text{Capacity}(\sigma)$
    \item Kosten-Nutzen-Analyse: Interventionen müssen Nutzen > Kosten haben
    \item Reversibilität: Reversible Interventionen bevorzugen
    \item Crowding-Out vermeiden: $\partial U_{\text{intrinsic}}/\partial \text{Incentive} < 0$ möglich
    \item Spillover beachten: Nicht-intendierte Effekte berücksichtigen
    \item Ausweichverhalten beachten: $P(\text{Evasion}|\text{REGULATE}) > 0$
    \item Opt-Out Prinzip: $\forall I \neq \text{REGULATE}: \text{OptOut}(I) = \text{possible}$
    \item Informed Consent: $\text{Legitimate}(I) \rightarrow \text{Informed}(\sigma) \wedge \text{Consenting}(\sigma)$
    \item Interventionen habituieren: $\partial \text{Effect}(I)/\partial \tau < 0$ für $\tau \rightarrow \infty$
    \item Iteration notwendig: $\text{Optimal}(I)$ requires Feedback-Loop
    \item Evaluation notwendig: $\forall I: \text{Measure}(\text{Effect}(I))$
    \item A/B-Testing bevorzugen: RCT > Observational für Kausalität
    \item Spillover messen: $\text{Measure}(\text{Effect}_{\text{unintended}})$ alongside $\text{Effect}_{\text{intended}}$
    \item Langfristeffekte beachten: $\text{Effect}(\tau=\text{short}) \neq \text{Effect}(\tau=\text{long})$
    \item Multi-Level Interventionen: Interventionen auf mehreren Ebenen kombinieren
\end{enumerate}

\subsection{Dynamik (DYN): 25 Axiome}
\begin{enumerate}[label=MA-DYN-\arabic*:, leftmargin=3cm]
    \item Zeit existiert: $\tau \in T, T \subseteq \mathbb{R}^+$
    \item Diskrete Events: $V(t) = \sum_i V_i \cdot \delta(t - t_i)$
    \item Pfadabhängigkeit: $U(V, t) = f(V, \{V_\tau : \tau < t\})$
    \item Multiple Equilibria: Es kann mehrere stabile Zustände geben
    \item Hysterese möglich: $\text{State}(\uparrow) \neq \text{State}(\downarrow)$ bei gleichem Trigger
    \item Tipping Points: Kritische Schwellen existieren
    \item Gewohnheitsformation: $P(V_{\text{auto}}|\text{Repetition}) \uparrow$ mit $n$
    \item Habit Loop: Cue $\rightarrow$ Routine $\rightarrow$ Reward
    \item Commitment Devices: $U(\text{Commit}) > U(\text{NoCommit})$ wenn $\delta_{\text{now}} > \delta_{\text{future}}$
    \item Soziale Diffusion: $\partial P(V)/\partial t = \beta \cdot P(V) \cdot (1 - P(V)) \cdot f(\text{Network})$
    \item Kritische Masse: Ab $P^* \approx 25\%$ wird Verhalten selbstverstärkend
    \item Norm-Kaskaden: Normen können rapid kippen
    \item Discounting: $U_{\text{present}}(U_{\text{future}}) = \delta^t \cdot U_{\text{future}}$
    \item Hyperbolisches Discounting: $D(t) = 1/(1 + k \cdot t)$ statt $D(t) = e^{-\delta t}$
    \item Projection Bias: $\mathbb{E}[U_{\text{future}}|\text{State}_{\text{now}}]$ biased toward $\text{State}_{\text{now}}$
    \item Hot-Cold Empathy Gap: $\mathbb{E}[U_{\text{hot}}|\text{State}_{\text{cold}}] < U_{\text{hot,actual}}$
    \item Ego Depletion: $\text{Willpower}(t) \downarrow$ mit Nutzung
    \item Implementation Intentions: $P(V|\text{If-Then Plan}) > P(V|\text{Goal only})$
    \item Adaptation: $U(\text{State}) \rightarrow U_{\text{baseline}}$ as $\tau \rightarrow \infty$
    \item Hedonic Treadmill: $\lim_{t \rightarrow \infty} U(t|\text{Change}) = U_{\text{baseline}}$
    \item Peak-End Rule: $\text{Recalled}_U = f(\text{Max}_U, \text{End}_U)$
    \item Identitäts-Shifts: $I(\tau + \Delta\tau) \not\subset I(\tau)$ möglich
    \item $\lambda$-Dynamik: $\lambda(\tau + \Delta\tau) = f(\lambda(\tau), \text{Experience}, \text{Context})$
    \item System-Evolution: $\Phi(\tau + \Delta\tau) = \Phi(\tau) + \mu \cdot \Delta\Phi$
    \item Feedback-Lags: Verzögerungen zwischen Aktion und Feedback
\end{enumerate}

%==============================================================================
% Detaillierte Axiom-Dokumentation
%==============================================================================



%==============================================================================
% DETAILLIERTE AXIOM-DOKUMENTATION
%==============================================================================

\section{Detaillierte Axiom-Dokumentation}

Die folgenden Abschnitte enthalten die vollständige 9-Felder-Dokumentation aller 115 BCM META-Axiome. Jedes Axiom wird systematisch beschrieben mit:

\begin{enumerate}
    \item \textbf{Statement}: Kernaussage des Axioms
    \item \textbf{Formel}: Mathematische Formalisierung
    \item \textbf{Beschreibung}: Ausführliche Erläuterung
    \item \textbf{BCM-Relevanz}: Bedeutung für das Behavioral Change Model
    \item \textbf{BEATRIX-Relevanz}: Implementierung im LLM-basierten System
    \item \textbf{Konsequenz bei Fehlen}: Was passiert, wenn dieses Axiom nicht gilt
    \item \textbf{Validiert}: Welche Module durch dieses Axiom validiert werden
    \item \textbf{Abhängigkeiten}: Vorausgesetzte Axiome
    \item \textbf{Literatur}: Wissenschaftliche Referenzen
\end{enumerate}

% Einbinden der detaillierten Axiom-Dokumentation
% =============================================================================
% ONT-AXIOME (ONTOLOGIE) - Vollständige Ausformulierung
% BCM 2.2 META-Modul
% =============================================================================

\section{Ontologie-Axiome (MA-ONT)}

Die ONT-Axiome definieren die fundamentalen Entitäten und Beziehungen des BCM.
Sie beantworten die Frage: \textit{Was existiert?}

%==============================================================================
% MA-ONT-01: Bereits in axiom_beispiele.tex
%==============================================================================

%==============================================================================
% MA-ONT-02: Existenz von Verhalten
%==============================================================================

\subsubsection*{MA-ONT-02: Existenz von Verhalten}

\textbf{Statement:} Es existieren beobachtbare Verhaltensweisen $V$.

\textbf{Formel:}
\begin{equation}
V \in \mathcal{V}, \quad |\mathcal{V}| \geq 1
\end{equation}

\textbf{Beschreibung:}
Verhalten umfasst alle beobachtbaren Handlungen eines Subjekts. Dazu gehören manifestes Verhalten (Kauf, Wahl, Bewegung), aber auch Unterlassungen (Nicht-Handeln als Handlung). Verhalten ist die abhängige Variable des BCM -- das, was erklärt und beeinflusst werden soll.

\textbf{BCM-Relevanz:}
Ohne beobachtbares Verhalten wäre das BCM nicht falsifizierbar. Alle Module (INU, KNU, IDN) haben Verhalten als Output. Die WTX-Komponente (Willingness to X) beschreibt die Transition von Intention zu Verhalten.

\textbf{BEATRIX-Relevanz:}
BEATRIX sagt Verhaltenswahrscheinlichkeiten vorher. Das LLM analysiert die Persona und schätzt $P(V_i)$ für alternative Verhaltensweisen. Die Verhaltensoptionen müssen dem LLM explizit präsentiert werden (Choice Set).

\textbf{Konsequenz bei Fehlen:}
Das BCM wäre ein reines Präferenzmodell ohne Handlungsbezug. Es könnte nur erklären, was Menschen wollen, nicht was sie tun. Die Intention-Action-Gap (MA-EPI-11) wäre nicht modellierbar.

\textbf{Validiert:} WTX

\textbf{Abhängigkeiten:} MA-ONT-01 (Subjekt existiert)

\textbf{Literatur:} Ajzen (1991); Fishbein \& Ajzen (1975)

%==============================================================================
% MA-ONT-03: Existenz von Nutzen
%==============================================================================

\subsubsection*{MA-ONT-03: Existenz von Nutzen}

\textbf{Statement:} Subjekte haben Nutzenfunktionen $U$.

\textbf{Formel:}
\begin{equation}
U: \mathcal{V} \times \Sigma \rightarrow \mathbb{R}
\end{equation}

\textbf{Beschreibung:}
Jedes Subjekt bewertet jedes mögliche Verhalten mit einem skalaren Nutzenwert. Diese Bewertung ist subjektiv und kann von objektiven Ergebnissen abweichen. Der Nutzen ist nicht direkt beobachtbar, sondern wird aus Verhalten erschlossen (revealed preferences).

\textbf{BCM-Relevanz:}
Nutzen ist die zentrale erklärende Variable. Das BCM postuliert (schwache) Nutzenmaximierung: Subjekte wählen tendenziell Verhaltensweisen mit höherem erwarteten Nutzen. Die Stärke der Präferenz wird durch Nutzendifferenzen ausgedrückt.

\textbf{BEATRIX-Relevanz:}
BEATRIX schätzt Nutzenparameter aus Persona-Beschreibungen. Die FEPSDE-Dimensionen (MA-ONT-08) strukturieren den individuellen Nutzen. Das LLM verwendet die geschätzten Parameter zur Verhaltensvorhersage.

\textbf{Konsequenz bei Fehlen:}
Kein Optimierungsmodell möglich. Verhalten wäre zufällig oder rein reaktiv. Keine Grundlage für Präferenzanalyse, Trade-off-Modellierung oder Interventionsdesign.

\textbf{Validiert:} INU

\textbf{Abhängigkeiten:} MA-ONT-01, MA-ONT-02

\textbf{Literatur:} von Neumann \& Morgenstern (1944); Samuelson (1938)

%==============================================================================
% MA-ONT-04: Existenz von Kontext
%==============================================================================

\subsubsection*{MA-ONT-04: Existenz von Kontext}

\textbf{Statement:} Verhalten findet in Kontexten $\rho$ statt.

\textbf{Formel:}
\begin{equation}
\rho \in \mathcal{P}, \quad |\mathcal{P}| \geq 1
\end{equation}

\textbf{Beschreibung:}
Kontext umfasst alle situativen Faktoren: physische Umgebung, soziales Setting, zeitlicher Rahmen, institutioneller Kontext. Derselbe Mensch verhält sich in verschiedenen Kontexten unterschiedlich. Kontext aktiviert spezifische Identitäten und Normen.

\textbf{BCM-Relevanz:}
Kontextabhängigkeit ist fundamental. Der $\lambda$-Parameter, Identitätsaktivierung und Normensalienz sind alle kontextabhängig. Das 5-Ebenen-Modell (MA-ONT-26) strukturiert Kontexte hierarchisch.

\textbf{BEATRIX-Relevanz:}
BEATRIX erhält Kontextinformationen als Input. Das LLM muss wissen: Wo? Wann? Mit wem? Unter welchen Regeln? Ohne Kontext sind Parameter-Schätzungen unspezifisch und ungenau.

\textbf{Konsequenz bei Fehlen:}
Trait-only-Modell. Verhalten wäre rein personenabhängig, situative Faktoren ignoriert. Ecological validity würde massiv sinken. Interventionen könnten nicht kontextspezifisch designt werden.

\textbf{Validiert:} INU, KNU, IDN

\textbf{Abhängigkeiten:} MA-ONT-02

\textbf{Literatur:} Mischel \& Shoda (1995); Ross \& Nisbett (1991)

%==============================================================================
% MA-ONT-05: Existenz von Zeit
%==============================================================================

\subsubsection*{MA-ONT-05: Existenz von Zeit}

\textbf{Statement:} Verhalten hat eine temporale Dimension $\tau$.

\textbf{Formel:}
\begin{equation}
\tau \in T, \quad T \subseteq \mathbb{R}^+
\end{equation}

\textbf{Beschreibung:}
Zeit ist nicht nur Hintergrund, sondern aktiver Faktor. Entscheidungen haben zeitliche Konsequenzen (sofort vs. später). Präferenzen können sich über Zeit ändern (dynamische Inkonsistenz). Gewohnheiten bilden sich über Wiederholung.

\textbf{BCM-Relevanz:}
Ermöglicht das DYN-Modul. Ohne Zeit keine Dynamik: keine Habituation (MA-ONT-20), keine Diskontierung (MA-DYN-13), keine Evolution von Normen (MA-ONT-24). Zeit ist die unabhängige Variable der Veränderung.

\textbf{BEATRIX-Relevanz:}
BEATRIX berücksichtigt temporale Aspekte:
\begin{itemize}
    \item Zeithorizont der Entscheidung (kurzfristig vs. langfristig)
    \item Zeitpunkt der Intervention
    \item Historische Verhaltensmuster der Persona
\end{itemize}

\textbf{Konsequenz bei Fehlen:}
Statisches Modell. Keine Verhaltensänderung modellierbar. Keine Gewohnheitsbildung, keine Lerneffekte, keine Trendvorhersagen. Das BCM wäre ein Snapshot ohne Dynamik.

\textbf{Validiert:} WTX

\textbf{Abhängigkeiten:} MA-ONT-02

\textbf{Literatur:} Strotz (1955); Frederick et al. (2002)

%==============================================================================
% MA-ONT-06: Existenz von Constraints
%==============================================================================

\subsubsection*{MA-ONT-06: Existenz von Constraints}

\textbf{Statement:} Es gibt Beschränkungen für Verhalten $\Gamma$.

\textbf{Formel:}
\begin{equation}
\Gamma: \mathcal{V} \rightarrow \{0, 1\}
\end{equation}

\textbf{Beschreibung:}
Nicht alles, was gewollt wird, kann getan werden. Constraints umfassen: Budget (finanzielle Grenzen), Zeit (24h/Tag), physische Grenzen (Körper), rechtliche Grenzen (Verbote), soziale Grenzen (Sanktionen). Verhalten wird aus der Menge der feasible options gewählt.

\textbf{BCM-Relevanz:}
Unterscheidet Präferenzen von Verhalten. Das BCM modelliert Optimierung unter Nebenbedingungen. Das GOV-Modul (Governance) operiert primär durch Constraint-Manipulation: Erweiterung (Ermöglichung) oder Einschränkung (Regulierung).

\textbf{BEATRIX-Relevanz:}
BEATRIX muss Constraints kennen, um realistische Vorhersagen zu machen. Das LLM prüft: Ist das Verhalten überhaupt möglich? Welche Barrieren existieren? Interventionen können auf Constraint-Reduktion zielen.

\textbf{Konsequenz bei Fehlen:}
Wunschmodell statt Verhaltensmodell. Alle gewünschten Verhaltensweisen würden als möglich angenommen. Interventionen würden Barrieren ignorieren. Feasibility-Checks wären nicht möglich.

\textbf{Validiert:} GOV

\textbf{Abhängigkeiten:} MA-ONT-02

\textbf{Literatur:} Becker (1965); Lancaster (1966)

%==============================================================================
% MA-ONT-07: Bereits in axiom_beispiele.tex
%==============================================================================

%==============================================================================
% MA-ONT-08: FEPSDE-Dimensionen
%==============================================================================

\subsubsection*{MA-ONT-08: FEPSDE-Dimensionen}

\textbf{Statement:} Der individuelle Nutzen hat sechs Dimensionen: Financial, Emotional, Physical, Social, Digital, Ecological.

\textbf{Formel:}
\begin{equation}
U_{\text{INU}} = \sum_{d \in \text{FEPSDE}} \omega_d \cdot u_d, \quad \sum_d \omega_d = 1
\end{equation}

\textbf{Beschreibung:}
FEPSDE ist ein erschöpfendes Kategoriensystem für individuelle Nutzendimensionen:
\begin{itemize}
    \item \textbf{F}inancial: Geld, Vermögen, materielle Ressourcen
    \item \textbf{E}motional: Gefühle, Wohlbefinden, Stressreduktion
    \item \textbf{P}hysical: Gesundheit, körperliche Unversehrtheit, Energie
    \item \textbf{S}ocial: Beziehungen, Status, Anerkennung
    \item \textbf{D}igital: Online-Präsenz, digitale Assets, Reputation
    \item \textbf{E}cological: Umwelt, Nachhaltigkeit, Naturerlebnis
\end{itemize}

\textbf{BCM-Relevanz:}
Strukturiert das INU-Modul. Jede Verhaltensalternative hat Konsequenzen in allen sechs Dimensionen. Die Gewichte $\omega_d$ sind personenspezifisch und kontextabhängig. Trade-offs zwischen Dimensionen sind explizit modellierbar.

\textbf{BEATRIX-Relevanz:}
BEATRIX schätzt die $\omega_d$-Gewichte aus Persona-Beschreibungen. Das LLM analysiert: Was ist dieser Person wichtig? Welche Dimensionen dominieren? Die Schätzung erfolgt als 6-dimensionaler Vektor mit Summe = 1.

\textbf{Konsequenz bei Fehlen:}
Eindimensionaler Nutzen (meist: Geld). Emotionale, soziale und ökologische Motive wären nicht modellierbar. Viele Verhaltensweisen (Ehrenamt, Umweltschutz) wären unerklärlich.

\textbf{Validiert:} INU

\textbf{Abhängigkeiten:} MA-ONT-07

\textbf{Literatur:} Eigenes Konstrukt, aufbauend auf Maslow (1943); Max-Neef (1991)

%==============================================================================
% MA-ONT-09: Multiple Identitäten
%==============================================================================

\subsubsection*{MA-ONT-09: Multiple Identitäten}

\textbf{Statement:} Subjekte haben multiple Identitäten.

\textbf{Formel:}
\begin{equation}
I_\sigma = \{I_1, I_2, ..., I_n\}, \quad n \geq 1
\end{equation}

\textbf{Beschreibung:}
Ein Mensch ist nicht ``eine'' Identität, sondern ein Portfolio. Die Identitäten können sein: beruflich (Arzt), familiär (Vater), kulturell (Europäer), politisch (Grüne), religiös (Muslim), Hobby (Läufer). Verschiedene Identitäten können konfligierende Normen haben.

\textbf{BCM-Relevanz:}
Fundament des IDN-Moduls. Identitäten haben Präskriptionen (was sollte ich tun?). Verhalten, das zur aktivierten Identität passt, gibt Identitätsnutzen. Konflikte zwischen Identitäten erzeugen kognitives Unbehagen.

\textbf{BEATRIX-Relevanz:}
BEATRIX extrahiert das Identitätsportfolio aus Persona-Beschreibungen. Das LLM identifiziert: Welche Identitäten hat diese Person? Welche sind zentral? Dies ermöglicht die Berechnung von $U_{\text{IDN}}$.

\textbf{Konsequenz bei Fehlen:}
Monolithisches Selbstbild. Identitätsbasierte Interventionen (IDENTITY-Typ) wären nicht begründbar. Konflikte zwischen Rollen (Arbeit vs. Familie) wären nicht modellierbar.

\textbf{Validiert:} IDN

\textbf{Abhängigkeiten:} MA-ONT-01

\textbf{Literatur:} Akerlof \& Kranton (2000); Burke \& Stets (2009)

%==============================================================================
% MA-ONT-10: Existenz von Gruppen
%==============================================================================

\subsubsection*{MA-ONT-10: Existenz von Gruppen}

\textbf{Statement:} Subjekte gehören zu Gruppen $G$.

\textbf{Formel:}
\begin{equation}
G \subseteq \mathcal{P}(\Sigma), \quad \sigma \in G_i
\end{equation}

\textbf{Beschreibung:}
Menschen sind soziale Wesen. Sie gehören zu Familien, Teams, Unternehmen, Vereinen, Nationen. Gruppenzugehörigkeit beeinflusst Normen, Identität und Kooperationsbereitschaft. Ein Subjekt kann mehreren Gruppen angehören (nested groups).

\textbf{BCM-Relevanz:}
Fundament des KNU-Moduls (Kollektiver Nutzen). Gruppen definieren, wer ``wir'' ist. Der $\lambda$-Parameter regelt, wie stark Gruppeninteressen vs. Eigeninteressen gewichtet werden. In-group/Out-group-Differenzierung ist zentral.

\textbf{BEATRIX-Relevanz:}
BEATRIX identifiziert Gruppenzugehörigkeiten aus Persona-Beschreibungen:
\begin{itemize}
    \item Primärgruppen (Familie, enge Freunde)
    \item Sekundärgruppen (Arbeitgeber, Vereine)
    \item Kategoriale Gruppen (Nation, Geschlecht)
\end{itemize}

\textbf{Konsequenz bei Fehlen:}
Atomistisches Menschenbild. Soziale Präferenzen, Kooperation und Normenkonformität wären nicht modellierbar. Das BCM würde auf homo oeconomicus reduziert.

\textbf{Validiert:} KNU

\textbf{Abhängigkeiten:} MA-ONT-01

\textbf{Literatur:} Tajfel \& Turner (1979); Brewer (1991)

%==============================================================================
% MA-ONT-11: Existenz von Normen
%==============================================================================

\subsubsection*{MA-ONT-11: Existenz von Normen}

\textbf{Statement:} Gruppen haben soziale Normen $N$.

\textbf{Formel:}
\begin{equation}
N: G \times \mathcal{V} \rightarrow [0, 1]
\end{equation}

\textbf{Beschreibung:}
Normen sind informelle Regeln: ``Was tut man in dieser Gruppe?'' Sie können deskriptiv (was andere tun) oder injunktiv (was andere erwarten) sein. Normen werden durch Sozialisation gelernt und durch Sanktionen (Lob, Tadel) durchgesetzt.

\textbf{BCM-Relevanz:}
Normen sind der Hauptmechanismus für KNU-Effekte. Die SOCIAL-Intervention (MA-GOV-07) wirkt durch Normenaktivierung. Conditional Cooperators orientieren sich an wahrgenommenen Normen. $\lambda$ steigt, wenn Kooperationsnormen salient sind.

\textbf{BEATRIX-Relevanz:}
BEATRIX schätzt Normensalienz und -stärke:
\begin{itemize}
    \item Welche Normen gelten in den Gruppen der Persona?
    \item Wie norm-konform ist die Persona typischerweise?
    \item Welche Sanktionen drohen bei Normverletzung?
\end{itemize}

\textbf{Konsequenz bei Fehlen:}
Nur individuelle Präferenzen, keine sozialen. Konformität wäre unerklärlich. SOCIAL-Interventionen hätten keine theoretische Basis.

\textbf{Validiert:} KNU

\textbf{Abhängigkeiten:} MA-ONT-10

\textbf{Literatur:} Cialdini et al. (1990); Bicchieri (2006)

%==============================================================================
% MA-ONT-12: Existenz von Komplementarität
%==============================================================================

\subsubsection*{MA-ONT-12: Existenz von Komplementarität}

\textbf{Statement:} Gegensätze bedingen sich gegenseitig ($\kappa$).

\textbf{Formel:}
\begin{equation}
\kappa: (A, B) \rightarrow \mathbb{R}, \quad \kappa(A, B) = \kappa(B, A)
\end{equation}

\textbf{Beschreibung:}
Individuum und Kollektiv sind keine Gegensätze, sondern komplementär. Eigennutz kann Kollektivnutzen fördern (unsichtbare Hand), Kooperation kann Eigennutz steigern (Reziprozität). Der $\kappa$-Operator misst diese Synergie.

\textbf{BCM-Relevanz:}
Überwindet falsche Dichotomien. Das BCM modelliert nicht ``Egoismus vs. Altruismus'', sondern Synergien. Der $\lambda$-Parameter regelt nicht ``gut vs. böse'', sondern den optimalen Mix für die jeweilige Situation.

\textbf{BEATRIX-Relevanz:}
BEATRIX erkennt Win-Win-Situationen:
\begin{itemize}
    \item Wann fördert Kooperation auch den Eigennutz?
    \item Wann ist Defektieren selbstschädigend (Tragedy of Commons)?
    \item Synergien werden positiv, Trade-offs negativ bewertet.
\end{itemize}

\textbf{Konsequenz bei Fehlen:}
Zero-Sum-Denken. Jeder Gewinn des Kollektivs wäre ein Verlust für das Individuum. Kooperation wäre nur durch Altruismus erklärbar. Win-Win-Interventionen wären nicht identifizierbar.

\textbf{Validiert:} META

\textbf{Abhängigkeiten:} MA-ONT-07

\textbf{Literatur:} Smith (1776); Ostrom (1990); Nowak (2006)

%==============================================================================
% MA-ONT-13: Kohärenz-Operator Ψ
%==============================================================================

\subsubsection*{MA-ONT-13: Kohärenz-Operator $\Psi$}

\textbf{Statement:} Die Kombination von Nutzen ist nicht-additiv.

\textbf{Formel:}
\begin{equation}
\Psi(U) \neq \sum_i U_i \quad \text{(im Allgemeinen)}
\end{equation}

\textbf{Beschreibung:}
INU, KNU und IDN werden nicht einfach addiert. Der $\Psi$-Operator modelliert Interaktionen: Ein Verhalten kann hohen INU haben, aber bei negativem IDN-Alignment trotzdem vermieden werden. Es gibt Schwelleneffekte und Vetos.

\textbf{BCM-Relevanz:}
Ermöglicht die Modellierung von:
\begin{itemize}
    \item Identitätsvetos: Trotz hohem INU wird nicht gehandelt
    \item Kooperationsbonus: KNU potenziert INU bei Reziprozität
    \item Kohärenzdruck: Inkonsistenz zwischen Komponenten erzeugt Unbehagen
\end{itemize}

\textbf{BEATRIX-Relevanz:}
BEATRIX verwendet nicht-lineare Kombinationsformeln. Das LLM wird instruiert, auch ``unerwartete'' Entscheidungen zu generieren, wenn Komponentenkonflikte vorliegen.

\textbf{Konsequenz bei Fehlen:}
Einfache Addition würde systematisch falsche Vorhersagen generieren. Identitätsbasiertes Verhalten (gegen Eigeninteresse) wäre nicht modellierbar.

\textbf{Validiert:} META, WAX

\textbf{Abhängigkeiten:} MA-ONT-07

\textbf{Literatur:} Arrow (1951); Sen (1977)

%==============================================================================
% MA-ONT-14: Kooperations-Regulator λ
%==============================================================================

\subsubsection*{MA-ONT-14: Kooperations-Regulator $\lambda$}

\textbf{Statement:} $\lambda$ reguliert die lokale vs. globale Orientierung.

\textbf{Formel:}
\begin{equation}
U_{\text{eff}} = (1-\lambda) \cdot U_{\text{lok}} + \lambda \cdot \kappa(U_{\text{lok}}, U_{\text{komp}})
\end{equation}

\textbf{Beschreibung:}
$\lambda$ ist der zentrale Parameter für soziale Präferenzen:
\begin{itemize}
    \item $\lambda = 0$: Pure Eigennutz-Maximierung
    \item $\lambda = 0.5$: Balancierte Berücksichtigung
    \item $\lambda = 1$: Reine Kollektivorientierung
\end{itemize}
$\lambda$ ist nicht fix, sondern kontextabhängig (MA-ONT-15).

\textbf{BCM-Relevanz:}
$\lambda$ ist DER Parameter für individuelle Unterschiede in Kooperationsbereitschaft. Er klassifiziert Verhaltenstypen: Self-Interested ($\lambda < 0.3$), Conditional Cooperators ($\lambda \approx 0.5$), Altruisten ($\lambda > 0.7$).

\textbf{BEATRIX-Relevanz:}
$\lambda$-Schätzung ist eine Kernaufgabe. Das LLM analysiert:
\begin{itemize}
    \item Hinweise auf Kooperationsbereitschaft in Persona
    \item Gruppenbezug und Kollektividentität
    \item Historisches Kooperationsverhalten
\end{itemize}

\textbf{Konsequenz bei Fehlen:}
Keine Differenzierung von Kooperationstypen. Alle Menschen würden gleich modelliert. Die empirisch robuste Heterogenität (CC vs. SI vs. ALT) wäre nicht abbildbar.

\textbf{Validiert:} WAX

\textbf{Abhängigkeiten:} MA-ONT-12, MA-ONT-13

\textbf{Literatur:} Fehr \& Schmidt (1999); Fischbacher et al. (2001)

%==============================================================================
% MA-ONT-15: Hybrid-Kontinuum
%==============================================================================

\subsubsection*{MA-ONT-15: Hybrid-Kontinuum}

\textbf{Statement:} $\lambda \in [0, 1]$ und ist kontextabhängig.

\textbf{Formel:}
\begin{equation}
\lambda = f(\rho, \sigma, \tau)
\end{equation}

\textbf{Beschreibung:}
$\lambda$ ist keine Persönlichkeitskonstante. Dieselbe Person kann in verschiedenen Situationen unterschiedliche $\lambda$-Werte zeigen:
\begin{itemize}
    \item Anonymität senkt $\lambda$
    \item Beobachtung erhöht $\lambda$
    \item In-Group-Kontext erhöht $\lambda$
    \item Stress kann $\lambda$ senken oder erhöhen
\end{itemize}

\textbf{BCM-Relevanz:}
Überwindet die Trait-vs-State-Debatte. $\lambda$ hat eine Trait-Komponente (Basis-Kooperativität) und eine State-Komponente (situative Modulation). Das LAMBDA-Modul modelliert diese Interaktion.

\textbf{BEATRIX-Relevanz:}
BEATRIX schätzt kontextspezifisches $\lambda$:
\begin{itemize}
    \item Basis-$\lambda$ aus Persona-Traits
    \item Kontext-Adjustment aus Situationsbeschreibung
    \item Finales $\lambda = \lambda_{\text{basis}} + \Delta\lambda(\rho)$
\end{itemize}

\textbf{Konsequenz bei Fehlen:}
Fixes $\lambda$ würde situative Variabilität ignorieren. Kontextinterventionen (Beobachtung, Feedback) hätten keine Wirkung.

\textbf{Validiert:} WAX

\textbf{Abhängigkeiten:} MA-ONT-14

\textbf{Literatur:} Haley \& Fessler (2005); List (2006)

%==============================================================================
% MA-ONT-16: Gewichtungsfunktion ω
%==============================================================================

\subsubsection*{MA-ONT-16: Gewichtungsfunktion $\omega$}

\textbf{Statement:} FEPSDE-Dimensionen haben Gewichte.

\textbf{Formel:}
\begin{equation}
\omega: D \rightarrow [0, 1], \quad \sum_d \omega_d = 1
\end{equation}

\textbf{Beschreibung:}
Nicht alle Nutzendimensionen sind gleich wichtig. Für manche dominiert Financial (Unternehmer), für andere Emotional (Künstler) oder Ecological (Aktivist). Die Gewichte sind personenspezifisch und können sich über Zeit ändern.

\textbf{BCM-Relevanz:}
Personalisierung des INU-Moduls. Zwei Personen können bei gleichem Verhalten verschiedene Nutzen empfinden, weil sie die Dimensionen unterschiedlich gewichten.

\textbf{BEATRIX-Relevanz:}
BEATRIX schätzt den $\omega$-Vektor als 6 Werte, die sich zu 1 summieren. Das LLM analysiert Wertvorstellungen, Lebensstil und explizite Prioritäten in der Persona.

\textbf{Konsequenz bei Fehlen:}
Gleiche Gewichtung für alle würde individuelle Unterschiede ignorieren. Interventionen könnten nicht personalisiert werden.

\textbf{Validiert:} INU

\textbf{Abhängigkeiten:} MA-ONT-08

\textbf{Literatur:} Schwartz (1992); Inglehart (1997)

%==============================================================================
% MA-ONT-17: Identitäts-Aktivierung π
%==============================================================================

\subsubsection*{MA-ONT-17: Identitäts-Aktivierung $\pi$}

\textbf{Statement:} Identitäten haben kontextabhängige Aktivierungswahrscheinlichkeiten.

\textbf{Formel:}
\begin{equation}
\pi: I \times \rho \rightarrow [0, 1], \quad \sum_j \pi_j = 1
\end{equation}

\textbf{Beschreibung:}
Nicht alle Identitäten sind gleichzeitig aktiv. Der Kontext bestimmt, welche Identität salient ist: Im Büro ``Manager'', zu Hause ``Vater'', im Verein ``Trainer''. Die aktivierte Identität bestimmt relevante Normen und Präskriptionen.

\textbf{BCM-Relevanz:}
Erklärt situative Verhaltensvariation ohne Persönlichkeitsänderung. Die Person bleibt gleich, aber verschiedene Identitäten mit verschiedenen Normen werden aktiviert.

\textbf{BEATRIX-Relevanz:}
BEATRIX berechnet $\pi_j$ für jede Identität gegeben den Kontext:
\begin{itemize}
    \item Welche Identitäts-Cues enthält die Situation?
    \item Welche Identität ist chronisch salient?
    \item Priming-Effekte durch vorherige Interaktionen?
\end{itemize}

\textbf{Konsequenz bei Fehlen:}
Alle Identitäten immer gleich aktiv -- unrealistisch. Identitätsbasierte Interventionen (Priming) wären nicht modellierbar.

\textbf{Validiert:} IDN

\textbf{Abhängigkeiten:} MA-ONT-09

\textbf{Literatur:} Stryker \& Burke (2000); Oyserman (2009)

%==============================================================================
% MA-ONT-18: Alignment-Funktion A
%==============================================================================

\subsubsection*{MA-ONT-18: Alignment-Funktion $A$}

\textbf{Statement:} Verhalten hat Alignment mit Identität.

\textbf{Formel:}
\begin{equation}
A: \mathcal{V} \times I \rightarrow [-1, 1]
\end{equation}

\textbf{Beschreibung:}
Jedes Verhalten passt mehr oder weniger zur aktivierten Identität:
\begin{itemize}
    \item $A = +1$: Perfekt kongruent (``Das bin ich!'')
    \item $A = 0$: Neutral (identitätsirrelevant)
    \item $A = -1$: Diametral entgegengesetzt (Identitätsverletzung)
\end{itemize}
$U_{\text{IDN}}$ ist proportional zu $A$.

\textbf{BCM-Relevanz:}
Operationalisiert Identitätsnutzen. Identitätskongruentes Verhalten gibt positiven Nutzen, inkongruentes negativen. Dies kann INU-Nutzen übertrumpfen (Identitätsveto).

\textbf{BEATRIX-Relevanz:}
BEATRIX berechnet $A(V, I_j)$ für jede Verhaltensalternative und aktivierte Identität. Das LLM fragt: ``Wie würde ein typischer [Identität] handeln?''

\textbf{Konsequenz bei Fehlen:}
Identität wäre nur Label ohne Verhaltenskonsequenz. Identitätsbasierte Motivation (``Ich tue das, weil ich Umweltschützer bin'') wäre nicht modellierbar.

\textbf{Validiert:} IDN

\textbf{Abhängigkeiten:} MA-ONT-09, MA-ONT-02

\textbf{Literatur:} Akerlof \& Kranton (2000); Oyserman (2015)

%==============================================================================
% MA-ONT-19: Zeitabhängigkeit des Nutzens
%==============================================================================

\subsubsection*{MA-ONT-19: Zeitabhängigkeit des Nutzens}

\textbf{Statement:} Nutzen verändert sich über Zeit.

\textbf{Formel:}
\begin{equation}
U(\tau + \Delta\tau) = U(\tau) + \frac{\partial U}{\partial \tau} \cdot \Delta\tau + \varepsilon
\end{equation}

\textbf{Beschreibung:}
Präferenzen sind nicht statisch. Was heute wichtig ist, kann morgen irrelevant sein. Verändernde Faktoren: Alter, Lebensereignisse, gesellschaftlicher Wandel, Erfahrungen. Die Veränderung hat systematische und stochastische Komponenten.

\textbf{BCM-Relevanz:}
Ermöglicht dynamische Modellierung. Erklärt Präferenzwandel, Werteverschiebung und life-stage Effekte. Ohne dies wäre das BCM ein statisches Snapshot-Modell.

\textbf{BEATRIX-Relevanz:}
BEATRIX berücksichtigt zeitliche Entwicklung:
\begin{itemize}
    \item Lebensphase der Persona
    \item Präferenzstabilität vs. -wandel
    \item Erwartete zukünftige Verschiebungen
\end{itemize}

\textbf{Konsequenz bei Fehlen:}
Keine Modellierung von Präferenzwandel. Langfristige Vorhersagen wären statisch extrapoliert. Life-Cycle-Effekte ignoriert.

\textbf{Validiert:} WTX

\textbf{Abhängigkeiten:} MA-ONT-03, MA-ONT-05

\textbf{Literatur:} Becker \& Murphy (1988); Loewenstein et al. (2003)

%==============================================================================
% MA-ONT-20: Habituation
%==============================================================================

\subsubsection*{MA-ONT-20: Habituation}

\textbf{Statement:} Wiederholung reduziert Nutzen.

\textbf{Formel:}
\begin{equation}
\frac{\partial U}{\partial n} < 0 \quad \text{für } n \rightarrow \infty
\end{equation}

\textbf{Beschreibung:}
Der erste Bissen schmeckt am besten. Wiederholte Exposition führt zu abnehmendem Grenznutzen. Dies gilt für Konsum, aber auch für Interventionen: Nudges verlieren Wirkung bei Wiederholung.

\textbf{BCM-Relevanz:}
Erklärt Variety Seeking und die Notwendigkeit von Interventions-Variation. Ein statischer Nudge wird habituiert und verliert Wirkung. Das DYN-Modul modelliert diese Adaptation.

\textbf{BEATRIX-Relevanz:}
BEATRIX berücksichtigt Habituationseffekte:
\begin{itemize}
    \item Wie oft wurde das Verhalten bereits gezeigt?
    \item Wie lange wurde die Intervention bereits eingesetzt?
    \item Wann ist Interventions-Refresh nötig?
\end{itemize}

\textbf{Konsequenz bei Fehlen:}
Überschätzung der Langzeitwirkung von Interventionen. Keine Erklärung für Novelty Seeking. Statische Nutzenmodelle würden dominieren.

\textbf{Validiert:} INU, WTX

\textbf{Abhängigkeiten:} MA-ONT-19

\textbf{Literatur:} McSweeney \& Swindell (1999); Frederick \& Loewenstein (1999)

%==============================================================================
% MA-ONT-21: Feedback-Schleifen
%==============================================================================

\subsubsection*{MA-ONT-21: Feedback-Schleifen}

\textbf{Statement:} Verhalten beeinflusst zukünftigen Nutzen.

\textbf{Formel:}
\begin{equation}
U(\tau + 1) = f(U(\tau), V(\tau))
\end{equation}

\textbf{Beschreibung:}
Verhalten hat Konsequenzen, die zukünftige Präferenzen beeinflussen:
\begin{itemize}
    \item Positive Erfahrung → Präferenzstärkung
    \item Negative Erfahrung → Aversion
    \item Gewohnheitsbildung durch Wiederholung
\end{itemize}
Das System ist rekursiv: Entscheidungen heute beeinflussen den Entscheider von morgen.

\textbf{BCM-Relevanz:}
Ermöglicht die Modellierung von Lerneffekten und Pfadabhängigkeiten. Interventionen können langfristige Verhaltensänderung auslösen, indem sie Feedback-Loops anstoßen.

\textbf{BEATRIX-Relevanz:}
BEATRIX modelliert Feedback-Effekte für Interventionsplanung:
\begin{itemize}
    \item Welche Verstärkungseffekte sind zu erwarten?
    \item Wie kann positives Feedback designt werden?
    \item Wann drohen Abwärtsspiralen?
\end{itemize}

\textbf{Konsequenz bei Fehlen:}
Keine Lerneffekte modellierbar. Verhalten wäre nur vom aktuellen Zustand abhängig. Dynamische Interventionsstrategien wären nicht planbar.

\textbf{Validiert:} WTX

\textbf{Abhängigkeiten:} MA-ONT-19

\textbf{Literatur:} Herrnstein (1997); Staddon \& Cerutti (2003)

%==============================================================================
% MA-ONT-22: Spillover
%==============================================================================

\subsubsection*{MA-ONT-22: Spillover}

\textbf{Statement:} Verhalten in einem Bereich beeinflusst andere.

\textbf{Formel:}
\begin{equation}
\frac{\partial U_j}{\partial V_i} \neq 0 \quad \text{für } i \neq j
\end{equation}

\textbf{Beschreibung:}
Verhaltensänderungen sind nicht isoliert:
\begin{itemize}
    \item Positiver Spillover: Recycling → andere Umweltverhaltensweisen
    \item Negativer Spillover (Moral Licensing): ``Ich habe ja schon...''
    \item Cross-Domain: Sport → bessere Ernährung
\end{itemize}

\textbf{BCM-Relevanz:}
Erklärt systemische Effekte von Interventionen. Eine gezielte Intervention kann unbeabsichtigte (positive oder negative) Effekte in anderen Bereichen haben. Das BCM modelliert diese Interdependenzen.

\textbf{BEATRIX-Relevanz:}
BEATRIX prüft auf Spillover-Potenzial:
\begin{itemize}
    \item Welche anderen Verhaltensweisen könnten beeinflusst werden?
    \item Risiko von Moral Licensing?
    \item Synergien mit anderen Interventionen?
\end{itemize}

\textbf{Konsequenz bei Fehlen:}
Isolierte Betrachtung von Verhaltensweisen. Systemische Effekte und unbeabsichtigte Konsequenzen würden übersehen.

\textbf{Validiert:} INU, KNU

\textbf{Abhängigkeiten:} MA-ONT-21

\textbf{Literatur:} Thøgersen \& Ölander (2003); Truelove et al. (2014)

%==============================================================================
% MA-ONT-23: Identitäten evolvieren
%==============================================================================

\subsubsection*{MA-ONT-23: Identitäten evolvieren}

\textbf{Statement:} Die Identitätslandschaft verändert sich über Zeit.

\textbf{Formel:}
\begin{equation}
I(\tau + \Delta\tau) \neq I(\tau)
\end{equation}

\textbf{Beschreibung:}
Identitäten sind nicht fix:
\begin{itemize}
    \item Neue Identitäten entstehen (``Ich bin jetzt Veganer'')
    \item Alte Identitäten verblassen (``Früher war ich Raucher'')
    \item Hierarchien verschieben sich (Karriere → Familie)
\end{itemize}
Lebensereignisse (Heirat, Kind, Jobwechsel) katalysieren Identitätswandel.

\textbf{BCM-Relevanz:}
Ermöglicht die Modellierung von Identitätsinterventionen. IDENTITY-Interventionen zielen auf die Modifikation der Identitätslandschaft, nicht nur auf Aktivierung bestehender Identitäten.

\textbf{BEATRIX-Relevanz:}
BEATRIX kann Identitätsentwicklung modellieren:
\begin{itemize}
    \item Welche Identitäten sind im Aufbau?
    \item Welche neigen zum Verblassen?
    \item Wie kann eine gewünschte Identität gestärkt werden?
\end{itemize}

\textbf{Konsequenz bei Fehlen:}
Statische Identität würde Lebenslauf-Veränderungen nicht abbilden. IDENTITY-Interventionen wären auf Priming beschränkt.

\textbf{Validiert:} IDN

\textbf{Abhängigkeiten:} MA-ONT-09, MA-ONT-05

\textbf{Literatur:} McAdams (2001); Ibarra \& Barbulescu (2010)

%==============================================================================
% MA-ONT-24: Normen evolvieren
%==============================================================================

\subsubsection*{MA-ONT-24: Normen evolvieren}

\textbf{Statement:} Soziale Normen verändern sich über Zeit.

\textbf{Formel:}
\begin{equation}
N(\tau + \Delta\tau) = N(\tau) + \mu_N \cdot \Delta N
\end{equation}

\textbf{Beschreibung:}
Was einmal normal war, kann inakzeptabel werden (Rauchen im Büro) und umgekehrt (gleichgeschlechtliche Ehe). Normenwandel kann langsam (Generationenwechsel) oder schnell (Tipping Points) sein. Normentrepreneure treiben Wandel voran.

\textbf{BCM-Relevanz:}
Ermöglicht die Modellierung von gesellschaftlichem Wandel. SOCIAL-Interventionen können auf Normenwandel zielen, nicht nur auf Konformität mit bestehenden Normen.

\textbf{BEATRIX-Relevanz:}
BEATRIX berücksichtigt Normendynamik:
\begin{itemize}
    \item In welche Richtung entwickeln sich relevante Normen?
    \item Wie resilient sind aktuelle Normen?
    \item Gibt es emergente Normen?
\end{itemize}

\textbf{Konsequenz bei Fehlen:}
Statische Normen würden gesellschaftlichen Wandel nicht abbilden. Langfristige Interventionsplanung wäre ahistorisch.

\textbf{Validiert:} KNU

\textbf{Abhängigkeiten:} MA-ONT-11, MA-ONT-05

\textbf{Literatur:} Centola et al. (2018); Bicchieri (2016)

%==============================================================================
% MA-ONT-25: Meta-Evolution
%==============================================================================

\subsubsection*{MA-ONT-25: Meta-Evolution}

\textbf{Statement:} Die Regeln der Veränderung verändern sich selbst.

\textbf{Formel:}
\begin{equation}
\mu(\tau + \Delta\tau) \neq \mu(\tau)
\end{equation}

\textbf{Beschreibung:}
Nicht nur Parameter ändern sich, sondern auch die Änderungsraten:
\begin{itemize}
    \item Beschleunigter Wandel in digitalen Kontexten
    \item Verlangsamung durch Institutionalisierung
    \item Phasenübergänge in der Veränderungsdynamik
\end{itemize}
Das System lernt zu lernen (Meta-Lernen).

\textbf{BCM-Relevanz:}
Höchstes Abstraktionsniveau. Ermöglicht die Modellierung von Systemtransformationen, nicht nur Parameterveränderungen. Relevant für langfristige Prognosen und fundamentale Interventionen.

\textbf{BEATRIX-Relevanz:}
Für die meisten Anwendungen irrelevant (zu abstrakt). Relevant nur bei:
\begin{itemize}
    \item Sehr langfristigen Szenarien
    \item Fundamentalen Systemtransformationen
    \item Meta-Interventionen (``Wie ändere ich die Änderungsrate?'')
\end{itemize}

\textbf{Konsequenz bei Fehlen:}
Statische Änderungsraten würden Phasenübergänge verpassen. Beschleunigungs- und Verlangsamungsmuster wären nicht modellierbar.

\textbf{Validiert:} META

\textbf{Abhängigkeiten:} MA-ONT-19

\textbf{Literatur:} Kauffman (1993); Holland (1995)

%==============================================================================
% MA-ONT-26: Bereits in axiom_beispiele.tex
%==============================================================================

\bigskip
\noindent\rule{\textwidth}{0.4pt}
\textit{Ende der ONT-Axiome. Alle 26 Ontologie-Axiome sind nun vollständig dokumentiert.}

% =============================================================================
% EPI-AXIOME (EPISTEMOLOGIE) - Vollständige Ausformulierung
% BCM 2.2 META-Modul
% =============================================================================

\section{Epistemologie-Axiome (MA-EPI)}

Die EPI-Axiome definieren die Grenzen des Wissens, Biases und kognitive Prozesse.
Sie beantworten die Frage: \textit{Was wissen wir?}

%==============================================================================
% MA-EPI-01: Unvollständiges Wissen
%==============================================================================

\subsubsection*{MA-EPI-01: Unvollständiges Wissen}

\textbf{Statement:} Wissen über Nutzen ist immer unvollständig.

\textbf{Formel:}
\begin{equation}
K(U) \subset U
\end{equation}

\textbf{Beschreibung:}
Menschen kennen nicht alle ihre Präferenzen, nicht alle Konsequenzen ihrer Handlungen und nicht alle verfügbaren Optionen. Das subjektive Wissen $K(U)$ ist eine echte Teilmenge des objektiven Nutzens $U$. Diese Lücke ist fundamental, nicht nur praktisch.

\textbf{BCM-Relevanz:}
Epistemologisches Fundament für alle EPI-Axiome. Ohne dieses Axiom wäre vollständige Rationalität möglich. Alle kognitiven Biases sind Konsequenzen dieser fundamentalen Wissenslücke.

\textbf{BEATRIX-Relevanz:}
BEATRIX arbeitet mit geschätzten, nicht wahren Parametern:
\begin{itemize}
    \item LLM-Schätzungen sind Approximationen
    \item Unsicherheitsintervalle müssen kommuniziert werden
    \item Persona-Beschreibungen sind unvollständig
\end{itemize}

\textbf{Konsequenz bei Fehlen:}
Perfektes Wissen würde perfekte Entscheidungen ermöglichen. Informationsinterventionen wären trivial. Die gesamte Verhaltensökonomie wäre irrelevant.

\textbf{Validiert:} AWX

\textbf{Abhängigkeiten:} MA-ONT-03 (Nutzen existiert)

\textbf{Literatur:} Knight (1921); Hayek (1945)

%==============================================================================
% MA-EPI-02: Bereits in axiom_beispiele.tex
%==============================================================================

%==============================================================================
% MA-EPI-03: Irreduzible Unsicherheit
%==============================================================================

\subsubsection*{MA-EPI-03: Irreduzible Unsicherheit}

\textbf{Statement:} Es gibt fundamentale, nicht eliminierbare Unsicherheit.

\textbf{Formel:}
\begin{equation}
\exists \varepsilon: P(\varepsilon) > 0 \quad \forall \text{ Modelle}
\end{equation}

\textbf{Beschreibung:}
Selbst mit unbegrenzten Ressourcen bleibt Unsicherheit:
\begin{itemize}
    \item Ontologische Unsicherheit: Die Zukunft ist nicht determiniert
    \item Epistemische Unsicherheit: Grenzen der Messbarkeit
    \item Modell-Unsicherheit: Jedes Modell ist vereinfachend
\end{itemize}
Diese Unsicherheit ist nicht ``noch nicht gelöst'', sondern prinzipiell irreduzibel.

\textbf{BCM-Relevanz:}
Begründet Demut gegenüber Vorhersagen. Das BCM kann Wahrscheinlichkeiten angeben, nie Gewissheit. Übermäßiges Vertrauen in Modelle wird strukturell verhindert.

\textbf{BEATRIX-Relevanz:}
BEATRIX-Outputs sind probabilistisch:
\begin{itemize}
    \item Keine deterministischen Vorhersagen
    \item Konfidenzintervalle immer angeben
    \item ``Unentscheidbar'' als legitimes Ergebnis
\end{itemize}

\textbf{Konsequenz bei Fehlen:}
Hybris der Vorhersagbarkeit. Modelle würden als deterministische Wahrheiten behandelt. Risiken würden unterschätzt.

\textbf{Validiert:} META

\textbf{Abhängigkeiten:} -- (fundamentales epistemisches Axiom)

\textbf{Literatur:} Knight (1921); Taleb (2007)

%==============================================================================
% MA-EPI-04: Adaptive Heuristiken
%==============================================================================

\subsubsection*{MA-EPI-04: Adaptive Heuristiken}

\textbf{Statement:} Heuristiken können ökologisch rational sein.

\textbf{Formel:}
\begin{equation}
E[U|\text{Heuristik}] \geq E[U|\text{Optimierung}] \quad \text{in } \rho
\end{equation}

\textbf{Beschreibung:}
Heuristiken sind nicht per se Fehler. In bestimmten Umgebungen (hohe Unsicherheit, Zeitdruck) können einfache Regeln komplexe Algorithmen übertreffen. ``Fast and frugal'' Heuristiken nutzen die Struktur der Umwelt aus.

\textbf{BCM-Relevanz:}
Differenziert zwischen adaptiven und maladaptiven Abweichungen von Rationalität. Nicht jede Vereinfachung ist ein Bias. BOOST-Interventionen stärken adaptive Heuristiken.

\textbf{BEATRIX-Relevanz:}
BEATRIX unterscheidet:
\begin{itemize}
    \item Adaptive Heuristik → Nicht korrigieren
    \item Maladaptiver Bias → Intervention möglich
    \item Kontextabhängigkeit beachten
\end{itemize}

\textbf{Konsequenz bei Fehlen:}
Alle Heuristiken wären Fehler. Debiasing wäre immer angezeigt. Die ökologische Intelligenz einfacher Regeln würde übersehen.

\textbf{Validiert:} INU

\textbf{Abhängigkeiten:} MA-EPI-02

\textbf{Literatur:} Gigerenzer et al. (1999); Todd \& Gigerenzer (2012)

%==============================================================================
% MA-EPI-05: Zwei Systeme
%==============================================================================

\subsubsection*{MA-EPI-05: Zwei-Systeme-Kognition}

\textbf{Statement:} Kognition operiert in zwei Modi: System 1 (schnell, automatisch) und System 2 (langsam, deliberativ).

\textbf{Formel:}
\begin{equation}
K = K_{S1} \cup K_{S2}
\end{equation}

\textbf{Beschreibung:}
System 1: Intuitiv, assoziativ, parallelverarbeitend, mühelos, emotional getrieben.
System 2: Analytisch, regelbasiert, seriell, anstrengend, rational kontrolliert.
Die meisten Alltagsentscheidungen sind System-1-getrieben. System 2 kann aktiviert werden, ist aber ressourcenbeschränkt.

\textbf{BCM-Relevanz:}
Fundamentale Unterscheidung für Interventionsdesign:
\begin{itemize}
    \item NUDGE: Wirkt auf System 1
    \item BOOST: Stärkt System 1 oder System 2
    \item THINK: Aktiviert System 2
\end{itemize}

\textbf{BEATRIX-Relevanz:}
BEATRIX bewertet Entscheidungsmodus:
\begin{itemize}
    \item Welches System dominiert bei dieser Persona/Situation?
    \item Interventionen auf richtigen Modus abstimmen
    \item Cognitive Load beachten
\end{itemize}

\textbf{Konsequenz bei Fehlen:}
Einheitliches Kognitionsmodell. Unterschied zwischen automatischen und deliberativen Entscheidungen nicht modellierbar. Interventionstypen nicht differenzierbar.

\textbf{Validiert:} AWX, WAX

\textbf{Abhängigkeiten:} MA-EPI-02

\textbf{Literatur:} Kahneman (2011); Stanovich \& West (2000)

%==============================================================================
% MA-EPI-06: Systematische Biases
%==============================================================================

\subsubsection*{MA-EPI-06: Systematische Biases}

\textbf{Statement:} Kognitive Verzerrungen sind systematisch und vorhersagbar.

\textbf{Formel:}
\begin{equation}
E[\text{Bias}] \neq 0
\end{equation}

\textbf{Beschreibung:}
Biases sind nicht zufällige Fehler, sondern haben Richtung:
\begin{itemize}
    \item Überschätzung (Overconfidence)
    \item Unterschätzung (Neglect)
    \item Systematische Verzerrung (Framing)
\end{itemize}
Weil sie systematisch sind, sind sie vorhersagbar und potenziell korrigierbar.

\textbf{BCM-Relevanz:}
Fundament für alle spezifischen Bias-Axiome (EPI-12 bis EPI-37). Biases sind keine ``Fehler'', sondern Eigenschaften des kognitiven Systems. Das BCM integriert sie als Parameter, nicht als Störterme.

\textbf{BEATRIX-Relevanz:}
BEATRIX modelliert Biases explizit:
\begin{itemize}
    \item Bias-Anfälligkeit als Persona-Parameter
    \item Kontextuelle Bias-Aktivierung
    \item Debiasing-Potenzial bewerten
\end{itemize}

\textbf{Konsequenz bei Fehlen:}
Biases wären zufällig und nicht modellierbar. Nudges und Debiasing-Strategien hätten keine theoretische Grundlage.

\textbf{Validiert:} INU, AWX

\textbf{Abhängigkeiten:} MA-EPI-05

\textbf{Literatur:} Kahneman \& Tversky (1974); Ariely (2008)

%==============================================================================
% MA-EPI-07: AWX (Awareness)
%==============================================================================

\subsubsection*{MA-EPI-07: AWX (Awareness)}

\textbf{Statement:} Bewusstsein über Optionen und Konsequenzen variiert.

\textbf{Formel:}
\begin{equation}
AWX = f(K_V, K_C)
\end{equation}

\textbf{Beschreibung:}
AWX (Awareness) ist die erste Stufe im Verhaltenstrichter:
\begin{itemize}
    \item $K_V$: Wissen über Verhaltensoptionen (Was kann ich tun?)
    \item $K_C$: Wissen über Konsequenzen (Was passiert dann?)
\end{itemize}
Ohne Awareness keine Entscheidung. Viele Verhaltensdefizite sind Awareness-Defizite.

\textbf{BCM-Relevanz:}
Erste Variable im AWX → WAX → WTX → V Trichter. Interventionen können auf jeder Stufe ansetzen. Information und Bildung erhöhen AWX.

\textbf{BEATRIX-Relevanz:}
BEATRIX schätzt AWX-Level:
\begin{itemize}
    \item Wie informiert ist die Persona?
    \item Welche Optionen sind ihr bekannt?
    \item Werden Konsequenzen korrekt antizipiert?
\end{itemize}

\textbf{Konsequenz bei Fehlen:}
Kein Unterschied zwischen ``nicht wissen'' und ``nicht wollen''. Informationsinterventionen wären nicht gezielt einsetzbar.

\textbf{Validiert:} AWX

\textbf{Abhängigkeiten:} MA-EPI-01

\textbf{Literatur:} Prochaska \& DiClemente (1983)

%==============================================================================
% MA-EPI-08: WAX (Willingness)
%==============================================================================

\subsubsection*{MA-EPI-08: WAX (Willingness)}

\textbf{Statement:} Bereitschaft zu handeln ist eine eigenständige Variable.

\textbf{Formel:}
\begin{equation}
WAX = f(AWX, U, \Gamma)
\end{equation}

\textbf{Beschreibung:}
WAX (Willingness to Act) hängt ab von:
\begin{itemize}
    \item AWX: Muss wissen, dass Option existiert
    \item U: Muss Nutzen sehen
    \item $\Gamma$: Constraints müssen erfüllt sein
\end{itemize}
WAX ist Verhaltensintention, noch nicht Verhalten selbst.

\textbf{BCM-Relevanz:}
Zweite Stufe im Trichter. WAX kann hoch sein ohne Verhalten (Intention-Action Gap). Motivationsinterventionen erhöhen WAX.

\textbf{BEATRIX-Relevanz:}
BEATRIX bewertet WAX:
\begin{itemize}
    \item Ist die Person motiviert?
    \item Welche Barrieren senken WAX?
    \item Interventionen zur WAX-Steigerung identifizieren
\end{itemize}

\textbf{Konsequenz bei Fehlen:}
Motivation und Verhalten würden nicht unterschieden. Die wichtige Lücke zwischen ``wollen'' und ``tun'' wäre unsichtbar.

\textbf{Validiert:} WAX

\textbf{Abhängigkeiten:} MA-EPI-07, MA-ONT-03

\textbf{Literatur:} Ajzen (1991); Schwarzer (2008)

%==============================================================================
% MA-EPI-09: WTX (Transition)
%==============================================================================

\subsubsection*{MA-EPI-09: WTX (Willingness to Transition)}

\textbf{Statement:} Der Übergang von Absicht zu Verhalten ist probabilistisch.

\textbf{Formel:}
\begin{equation}
P(V|WAX) = WTX
\end{equation}

\textbf{Beschreibung:}
WTX modelliert die ``letzte Meile'': Von der Intention zur Handlung. Selbst bei hohem WAX kann WTX niedrig sein:
\begin{itemize}
    \item Prokrastination
    \item Selbstkontrollversagen
    \item Situative Barrieren
\end{itemize}

\textbf{BCM-Relevanz:}
Dritte Stufe im Trichter. WTX ist der kritische Übergangspunkt. Commitment Devices (MA-DYN-09) und Implementation Intentions (MA-DYN-18) erhöhen WTX.

\textbf{BEATRIX-Relevanz:}
BEATRIX schätzt WTX:
\begin{itemize}
    \item Wie wahrscheinlich ist die Umsetzung?
    \item Welche Transitonsbarrieren existieren?
    \item Braucht es einen ``letzten Schubs''?
\end{itemize}

\textbf{Konsequenz bei Fehlen:}
Intention würde mit Verhalten gleichgesetzt. Die Intention-Action Gap (MA-EPI-11) wäre nicht modellierbar.

\textbf{Validiert:} WTX

\textbf{Abhängigkeiten:} MA-EPI-08

\textbf{Literatur:} Sheeran (2002); Gollwitzer \& Sheeran (2006)

%==============================================================================
% MA-EPI-10: AWX-WAX-WTX Sequenz
%==============================================================================

\subsubsection*{MA-EPI-10: Verhaltenssequenz AWX → WAX → WTX → V}

\textbf{Statement:} Verhalten folgt einer sequenziellen Kette.

\textbf{Formel:}
\begin{equation}
V = WTX(WAX(AWX(K)))
\end{equation}

\textbf{Beschreibung:}
Die Kette ist kausal geordnet:
\begin{enumerate}
    \item Ohne AWX kein WAX (Was ich nicht kenne, kann ich nicht wollen)
    \item Ohne WAX kein WTX (Ohne Motivation keine Handlung)
    \item Ohne WTX kein V (Ohne Umsetzung kein Verhalten)
\end{enumerate}
Jede Stufe ist notwendig, keine ist hinreichend.

\textbf{BCM-Relevanz:}
Ermöglicht Diagnose: Wo liegt das Problem? Interventionen müssen auf die richtige Stufe zielen:
\begin{itemize}
    \item AWX-Defizit → Information
    \item WAX-Defizit → Motivation
    \item WTX-Defizit → Facilitation
\end{itemize}

\textbf{BEATRIX-Relevanz:}
BEATRIX identifiziert den Engpass:
\begin{itemize}
    \item Welche Stufe ist der Bottleneck?
    \item Interventionen entsprechend priorisieren
    \item Kaskadeneffekte modellieren
\end{itemize}

\textbf{Konsequenz bei Fehlen:}
Alle Verhaltensdefizite würden gleich behandelt. Intervention ohne Diagnose.

\textbf{Validiert:} AWX, WAX, WTX

\textbf{Abhängigkeiten:} MA-EPI-07, MA-EPI-08, MA-EPI-09

\textbf{Literatur:} Michie et al. (2011)

%==============================================================================
% MA-EPI-11: Intention-Action Gap
%==============================================================================

\subsubsection*{MA-EPI-11: Intention-Action Gap}

\textbf{Statement:} Absicht führt nicht immer zu Verhalten.

\textbf{Formel:}
\begin{equation}
P(V|WAX=1) < 1
\end{equation}

\textbf{Beschreibung:}
Selbst bei maximaler Motivation wird das Verhalten nicht immer umgesetzt:
\begin{itemize}
    \item 47\% der Intentions werden nicht umgesetzt (Sheeran, 2002)
    \item Gründe: Vergessen, Aufschub, situative Hindernisse
    \item Chronische Gap bei Gesundheits- und Umweltverhalten
\end{itemize}

\textbf{BCM-Relevanz:}
Zentral für WTX-Forschung. Die Gap ist nicht Motivationsdefizit, sondern Ausführungsdefizit. Andere Interventionen nötig: Commitment, Defaults, Erinnerungen.

\textbf{BEATRIX-Relevanz:}
BEATRIX quantifiziert die Gap:
\begin{itemize}
    \item Wie groß ist die typische Gap bei dieser Persona/Verhalten?
    \item Welche Interventionen schließen die Gap?
    \item Modelliert als $P(V|WAX)$
\end{itemize}

\textbf{Konsequenz bei Fehlen:}
$WAX = 1$ würde $V = 1$ implizieren. Umsetzungsprobleme wären unsichtbar.

\textbf{Validiert:} WTX

\textbf{Abhängigkeiten:} MA-EPI-09

\textbf{Literatur:} Sheeran (2002); Webb \& Sheeran (2006)

%==============================================================================
% MA-EPI-12: Verfügbarkeitsheuristik
%==============================================================================

\subsubsection*{MA-EPI-12: Verfügbarkeitsheuristik}

\textbf{Statement:} Leicht Abrufbares wird als wahrscheinlicher eingeschätzt.

\textbf{Formel:}
\begin{equation}
P_{\text{subj}}(E) = f(\text{Accessibility}(E))
\end{equation}

\textbf{Beschreibung:}
Die Verfügbarkeitsheuristik ersetzt ``Wie wahrscheinlich?'' durch ``Wie leicht fällt mir ein Beispiel ein?'':
\begin{itemize}
    \item Medial präsente Risiken (Terror) überschätzt
    \item Alltägliche Risiken (Autounfall) unterschätzt
    \item Persönliche Erfahrungen übergewichtet
\end{itemize}

\textbf{BCM-Relevanz:}
Erklärt systematische Risikofehleinschätzungen. Kann für Interventionen genutzt werden: Saliente Beispiele erhöhen wahrgenommene Wahrscheinlichkeit.

\textbf{BEATRIX-Relevanz:}
BEATRIX berücksichtigt Verfügbarkeit:
\begin{itemize}
    \item Welche Beispiele sind für die Persona salient?
    \item Welche Medienexposition?
    \item Können saliente Counter-Examples präsentiert werden?
\end{itemize}

\textbf{Konsequenz bei Fehlen:}
Wahrscheinlichkeitsurteile wären Bayes-rational. Media-Effekte auf Risikowahrnehmung unerklärlich.

\textbf{Validiert:} AWX

\textbf{Abhängigkeiten:} MA-EPI-05

\textbf{Literatur:} Tversky \& Kahneman (1973)

%==============================================================================
% MA-EPI-13: Framing-Effekte
%==============================================================================

\subsubsection*{MA-EPI-13: Framing-Effekte}

\textbf{Statement:} Die Darstellung beeinflusst die Entscheidung.

\textbf{Formel:}
\begin{equation}
U(V|\text{Frame}_A) \neq U(V|\text{Frame}_B)
\end{equation}

\textbf{Beschreibung:}
Dieselbe Information, unterschiedlich präsentiert, führt zu unterschiedlichen Entscheidungen:
\begin{itemize}
    \item ``90\% Überlebensrate'' vs. ``10\% Sterberate''
    \item ``500€ Rabatt'' vs. ``500€ weniger zahlen''
    \item Gewinn-Frame vs. Verlust-Frame
\end{itemize}

\textbf{BCM-Relevanz:}
Mächtiges Interventionsinstrument. NUDGE-Interventionen nutzen häufig Framing. Ethisch: Transparenz über Frame-Wahl nötig.

\textbf{BEATRIX-Relevanz:}
BEATRIX analysiert Frame-Sensitivität:
\begin{itemize}
    \item Wie Frame-sensitiv ist die Persona?
    \item Welcher Frame ist für das Zielverhalten optimal?
    \item Ethische Bewertung des Framings
\end{itemize}

\textbf{Konsequenz bei Fehlen:}
Präsentation wäre irrelevant. Nur Inhalt zählt (homo oeconomicus). Kommunikationsdesign hätte keinen Wert.

\textbf{Validiert:} INU, AWX

\textbf{Abhängigkeiten:} MA-EPI-06

\textbf{Literatur:} Tversky \& Kahneman (1981)

%==============================================================================
% MA-EPI-14: Anchoring
%==============================================================================

\subsubsection*{MA-EPI-14: Anchoring}

\textbf{Statement:} Anker beeinflussen numerische Schätzungen.

\textbf{Formel:}
\begin{equation}
E[X|\text{Anchor}] = f(\text{Anchor})
\end{equation}

\textbf{Beschreibung:}
Beliebige Zahlen (auch irrelevante) ziehen Schätzungen in ihre Richtung:
\begin{itemize}
    \item Preisverhandlungen: Erstes Angebot dominiert
    \item Spendenaufrufe: Vorgeschlagene Beträge wirken als Anker
    \item Selbst offensichtlich irrelevante Zahlen haben Wirkung
\end{itemize}

\textbf{BCM-Relevanz:}
Wird in NUDGE-Interventionen genutzt. Default-Werte sind Anker. Ethik: Anker sollten nicht manipulativ sein.

\textbf{BEATRIX-Relevanz:}
BEATRIX berücksichtigt Ankereffekte:
\begin{itemize}
    \item Welche Anker sind in der Situation präsent?
    \item Wie kann Anchoring für positive Verhaltensänderung genutzt werden?
    \item Debiasing durch multiple Anker
\end{itemize}

\textbf{Konsequenz bei Fehlen:}
Schätzungen wären unbeeinflussbar durch Kontext. Defaults und Vorschläge hätten keine Wirkung.

\textbf{Validiert:} INU

\textbf{Abhängigkeiten:} MA-EPI-06

\textbf{Literatur:} Tversky \& Kahneman (1974); Ariely et al. (2003)

%==============================================================================
% MA-EPI-15: Bereits in axiom_beispiele.tex
%==============================================================================

%==============================================================================
% MA-EPI-16: Present Bias
%==============================================================================

\subsubsection*{MA-EPI-16: Present Bias}

\textbf{Statement:} Die Gegenwart wird übergewichtet.

\textbf{Formel:}
\begin{equation}
\delta_{\text{now}} > \delta_{\text{future}}
\end{equation}

\textbf{Beschreibung:}
Unmittelbare Belohnungen werden unverhältnismäßig stark gewichtet:
\begin{itemize}
    \item ``Jetzt 50€'' wird ``in einer Woche 55€'' vorgezogen
    \item Aber: ``In einem Jahr 50€'' vs. ``in einem Jahr + einer Woche 55€'' -- hier wird gewartet
    \item Dies ist zeitinkonsistent (dynamische Inkonsistenz)
\end{itemize}

\textbf{BCM-Relevanz:}
Erklärt Prokrastination, Suchtverhalten und Spar-Defizite. Present Bias ist der Kern von MA-DYN-14 (hyperbolisches Discounting).

\textbf{BEATRIX-Relevanz:}
BEATRIX schätzt Present-Bias-Stärke:
\begin{itemize}
    \item Hinweise auf Impulsivität in Persona
    \item Interventionen mit sofortiger vs. verzögerter Belohnung
    \item Commitment Devices für Hochimpulsive
\end{itemize}

\textbf{Konsequenz bei Fehlen:}
Exponentielles Discounting (zeitkonsistent). Prokrastination wäre unerklärlich. Commitment Devices nutzlos.

\textbf{Validiert:} INU, WTX

\textbf{Abhängigkeiten:} MA-ONT-05

\textbf{Literatur:} Laibson (1997); O'Donoghue \& Rabin (1999)

%==============================================================================
% MA-EPI-17: Status Quo Bias
%==============================================================================

\subsubsection*{MA-EPI-17: Status Quo Bias}

\textbf{Statement:} Der aktuelle Zustand wird bevorzugt.

\textbf{Formel:}
\begin{equation}
U(\text{Status Quo}) + \text{Bias} > U(\text{Alternative})
\end{equation}

\textbf{Beschreibung:}
Menschen bleiben bei dem, was sie haben, selbst wenn Alternativen objektiv besser sind:
\begin{itemize}
    \item Nichtwechseln von Stromanbietern
    \item Beibehalten suboptimaler Voreinstellungen
    \item Festhalten an Investitionsentscheidungen
\end{itemize}

\textbf{BCM-Relevanz:}
Erklärt die Macht von Defaults. Wenn der Default ``Organspende'' ist, spenden mehr. Status Quo Bias + Loss Aversion verstärken sich.

\textbf{BEATRIX-Relevanz:}
BEATRIX nutzt Status Quo Bias:
\begin{itemize}
    \item Was ist der aktuelle Default?
    \item Kann der Default geändert werden?
    \item Bei SQ-Bias-anfälliger Persona: Opt-Out statt Opt-In
\end{itemize}

\textbf{Konsequenz bei Fehlen:}
Defaults wären irrelevant. Alle Optionen würden neutral bewertet.

\textbf{Validiert:} INU, WAX

\textbf{Abhängigkeiten:} MA-EPI-15

\textbf{Literatur:} Samuelson \& Zeckhauser (1988); Kahneman et al. (1991)

%==============================================================================
% MA-EPI-18: Overconfidence
%==============================================================================

\subsubsection*{MA-EPI-18: Overconfidence}

\textbf{Statement:} Eigene Fähigkeiten werden überschätzt.

\textbf{Formel:}
\begin{equation}
P_{\text{subj}}(\text{Success}) > P_{\text{obj}}(\text{Success})
\end{equation}

\textbf{Beschreibung:}
Die meisten Menschen glauben, überdurchschnittlich zu sein:
\begin{itemize}
    \item 93\% der US-Fahrer halten sich für ``überdurchschnittlich''
    \item Überschätzung eigener Prognosegenauigkeit
    \item ``Planning Fallacy'': Projekte dauern länger als geplant
\end{itemize}

\textbf{BCM-Relevanz:}
Overconfidence beeinflusst Risikobereitschaft und Planungsverhalten. Kann zu suboptimalen Entscheidungen führen (Unternehmensgründungen, Investitionen).

\textbf{BEATRIX-Relevanz:}
BEATRIX korrigiert für Overconfidence:
\begin{itemize}
    \item Persona mit hohem Overconfidence → Risikowarnung
    \item Zeitschätzungen nach oben korrigieren
    \item Objektive Benchmarks präsentieren
\end{itemize}

\textbf{Konsequenz bei Fehlen:}
Selbsteinschätzungen würden als akkurat behandelt. Unrealistische Pläne blieben unkorrigiert.

\textbf{Validiert:} INU, AWX

\textbf{Abhängigkeiten:} MA-EPI-06

\textbf{Literatur:} Moore \& Healy (2008); Buehler et al. (1994)

%==============================================================================
% MA-EPI-19: Lernen ist möglich
%==============================================================================

\subsubsection*{MA-EPI-19: Lernen ist möglich}

\textbf{Statement:} Wissen kann akkumuliert werden.

\textbf{Formel:}
\begin{equation}
K(\tau + \Delta\tau) \supseteq K(\tau)
\end{equation}

\textbf{Beschreibung:}
Trotz fundamentaler Wissensgrenzen (MA-EPI-01) kann Wissen erweitert werden:
\begin{itemize}
    \item Erfahrungslernen
    \item Instruktion und Bildung
    \item Feedback-basiertes Lernen
\end{itemize}
Die Wissensmenge ist monoton nicht-abnehmend (Vergessen ignoriert).

\textbf{BCM-Relevanz:}
Begründet BOOST- und THINK-Interventionen. Wenn Lernen unmöglich wäre, könnten nur Constraints und Nudges wirken.

\textbf{BEATRIX-Relevanz:}
BEATRIX modelliert Lernpotenzial:
\begin{itemize}
    \item Wie lernfähig ist die Persona?
    \item Welche Lerninterventionen sind effektiv?
    \item Erwartete Wissensakkumulation über Zeit
\end{itemize}

\textbf{Konsequenz bei Fehlen:}
Statisches Wissen. Bildungsinterventionen wären wirkungslos. Keine Kompetenzentwicklung modellierbar.

\textbf{Validiert:} AWX

\textbf{Abhängigkeiten:} MA-EPI-01, MA-ONT-05

\textbf{Literatur:} Anderson (1982); Ericsson \& Charness (1994)

%==============================================================================
% MA-EPI-20: Feedback verbessert Schätzungen
%==============================================================================

\subsubsection*{MA-EPI-20: Feedback verbessert Schätzungen}

\textbf{Statement:} Mit Feedback werden Prognosen akkurater.

\textbf{Formel:}
\begin{equation}
|K(\tau+1) - U| < |K(\tau) - U| \quad \text{mit Feedback}
\end{equation}

\textbf{Beschreibung:}
Feedback ermöglicht Kalibrierung:
\begin{itemize}
    \item Überprüfung von Vorhersagen an der Realität
    \item Korrektur systematischer Biases
    \item Verbesserung der Metakognition
\end{itemize}
Voraussetzung: Feedback muss zeitnah, spezifisch und glaubwürdig sein.

\textbf{BCM-Relevanz:}
Begründet Feedback-Interventionen. Real-time Feedback (Stromverbrauch, Fitness) kann Verhalten ändern.

\textbf{BEATRIX-Relevanz:}
BEATRIX empfiehlt Feedback-Mechanismen:
\begin{itemize}
    \item Welches Feedback wäre hilfreich?
    \item Wie schnell und spezifisch sollte es sein?
    \item Feedback-Loop in Interventionsdesign integrieren
\end{itemize}

\textbf{Konsequenz bei Fehlen:}
Feedback wäre wirkungslos. Keine Lernkurven. Statische Biases.

\textbf{Validiert:} AWX

\textbf{Abhängigkeiten:} MA-EPI-19

\textbf{Literatur:} Kluger \& DeNisi (1996); Hattie \& Timperley (2007)

%==============================================================================
% MA-EPI-21: Gewohnheiten reduzieren Kognition
%==============================================================================

\subsubsection*{MA-EPI-21: Gewohnheiten reduzieren kognitive Last}

\textbf{Statement:} Automatisierung spart kognitive Ressourcen.

\textbf{Formel:}
\begin{equation}
K_{\text{needed}}(V_{\text{habit}}) < K_{\text{needed}}(V_{\text{new}})
\end{equation}

\textbf{Beschreibung:}
Gewohnheiten sind kognitiv effizient:
\begin{itemize}
    \item Automatische Ausführung ohne bewusste Entscheidung
    \item Freisetzung kognitiver Ressourcen für andere Aufgaben
    \item Stabilität auch unter Stress
\end{itemize}
Nachteile: Schwer zu ändern, können maladaptiv sein.

\textbf{BCM-Relevanz:}
Erklärt Persistenz von Verhalten. Gewohnheitsänderung erfordert temporär erhöhte kognitive Ressourcen. DESIGN-Interventionen können Habit-Cues unterbrechen.

\textbf{BEATRIX-Relevanz:}
BEATRIX analysiert Gewohnheitsgrad:
\begin{itemize}
    \item Ist das Verhalten habituiert?
    \item Welche Cues triggern die Gewohnheit?
    \item Interventionsstrategie: Habit Disruption oder Habit Formation
\end{itemize}

\textbf{Konsequenz bei Fehlen:}
Alle Entscheidungen würden deliberativ behandelt. Automatisches Verhalten nicht modellierbar.

\textbf{Validiert:} WTX

\textbf{Abhängigkeiten:} MA-EPI-05

\textbf{Literatur:} Wood \& Neal (2007); Verplanken \& Wood (2006)

%==============================================================================
% MA-EPI-22: Debiasing ist schwierig
%==============================================================================

\subsubsection*{MA-EPI-22: Debiasing ist schwierig}

\textbf{Statement:} Kognitive Biases sind persistent.

\textbf{Formel:}
\begin{equation}
\frac{\partial \text{Bias}}{\partial \text{Intervention}} < \text{expected}
\end{equation}

\textbf{Beschreibung:}
Biases lassen sich nicht einfach ``wegtrainieren'':
\begin{itemize}
    \item Wissen über einen Bias schützt nicht davor
    \item Effekte von Debiasing-Training sind kurzlebig
    \item Biases kehren unter Stress zurück
\end{itemize}
Besser: Umgebungen designen, die Biases irrelevant machen.

\textbf{BCM-Relevanz:}
Erklärt Grenzen von THINK-Interventionen. Reine Aufklärung reicht nicht. DESIGN-Interventionen, die Bias-Effekte neutralisieren, sind oft wirksamer.

\textbf{BEATRIX-Relevanz:}
BEATRIX warnt vor Überoptimismus bezüglich Debiasing:
\begin{itemize}
    \item Debiasing-Interventionen mit Vorsicht empfehlen
    \item Alternative: Choice Architecture
    \item Langfristige Wirkung skeptisch bewerten
\end{itemize}

\textbf{Konsequenz bei Fehlen:}
Biases würden als leicht korrigierbar angenommen. Überinvestition in Aufklärung statt Umgebungsdesign.

\textbf{Validiert:} AWX

\textbf{Abhängigkeiten:} MA-EPI-06

\textbf{Literatur:} Fischhoff (1982); Lilienfeld et al. (2009)

%==============================================================================
% MA-EPI-23: Exponentielle Sättigung
%==============================================================================

\subsubsection*{MA-EPI-23: Exponentielle Sättigung (MA\_0.ES)}

\textbf{Statement:} Systeme mit proportionaler Unsicherheitsreduktion konvergieren exponentiell.

\textbf{Formel:}
\begin{equation}
\frac{dU}{dt} = -k \cdot U \quad \Rightarrow \quad U(t) = U_0 \cdot e^{-kt}
\end{equation}

\textbf{Beschreibung:}
Wenn Unsicherheit proportional zu sich selbst reduziert wird:
\begin{itemize}
    \item Schneller Anfangsprogress
    \item Asymptotische Annäherung an Null
    \item Übergangspunkt bei ~63\% Abbau (Rest $\approx 1/e$)
\end{itemize}
Anwendung: Lernkurven, Informationsakkumulation, Habituation.

\textbf{BCM-Relevanz:}
Mathematische Grundlage für viele dynamische Prozesse im BCM. Erklärt, warum initiale Interventionseffekte groß sind, dann abnehmen.

\textbf{BEATRIX-Relevanz:}
BEATRIX modelliert Sättigungseffekte:
\begin{itemize}
    \item Wann ist der Übergangspunkt erreicht?
    \item Ab wann lohnen weitere Interventionen nicht mehr?
    \item Optimaler Interventionszeitpunkt
\end{itemize}

\textbf{Konsequenz bei Fehlen:}
Lineare Extrapolation statt Sättigung. Langfristprognosen wären falsch.

\textbf{Validiert:} META

\textbf{Abhängigkeiten:} MA-EPI-19

\textbf{Literatur:} Eigenes Konstrukt; vgl. Newell \& Rosenbloom (1981)

%==============================================================================
% MA-EPI-24: Bayes'sches Updating
%==============================================================================

\subsubsection*{MA-EPI-24: Bayes'sches Updating}

\textbf{Statement:} Rationale Belief-Updates folgen dem Bayes-Theorem.

\textbf{Formel:}
\begin{equation}
P(H|E) = \frac{P(E|H) \cdot P(H)}{P(E)}
\end{equation}

\textbf{Beschreibung:}
Bayes beschreibt, wie Überzeugungen angesichts neuer Evidenz aktualisiert werden sollten:
\begin{itemize}
    \item Prior $P(H)$: Vorherige Überzeugung
    \item Likelihood $P(E|H)$: Wie wahrscheinlich ist die Evidenz unter der Hypothese?
    \item Posterior $P(H|E)$: Aktualisierte Überzeugung
\end{itemize}

\textbf{BCM-Relevanz:}
Normativer Benchmark für Lernprozesse. Abweichungen von Bayes (Base Rate Neglect, Confirmation Bias) sind als Biases modelliert.

\textbf{BEATRIX-Relevanz:}
BEATRIX verwendet Bayes für:
\begin{itemize}
    \item Eigene Inferenz über Persona-Parameter
    \item Modellierung, wie die Persona lernt
    \item Benchmark für Bias-Identifikation
\end{itemize}

\textbf{Konsequenz bei Fehlen:}
Kein normativer Standard für Lernen. Biases wären nicht als Abweichungen identifizierbar.

\textbf{Validiert:} AWX

\textbf{Abhängigkeiten:} MA-EPI-19

\textbf{Literatur:} Bayes (1763); Edwards et al. (1963)

%==============================================================================
% MA-EPI-25: Confirmation Bias
%==============================================================================

\subsubsection*{MA-EPI-25: Confirmation Bias}

\textbf{Statement:} Bestätigende Information wird bevorzugt gesucht und verarbeitet.

\textbf{Formel:}
\begin{equation}
P(\text{Search}|\text{Confirms}) > P(\text{Search}|\text{Disconfirms})
\end{equation}

\textbf{Beschreibung:}
Menschen suchen aktiv nach Information, die bestehende Überzeugungen bestätigt:
\begin{itemize}
    \item Selektive Mediennutzung (Echo Chambers)
    \item Asymmetrische Evidenzbewertung
    \item Immunisierung gegen Gegenargumente
\end{itemize}

\textbf{BCM-Relevanz:}
Erklärt Persistenz von Fehlüberzeugungen. Information allein ändert Beliefs nicht. Dialogische Interventionen nötig.

\textbf{BEATRIX-Relevanz:}
BEATRIX berücksichtigt Confirmation Bias:
\begin{itemize}
    \item Welche Überzeugungen hat die Persona?
    \item Wird neue Information abgelehnt werden?
    \item Framing so, dass es nicht als Angriff wirkt
\end{itemize}

\textbf{Konsequenz bei Fehlen:}
Information würde neutral verarbeitet. Fake News und Polarisierung unerklärlich.

\textbf{Validiert:} AWX

\textbf{Abhängigkeiten:} MA-EPI-06

\textbf{Literatur:} Nickerson (1998); Klayman (1995)

%==============================================================================
% MA-EPI-26: Bereits in axiom_beispiele.tex
%==============================================================================

%==============================================================================
% MA-EPI-27: Representativeness Heuristic
%==============================================================================

\subsubsection*{MA-EPI-27: Representativeness Heuristic}

\textbf{Statement:} Urteile basieren auf Ähnlichkeit zu Prototypen.

\textbf{Formel:}
\begin{equation}
\hat{P}(A \in \text{Category}) = f(\text{Similarity}(A, \text{Prototype}))
\end{equation}

\textbf{Beschreibung:}
Die Repräsentativitätsheuristik ersetzt Wahrscheinlichkeit durch Ähnlichkeit:
\begin{itemize}
    \item ``Linda Problem'': Konjunktionsfehlschluss
    \item Stereotypen-basierte Urteile
    \item Ignorieren von Base Rates (MA-EPI-26)
\end{itemize}

\textbf{BCM-Relevanz:}
Fundamental für Typenklassifikation. Menschen kategorisieren andere (und sich selbst) nach Prototypen. Kann zu systematischen Fehlklassifikationen führen.

\textbf{BEATRIX-Relevanz:}
BEATRIX korrigiert für Representativeness:
\begin{itemize}
    \item Prototypen explizit definieren
    \item Base Rates gewichten
    \item Vorsicht bei ``passt perfekt''-Urteilen
\end{itemize}

\textbf{Konsequenz bei Fehlen:}
Kategorisierung wäre rein probabilistisch. Intuitive Typisierung nicht modellierbar.

\textbf{Validiert:} AWX, ESTIMATION

\textbf{Abhängigkeiten:} MA-EPI-04

\textbf{Literatur:} Kahneman \& Tversky (1972)

%==============================================================================
% MA-EPI-28: Sunk Cost Fallacy
%==============================================================================

\subsubsection*{MA-EPI-28: Sunk Cost Fallacy}

\textbf{Statement:} Versunkene Kosten beeinflussen zukünftige Entscheidungen.

\textbf{Formel:}
\begin{equation}
U(\text{Continue}|\text{SunkCost}) > U(\text{Continue}|\text{NoSunkCost})
\end{equation}

\textbf{Beschreibung:}
Rationale Entscheidungen sollten nur zukünftige Kosten und Nutzen berücksichtigen. Aber:
\begin{itemize}
    \item ``Ich habe schon so viel investiert''
    \item Weitermachen bei gescheiterten Projekten
    \item Essen aufessen ``weil bezahlt''
\end{itemize}

\textbf{BCM-Relevanz:}
Erklärt irrationale Persistenz. Relevant für Commitment und Eskalation. Kann Verhaltensänderung blockieren (``habe schon so lange geraucht'').

\textbf{BEATRIX-Relevanz:}
BEATRIX identifiziert Sunk-Cost-Fallen:
\begin{itemize}
    \item Welche vergangenen Investitionen blockieren Änderung?
    \item Reframing: Zukunftsorientierte Perspektive fördern
    \item ``Sunk Costs sind sunk'' -- Debiasing-Intervention
\end{itemize}

\textbf{Konsequenz bei Fehlen:}
Vergangene Investitionen wären irrelevant. Eskalation von Commitment unerklärlich.

\textbf{Validiert:} INU

\textbf{Abhängigkeiten:} MA-EPI-15

\textbf{Literatur:} Arkes \& Blumer (1985); Thaler (1980)

%==============================================================================
% MA-EPI-29: Endowment Effect
%==============================================================================

\subsubsection*{MA-EPI-29: Endowment Effect}

\textbf{Statement:} Eigene Güter werden höher bewertet als identische fremde.

\textbf{Formel:}
\begin{equation}
WTA > WTP, \quad \text{typisch: } WTA/WTP \approx 2
\end{equation}

\textbf{Beschreibung:}
Die Verkaufsbereitschaft (WTA = Willingness to Accept) ist systematisch höher als die Kaufbereitschaft (WTP = Willingness to Pay) für dasselbe Gut:
\begin{itemize}
    \item Kaffeebecherexperiment: WTA 2x so hoch wie WTP
    \item ``Mein Kind ist das klügste''
    \item Überbewertung eigener Ideen/Arbeit
\end{itemize}

\textbf{BCM-Relevanz:}
Konsequenz von Loss Aversion (MA-EPI-15). Abgeben ist ein Verlust. Erklärt Marktfriktionen und suboptimalen Handel.

\textbf{BEATRIX-Relevanz:}
BEATRIX berücksichtigt Endowment:
\begin{itemize}
    \item Was ``besitzt'' die Persona (materiell und immateriell)?
    \item Wie stark ist der Endowment Effect bei ihr?
    \item Interventionen: Probebesitz vs. direkter Kauf
\end{itemize}

\textbf{Konsequenz bei Fehlen:}
WTA = WTP (ökonomisches Standardmodell). Besitzeffekte unerklärlich.

\textbf{Validiert:} INU

\textbf{Abhängigkeiten:} MA-EPI-15

\textbf{Literatur:} Kahneman, Knetsch \& Thaler (1990)

%==============================================================================
% MA-EPI-30: Optimism Bias
%==============================================================================

\subsubsection*{MA-EPI-30: Optimism Bias}

\textbf{Statement:} Positive Outcomes werden für sich selbst überschätzt.

\textbf{Formel:}
\begin{equation}
\hat{P}(\text{Positive}|\text{Self}) > P(\text{Positive})
\end{equation}

\textbf{Beschreibung:}
Unrealistischer Optimismus für eigene Zukunft:
\begin{itemize}
    \item ``Mir passiert das nicht'' (Unfälle, Krankheiten)
    \item Überschätzung von Karrierechancen
    \item Unterschätzung eigener Risiken
\end{itemize}

\textbf{BCM-Relevanz:}
Erklärt Unterversicherung, Risikoignoranz und mangelnde Prävention. Optimism Bias ist adaptiv (Wohlbefinden), aber kann zu Fehlentscheidungen führen.

\textbf{BEATRIX-Relevanz:}
BEATRIX korrigiert für Optimism Bias:
\begin{itemize}
    \item Risikowahrnehmung nach oben korrigieren
    \item Personalisierte Risikostatistiken präsentieren
    \item ``Das kann mir passieren''-Interventionen
\end{itemize}

\textbf{Konsequenz bei Fehlen:}
Selbst-Risiken würden korrekt geschätzt. Präventionslücken unerklärlich.

\textbf{Validiert:} AWX, IDN

\textbf{Abhängigkeiten:} MA-EPI-06

\textbf{Literatur:} Weinstein (1980); Sharot (2011)

%==============================================================================
% MA-EPI-31: Self-Serving Bias
%==============================================================================

\subsubsection*{MA-EPI-31: Self-Serving Bias}

\textbf{Statement:} Erfolge werden internal, Misserfolge external attribuiert.

\textbf{Formel:}
\begin{equation}
P(\text{Internal}|\text{Success}) > P(\text{Internal}|\text{Failure})
\end{equation}

\textbf{Beschreibung:}
Attributionsmuster zum Selbstwertschutz:
\begin{itemize}
    \item Erfolg: ``Ich bin gut''
    \item Misserfolg: ``Pech gehabt'' oder ``unfair''
\end{itemize}
Dient der Selbstwerterhaltung, behindert aber Lernen aus Fehlern.

\textbf{BCM-Relevanz:}
Kern-Mechanismus für Motivated Inference im IDN-Modul. Identitätsschutz durch selektive Attribution. Relevant für Feedback-Interventionen.

\textbf{BEATRIX-Relevanz:}
BEATRIX berücksichtigt Self-Serving Bias:
\begin{itemize}
    \item Wie wird die Persona Misserfolge attribuieren?
    \item Feedback so gestalten, dass Selbstwert nicht bedroht wird
    \item Lernen aus Fehlern ermöglichen
\end{itemize}

\textbf{Konsequenz bei Fehlen:}
Symmetrische Attribution. Selbstwertschutz-Dynamik nicht modellierbar.

\textbf{Validiert:} IDN

\textbf{Abhängigkeiten:} MA-EPI-06

\textbf{Literatur:} Miller \& Ross (1975); Mezulis et al. (2004)

%==============================================================================
% MA-EPI-32: Social Desirability Bias
%==============================================================================

\subsubsection*{MA-EPI-32: Social Desirability Bias}

\textbf{Statement:} Selbstberichte sind Richtung sozialer Erwünschtheit verzerrt.

\textbf{Formel:}
\begin{equation}
\text{Report}(V) = V_{\text{true}} + \varepsilon_{SD}, \quad \varepsilon_{SD} = f(\text{Visibility}, \text{Norm})
\end{equation}

\textbf{Beschreibung:}
Menschen berichten nicht wahrheitsgemäß, wenn die Wahrheit sozial unerwünscht ist:
\begin{itemize}
    \item Übertreibung von prosozialem Verhalten
    \item Untertreibung von Tabus (Alkohol, Rassismus)
    \item Stärker bei beobachteten Antworten
\end{itemize}

\textbf{BCM-Relevanz:}
\textbf{Kritisch für BEATRIX!} Persona-Beschreibungen sind systematisch verzerrt. Das LLM sieht eine ``aufgehübschte'' Version.

\textbf{BEATRIX-Relevanz:}
BEATRIX korrigiert für Social Desirability:
\begin{itemize}
    \item Persona-Angaben diskontieren
    \item Verhalten > Selbstbericht gewichten
    \item ``True self'' liegt näher am Mittelwert
\end{itemize}

\textbf{Konsequenz bei Fehlen:}
Persona-Beschreibungen würden für bare Münze genommen. Systematische Überschätzung von ``guten'' Eigenschaften.

\textbf{Validiert:} AWX, ESTIMATION

\textbf{Abhängigkeiten:} MA-EPI-06, MA-ONT-11

\textbf{Literatur:} Crowne \& Marlowe (1960); Paulhus (1984)

%==============================================================================
% MA-EPI-33: Omission Bias
%==============================================================================

\subsubsection*{MA-EPI-33: Omission Bias}

\textbf{Statement:} Unterlassung wird milder bewertet als aktive Handlung.

\textbf{Formel:}
\begin{equation}
U(\text{Outcome}|\text{Omission}) > U(\text{Outcome}|\text{Commission})
\end{equation}

\textbf{Beschreibung:}
Bei gleichem negativen Outcome wird Nichtstun weniger negativ bewertet als aktives Handeln:
\begin{itemize}
    \item Trolley-Problem: Nicht-Umleiten vs. aktives Stoßen
    \item Impfverweigerung: ``Nebenwirkung durch Impfung'' > ``Krankheit durch Nicht-Impfen''
\end{itemize}

\textbf{BCM-Relevanz:}
Relevant für ethische Bewertung von Interventionen. Opt-Out wird anders bewertet als Opt-In, selbst bei gleichem Ergebnis.

\textbf{BEATRIX-Relevanz:}
BEATRIX nutzt Omission Bias:
\begin{itemize}
    \item Default-Interventionen sind ``Omission'' für den Nutzer
    \item Aktive Entscheidung fühlt sich ``riskanter'' an
    \item Framing: ``Nichtstun hat Konsequenzen''
\end{itemize}

\textbf{Konsequenz bei Fehlen:}
Handlung und Unterlassung würden gleich bewertet. Default-Effekte wären rein durch Effort erklärbar.

\textbf{Validiert:} GOV, IDN

\textbf{Abhängigkeiten:} MA-EPI-06

\textbf{Literatur:} Spranca, Minsk \& Baron (1991); Ritov \& Baron (1990)

%==============================================================================
% MA-EPI-34: Moral Licensing
%==============================================================================

\subsubsection*{MA-EPI-34: Moral Licensing}

\textbf{Statement:} Gute Taten lizenzieren nachfolgende schlechte.

\textbf{Formel:}
\begin{equation}
P(\text{BadAction}|\text{PriorGoodAction}) > P(\text{BadAction}|\text{NoPriorAction})
\end{equation}

\textbf{Beschreibung:}
Moral Licensing ist ein negativer Spillover-Effekt:
\begin{itemize}
    \item ``Ich habe Sport gemacht, also kann ich jetzt Kuchen essen''
    \item ``Ich habe gespendet, also muss ich nicht recyceln''
    \item Selbstbild-Kredit wird aufgebaut und ausgegeben
\end{itemize}

\textbf{BCM-Relevanz:}
Identitäts-Dynamik: Gute Taten stärken das Selbstbild als ``guter Mensch'', was temporär schlechtes Verhalten erlaubt. Wichtig für Interventionskombinationen.

\textbf{BEATRIX-Relevanz:}
BEATRIX warnt vor Moral Licensing:
\begin{itemize}
    \item Nach erfolgreicher Intervention 1: Risiko für Backlash in Bereich 2
    \item Nicht mehrere ``gute'' Verhaltensweisen simultan einfordern
    \item Identitätsbasierte Interventionen können Licensing auslösen
\end{itemize}

\textbf{Konsequenz bei Fehlen:}
Positiver Spillover würde überschätzt. Interventionspakete könnten kontraproduktiv sein.

\textbf{Validiert:} IDN, KNU

\textbf{Abhängigkeiten:} MA-EPI-06

\textbf{Literatur:} Merritt, Effron \& Monin (2010)

%==============================================================================
% MA-EPI-35: Identifiable Victim Effect
%==============================================================================

\subsubsection*{MA-EPI-35: Identifiable Victim Effect}

\textbf{Statement:} Ein identifizierbares Opfer erzeugt mehr Empathie als statistische Opfer.

\textbf{Formel:}
\begin{equation}
U_{\text{empathy}}(1_{\text{identified}}) > U_{\text{empathy}}(N_{\text{statistical}})
\end{equation}

\textbf{Beschreibung:}
``Ein Toter ist eine Tragödie, eine Million Tote sind eine Statistik'' (Stalin zugeschrieben):
\begin{itemize}
    \item Spendenbereitschaft höher für Einzelschicksale
    \item Namen und Gesichter aktivieren Empathie
    \item Abstrakte Zahlen lassen kalt
\end{itemize}

\textbf{BCM-Relevanz:}
Erklärt asymmetrische Ressourcenallokation. Medial präsente Einzelfälle bekommen mehr als statistisch größere Probleme.

\textbf{BEATRIX-Relevanz:}
BEATRIX nutzt Identifiable Victim Effect:
\begin{itemize}
    \item Personalisierte Beispiele statt Statistiken
    \item ``Das könnte [Name] sein'' statt ``X\% Betroffene''
    \item Aber: Ethische Grenzen der Emotionalisierung
\end{itemize}

\textbf{Konsequenz bei Fehlen:}
Empathie wäre proportional zu Anzahl Betroffener. Fundraising-Strategien wären anders.

\textbf{Validiert:} KNU, IDN

\textbf{Abhängigkeiten:} MA-EPI-06

\textbf{Literatur:} Small \& Loewenstein (2003); Slovic (2007)

%==============================================================================
% MA-EPI-36: Scope Insensitivity
%==============================================================================

\subsubsection*{MA-EPI-36: Scope Insensitivity}

\textbf{Statement:} Bewertungen skalieren nicht proportional mit der Anzahl Betroffener.

\textbf{Formel:}
\begin{equation}
\frac{\partial U}{\partial N} \rightarrow 0 \quad \text{für } N \rightarrow \infty
\end{equation}

\textbf{Beschreibung:}
Die Zahlungsbereitschaft für Rettung von 2000 Vögeln ist kaum höher als für 200:
\begin{itemize}
    \item WTP für 2k, 20k, 200k Vögel: \$80, \$78, \$88
    \item ``Psychic numbing'': Große Zahlen werden nicht gefühlt
    \item 1 \approx 2 \approx viele
\end{itemize}

\textbf{BCM-Relevanz:}
Kollektiver Nutzen (KNU) wird systematisch unterschätzt. Je größer das Problem, desto mehr wird ignoriert (Klimawandel, Pandemie).

\textbf{BEATRIX-Relevanz:}
BEATRIX korrigiert für Scope Insensitivity:
\begin{itemize}
    \item Große Zahlen herunterbrechen (``pro Kopf'')
    \item Visualisierungen statt abstrakte Zahlen
    \item KNU explizit hochrechnen
\end{itemize}

\textbf{Konsequenz bei Fehlen:}
Nutzen wäre linear in Anzahl Betroffener. Kollektive Probleme würden angemessen gewichtet.

\textbf{Validiert:} KNU

\textbf{Abhängigkeiten:} MA-EPI-35

\textbf{Literatur:} Kahneman (1986); Desvousges et al. (1993)

%==============================================================================
% MA-EPI-37: Affect Heuristic
%==============================================================================

\subsubsection*{MA-EPI-37: Affect Heuristic}

\textbf{Statement:} Emotionale Reaktionen bestimmen Risiko- und Nutzenurteile.

\textbf{Formel:}
\begin{equation}
\hat{U}_{\text{risk}} = f(\text{Affect}), \quad \hat{U}_{\text{benefit}} = g(\text{Affect})
\end{equation}
\begin{equation}
\text{Corr}(\hat{U}_{\text{risk}}, \hat{U}_{\text{benefit}}) < 0
\end{equation}

\textbf{Beschreibung:}
Die Affect Heuristic ersetzt analytische Bewertung durch Bauchgefühl:
\begin{itemize}
    \item Positives Gefühl → niedriges Risiko, hoher Nutzen
    \item Negatives Gefühl → hohes Risiko, niedriger Nutzen
    \item Dies erzeugt negative Korrelation, die objektiv nicht existiert
\end{itemize}

\textbf{BCM-Relevanz:}
Erklärt emotionale Bewertung von Technologien (Kernenergie: emotional negativ → Risiko hoch, Nutzen niedrig; Solar: umgekehrt). Emotions-basierte Interventionen nutzen dies.

\textbf{BEATRIX-Relevanz:}
BEATRIX modelliert Affect:
\begin{itemize}
    \item Welche Emotionen sind mit dem Verhalten assoziiert?
    \item Können positive Affekte aktiviert werden?
    \item Emotionale vs. analytische Persona-Typen
\end{itemize}

\textbf{Konsequenz bei Fehlen:}
Risiko und Nutzen wären unabhängig bewertet. Emotionale Komponente von Entscheidungen unsichtbar.

\textbf{Validiert:} AWX, INU

\textbf{Abhängigkeiten:} MA-EPI-05

\textbf{Literatur:} Slovic et al. (2002); Finucane et al. (2000)

\bigskip
\noindent\rule{\textwidth}{0.4pt}
\textit{Ende der EPI-Axiome. Alle 37 Epistemologie-Axiome sind nun vollständig dokumentiert.}

% =============================================================================
% GOV-AXIOME (GOVERNANCE) - Vollständige Ausformulierung
% BCM 2.2 META-Modul
% =============================================================================

\section{Governance-Axiome (MA-GOV)}

Die GOV-Axiome definieren die Regeln, Constraints und ethischen Leitplanken.
Sie beantworten die Frage: \textit{Was ist erlaubt?}

%==============================================================================
% MA-GOV-01: Constraints existieren
%==============================================================================

\subsubsection*{MA-GOV-01: Constraints existieren}

\textbf{Statement:} Es gibt Beschränkungen für Verhalten.

\textbf{Formel:}
\begin{equation}
\Gamma \neq \emptyset
\end{equation}

\textbf{Beschreibung:}
Nicht alles ist erlaubt oder möglich. Constraints umfassen:
\begin{itemize}
    \item Physische Grenzen (Zeit, Raum, Körper)
    \item Rechtliche Grenzen (Gesetze, Verbote)
    \item Soziale Grenzen (Normen, Sanktionen)
    \item Ökonomische Grenzen (Budget, Ressourcen)
\end{itemize}

\textbf{BCM-Relevanz:}
Fundament des GOV-Moduls. Das BCM modelliert Verhalten unter Nebenbedingungen. REGULATE-Interventionen wirken durch Constraint-Manipulation.

\textbf{BEATRIX-Relevanz:}
BEATRIX muss Constraints kennen:
\begin{itemize}
    \item Welche Constraints limitieren die Persona?
    \item Sind Verhaltensoptionen feasible?
    \item Können Constraints gelockert werden?
\end{itemize}

\textbf{Konsequenz bei Fehlen:}
Alle Verhaltensweisen wären möglich. Keine Beschränkungen modellierbar. Regulierung hätte keine Grundlage.

\textbf{Validiert:} META

\textbf{Abhängigkeiten:} MA-ONT-06

\textbf{Literatur:} Becker (1968); North (1990)

%==============================================================================
% MA-GOV-02: Kontextabhängigkeit von Constraints
%==============================================================================

\subsubsection*{MA-GOV-02: Constraints sind kontextabhängig}

\textbf{Statement:} Erlaubtes hängt vom Kontext ab.

\textbf{Formel:}
\begin{equation}
\Gamma = f(\rho)
\end{equation}

\textbf{Beschreibung:}
Was in einem Kontext erlaubt ist, kann in einem anderen verboten sein:
\begin{itemize}
    \item Alkohol: erlaubt privat, verboten am Steuer
    \item Kleidung: angemessen im Büro ≠ am Strand
    \item Kommunikation: informell unter Freunden ≠ formell im Gericht
\end{itemize}

\textbf{BCM-Relevanz:}
Ermöglicht kontextspezifische Governance. Das 5-Ebenen-Modell (MA-ONT-26) strukturiert Kontexte hierarchisch.

\textbf{BEATRIX-Relevanz:}
BEATRIX prüft kontextspezifische Constraints:
\begin{itemize}
    \item Welche Regeln gelten in diesem spezifischen Kontext?
    \item Gibt es Kontextwechsel, die Constraints ändern?
    \item Kontextspezifische Interventionen designen
\end{itemize}

\textbf{Konsequenz bei Fehlen:}
Universelle Constraints. Kontextsensitivität von Regeln nicht abbildbar.

\textbf{Validiert:} META

\textbf{Abhängigkeiten:} MA-GOV-01, MA-ONT-04

\textbf{Literatur:} Ostrom (2005); Crawford \& Ostrom (1995)

%==============================================================================
% MA-GOV-03: Autonomie-Primat
%==============================================================================

\subsubsection*{MA-GOV-03: Autonomie-Primat}

\textbf{Statement:} Individuelle Autonomie hat Vorrang.

\textbf{Formel:}
\begin{equation}
\forall \Gamma: \min(\text{Einschränkung}) \quad \text{s.t. } \max(\text{Zielerreichung})
\end{equation}

\textbf{Beschreibung:}
Das BCM ist nicht paternalistisch. Autonomie ist ein Wert an sich:
\begin{itemize}
    \item Menschen haben das Recht auf eigene Entscheidungen
    \item Auch ``schlechte'' Entscheidungen sind zu respektieren
    \item Interventionen nur bei klarem Mehrwert und minimal invasiv
\end{itemize}

\textbf{BCM-Relevanz:}
Ethisches Leitprinzip. Begrenzt den Interventionsraum. Begründet die Präferenz für Nudges über Verbote.

\textbf{BEATRIX-Relevanz:}
BEATRIX respektiert Autonomie:
\begin{itemize}
    \item Keine manipulativen Interventionen empfehlen
    \item Transparenz über Interventionsziele
    \item Opt-Out-Möglichkeiten vorsehen
\end{itemize}

\textbf{Konsequenz bei Fehlen:}
Keine ethische Grenze für Interventionen. Paternalismus würde legitimiert.

\textbf{Validiert:} META

\textbf{Abhängigkeiten:} MA-GOV-01

\textbf{Literatur:} Mill (1859); Sunstein (2014)

%==============================================================================
% MA-GOV-04: Nicht-Normativität
%==============================================================================

\subsubsection*{MA-GOV-04: Nicht-Normativität}

\textbf{Statement:} Das BCM beschreibt, schreibt nicht vor.

\textbf{Formel:}
\begin{equation}
\text{BCM}: \text{is} \rightarrow \text{is}, \quad \text{nicht: is} \rightarrow \text{ought}
\end{equation}

\textbf{Beschreibung:}
Das BCM ist ein deskriptives, kein normatives System:
\begin{itemize}
    \item Es sagt, wie Menschen sich verhalten, nicht wie sie sich verhalten sollen
    \item Es ist wertfrei bezüglich der Inhalte
    \item Die Ziele der Intervention werden extern gesetzt
\end{itemize}

\textbf{BCM-Relevanz:}
Wichtige Abgrenzung. Das BCM kann für viele Ziele eingesetzt werden. Es ist ein Werkzeug, keine Weltanschauung.

\textbf{BEATRIX-Relevanz:}
BEATRIX ist agnostisch bezüglich Zielen:
\begin{itemize}
    \item Das Zielverhalten wird vom Nutzer vorgegeben
    \item BEATRIX optimiert für das vorgegebene Ziel
    \item Ethische Bewertung des Ziels ist Nutzer-Verantwortung
\end{itemize}

\textbf{Konsequenz bei Fehlen:}
Das BCM würde bestimmte Verhaltensweisen als ``richtig'' definieren. Ideologische Aufladung.

\textbf{Validiert:} META

\textbf{Abhängigkeiten:} -- (meta-ethisches Axiom)

\textbf{Literatur:} Weber (1917); Hume (1739)

%==============================================================================
% MA-GOV-05: Transparenz-Gebot
%==============================================================================

\subsubsection*{MA-GOV-05: Transparenz-Gebot}

\textbf{Statement:} Interventionen sollen transparent sein.

\textbf{Formel:}
\begin{equation}
\text{Transparency}(I) \rightarrow \max
\end{equation}

\textbf{Beschreibung:}
Menschen haben ein Recht zu wissen, dass und wie sie beeinflusst werden:
\begin{itemize}
    \item Offenlegung von Nudges
    \item Erklärung von Default-Setzungen
    \item Keine verdeckte Manipulation
\end{itemize}
Transparenz kann Wirkung reduzieren, ist aber ethisch geboten.

\textbf{BCM-Relevanz:}
Ethische Grenze für Interventionsdesign. ``Libertarian Paternalism'' verlangt Transparenz als Legitimationsbedingung.

\textbf{BEATRIX-Relevanz:}
BEATRIX empfiehlt transparente Interventionen:
\begin{itemize}
    \item Interventionen, die bei Offenlegung noch wirken, bevorzugen
    \item Warnung bei verdeckten Interventionen
    \item Transparenz-Score als Bewertungskriterium
\end{itemize}

\textbf{Konsequenz bei Fehlen:}
Verdeckte Manipulation wäre legitimiert. Vertrauensverlust bei Bekanntwerden.

\textbf{Validiert:} META

\textbf{Abhängigkeiten:} MA-GOV-03

\textbf{Literatur:} Thaler \& Sunstein (2008); Hansen \& Jespersen (2013)

%==============================================================================
% MA-GOV-06: Bereits in axiom_beispiele.tex
%==============================================================================

%==============================================================================
% MA-GOV-07: 8 Interventionstypen
%==============================================================================

\subsubsection*{MA-GOV-07: Acht Interventionstypen}

\textbf{Statement:} Es gibt acht grundlegende Interventionstypen.

\textbf{Formel:}
\begin{equation}
I \in \{\text{NUDGE, BOOST, THINK, INCENTIVE, REGULATE, DESIGN, SOCIAL, IDENTITY}\}
\end{equation}

\textbf{Beschreibung:}
Die acht Typen decken das Interventionsspektrum ab:
\begin{itemize}
    \item \textbf{NUDGE}: Entscheidungsarchitektur ändern
    \item \textbf{BOOST}: Kompetenz stärken
    \item \textbf{THINK}: Reflexion anregen
    \item \textbf{INCENTIVE}: Anreize setzen
    \item \textbf{REGULATE}: Regeln setzen/durchsetzen
    \item \textbf{DESIGN}: Physische Umgebung ändern
    \item \textbf{SOCIAL}: Soziale Normen aktivieren
    \item \textbf{IDENTITY}: Identität ansprechen
\end{itemize}

\textbf{BCM-Relevanz:}
Taxonomie für Interventionsplanung. Jeder Typ hat charakteristische Stärken und Grenzen.

\textbf{BEATRIX-Relevanz:}
BEATRIX klassifiziert und empfiehlt Interventionstypen:
\begin{itemize}
    \item Welcher Typ passt zur Persona und zum Kontext?
    \item Kombination verschiedener Typen planen
    \item Trade-offs zwischen Typen explizit machen
\end{itemize}

\textbf{Konsequenz bei Fehlen:}
Keine systematische Klassifikation von Interventionen. Ad-hoc-Interventionsdesign.

\textbf{Validiert:} META

\textbf{Abhängigkeiten:} MA-GOV-06

\textbf{Literatur:} Eigene Taxonomie, aufbauend auf Thaler \& Sunstein (2008); Hertwig \& Grüne-Yanoff (2017)

%==============================================================================
% MA-GOV-08: Freiheits-Hierarchie
%==============================================================================

\subsubsection*{MA-GOV-08: Freiheits-Hierarchie}

\textbf{Statement:} Interventionen haben unterschiedliche Freiheitsgrade.

\textbf{Formel:}
\begin{equation}
\text{Freedom}: \text{BOOST} > \text{THINK} > \text{NUDGE} > \text{DESIGN} > \text{SOCIAL} > \text{IDENTITY} > \text{INCENTIVE} > \text{REGULATE}
\end{equation}

\textbf{Beschreibung:}
Die Ranking ordnet Interventionen nach Autonomieeinschränkung:
\begin{itemize}
    \item BOOST: Erweitert Fähigkeiten, maximale Freiheit
    \item THINK: Aktiviert eigene Reflexion
    \item NUDGE: Sanfter Schubs, leicht überwindbar
    \item ...
    \item REGULATE: Verbot/Gebot, minimale Freiheit
\end{itemize}

\textbf{BCM-Relevanz:}
Operationalisierung des Subsidiaritätsprinzips (MA-GOV-06). Gibt klare Eskalationsreihenfolge vor.

\textbf{BEATRIX-Relevanz:}
BEATRIX verwendet die Hierarchie:
\begin{itemize}
    \item Empfehlungen mit Freedom-Score versehen
    \item Bei mehreren wirksamen Optionen: höhere Freiheit bevorzugen
    \item Trade-off Freiheit vs. Effektivität transparent machen
\end{itemize}

\textbf{Konsequenz bei Fehlen:}
Keine Ordnung von Interventionen nach Freiheit. Subsidiarität wäre nicht operationalisierbar.

\textbf{Validiert:} META

\textbf{Abhängigkeiten:} MA-GOV-07

\textbf{Literatur:} Eigenes Konstrukt; vgl. Sunstein (2016)

%==============================================================================
% MA-GOV-09: Eskalations-Prinzip
%==============================================================================

\subsubsection*{MA-GOV-09: Eskalations-Prinzip}

\textbf{Statement:} Von mild nach stark eskalieren.

\textbf{Formel:}
\begin{equation}
\text{Try}(I_n) \text{ before } \text{Try}(I_{n+1}) \quad \text{where } \text{Freedom}(I_n) > \text{Freedom}(I_{n+1})
\end{equation}

\textbf{Beschreibung:}
Die Interventionsleiter wird von oben nach unten abgearbeitet:
\begin{enumerate}
    \item Erst BOOST versuchen
    \item Wenn unwirksam: NUDGE
    \item Wenn unwirksam: INCENTIVE
    \item Nur als letztes: REGULATE
\end{enumerate}

\textbf{BCM-Relevanz:}
Prozedurale Norm für ethisches Interventionsdesign. Verhindert vorschnelle Regulierung.

\textbf{BEATRIX-Relevanz:}
BEATRIX empfiehlt Eskalationspfade:
\begin{itemize}
    \item Startpunkt: mildeste wirksame Intervention
    \item Fallback-Optionen bei Unwirksamkeit definieren
    \item Monitoring für Eskalationsentscheidung
\end{itemize}

\textbf{Konsequenz bei Fehlen:}
Interventionen könnten mit der härtesten Stufe beginnen. Unverhältnismäßigkeit.

\textbf{Validiert:} META

\textbf{Abhängigkeiten:} MA-GOV-08, MA-GOV-06

\textbf{Literatur:} Sunstein (2014)

%==============================================================================
% MA-GOV-10: 28 Interventions-Dimensionen
%==============================================================================

\subsubsection*{MA-GOV-10: 28 Interventions-Dimensionen}

\textbf{Statement:} Interventionen werden über 28 Dimensionen in 9 Kategorien charakterisiert.

\textbf{Formel:}
\begin{equation}
\text{Dim} = \{\text{Freiheit}, \text{Psychologie}, \text{Zeit}, \text{Reichweite}, \text{Kontext}, \text{Ethik}, \text{Implementation}, \text{Effekt}, \text{Dynamik}\}
\end{equation}

\textbf{Beschreibung:}
Jede Intervention kann auf 28 Dimensionen bewertet werden:
\begin{itemize}
    \item \textbf{Freiheit}: Restrictiveness, Opt-Out, Reversibility
    \item \textbf{Psychologie}: Target System (1/2), Mechanism, Awareness
    \item \textbf{Zeit}: Duration, Onset, Decay
    \item ...
\end{itemize}

\textbf{BCM-Relevanz:}
Ermöglicht präzise Charakterisierung und Vergleich von Interventionen. Systematische Planung statt Intuition.

\textbf{BEATRIX-Relevanz:}
BEATRIX verwendet die Dimensionen für:
\begin{itemize}
    \item Detaillierte Beschreibung empfohlener Interventionen
    \item Matching Intervention-Persona-Kontext
    \item Trade-off-Analyse zwischen Dimensionen
\end{itemize}

\textbf{Konsequenz bei Fehlen:}
Grobe Klassifikation nur nach Typ. Nuancen zwischen Interventionen nicht erfassbar.

\textbf{Validiert:} META

\textbf{Abhängigkeiten:} MA-GOV-07

\textbf{Literatur:} Eigenes Konstrukt

%==============================================================================
% MA-GOV-11: Distributive Gerechtigkeit
%==============================================================================

\subsubsection*{MA-GOV-11: Distributive Gerechtigkeit}

\textbf{Statement:} Kosten und Nutzen fair verteilen.

\textbf{Formel:}
\begin{equation}
\forall \sigma: \text{Burden}(\sigma) \propto \text{Capacity}(\sigma)
\end{equation}

\textbf{Beschreibung:}
Interventionen sollen nicht systematisch bestimmte Gruppen benachteiligen:
\begin{itemize}
    \item Wer profitiert, wer trägt die Kosten?
    \item Vulnerable Gruppen besonders schützen
    \item Progressivität: Stärkere Schultern tragen mehr
\end{itemize}

\textbf{BCM-Relevanz:}
Gerechtigkeitsprinzip für Interventionsdesign. Verhindert, dass Interventionen bestehende Ungleichheiten verstärken.

\textbf{BEATRIX-Relevanz:}
BEATRIX prüft distributive Effekte:
\begin{itemize}
    \item Wer ist von der Intervention betroffen?
    \item Sind vulnerable Gruppen überproportional belastet?
    \item Alternative Designs mit besserer Verteilung?
\end{itemize}

\textbf{Konsequenz bei Fehlen:}
Interventionen könnten systematisch ungerecht sein. Verstärkung von Ungleichheit.

\textbf{Validiert:} META

\textbf{Abhängigkeiten:} MA-GOV-03

\textbf{Literatur:} Rawls (1971); Sen (2009)

%==============================================================================
% MA-GOV-12: Reversibilität bevorzugen
%==============================================================================

\subsubsection*{MA-GOV-12: Reversibilität bevorzugen}

\textbf{Statement:} Reversible Interventionen sind vorzuziehen.

\textbf{Formel:}
\begin{equation}
\text{Prefer}(I_{\text{reversible}}) \text{ over } I_{\text{irreversible}}
\end{equation}

\textbf{Beschreibung:}
Interventionen sollten rückgängig machbar sein:
\begin{itemize}
    \item Wenn sie nicht wirken: Abbruch möglich
    \item Wenn unerwünschte Effekte: Korrektur möglich
    \item Experimentation wird ermöglicht
\end{itemize}
Irreversible Interventionen erfordern höhere Evidenzschwellen.

\textbf{BCM-Relevanz:}
Risikominimierung bei Unsicherheit. Ermöglicht iteratives Vorgehen (MA-GOV-20).

\textbf{BEATRIX-Relevanz:}
BEATRIX bewertet Reversibilität:
\begin{itemize}
    \item Reversibility-Score für jede Intervention
    \item Bei Unsicherheit: reversible Option wählen
    \item Warnung bei irreversiblen Empfehlungen
\end{itemize}

\textbf{Konsequenz bei Fehlen:}
Irreversible Interventionen ohne besondere Vorsicht. Keine Korrekturmöglichkeit bei Fehlern.

\textbf{Validiert:} META

\textbf{Abhängigkeiten:} MA-GOV-06

\textbf{Literatur:} Arrow \& Fisher (1974); Sunstein (2005)

%==============================================================================
% MA-GOV-13: Manipulation vermeiden
%==============================================================================

\subsubsection*{MA-GOV-13: Manipulation vermeiden}

\textbf{Statement:} Keine verdeckte, täuschende Beeinflussung.

\textbf{Formel:}
\begin{equation}
\text{If } \text{Hidden}(I) \land \text{Deceptive}(I) \rightarrow \text{Reject}(I)
\end{equation}

\textbf{Beschreibung:}
Manipulation = verdeckt + täuschend. Nicht jeder Nudge ist Manipulation:
\begin{itemize}
    \item Offener Nudge: OK
    \item Verdeckter, aber wahrhaftiger Nudge: Grauzone
    \item Verdeckter + täuschender Nudge: Manipulation = NICHT OK
\end{itemize}

\textbf{BCM-Relevanz:}
Ethische Grenze. Das BCM lehnt Manipulation ab, nicht Beeinflussung per se. Der Unterschied ist Transparenz und Wahrhaftigkeit.

\textbf{BEATRIX-Relevanz:}
BEATRIX prüft auf Manipulation:
\begin{itemize}
    \item Ist die Intervention bei Offenlegung noch vertretbar?
    \item Werden falsche Informationen verwendet?
    \item Manipulation-Flag bei Grenzfällen
\end{itemize}

\textbf{Konsequenz bei Fehlen:}
Keine Grenze zwischen Nudge und Manipulation. Ethische Kritik berechtigt.

\textbf{Validiert:} META

\textbf{Abhängigkeiten:} MA-GOV-05

\textbf{Literatur:} Sunstein (2016); Bovens (2009)

%==============================================================================
% MA-GOV-14: Bereits in axiom_beispiele.tex
%==============================================================================

%==============================================================================
% MA-GOV-15: Reaktanz beachten
%==============================================================================

\subsubsection*{MA-GOV-15: Reaktanz beachten}

\textbf{Statement:} Einschränkungen können psychologische Reaktanz auslösen.

\textbf{Formel:}
\begin{equation}
P(\text{Resistance}|\text{Restriction}) > 0
\end{equation}

\textbf{Beschreibung:}
Reaktanz ist die psychologische Gegenwehr gegen wahrgenommene Freiheitseinschränkung:
\begin{itemize}
    \item ``Verbotene Frucht'' wird attraktiver
    \item ``Trotzreaktion'': das Verbotene tun
    \item Besonders stark bei autonomieorientierten Personen
\end{itemize}

\textbf{BCM-Relevanz:}
Erklärt kontraproduktive Effekte von Regulierung. REGULATE kann Backlash erzeugen. Reaktanz ist besonders relevant bei IDN-starken Personen.

\textbf{BEATRIX-Relevanz:}
BEATRIX schätzt Reaktanz-Risiko:
\begin{itemize}
    \item Wie reaktanz-anfällig ist die Persona?
    \item Wie freiheitsbeschränkend ist die Intervention?
    \item Alternativen mit weniger Reaktanz-Risiko?
\end{itemize}

\textbf{Konsequenz bei Fehlen:}
Nur positive Interventionseffekte modelliert. Backlash-Risiko ignoriert.

\textbf{Validiert:} KNU

\textbf{Abhängigkeiten:} MA-GOV-01

\textbf{Literatur:} Brehm (1966); Steindl et al. (2015)

%==============================================================================
% MA-GOV-16: Ausweichverhalten beachten
%==============================================================================

\subsubsection*{MA-GOV-16: Ausweichverhalten beachten}

\textbf{Statement:} Regulierung kann zu Umgehung führen.

\textbf{Formel:}
\begin{equation}
P(\text{Evasion}|\text{REGULATE}) > 0
\end{equation}

\textbf{Beschreibung:}
Verbote werden nicht immer befolgt -- oft wird nach Schlupflöchern gesucht:
\begin{itemize}
    \item Steuerregeln → Steuervermeidung
    \item Geschwindigkeitsbegrenzung → Radarwarner
    \item Drogenverbot → illegaler Markt
\end{itemize}

\textbf{BCM-Relevanz:}
Begrenzt die Wirksamkeit von REGULATE. Bei hoher Evasion-Wahrscheinlichkeit sind andere Interventionstypen vorzuziehen.

\textbf{BEATRIX-Relevanz:}
BEATRIX schätzt Evasion-Wahrscheinlichkeit:
\begin{itemize}
    \item Wie einfach ist Umgehung?
    \item Wie hoch ist die Entdeckungswahrscheinlichkeit?
    \item Ist die Regulierung durchsetzbar?
\end{itemize}

\textbf{Konsequenz bei Fehlen:}
Perfekte Compliance angenommen. Regulierungs-Illusionen.

\textbf{Validiert:} WTX

\textbf{Abhängigkeiten:} MA-GOV-07

\textbf{Literatur:} Becker (1968); Slemrod (2007)

%==============================================================================
% MA-GOV-17: Opt-Out-Prinzip
%==============================================================================

\subsubsection*{MA-GOV-17: Opt-Out-Prinzip}

\textbf{Statement:} Austritt muss möglich sein (außer bei REGULATE).

\textbf{Formel:}
\begin{equation}
\forall I \neq \text{REGULATE}: \text{OptOut}(I) = \text{possible}
\end{equation}

\textbf{Beschreibung:}
Alle Interventionen außer REGULATE müssen Opt-Out erlauben:
\begin{itemize}
    \item Default kann überschrieben werden
    \item Nudge kann ignoriert werden
    \item Incentive kann abgelehnt werden
\end{itemize}
REGULATE ist die Ausnahme -- Gesetze gelten für alle.

\textbf{BCM-Relevanz:}
Operationalisiert Autonomie-Primat (MA-GOV-03). Unterscheidet ``libertarian paternalism'' von echtem Paternalismus.

\textbf{BEATRIX-Relevanz:}
BEATRIX stellt Opt-Out sicher:
\begin{itemize}
    \item Jede Intervention außer REGULATE hat Opt-Out-Weg
    \item Opt-Out muss praktisch möglich sein (nicht nur theoretisch)
    \item Warnung bei de-facto nicht-überwindbaren Defaults
\end{itemize}

\textbf{Konsequenz bei Fehlen:}
Alle Interventionen wären de facto Zwang. Verlust des libertären Elements.

\textbf{Validiert:} META

\textbf{Abhängigkeiten:} MA-GOV-03

\textbf{Literatur:} Thaler \& Sunstein (2008)

%==============================================================================
% MA-GOV-18: Informed Consent
%==============================================================================

\subsubsection*{MA-GOV-18: Informed Consent}

\textbf{Statement:} Informierte Zustimmung bei Interventionen.

\textbf{Formel:}
\begin{equation}
\text{Legitimate}(I) \rightarrow \text{Informed}(\sigma) \land \text{Consenting}(\sigma)
\end{equation}

\textbf{Beschreibung:}
Ideale Intervention: Die Zielperson versteht und akzeptiert sie:
\begin{itemize}
    \item Informed: Weiß, dass und wie sie beeinflusst wird
    \item Consenting: Stimmt der Beeinflussung zu
\end{itemize}
Bei öffentlichen Interventionen: implizite Zustimmung durch demokratische Legitimation.

\textbf{BCM-Relevanz:}
Höchster ethischer Standard. In der Praxis oft abgeschwächt, aber als Ideal wichtig.

\textbf{BEATRIX-Relevanz:}
BEATRIX bewertet Consent-Level:
\begin{itemize}
    \item Ist explizite Zustimmung möglich/nötig?
    \item Ist implizite Zustimmung ausreichend?
    \item Wie kann Informiertheit erhöht werden?
\end{itemize}

\textbf{Konsequenz bei Fehlen:}
Interventionen ohne jede Zustimmung legitim. Paternalismus pur.

\textbf{Validiert:} META

\textbf{Abhängigkeiten:} MA-GOV-05

\textbf{Literatur:} Faden \& Beauchamp (1986)

%==============================================================================
% MA-GOV-19: Interventionen habituieren
%==============================================================================

\subsubsection*{MA-GOV-19: Interventionen habituieren}

\textbf{Statement:} Interventionswirkung lässt über Zeit nach.

\textbf{Formel:}
\begin{equation}
\frac{\partial \text{Effect}(I)}{\partial \tau} < 0 \quad \text{für } \tau \rightarrow \infty
\end{equation}

\textbf{Beschreibung:}
Die meisten Interventionen verlieren Wirkung:
\begin{itemize}
    \item Nudges werden habituiert (``Blindheit'' für Hinweise)
    \item Incentives verlieren Neuheit
    \item Nur REGULATE und DESIGN bleiben stabil
\end{itemize}

\textbf{BCM-Relevanz:}
Erklärt Langzeit-Ineffektivität von Interventionen. Begründet Notwendigkeit von Iteration (MA-GOV-20).

\textbf{BEATRIX-Relevanz:}
BEATRIX modelliert Decay:
\begin{itemize}
    \item Erwartete Wirkdauer der Intervention
    \item Wann ist Refresh nötig?
    \item Decay-resistente Interventionstypen bevorzugen für Langzeiteffekte
\end{itemize}

\textbf{Konsequenz bei Fehlen:}
Konstante Interventionswirkung angenommen. Überschätzung von Langzeiteffekten.

\textbf{Validiert:} WTX

\textbf{Abhängigkeiten:} MA-ONT-20

\textbf{Literatur:} Allcott \& Rogers (2014)

%==============================================================================
% MA-GOV-20: Iteration notwendig
%==============================================================================

\subsubsection*{MA-GOV-20: Iteration notwendig}

\textbf{Statement:} Interventionen müssen iterativ angepasst werden.

\textbf{Formel:}
\begin{equation}
\text{Optimal}(I) \text{ requires Feedback-Loop}
\end{equation}

\textbf{Beschreibung:}
``Set and forget'' funktioniert nicht:
\begin{itemize}
    \item Erstimplementierung ist selten optimal
    \item Kontexte ändern sich
    \item Habituation erfordert Variation
\end{itemize}
PDCA-Zyklus: Plan-Do-Check-Act.

\textbf{BCM-Relevanz:}
Prozessprinzip. Interventionen sind keine einmaligen Events, sondern laufende Optimierung.

\textbf{BEATRIX-Relevanz:}
BEATRIX empfiehlt iteratives Design:
\begin{itemize}
    \item Monitoring-Metriken definieren
    \item Anpassungszeitpunkte festlegen
    \item A/B-Test-Optionen vorsehen
\end{itemize}

\textbf{Konsequenz bei Fehlen:}
Statische Interventionen. Keine Lernschleifen. Suboptimale Langzeitergebnisse.

\textbf{Validiert:} META

\textbf{Abhängigkeiten:} MA-GOV-19

\textbf{Literatur:} Deming (1986); Haynes et al. (2012)

%==============================================================================
% MA-GOV-21: Evaluation notwendig
%==============================================================================

\subsubsection*{MA-GOV-21: Evaluation notwendig}

\textbf{Statement:} Wirksamkeit muss gemessen werden.

\textbf{Formel:}
\begin{equation}
\forall I: \text{Measure}(\text{Effect}(I))
\end{equation}

\textbf{Beschreibung:}
Keine Intervention ohne Evaluation:
\begin{itemize}
    \item Wurde das Zielverhalten erreicht?
    \item War die Intervention ursächlich?
    \item Gibt es unbeabsichtigte Effekte?
\end{itemize}

\textbf{BCM-Relevanz:}
Wissenschaftliches Prinzip. Das BCM verlangt Evidenzbasierung, nicht nur Plausibilität.

\textbf{BEATRIX-Relevanz:}
BEATRIX integriert Evaluation:
\begin{itemize}
    \item KPIs für jede Intervention definieren
    \item Evaluationsdesign (RCT, Quasi-Experiment) vorschlagen
    \item Datenquellen identifizieren
\end{itemize}

\textbf{Konsequenz bei Fehlen:}
Intuitionsbasierte Interventionen. Keine Lernfähigkeit. Persistenz unwirksamer Maßnahmen.

\textbf{Validiert:} META

\textbf{Abhängigkeiten:} MA-GOV-20

\textbf{Literatur:} Campbell \& Stanley (1966); Shadish et al. (2002)

%==============================================================================
% MA-GOV-22: A/B-Testing bevorzugen
%==============================================================================

\subsubsection*{MA-GOV-22: A/B-Testing bevorzugen}

\textbf{Statement:} Experimentelle Evaluation wenn möglich.

\textbf{Formel:}
\begin{equation}
\text{RCT} > \text{Observational} \quad \text{für Kausalität}
\end{equation}

\textbf{Beschreibung:}
Randomisierte Kontrollversuche (RCT) liefern kausale Evidenz:
\begin{itemize}
    \item Zufällige Zuweisung eliminiert Selektionsbias
    \item Vergleich Treatment vs. Control
    \item Gold-Standard für Wirksamkeitsnachweis
\end{itemize}

\textbf{BCM-Relevanz:}
Methodisches Ideal. Wo RCTs nicht möglich: Quasi-experimentelle Designs, natürliche Experimente.

\textbf{BEATRIX-Relevanz:}
BEATRIX empfiehlt Testdesigns:
\begin{itemize}
    \item Ist ein RCT möglich?
    \item Sample-Size-Berechnung
    \item Alternative Designs bei RCT-Unmöglichkeit
\end{itemize}

\textbf{Konsequenz bei Fehlen:}
Korrelation = Kausalität-Fehler. Überschätzung von Interventionseffekten.

\textbf{Validiert:} META

\textbf{Abhängigkeiten:} MA-GOV-21

\textbf{Literatur:} Fisher (1935); Duflo et al. (2008)

%==============================================================================
% MA-GOV-23: Spillover messen
%==============================================================================

\subsubsection*{MA-GOV-23: Spillover messen}

\textbf{Statement:} Nicht-intendierte Effekte müssen erfasst werden.

\textbf{Formel:}
\begin{equation}
\text{Measure}(\text{Effect}_{\text{unintended}}) \text{ alongside } \text{Effect}_{\text{intended}}
\end{equation}

\textbf{Beschreibung:}
Interventionen haben oft unbeabsichtigte Nebenwirkungen:
\begin{itemize}
    \item Positiver Spillover: Andere Verhaltensweisen verbessern sich
    \item Negativer Spillover: Moral Licensing, Ausweichverhalten
    \item Spillover auf Nicht-Zielpersonen
\end{itemize}

\textbf{BCM-Relevanz:}
Ganzheitliche Evaluation. Eine Intervention, die das Zielverhalten verbessert, aber Schlimmeres verursacht, ist netto negativ.

\textbf{BEATRIX-Relevanz:}
BEATRIX prüft auf Spillover:
\begin{itemize}
    \item Welche anderen Verhaltensweisen könnten beeinflusst werden?
    \item Spillover-Metriken in Evaluation aufnehmen
    \item Warnung bei hohem negativen Spillover-Risiko
\end{itemize}

\textbf{Konsequenz bei Fehlen:}
Nur Zieleffekt betrachtet. Systemische Auswirkungen ignoriert.

\textbf{Validiert:} META

\textbf{Abhängigkeiten:} MA-ONT-22, MA-GOV-21

\textbf{Literatur:} Dolan \& Galizzi (2015)

%==============================================================================
% MA-GOV-24: Langfristeffekte beachten
%==============================================================================

\subsubsection*{MA-GOV-24: Langfristeffekte beachten}

\textbf{Statement:} Kurzfristige Wirkung $\neq$ Langfristige Wirkung.

\textbf{Formel:}
\begin{equation}
\text{Effect}(\tau = \text{short}) \neq \text{Effect}(\tau = \text{long})
\end{equation}

\textbf{Beschreibung:}
Interventionen können unterschiedliche Kurz- und Langzeiteffekte haben:
\begin{itemize}
    \item Initial stark, dann Decay (NUDGE)
    \item Initial schwach, dann Wachstum (BOOST)
    \item Langfristige Backlash-Risiken
\end{itemize}

\textbf{BCM-Relevanz:}
Temporale Differenzierung. Kurze Studien überschätzen möglicherweise langfristige Wirkung.

\textbf{BEATRIX-Relevanz:}
BEATRIX modelliert zeitliche Dynamik:
\begin{itemize}
    \item Erwartete Effektkurve über Zeit
    \item Langzeit-Follow-Up empfehlen
    \item Warnung bei bekannten Decay-Patterns
\end{itemize}

\textbf{Konsequenz bei Fehlen:}
Kurzfristeffekte extrapoliert. Systematische Überschätzung von Langzeitwirkung.

\textbf{Validiert:} META

\textbf{Abhängigkeiten:} MA-ONT-05, MA-GOV-21

\textbf{Literatur:} Allcott \& Rogers (2014); Brandon et al. (2017)

%==============================================================================
% MA-GOV-25: Governance evolviert
%==============================================================================

\subsubsection*{MA-GOV-25: Governance evolviert}

\textbf{Statement:} Regeln und Constraints passen sich an.

\textbf{Formel:}
\begin{equation}
\Gamma(\tau + \Delta\tau) = f(\Gamma(\tau), \text{Feedback})
\end{equation}

\textbf{Beschreibung:}
Governance ist nicht statisch:
\begin{itemize}
    \item Gesetze werden reformiert
    \item Normen wandeln sich
    \item Best Practices entwickeln sich
\end{itemize}
Das GOV-Modul selbst ist Gegenstand von Evolution.

\textbf{BCM-Relevanz:}
Meta-Dynamik. Die Regeln für Interventionen ändern sich. Ethische Standards entwickeln sich.

\textbf{BEATRIX-Relevanz:}
BEATRIX berücksichtigt Governance-Evolution:
\begin{itemize}
    \item Sind aktuelle Best Practices noch aktuell?
    \item Trends in der Interventionsforschung beobachten
    \item Empfehlungen können veralten
\end{itemize}

\textbf{Konsequenz bei Fehlen:}
Statische Governance-Annahmen. Keine Anpassung an veränderte Rahmenbedingungen.

\textbf{Validiert:} META

\textbf{Abhängigkeiten:} MA-ONT-25

\textbf{Literatur:} Ostrom (1990); North (1990)

\bigskip
\noindent\rule{\textwidth}{0.4pt}
\textit{Ende der GOV-Axiome. Alle 25 Governance-Axiome sind nun vollständig dokumentiert.}

% =============================================================================
% DYN-AXIOME (DYNAMIK) - Vollständige Ausformulierung
% BCM 2.2 META-Modul
% =============================================================================

\section{Dynamik-Axiome (MA-DYN)}

Die DYN-Axiome definieren zeitliche Prozesse, Veränderungen und Evolution.
Sie beantworten die Frage: \textit{Wie verändert es sich?}

%==============================================================================
% MA-DYN-01: Zeit ist kontinuierlich
%==============================================================================

\subsubsection*{MA-DYN-01: Zeit ist kontinuierlich}

\textbf{Statement:} Verhalten findet in kontinuierlicher Zeit statt.

\textbf{Formel:}
\begin{equation}
\tau \in \mathbb{R}^+
\end{equation}

\textbf{Beschreibung:}
Zeit ist nicht diskret, sondern fließend:
\begin{itemize}
    \item Zustände ändern sich stetig
    \item Übergänge sind graduell
    \item Differentialgleichungen sind angemessene Modelle
\end{itemize}

\textbf{BCM-Relevanz:}
Ermöglicht dynamische Modellierung mit kontinuierlichen Funktionen. Habituation, Lernen und Evolution sind als stetige Prozesse modellierbar.

\textbf{BEATRIX-Relevanz:}
BEATRIX arbeitet mit diskreten Zeitpunkten (Snapshot), aber:
\begin{itemize}
    \item Interpolation zwischen Zeitpunkten möglich
    \item Trajektorien statt Zustände modellieren
    \item Zeitliche Glättung bei Messungen
\end{itemize}

\textbf{Konsequenz bei Fehlen:}
Nur diskrete Zustände. Keine kontinuierliche Dynamik. Sprunghafte Übergänge statt gradueller.

\textbf{Validiert:} WTX

\textbf{Abhängigkeiten:} MA-ONT-05

\textbf{Literatur:} Newton (1687); Strogatz (2015)

%==============================================================================
% MA-DYN-02: Diskrete Events
%==============================================================================

\subsubsection*{MA-DYN-02: Diskrete Events}

\textbf{Statement:} Entscheidungen sind diskrete Ereignisse in kontinuierlicher Zeit.

\textbf{Formel:}
\begin{equation}
\text{Decision}_i \text{ at } \tau_i \in T
\end{equation}

\textbf{Beschreibung:}
Obwohl Zeit kontinuierlich ist, sind Entscheidungen punktuell:
\begin{itemize}
    \item Kaufentscheidung: Zeitpunkt $\tau_k$
    \item Wahlentscheidung: Zeitpunkt $\tau_w$
    \item Sequenz von Decisions: $\{\tau_1, \tau_2, ..., \tau_n\}$
\end{itemize}

\textbf{BCM-Relevanz:}
Verbindet kontinuierliche Dynamik mit diskreten Handlungen. Interventionen zielen auf Entscheidungszeitpunkte.

\textbf{BEATRIX-Relevanz:}
BEATRIX fokussiert auf Decision Points:
\begin{itemize}
    \item Wann fällt die Entscheidung?
    \item Interventions-Timing optimieren
    \item Sequenz von Entscheidungen modellieren
\end{itemize}

\textbf{Konsequenz bei Fehlen:}
Keine Unterscheidung Prozess/Event. Entscheidungen wären stetig statt punktuell.

\textbf{Validiert:} WTX

\textbf{Abhängigkeiten:} MA-DYN-01

\textbf{Literatur:} Ross (1983)

%==============================================================================
% MA-DYN-03: Pfadabhängigkeit
%==============================================================================

\subsubsection*{MA-DYN-03: Pfadabhängigkeit}

\textbf{Statement:} Geschichte beeinflusst Gegenwart.

\textbf{Formel:}
\begin{equation}
\text{State}(\tau) = f(\text{History}(0..\tau))
\end{equation}

\textbf{Beschreibung:}
Der aktuelle Zustand ist nicht nur von aktuellen Bedingungen abhängig:
\begin{itemize}
    \item Frühere Entscheidungen schaffen Pfade
    \item Lock-In-Effekte: QWERTY-Tastatur
    \item Gewohnheiten sind akkumulierte Geschichte
\end{itemize}

\textbf{BCM-Relevanz:}
Erklärt Persistenz von Verhalten. Verhaltensänderung erfordert ``Pfadbruch''. Frühe Intervention hat höheren Hebel.

\textbf{BEATRIX-Relevanz:}
BEATRIX berücksichtigt Verhaltensgeschichte:
\begin{itemize}
    \item Wie lange besteht das Verhalten schon?
    \item Welche Lock-Ins existieren?
    \item Historische Daten für bessere Vorhersage nutzen
\end{itemize}

\textbf{Konsequenz bei Fehlen:}
Zustand nur von aktuellen Bedingungen abhängig. Markov-Annahme. Geschichte irrelevant.

\textbf{Validiert:} WTX

\textbf{Abhängigkeiten:} MA-DYN-01

\textbf{Literatur:} Arthur (1989); David (1985)

%==============================================================================
% MA-DYN-04: Multiple Equilibria
%==============================================================================

\subsubsection*{MA-DYN-04: Multiple Equilibria}

\textbf{Statement:} Es kann mehrere stabile Zustände geben.

\textbf{Formel:}
\begin{equation}
|\{\text{State}: \frac{d\text{State}}{dt} = 0\}| \geq 1
\end{equation}

\textbf{Beschreibung:}
Systeme können mehrere Gleichgewichte haben:
\begin{itemize}
    \item Kooperation vs. Defektion als stabile Zustände
    \item ``Gute'' vs. ``schlechte'' Gewohnheit
    \item Gesellschaftliche Normen A vs. B
\end{itemize}
Welches Gleichgewicht erreicht wird, hängt von Anfangsbedingungen ab.

\textbf{BCM-Relevanz:}
Erklärt, warum ähnliche Populationen unterschiedliche Normen haben können. Interventionen können Gleichgewichtswechsel auslösen.

\textbf{BEATRIX-Relevanz:}
BEATRIX analysiert Gleichgewichte:
\begin{itemize}
    \item In welchem Gleichgewicht befindet sich die Persona/Gruppe?
    \item Ist ein Wechsel zu einem besseren Gleichgewicht möglich?
    \item Welche Intervention löst den Wechsel aus?
\end{itemize}

\textbf{Konsequenz bei Fehlen:}
Einziges Gleichgewicht. Keine alternativen Attraktoren. Deterministische Konvergenz.

\textbf{Validiert:} WTX

\textbf{Abhängigkeiten:} MA-DYN-03

\textbf{Literatur:} Schelling (1978); Young (1998)

%==============================================================================
% MA-DYN-05: Hysterese
%==============================================================================

\subsubsection*{MA-DYN-05: Hysterese möglich}

\textbf{Statement:} Systemzustand hängt von der Richtung der Veränderung ab.

\textbf{Formel:}
\begin{equation}
\text{State}(\uparrow) \neq \text{State}(\downarrow) \quad \text{bei gleichem Trigger}
\end{equation}

\textbf{Beschreibung:}
Der Weg zählt:
\begin{itemize}
    \item Gewohnheit aufbauen ist leichter als abbauen
    \item Vertrauen aufbauen dauert, zerstören geht schnell
    \item Norm-Etablierung vs. Norm-Erosion asymmetrisch
\end{itemize}

\textbf{BCM-Relevanz:}
Asymmetrie von Aufbau und Abbau. Positive Verhaltensänderung kann andere Schwellen haben als Rückfall.

\textbf{BEATRIX-Relevanz:}
BEATRIX modelliert Hysterese:
\begin{itemize}
    \item Unterschiedliche Schwellen für Veränderung vs. Rückfall
    \item Stabilisierung nach Verhaltensänderung nötig
    \item Warnung vor ``zu schnellem'' Absetzen von Interventionen
\end{itemize}

\textbf{Konsequenz bei Fehlen:}
Symmetrische Dynamik angenommen. Richtungsunabhängigkeit.

\textbf{Validiert:} WTX

\textbf{Abhängigkeiten:} MA-DYN-03

\textbf{Literatur:} Scheffer (2009)

%==============================================================================
% MA-DYN-06: Tipping Points
%==============================================================================

\subsubsection*{MA-DYN-06: Tipping Points existieren}

\textbf{Statement:} Es gibt kritische Schwellen für qualitative Zustandswechsel.

\textbf{Formel:}
\begin{equation}
\exists \theta: \text{State}(x < \theta) \neq \text{State}(x > \theta) \quad \text{qualitativ}
\end{equation}

\textbf{Beschreibung:}
Kleine Änderungen können große Wirkung haben -- wenn sie eine Schwelle überschreiten:
\begin{itemize}
    \item Epidemiologischer Tipping Point: $R > 1$
    \item Soziale Tipping Points: Norm-Kaskaden
    \item Individuelle Tipping Points: Verhaltensumbruch
\end{itemize}

\textbf{BCM-Relevanz:}
Erklärt Nicht-Linearität. Kleine Interventionen können große Effekte haben -- wenn sie den Tipping Point erreichen.

\textbf{BEATRIX-Relevanz:}
BEATRIX identifiziert Tipping Points:
\begin{itemize}
    \item Wie nah ist das System am Tipping Point?
    \item Welche Intervention erreicht den Kipppunkt?
    \item Warnung vor instabilen Zuständen
\end{itemize}

\textbf{Konsequenz bei Fehlen:}
Lineare Extrapolation. Sprunghafte Veränderungen unerwartet.

\textbf{Validiert:} WTX

\textbf{Abhängigkeiten:} MA-DYN-04

\textbf{Literatur:} Gladwell (2000); Scheffer et al. (2012)

%==============================================================================
% MA-DYN-07: Gewohnheitsformation
%==============================================================================

\subsubsection*{MA-DYN-07: Gewohnheitsformation}

\textbf{Statement:} Wiederholung führt zu Automatisierung.

\textbf{Formel:}
\begin{equation}
P(V_{\text{auto}}|\text{Repetition}) \uparrow \text{ mit } n
\end{equation}

\textbf{Beschreibung:}
Gewohnheiten bilden sich durch Wiederholung:
\begin{itemize}
    \item Ca. 66 Tage für automatisches Verhalten (Lally et al.)
    \item Exponentieller Anstieg der Automatisierung
    \item Plateau nach Stabilisierung
\end{itemize}

\textbf{BCM-Relevanz:}
Ziel vieler Interventionen: Gewohnheitsbildung. Einmalige Verhaltensänderung ist weniger wertvoll als dauerhafte.

\textbf{BEATRIX-Relevanz:}
BEATRIX plant für Gewohnheitsbildung:
\begin{itemize}
    \item Wie oft muss das Verhalten wiederholt werden?
    \item Welche Cues etablieren die Gewohnheit?
    \item Tracking der Gewohnheitsstärke
\end{itemize}

\textbf{Konsequenz bei Fehlen:}
Keine Automatisierung. Jede Entscheidung deliberativ.

\textbf{Validiert:} WTX

\textbf{Abhängigkeiten:} MA-EPI-21

\textbf{Literatur:} Lally et al. (2010); Wood \& Rünger (2016)

%==============================================================================
% MA-DYN-08: Habit Loop
%==============================================================================

\subsubsection*{MA-DYN-08: Habit Loop}

\textbf{Statement:} Gewohnheiten folgen dem Cue-Routine-Reward-Muster.

\textbf{Formel:}
\begin{equation}
\text{Habit} = f(\text{Cue}, \text{Routine}, \text{Reward})
\end{equation}

\textbf{Beschreibung:}
Der Habit Loop nach Duhigg:
\begin{enumerate}
    \item \textbf{Cue}: Trigger (Zeit, Ort, Emotion, andere Personen)
    \item \textbf{Routine}: Das automatische Verhalten
    \item \textbf{Reward}: Belohnung, die den Loop verstärkt
\end{enumerate}

\textbf{BCM-Relevanz:}
Interventionsansatz: Bestehenden Cue mit neuer Routine verbinden. ``Habit stacking''. Reward engineering.

\textbf{BEATRIX-Relevanz:}
BEATRIX analysiert Habit Loops:
\begin{itemize}
    \item Welche Cues triggern welches Verhalten?
    \item Welche Rewards verstärken es?
    \item Wie kann der Loop umgestaltet werden?
\end{itemize}

\textbf{Konsequenz bei Fehlen:}
Keine Struktur für Gewohnheitsänderung. Nur globale Verhaltensänderung statt gezielter Manipulation.

\textbf{Validiert:} WTX

\textbf{Abhängigkeiten:} MA-DYN-07

\textbf{Literatur:} Duhigg (2012); Wood \& Neal (2007)

%==============================================================================
% MA-DYN-09: Commitment Devices
%==============================================================================

\subsubsection*{MA-DYN-09: Commitment Devices}

\textbf{Statement:} Selbstbindung ist möglich und kann rational sein.

\textbf{Formel:}
\begin{equation}
U(\text{Commit}) > U(\text{NoCommit}) \quad \text{wenn } \delta_{\text{now}} > \delta_{\text{future}}
\end{equation}

\textbf{Beschreibung:}
Menschen können sich selbst binden, um zukünftige Schwäche zu überwinden:
\begin{itemize}
    \item Ulysses und die Sirenen
    \item Sparverträge mit Strafe bei vorzeitiger Auflösung
    \item Öffentliche Ankündigungen
\end{itemize}

\textbf{BCM-Relevanz:}
Lösung für zeitinkonsistente Präferenzen (MA-DYN-14). Commitment ist keine Irrationalität, sondern Metarationalität.

\textbf{BEATRIX-Relevanz:}
BEATRIX empfiehlt Commitment Devices:
\begin{itemize}
    \item Welche Form der Selbstbindung passt zur Persona?
    \item Soziale vs. finanzielle Commitments
    \item Stärke des Commitments kalibrieren
\end{itemize}

\textbf{Konsequenz bei Fehlen:}
Keine Lösung für Present Bias. Zeitinkonsistenz bleibt unaddressiert.

\textbf{Validiert:} WAX

\textbf{Abhängigkeiten:} MA-EPI-16

\textbf{Literatur:} Elster (1979); Bryan et al. (2010)

%==============================================================================
% MA-DYN-10: Soziale Diffusion
%==============================================================================

\subsubsection*{MA-DYN-10: Soziale Diffusion}

\textbf{Statement:} Verhalten breitet sich sozial aus.

\textbf{Formel:}
\begin{equation}
P(V_i|V_{\text{neighbors}}) = f(\text{Density}(V_{\text{neighbors}}))
\end{equation}

\textbf{Beschreibung:}
Verhalten ist ansteckend:
\begin{itemize}
    \item Peer Effects: Freunde beeinflussen Verhalten
    \item Netzwerk-Diffusion: Von Early Adopters zur Mehrheit
    \item Kaskaden: Kleine Anfänge, große Wirkung
\end{itemize}

\textbf{BCM-Relevanz:}
Grundlage für SOCIAL-Interventionen. Der $\lambda$-Parameter hängt von wahrgenommener Kooperation anderer ab.

\textbf{BEATRIX-Relevanz:}
BEATRIX modelliert Netzwerkeffekte:
\begin{itemize}
    \item Wer sind die relevanten Peers der Persona?
    \item Welche Diffusionsstrategien sind effektiv?
    \item Seed-Auswahl für maximale Verbreitung
\end{itemize}

\textbf{Konsequenz bei Fehlen:}
Individualistisches Modell. Soziale Ansteckung nicht abbildbar.

\textbf{Validiert:} KNU

\textbf{Abhängigkeiten:} MA-ONT-10

\textbf{Literatur:} Rogers (2003); Centola \& Macy (2007)

%==============================================================================
% MA-DYN-11: Bereits in axiom_beispiele.tex
%==============================================================================

%==============================================================================
% MA-DYN-12: Norm-Kaskaden
%==============================================================================

\subsubsection*{MA-DYN-12: Norm-Kaskaden}

\textbf{Statement:} Normen können rapid kippen.

\textbf{Formel:}
\begin{equation}
\frac{dN}{dt} \text{ kann sprunghaft sein}
\end{equation}

\textbf{Beschreibung:}
Normenwandel ist oft nicht graduell:
\begin{itemize}
    \item Raucher-Norm: von ``cool'' zu ``uncool'' in einer Generation
    \item Gleichgeschlechtliche Ehe: rapide Akzeptanzsteigerung
    \item ``Preference falsification'' → plötzliches Offenlegen
\end{itemize}

\textbf{BCM-Relevanz:}
Dynamische Ergänzung zu MA-ONT-24 (Normen evolvieren). Nicht nur dass, sondern wie: Sprünge statt lineare Trends.

\textbf{BEATRIX-Relevanz:}
BEATRIX erkennt Kaskadenrisiken:
\begin{itemize}
    \item Ist eine Norm-Kaskade im Gange?
    \item Kann eine Intervention die Kaskade auslösen?
    \item Timing für maximalen Impact
\end{itemize}

\textbf{Konsequenz bei Fehlen:}
Nur gradueller Normenwandel modelliert. Plötzliche Umbrüche überraschen.

\textbf{Validiert:} KNU

\textbf{Abhängigkeiten:} MA-ONT-24, MA-DYN-11

\textbf{Literatur:} Kuran (1995); Bicchieri (2016)

%==============================================================================
% MA-DYN-13: Discounting
%==============================================================================

\subsubsection*{MA-DYN-13: Discounting}

\textbf{Statement:} Zukünftiger Nutzen wird diskontiert.

\textbf{Formel:}
\begin{equation}
U_{\text{present}}(U_{\text{future}}) = \delta^t \cdot U_{\text{future}}
\end{equation}

\textbf{Beschreibung:}
Ein Euro heute ist mehr wert als ein Euro morgen:
\begin{itemize}
    \item Zeitpräferenz: Gegenwart wird vorgezogen
    \item Unsicherheit: Zukunft ist ungewiss
    \item Opportunity Cost: Verzicht auf heutige Nutzung
\end{itemize}

\textbf{BCM-Relevanz:}
Fundamental für intertemporale Entscheidungen. Erklärt Unterinvestition in Zukunft (Sparen, Prävention, Klimaschutz).

\textbf{BEATRIX-Relevanz:}
BEATRIX schätzt Diskontrate:
\begin{itemize}
    \item Wie stark diskontiert die Persona?
    \item Kurzfristige Belohnungen vs. langfristiger Nutzen
    \item Interventionen mit optimiertem zeitlichen Profil
\end{itemize}

\textbf{Konsequenz bei Fehlen:}
Keine Zeitpräferenz. Gegenwart und Zukunft gleichwertig.

\textbf{Validiert:} INU

\textbf{Abhängigkeiten:} MA-EPI-16

\textbf{Literatur:} Samuelson (1937); Frederick et al. (2002)

%==============================================================================
% MA-DYN-14: Bereits in axiom_beispiele.tex
%==============================================================================

%==============================================================================
% MA-DYN-15: Projection Bias
%==============================================================================

\subsubsection*{MA-DYN-15: Projection Bias}

\textbf{Statement:} Aktuelle Zustände werden in die Zukunft projiziert.

\textbf{Formel:}
\begin{equation}
E[U_{\text{future}}|\text{State}_{\text{now}}] \text{ biased toward State}_{\text{now}}
\end{equation}

\textbf{Beschreibung:}
Menschen überschätzen die Persistenz aktueller Zustände:
\begin{itemize}
    \item Hungrig einkaufen: Zu viel kaufen
    \item Gute Laune: Optimistische Prognosen
    \item Schlechte Laune: Pessimistische Prognosen
\end{itemize}

\textbf{BCM-Relevanz:}
Erklärt systematische Fehlplanung. Zukünftige Zustände (Hunger, Motivation, Energie) werden falsch antizipiert.

\textbf{BEATRIX-Relevanz:}
BEATRIX korrigiert für Projection Bias:
\begin{itemize}
    \item In welchem Zustand befindet sich die Persona jetzt?
    \item Wie würde sie in anderem Zustand entscheiden?
    \item Interventionen nicht in extremen Zuständen planen
\end{itemize}

\textbf{Konsequenz bei Fehlen:}
Aktuelle Zustände würden als repräsentativ für zukünftige gelten.

\textbf{Validiert:} INU

\textbf{Abhängigkeiten:} MA-EPI-06

\textbf{Literatur:} Loewenstein et al. (2003)

%==============================================================================
% MA-DYN-16: Hot-Cold Empathy Gap
%==============================================================================

\subsubsection*{MA-DYN-16: Hot-Cold Empathy Gap}

\textbf{Statement:} In ``cold'' Zuständen werden ``hot'' Zustände unterschätzt.

\textbf{Formel:}
\begin{equation}
E[U_{\text{hot}}|\text{State}_{\text{cold}}] < U_{\text{hot, actual}}
\end{equation}

\textbf{Beschreibung:}
Die Intensität emotionaler Zustände wird nicht antizipiert:
\begin{itemize}
    \item Nüchtern: Unterschätzen von Alkohol-Verlangen
    \item Satt: Unterschätzen von Hunger
    \item Ruhig: Unterschätzen von Wut
\end{itemize}

\textbf{BCM-Relevanz:}
Erklärt Scheitern von ``cold'' geplanten Vorsätzen in ``hot'' Situationen. Commitment Devices helfen, wenn Cold-Self den Hot-Self bindet.

\textbf{BEATRIX-Relevanz:}
BEATRIX berücksichtigt Hot-Cold Gap:
\begin{itemize}
    \item Plane Interventionen für den ``hot'' Zustand
    \item Commitments in ``cold'' Zuständen eingehen lassen
    \item Warnung vor Hot-State-Entscheidungen
\end{itemize}

\textbf{Konsequenz bei Fehlen:}
Emotionale Zustände würden korrekt antizipiert. Vorsätze würden immer halten.

\textbf{Validiert:} INU, AWX

\textbf{Abhängigkeiten:} MA-DYN-15

\textbf{Literatur:} Loewenstein (1996, 2005)

%==============================================================================
% MA-DYN-17: Ego Depletion
%==============================================================================

\subsubsection*{MA-DYN-17: Ego Depletion}

\textbf{Statement:} Selbstkontrolle ist eine erschöpfliche Ressource.

\textbf{Formel:}
\begin{equation}
\text{Willpower}(t) \downarrow \text{ mit Nutzung}
\end{equation}

\textbf{Beschreibung:}
Selbstkontrolle funktioniert wie ein Muskel:
\begin{itemize}
    \item Ermüdet durch Nutzung
    \item Regeneriert durch Ruhe
    \item Kann durch Training gestärkt werden
\end{itemize}
Hinweis: Effekt ist umstritten (Replication Crisis), aber praktisch relevant.

\textbf{BCM-Relevanz:}
Erklärt Selbstkontrollversagen am Abend, nach Stress. Interventionen sollten Willpower-Status berücksichtigen.

\textbf{BEATRIX-Relevanz:}
BEATRIX berücksichtigt Ego Depletion:
\begin{itemize}
    \item Wann ist die Persona wahrscheinlich depleted?
    \item Interventionen mit geringem Willpower-Bedarf bevorzugen
    \item Timing für Willpower-intensive Entscheidungen
\end{itemize}

\textbf{Konsequenz bei Fehlen:}
Konstante Selbstkontrolle angenommen. Ermüdung ignoriert.

\textbf{Validiert:} WAX, WTX

\textbf{Abhängigkeiten:} MA-EPI-05

\textbf{Literatur:} Baumeister et al. (1998); (umstritten: Carter \& McCullough, 2014)

%==============================================================================
% MA-DYN-18: Implementation Intentions
%==============================================================================

\subsubsection*{MA-DYN-18: Implementation Intentions}

\textbf{Statement:} Wenn-Dann-Pläne erhöhen Umsetzungswahrscheinlichkeit.

\textbf{Formel:}
\begin{equation}
P(V|\text{If-Then Plan}) > P(V|\text{Goal only})
\end{equation}

\textbf{Beschreibung:}
Konkrete Wenn-Dann-Pläne schließen die Intention-Action-Gap:
\begin{itemize}
    \item ``Wenn X, dann tue ich Y''
    \item Automatische Aktivierung durch Cue
    \item Reduziert deliberativen Aufwand
\end{itemize}

\textbf{BCM-Relevanz:}
Effektive Intervention für WTX-Defizite. Kombiniert gut mit NUDGE (Default als Implementation Intention).

\textbf{BEATRIX-Relevanz:}
BEATRIX empfiehlt Implementation Intentions:
\begin{itemize}
    \item Konkrete Wenn-Dann-Formulierungen generieren
    \item Passende Cues für die Persona identifizieren
    \item Als Teil von BOOST-Interventionen
\end{itemize}

\textbf{Konsequenz bei Fehlen:}
Nur Ziele ohne Umsetzungspläne. Intention-Action-Gap bleibt.

\textbf{Validiert:} WTX

\textbf{Abhängigkeiten:} MA-EPI-11

\textbf{Literatur:} Gollwitzer (1999); Gollwitzer \& Sheeran (2006)

%==============================================================================
% MA-DYN-19: Adaptation
%==============================================================================

\subsubsection*{MA-DYN-19: Adaptation}

\textbf{Statement:} Nutzen adaptiert an neuen Status Quo.

\textbf{Formel:}
\begin{equation}
U(\text{State}) \rightarrow U_{\text{baseline}} \quad \text{as } \tau \rightarrow \infty
\end{equation}

\textbf{Beschreibung:}
Menschen gewöhnen sich an alles:
\begin{itemize}
    \item Positive Veränderungen: Freude verblasst
    \item Negative Veränderungen: Leid verblasst
    \item Neuer Normalzustand wird Referenzpunkt
\end{itemize}

\textbf{BCM-Relevanz:}
Erklärt das Hedonic Treadmill (MA-DYN-20). Dauerhafte Nutzenänderung ist schwierig zu erreichen.

\textbf{BEATRIX-Relevanz:}
BEATRIX berücksichtigt Adaptation:
\begin{itemize}
    \item Erwartete Adaptationsrate modellieren
    \item Nachhaltige vs. temporäre Nutzeneffekte
    \item Intervention gegen Adaptation designen
\end{itemize}

\textbf{Konsequenz bei Fehlen:}
Permanente Nutzenveränderungen angenommen. Überschätzung langfristiger Effekte.

\textbf{Validiert:} INU

\textbf{Abhängigkeiten:} MA-ONT-20

\textbf{Literatur:} Helson (1964); Wilson \& Gilbert (2005)

%==============================================================================
% MA-DYN-20: Hedonic Treadmill
%==============================================================================

\subsubsection*{MA-DYN-20: Hedonic Treadmill}

\textbf{Statement:} Glück kehrt zu einer individuellen Baseline zurück.

\textbf{Formel:}
\begin{equation}
\text{Happiness} \rightarrow \text{Baseline} \quad \text{nach positiven/negativen Events}
\end{equation}

\textbf{Beschreibung:}
Das hedonistische Laufband:
\begin{itemize}
    \item Lottogewinner: Nach 1 Jahr wieder auf Ausgangsniveau
    \item Querschnittslähmung: Nach 2 Jahren signifikante Erholung
    \item Set Point Theory: Genetisch determinierte Baseline
\end{itemize}

\textbf{BCM-Relevanz:}
Relativiert materielle Interventionen. Mehr Geld macht nicht dauerhaft glücklicher. Fokus auf prozessuale statt ergebnisbasierte Interventionen.

\textbf{BEATRIX-Relevanz:}
BEATRIX berücksichtigt Hedonic Adaptation:
\begin{itemize}
    \item Welche Interventionen haben nachhaltige Effekte?
    \item Prozess-Glück vs. Ergebnis-Glück
    \item Warnung vor Überschätzung materieller Anreize
\end{itemize}

\textbf{Konsequenz bei Fehlen:}
Lineare Beziehung Ressourcen → Glück. Materielle Interventionen überschätzt.

\textbf{Validiert:} INU

\textbf{Abhängigkeiten:} MA-DYN-19

\textbf{Literatur:} Brickman \& Campbell (1971); Diener et al. (2006)

%==============================================================================
% MA-DYN-21: Peak-End Rule
%==============================================================================

\subsubsection*{MA-DYN-21: Peak-End Rule}

\textbf{Statement:} Erinnerung = Peak + End.

\textbf{Formel:}
\begin{equation}
\text{Recalled}_U = f(\max_U, \text{End}_U)
\end{equation}

\textbf{Beschreibung:}
Die Erinnerung an eine Erfahrung wird dominiert von:
\begin{itemize}
    \item Dem intensivsten Moment (Peak)
    \item Dem Endmoment
    \item Dauer spielt kaum eine Rolle
\end{itemize}

\textbf{BCM-Relevanz:}
Erklärt, warum Erfahrungen anders erinnert werden als erlebt. Interventionen können das Ende optimieren für bessere Gesamtbewertung.

\textbf{BEATRIX-Relevanz:}
BEATRIX nutzt Peak-End Rule:
\begin{itemize}
    \item Positive Enden für Interventionserfahrungen designen
    \item Negative Peaks vermeiden
    \item Dauer ist weniger wichtig als Qualität der Momente
\end{itemize}

\textbf{Konsequenz bei Fehlen:}
Dauer-gewichtete Durchschnitte angenommen. Peak und End-Effekte ignoriert.

\textbf{Validiert:} INU

\textbf{Abhängigkeiten:} MA-EPI-12

\textbf{Literatur:} Kahneman et al. (1993); Kahneman (1999)

%==============================================================================
% MA-DYN-22: Identitäts-Shifts
%==============================================================================

\subsubsection*{MA-DYN-22: Identitäts-Shifts}

\textbf{Statement:} Identitäten können sich fundamental ändern.

\textbf{Formel:}
\begin{equation}
I(\tau + \Delta\tau) \not\subset I(\tau) \quad \text{möglich}
\end{equation}

\textbf{Beschreibung:}
Nicht nur Evolution, sondern Transformation:
\begin{itemize}
    \item Konversion: Religiöse, politische
    \item Life Events: Elternschaft, Scheidung, Krankheit
    \item ``Ich bin nicht mehr derselbe Mensch''
\end{itemize}

\textbf{BCM-Relevanz:}
Ermöglicht Modellierung fundamentaler Identitätsänderungen. IDENTITY-Interventionen können auf Shift zielen, nicht nur Salienz.

\textbf{BEATRIX-Relevanz:}
BEATRIX erkennt Identitäts-Shift-Potenzial:
\begin{itemize}
    \item Steht die Persona vor einem Life Event?
    \item Gibt es Anzeichen für Identitätskrise?
    \item Interventionen für Übergangszeiten optimieren
\end{itemize}

\textbf{Konsequenz bei Fehlen:}
Nur graduelle Identitätsänderung. Fundamentale Transformationen nicht modellierbar.

\textbf{Validiert:} IDN

\textbf{Abhängigkeiten:} MA-ONT-23

\textbf{Literatur:} James (1902); Paloutzian (2005)

%==============================================================================
% MA-DYN-23: λ-Dynamik
%==============================================================================

\subsubsection*{MA-DYN-23: $\lambda$-Dynamik}

\textbf{Statement:} Der Kooperationsparameter $\lambda$ verändert sich über Zeit.

\textbf{Formel:}
\begin{equation}
\lambda(\tau + \Delta\tau) = f(\lambda(\tau), \text{Experience}, \text{Context})
\end{equation}

\textbf{Beschreibung:}
$\lambda$ ist nicht statisch:
\begin{itemize}
    \item Positive Kooperationserfahrungen erhöhen $\lambda$
    \item Betrug und Enttäuschung senken $\lambda$
    \item Gesellschaftlicher Wandel beeinflusst Basis-$\lambda$
\end{itemize}

\textbf{BCM-Relevanz:}
Dynamische Erweiterung des $\lambda$-Konzepts. Erklärt, wie Kooperationsbereitschaft sich entwickelt.

\textbf{BEATRIX-Relevanz:}
BEATRIX modelliert $\lambda$-Entwicklung:
\begin{itemize}
    \item Wie hat sich $\lambda$ entwickelt?
    \item Welche Erfahrungen beeinflussen es?
    \item Kann $\lambda$ durch Intervention erhöht werden?
\end{itemize}

\textbf{Konsequenz bei Fehlen:}
Statisches $\lambda$ als Trait. Keine Kooperations-Entwicklung modellierbar.

\textbf{Validiert:} WAX

\textbf{Abhängigkeiten:} MA-ONT-14

\textbf{Literatur:} Ostrom (1998); Fehr \& Gächter (2000)

%==============================================================================
% MA-DYN-24: System-Evolution
%==============================================================================

\subsubsection*{MA-DYN-24: System-Evolution}

\textbf{Statement:} Das BCM-Gesamtsystem evolviert.

\textbf{Formel:}
\begin{equation}
\Phi(\tau + \Delta\tau) = \Phi(\tau) + \mu \cdot \Delta\Phi
\end{equation}

\textbf{Beschreibung:}
Nicht nur Parameter ändern sich, sondern das System als Ganzes:
\begin{itemize}
    \item Neue Verhaltenskategorien entstehen (Digital Behavior)
    \item Neue Interventionstypen werden entdeckt
    \item Systemgrenzen verschieben sich
\end{itemize}

\textbf{BCM-Relevanz:}
Meta-Ebene: Das BCM selbst ist nicht statisch. Neue Forschung erweitert das Modell. Version 2.2 ≠ Version 3.0.

\textbf{BEATRIX-Relevanz:}
BEATRIX wird kontinuierlich aktualisiert:
\begin{itemize}
    \item Neue Erkenntnisse integrieren
    \item Modellparameter rekalibrieren
    \item Systementwicklung dokumentieren
\end{itemize}

\textbf{Konsequenz bei Fehlen:}
Statisches Modell. Keine Weiterentwicklung vorgesehen.

\textbf{Validiert:} META

\textbf{Abhängigkeiten:} MA-ONT-25

\textbf{Literatur:} Kuhn (1962); Lakatos (1978)

%==============================================================================
% MA-DYN-25: Meta-Stabilität
%==============================================================================

\subsubsection*{MA-DYN-25: Meta-Stabilität}

\textbf{Statement:} Fundamentale Struktur bleibt erhalten, auch wenn Inhalte sich ändern.

\textbf{Formel:}
\begin{equation}
\frac{\partial \text{Structure}}{\partial \tau} \rightarrow 0 \quad \text{auch wenn } \frac{\partial \text{Content}}{\partial \tau} \neq 0
\end{equation}

\textbf{Beschreibung:}
Trotz aller Dynamik gibt es Invarianten:
\begin{itemize}
    \item Menschen haben Präferenzen (welche auch immer)
    \item Soziale Normen existieren (welche auch immer)
    \item Biases sind systematisch (welche auch immer)
\end{itemize}
Die Struktur des BCM bleibt stabil, die Inhalte variieren.

\textbf{BCM-Relevanz:}
Ontologisches Fundament. Das BCM ist nicht nur ein historisches Snapshot, sondern beschreibt tiefere Strukturen, die invariant sind.

\textbf{BEATRIX-Relevanz:}
BEATRIX verlässt sich auf Meta-Stabilität:
\begin{itemize}
    \item Strukturannahmen sind robust
    \item Parameter-Aktualisierung, nicht Modell-Austausch
    \item Vertrauen in Grundarchitektur
\end{itemize}

\textbf{Konsequenz bei Fehlen:}
Alles wäre in Fluss. Kein stabiler Rahmen für Modellierung. Beliebigkeit.

\textbf{Validiert:} META

\textbf{Abhängigkeiten:} MA-DYN-24

\textbf{Literatur:} Eigen (1971); Kauffman (1993)

\bigskip
\noindent\rule{\textwidth}{0.4pt}
\textit{Ende der DYN-Axiome. Alle 25 Dynamik-Axiome sind nun vollständig dokumentiert.}


\end{document}
